\documentclass[11pt,letterpaper]{article}
\usepackage[T1]{fontenc}
\usepackage[utf8]{inputenc}
\usepackage{amsmath,amsthm,mathtools}
\usepackage{newtxtext,newtxmath}
\usepackage[margin=1in,headheight=15pt]{geometry}
\usepackage{microtype,booktabs,array,longtable,enumitem,xurl}
\usepackage{xcolor,fancyhdr,needspace}
\usepackage[colorlinks=true,linkcolor=ink,citecolor=ink,urlcolor=ink]{hyperref}
\usepackage{aliascnt}
\usepackage[nameinlink,noabbrev]{cleveref}
\definecolor{ink}{HTML}{354E42}
\hypersetup{pdftitle={Surreal Fields Across Set-Theoretic Universes},
 pdfauthor={Merged research report prepared with ChatGPT},
 pdfsubject={External saturation of old surreal fields, fresh-sign gaps, first new birthdays, surcomplex real forms, a Boolean completion of reverse simplicity, and full gap spectra, relative grounds and classification (Part VIII)}}
\pagestyle{fancy}\fancyhf{}
\fancyhead[L]{\small\scshape Surreal fields across universes}
\fancyhead[R]{\small Set theory and model theory}
\fancyfoot[C]{\thepage}
\renewcommand{\headrulewidth}{0.3pt}
\setlength{\parskip}{3.5pt}
\setlength{\parindent}{1.15em}
\setlength{\emergencystretch}{2em}
\setlist{leftmargin=1.7em,itemsep=2pt,topsep=5pt}
\clubpenalty=10000\widowpenalty=10000
\allowdisplaybreaks[2]
\numberwithin{equation}{section}
\newtheorem{theorem}{Theorem}[section]
\newaliascnt{lemma}{theorem}
\newtheorem{lemma}[lemma]{Lemma}
\aliascntresetthe{lemma}
\crefname{lemma}{lemma}{lemmas}
\Crefname{lemma}{Lemma}{Lemmas}
\newaliascnt{proposition}{theorem}
\newtheorem{proposition}[proposition]{Proposition}
\aliascntresetthe{proposition}
\crefname{proposition}{proposition}{propositions}
\Crefname{proposition}{Proposition}{Propositions}
\newaliascnt{corollary}{theorem}
\newtheorem{corollary}[corollary]{Corollary}
\aliascntresetthe{corollary}
\crefname{corollary}{corollary}{corollaries}
\Crefname{corollary}{Corollary}{Corollaries}
\theoremstyle{definition}
\newaliascnt{definition}{theorem}
\newtheorem{definition}[definition]{Definition}
\aliascntresetthe{definition}
\crefname{definition}{definition}{definitions}
\Crefname{definition}{Definition}{Definitions}
\newaliascnt{example}{theorem}
\newtheorem{example}[example]{Example}
\aliascntresetthe{example}
\crefname{example}{example}{examples}
\Crefname{example}{Example}{Examples}
\theoremstyle{remark}
\newaliascnt{remark}{theorem}
\newtheorem{remark}[remark]{Remark}
\aliascntresetthe{remark}
\crefname{remark}{remark}{remarks}
\Crefname{remark}{Remark}{Remarks}
\newaliascnt{question}{theorem}
\newtheorem{question}[question]{Question}
\aliascntresetthe{question}
\crefname{question}{question}{questions}
\Crefname{question}{Question}{Questions}
\crefname{part}{part}{parts}
\Crefname{part}{Part}{Parts}
\newcommand{\No}{\mathbf{No}}
\newcommand{\SC}{\mathrm{SC}}
\newcommand{\NoM}{\No^{M}}
\newcommand{\NoN}{\No^{N}}
\newcommand{\SCM}{\SC^{M}}
\newcommand{\SCN}{\SC^{N}}
\newcommand{\Ord}{\mathrm{Ord}}
\newcommand{\R}{\mathbb R}
\newcommand{\C}{\mathbb C}
\newcommand{\Q}{\mathbb Q}
\newcommand{\Pow}{\mathcal P}
\newcommand{\bd}{\operatorname{b}}
\newcommand{\cf}{\operatorname{cf}}
\newcommand{\ci}{\operatorname{ci}}
\newcommand{\rc}{\operatorname{rc}}
\newcommand{\dom}{\operatorname{dom}}
\newcommand{\ran}{\operatorname{ran}}
\newcommand{\Fix}{\operatorname{Fix}}
\newcommand{\Aut}{\operatorname{Aut}}
\newcommand{\Cone}{\operatorname{Cone}}
\newcommand{\Fresh}{\operatorname{Fresh}}
\newcommand{\GapSpec}{\operatorname{GapSpec}}
\newcommand{\code}{\operatorname{code}}
\newcommand{\codeplus}{\operatorname{code}^{+}}
\newcommand{\Add}{\operatorname{Add}}
\newcommand{\res}{\mathbin{\upharpoonright}}
\newcommand{\cat}{\mathbin{{}^{\frown}}}
\newcommand{\NS}{\mathsf{NS}}
\newcommand{\RCF}{\mathsf{RCF}}
\newcommand{\ACF}{\mathsf{ACF}_0}
\newcommand{\GBC}{\mathsf{GBC}}
\newcommand{\ZFC}{\mathsf{ZFC}}
\newcommand{\Lang}{\mathcal L}
\newcommand{\abs}[1]{\lvert#1\rvert}
\newcommand{\cut}[2]{\left\{\,#1\;\middle|\;#2\,\right\}}
\newcommand{\seq}[1]{\langle#1\rangle}
\newcommand{\hatord}[1]{\widehat{#1}}
\newcommand{\shorthash}{\texttt{4f26455}}
\newcommand{\fullhash}{\texttt{4f264559\allowbreak 9121fa87\allowbreak 2c710499\allowbreak 5e47f86f\allowbreak 0382351f}}
\newcommand{\repo}{\texttt{VladimirReshetnikov/Surreal}}
% Part VIII (sources 09--14, added 28 September 2026)
\newcommand{\Char}{\operatorname{Char}}
\newcommand{\FF}{\mathfrak F}
\newcommand{\Fr}{\operatorname{Fr}}
\newcommand{\FRESH}{\operatorname{FRESH}}
\newcommand{\Reg}{\mathrm{Reg}}
\newcommand{\Fn}{\operatorname{Fn}}
\newcommand{\Coll}{\operatorname{Coll}}
\newcommand{\Ecl}{\operatorname{Ecl}}
\newcommand{\Tr}{\operatorname{Tr}_{\R}}
\newcommand{\st}{\operatorname{st}}
\newcommand{\dcode}{\operatorname{dcode}}
\newcommand{\bcode}{\operatorname{bcode}}
\newcommand{\GCH}{\mathsf{GCH}}
\newcommand{\CH}{\mathsf{CH}}
\newcommand{\forces}{\Vdash}
\newcommand{\gspin}{\texttt{e21766d}}
\newcommand{\gsfull}{\texttt{e21766d0\allowbreak 4c2b8a9b\allowbreak 2cdba0cd\allowbreak 43e56302\allowbreak 4b3bd1b9}}
\newcolumntype{P}[1]{>{\raggedright\arraybackslash}p{#1}}

\begin{document}
\hypersetup{pageanchor=false}
\begin{titlepage}
\centering
\vspace*{0.2in}
{\large\scshape Merged research report\par}
\vspace{0.35in}
{\Huge\bfseries\color{ink} Surreal Fields Across\par Set-Theoretic Universes\par}
\vspace{0.2in}
{\Large Omitted cuts, fresh-sign gaps and saturation spectra of old surreal
fields; first new birthdays; surcomplex real forms;\par
and a Boolean completion of reverse simplicity\par}
\vspace{0.25in}
{22 September 2026\par}
\vspace{0.25in}
\begin{minipage}{0.94\textwidth}
\begin{abstract}
Let $M\subseteq N$ be transitive universes of ZFC with the same ordinals, and
view the old surreal field $\NoM$ as a proper class of $N$. All sizes, types
and parameter sets are computed in $N$. For every uncountable cardinal
$\kappa$ of $N$, the ordered field $\NoM$ is $\kappa$-saturated exactly when
$N$ has no new ordinal-valued sequences of length less than $\kappa$. The
same condition is equivalent to interpolation of every set-sized cut with
just one side of size less than $\kappa$. The least length $\delta$ of a new
sequence is a regular cardinal of both universes. A fresh sequence of length
$\delta$ yields an omitted gap of exact cofinality pair $(\delta,\delta)$, by
three explicit constructions: a nested interval tree, an absolute interval
code with all ordinals as labels, and ordinal block codes of signs.

Every gap of $\NoM$ with set-sized cofinal and coinitial witnesses comes from a
unique fresh sign, a new sign sequence all of whose proper prefixes are old.
Its fillers are exactly the sign extensions of that fresh sign, and every
such gap is symmetric. The first new birthday $\beta$ satisfies
$\delta\le\cf^N\beta\le\beta$. At the countable boundary, $\omega$-saturation
detects exactly the new reals. The pure field $\NoM[i]$ stays saturated over
every $N$-set, while naming conjugation restores the real spectrum. The old
field is closed and nowhere dense in $\NoN$, and each $N$-set of old numbers is
uniformly discrete. We also treat bounded birthday fragments and the Cohen,
$\Add(\kappa,1)$ and Prikry extensions. Under Global Choice, we give
nonconjugate real forms of $\NoN[i]$ none of which has a set-sized cofinal
subset. Finally we give a set-complete class Boolean algebra in which reverse
simplicity is dense; class-completeness is impossible under Global Choice.
\end{abstract}
\end{minipage}
\vfill
\begin{minipage}{0.94\textwidth}\small
\textbf{Status and provenance.} An AI-assisted research report merged from
four manuscripts of 22 September 2026 (sources 01, 03, 04 and 08 of the batch
that placed them), each pinned to repository commit \shorthash; see
\cref{univ:sub:provenance}. The principal results are proposed contributions
with written proofs, not certified priority claims. The report has not been
refereed, and no statement in it is Lean-verified. Classical surreal facts,
quantifier elimination, standard forcing facts, and the class back-and-forth
and conjugacy criteria already in the collection are credited, not claimed.
The finite Python checks test finite sign and interval formulas only.

\textbf{Addition of 28 September 2026.} Part~VIII merges six further
manuscripts (sources 09--14, pinned to ProveIt commit \gspin): full gap
spectra of old surreal fields after forcing, the spectra of intermediate
grounds, and old fields with equal full gap spectra that are not isomorphic
(\cref{univ:gs:sub:provenance}). It is equally unrefereed and not Lean-verified.
\end{minipage}
\end{titlepage}
\hypersetup{pageanchor=true}
\pagenumbering{roman}
\setcounter{tocdepth}{2}
{\small\setlength{\parskip}{0pt}\tableofcontents}
\clearpage
\pagenumbering{arabic}

\section{The question, the answers, and how this report was assembled}\label{univ:sec:overview}

\subsection{The question}

Forcing preserves ordinals, but it can add new functions on old ordinals.
Consequently, it can add new surreal sign sequences without changing the
ordinal labels used for birthdays. This raises a more delicate question than
whether new surreal numbers exist:
\begin{quote}
How much of the old surreal field's ability to realize cuts and first-order
types survives when the surrounding universe acquires new sets?
\end{quote}
An old field can remain an elementary subfield of a larger surreal field while
losing saturation. Elementarity tests individual first-order formulas with
finitely many parameters. Saturation also tests whole sets of formulas, and
those sets can be new even when every formula and every parameter is old.

The familiar cut principle says that every pair of \emph{sets} $A<B$ of
surreals has a separator. The word ``set'' hides a parameter: which universe
decides what counts as a set? When $M$ is enlarged to $N$ without changing the
ordinals, every old surreal number remains available, but a new set of old
numbers need not have belonged to $M$. The old cut principle promises nothing
for such a collection. This report measures exactly what is lost.

\subsection{Standing hypotheses and the two levels of assumptions}\label{univ:sub:levels}

\begin{enumerate}[label=\textbf{(\Alph*)},leftmargin=2.2em]
\item[\textbf{(H)}] $N\models\ZFC$, and $M$ is a parameter-definable
transitive inner class of $N$ that contains every ordinal of $N$ and satisfies
$\ZFC$ (as a scheme). Equivalently for our purposes, $M$ may be an available
class predicate in a class theory that has the Separation and Replacement
instances for that predicate, or a two-universe metatheory with these closure
properties stipulated. Forcing extensions $N=M[G]$ are the intended examples,
but no proof uses forcing except in \cref{univ:sec:examples} and
\cref{univ:cor:finitetower}.
\item[\textbf{(GC)}] In addition, the ambient theory is $\GBC$ with Global
Choice, and $M$ is a class of that class universe; $M$ need not be definable.
\end{enumerate}
Everything in this report is proved under \textbf{(H)} except the following,
which also use \textbf{(GC)}: the class isomorphisms and real forms of
\cref{univ:part:realforms}, and the impossibility of a class-complete Boolean
algebra (\cref{univ:thm:booleanno}). Sources~03 and~08 were written with
\textbf{(GC)} as their ambient theory. Their own dependency statements confine
Global Choice to the class back-and-forth and real-form arguments (source~03,
its audit table; source~08, its foundations section). Every proof taken from
them for a statement outside \cref{univ:part:realforms} was re-read for this
merge. It uses only set-sized choices, Choice inside $M$, and Replacement for
the predicate $M$. One re-reading was needed: source~03 quantifies over gaps
of $\NoM$ as class partitions. Under \textbf{(H)} a gap is instead given by its
set presentation (\cref{univ:def:setgap}), and the statements of
\cref{univ:part:gaps} become schemes over such presentations.

External means computed in $N$, never in a larger metatheory. If $M$ and $N$
are represented by countable transitive \emph{set} models inside a larger
universe, all claims about sets, cardinalities and saturation must still be
read inside $N$. They are not claims of saturation with respect to every
subset in that larger universe. We use no unrestricted truth predicate for
class structures. In the languages used here, quantifier elimination reduces
formulas to finite algebraic conditions and so supplies the satisfaction
relation directly. Types are ordinary $N$-sets of finite syntactic
expressions over set-sized parameter sets. This avoids the general
class-satisfaction obstacle discussed by van den Dries--Ehrlich and by
Ehrlich \cite{vdDE,Ehrlich1989}.

The Prikry examples assume a measurable cardinal in $M$. The
$\Add(\kappa,1)$ examples need no cardinal arithmetic; source~04's extra
hypothesis $\kappa^{<\kappa}=\kappa$ is used only for its additional
preservation clause (\cref{univ:ex:addkappa}).

\subsection{Conventions and notation}\label{univ:sub:notation}

The four sources used different letters for the same objects, and several
letters for two objects. This report fixes one symbol per notion.

\begin{center}\small
\begin{longtable}{@{}P{0.33\textwidth}P{0.2\textwidth}P{0.38\textwidth}@{}}
\toprule
Notion & Symbol here & In the sources\\
\midrule
\endhead
Inner and outer universe & $M\subseteq N$ & $M\subseteq N$ (01, 03); $M\subseteq W$ (04); $M\subseteq V$ (08)\\
Old and new surreal fields & $\NoM$, $\NoN$ & $K,K^+$ (01); $F,E$ (03); $S,T$ (04); $F_M,F_V$ (08)\\
Old and new surcomplex fields & $\SCM=\NoM[i]$, $\SCN=\NoN[i]$ & $C,C^+$ (01); $K_M,K_N$ (03); $K_M,K_W$ (04); $K_M,K_V$ (08)\\
Canonical conjugations & $c_M$, $c_N$ & $c$ (01); $c_M,c_N$ (03); $c_M,c_W$ (04); $c_M,c_V$ (08)\\
Transported conjugation & $j=\Phi c_M\Phi^{-1}$ & $j_M$ (03); $d$ (04); $c'$ (08)\\
First new length of an ordinal-valued sequence & $\delta=\delta(M,N)$ & $\delta$ (01, 08); $\nu$ (03); $d(M,W)$ (04)\\
First new birthday & $\beta=\beta(M,N)$ & $\beta$ (08)\\
No new short sequences & $\NS_{<\kappa}(M,N)$ & same (01, 08); $\mathsf D_{<\kappa}(M,W)$ (04)\\
Birthday (length) of a sign & $\bd(s)$ & $\bd$ (01, 04); $\ell(s)$ (03); $\operatorname{lh}(s)$ (08)\\
Sign cone of $s$ & $\Cone(s)$ & $(L_s,R_s)$ separators (01); $\Cone(s)$ (03, 08); $C_s$ (04)\\
Ordinal block codes & $\code(f)$, $\codeplus(f)$ & $C(f)$ (08); $s_f$ (03)\\
Measurable cardinal (Prikry) & $\rho$ & $\rho$ (01); $\lambda$ (03); $\nu$ (04); $\theta$ (08)\\
Regular cardinal of $\Add(\kappa,1)$ & $\kappa$ & $\rho$ in $\operatorname{Fn}_{<\rho}(\rho,2)$ (01); $\kappa$ (03, 04, 08)\\
Gap with set-sized witnesses & set-presented gap & set-character gap (01); set-presented gap (03); Dedekind gap with character (08)\\
Cut property & $\eta_\kappa$ & $\eta_\kappa$ (01, 03, 04); $\mathrm{Cut}_{<\kappa}$ (08)\\
\bottomrule
\end{longtable}
\end{center}

The letter $K$ is not used for the proper-class fields $\NoM$ and $\SCM$:
most reports of the collection reserve $K$ for a \emph{set-sized} field, and
the collection's notation guide recommends $\SC=\No[i]$ for the surcomplex
class (\path{docs/NOTATION.md}). The letter $\delta$ always denotes the first
new length of an ordinal-valued sequence, never a birthday; the first new
birthday is $\beta$, as in source~08. The sibling report
\path{docs/foundations-and-computation/birthday-cutoffs-and-hereditary-sets}
uses the same letter $\beta$ for the first new birthday. Several of that
report's source manuscripts wrote $\delta$ for the first new birthday; that
$\delta$ is the $\beta$ of this report, not its $\delta$. The
letter $\theta$ denotes an ordinal bound, on the range of a sequence or on a
set of birthdays, and $\rho$ is reserved for the measurable cardinal of the
Prikry examples. The letter $\nu$ is not used.

\emph{Tempting false readings.} ``Gap'' here always means a gap of the class
order $\NoM$ seen from $N$, with set-sized witnesses; it is not the
proper-class ``finite--infinite gap'' of $\No$ itself studied in
\path{docs/surreal/genetic-gaps-and-primitives}. ``External'' means computed
in the outer universe $N$; in the foundations report the same word refers to
a stronger metatheory (\cref{univ:rem:foundations}). $\eta_\kappa$ is indexed
by a size bound, not by an aleph subscript. $\beta(M,N)$ is an ordinal
attached to a pair of universes, not the birthday of a particular number.

\subsection{The main results}\label{univ:sub:main}

For an infinite cardinal $\kappa$ of $N$, write
\begin{equation}\label{univ:eq:NSdef}
 \NS_{<\kappa}(M,N)
 \quad\Longleftrightarrow\quad
 ({}^\lambda\theta)^N\subseteq M
 \quad\text{for all ordinals }\lambda<\kappa\text{ and }\theta.
\end{equation}
Every ordinal-valued sequence has bounded range, so this is exactly the
absence of new ordinal-valued sequences of length less than $\kappa$.
The range bound $\theta$ is \emph{not} required to be below $\kappa$.

\begin{theorem}[Main spectrum theorem; sources 01, 03, 04, 08]\label{univ:thm:main-summary}
For every infinite cardinal $\kappa$ of $N$, the following are equivalent.
\begin{enumerate}[label=\textnormal{(\roman*)}]
\item $\NS_{<\kappa}(M,N)$ holds.
\item Every pair of $N$-sets $L,R\subseteq\NoM$ with $L<R$ and
$\min(\abs L,\abs R)<\kappa$ has a separator in $\NoM$.
\item Every such pair with $\abs{L\cup R}<\kappa$ has a separator in $\NoM$.
\end{enumerate}
For uncountable $\kappa$, these are also equivalent to $\kappa$-saturation of
$\NoM$ in the ordered-field language, and to $\kappa$-saturation of the
order $(\NoM,<)$ as a dense linear order, both evaluated in $N$. When they
fail, an omitted pair as in \textnormal{(iii)} can be chosen inside $(0,1)$.
\end{theorem}

Condition (ii), from source~01, is stronger in appearance than the usual
small-cut condition: the other side can be an arbitrarily large set. Neither
side may be a proper class. In particular a cut is always filled if one side
is an actual ground-model set (\cref{univ:thm:oneside}). The two-sided form
(iii) and the saturation clause were proved independently by all four
sources; the dense-linear-order clause is source~08's and the $(0,1)$ clause
source~04's.

\begin{theorem}[First loss and language sensitivity; sources 01, 03, 04, 08]\label{univ:thm:package-summary}
Suppose $M\ne N$, and let $\delta$ be the least ordinal length of a new
ordinal-valued sequence. Then:
\begin{enumerate}[label=\textnormal{(\roman*)}]
\item $\delta$ is an infinite regular cardinal in both $M$ and $N$.
\item $\NoM$ has a set-presented gap of exact cofinality pair
$(\delta,\delta)$, and no set-presented gap has either cofinality smaller
than $\delta$.
\item For uncountable $N$-cardinals $\kappa$, $\NoM$ is $\kappa$-saturated
exactly when $\kappa\leq\delta$.
\item $\NoM$ is $\omega$-saturated exactly when $\R^M=\R^N$.
\item $\SCM$ in the pure field language is saturated over every $N$-set of
parameters. The expansion $(\SCM,c_M)$ by $c_M(a+bi)=a-bi$ has exactly the
saturation spectrum of $\NoM$, including at $\omega$.
\end{enumerate}
\end{theorem}

Beyond this common spine, each source contributed results that the others do
not have.
\begin{itemize}
\item \textbf{Fresh signs classify the gaps} (source~03,
\cref{univ:part:gaps}). The set-presented gaps of $\NoM$ correspond
bijectively to the \emph{fresh signs}, new sign sequences all of whose proper
prefixes are old. The gap of a fresh sign $s$ has both cofinalities equal to
$\cf^N\bd(s)$, and its fillers in $\NoN$ are exactly the extensions of $s$.
Hence every set-presented gap is symmetric, and $\NoN\setminus\NoM$ is the
disjoint union of the fresh cones.
\item \textbf{The first new birthday} (source~08, \cref{univ:part:beta}). The
least birthday $\beta$ of a new surreal is a cardinal of both universes, and
$\delta\le\cf^N\beta\le\beta$.
\item \textbf{Bounded birthday fragments} (source~01,
\cref{univ:sec:fragments}). For an uncountable regular cardinal $\kappa$ of
$M$ that remains a cardinal in $N$,
\begin{equation}\label{univ:eq:bounded-summary}
\begin{split}
 (\No_{<\kappa})^M\text{ is }\kappa\text{-saturated in }N
 \quad\Longleftrightarrow\quad&
 \cf^N(\kappa)=\kappa\ \text{ and}\\[-2pt]
 &({}^\alpha 2)^N=({}^\alpha 2)^M\quad(\alpha<\kappa).
\end{split}
\end{equation}
Thus the full-class theorem and the bounded-fragment theorem test different
preservation properties.
\item \textbf{Topology} (sources 01, 03, 04, \cref{univ:sec:topology}). Every
$N$-set of old numbers is uniformly discrete and closed, and set-indexed Cauchy
nets of old numbers are eventually constant. If $M\ne N$, the old field is
closed and nowhere dense in $\NoN$, and its subspace topology is its own order
topology.
\item \textbf{Real forms under Global Choice} (sources 03, 04, 08,
\cref{univ:part:realforms}). $\SCM\cong\SCN$ as pure fields, over any $N$-set
of common parameters. Transporting $c_M$ gives involutions of $\SCN$ that are
not conjugate to $c_N$, although neither fixed field has a set-sized cofinal
subset. Several such involutions can be separated by the gap invariant, and
they can agree with $c_N$ on every ordinary complex number.
\item \textbf{A Boolean completion} (source~01, \cref{univ:sec:boolean}). The
reverse-simplicity order on signs embeds densely in an explicit set-complete
class Boolean algebra. Under Global Choice no class-complete such algebra
exists.
\end{itemize}

\subsection{What is offered as new, and what is not}\label{univ:sub:new}

The proposed contributions are the exact external spectrum of
\cref{univ:thm:main-summary} with its one-sided form; the three explicit
codings of a fresh sequence by an omitted cut (the arbitrary-alphabet interval
tree of source~01, the absolute all-ordinal interval code of source~04, and
the ordinal block codes of sources~03 and~08); the exact gap threshold; the
fresh-sign classification of all set-presented gaps; the separation of the
first new birthday from the first saturation obstruction; the distinction
between $\omega$- and $\omega_1$-saturation under no-real extensions; the
resulting surcomplex language comparison; the bounded-fragment criterion; the
external topological statements; the real-form invariants; and the Boolean
completion. Each source stated its own contributions as proposed, with
priority not certified.

Their ingredients are not claimed as new: Conway's cut theorem and sign
expansions, real closedness of $\No$ and algebraic closedness of $\No[i]$
\cite{Conway,Gonshor}; quantifier elimination for real closed and
algebraically closed fields \cite{Marker}; saturation and homogeneous
universality of surreal class models \cite{Ehrlich1989,vdDE2018}; standard
forcing facts \cite{Jech,Prikry,Gitik,Benhamou}; the forcing notion of a fresh
function \cite{FKW}; and, from the collection, class back-and-forth for
proper-class algebraically closed fields, the conjugacy criteria for real
forms, and a real form with a countable cofinal subset
\cite{RepoAutomorphisms}. Several statements are deliberately labelled
consequences rather than independent discoveries. Regularity of the first new
length is an elementary set-theoretic fact, not a new invariant.

The Boolean section answers a concrete question in Berenbeim's 2020
\emph{notes}, not a conjecture of a peer-reviewed paper. Its construction is
a natural direct limit and may be folklore in class forcing; its value here is
an explicit, size-correct answer, not a claim that Boolean completions were
previously unavailable. General class-forcing completion theory predates
those notes \cite{HolyEtAl}. Berenbeim also proposes studying surreals under
forcing in the final section of the same notes \cite[pp.~32, 34]{Berenbeim};
this report develops one concrete direction of that program and does not
settle the notes' broader suggestions about embeddings of universes.

\subsection{Four manuscripts, one report}\label{univ:sub:provenance}

This report is merged from four manuscripts, all dated 22 September 2026 and
all written against commit \fullhash{} of \repo. The source numbers are those
of the batch in which they arrived. Delivered code and data files keep those
numbers as prefixes. No source manuscript is shipped.

{\small
\begin{longtable}{@{}P{0.07\textwidth}P{0.33\textwidth}P{0.52\textwidth}@{}}
\toprule
Source & Manuscript & Contribution here\\
\midrule
\endhead
01 (base) & \emph{Forcing, Omitted Cuts, and Old Surreal Fields: ordinal-sequence
spectra, surcomplex conjugation, and a Boolean completion of reverse
simplicity} (25 pages) & The foundations and spine
(\cref{univ:sec:foundations,univ:sec:bounds,univ:sec:delta,univ:sec:decoder,univ:sec:spectrum,univ:sec:omega}),
the one-old-side theorem, the interval tree, the surcomplex dichotomy
(\cref{univ:sec:complex}), bounded fragments (\cref{univ:sec:fragments}), the
forcing framework (\cref{univ:sec:examples}), external Cauchy rigidity, the
Boolean completion (\cref{univ:sec:boolean}), and most of \cref{univ:sec:scope}.
LaTeX \texttt{Makefile}, as \path{code/01-forcing-omitted-cuts-Makefile}.\\
03 & \emph{Surreal Fields Across Set-Theoretic Universes: fresh gaps, exact
saturation thresholds, and surcomplex real forms} (21 pages) & All of
\cref{univ:part:gaps} (fresh-sign gap classification, filling cones, the
terminal-plus block code, the least-gap theorem); old ordinal bounds
(\cref{univ:lem:oldbounds}); the ordinal-bound form of the cut criterion; the
$\Add(\kappa,1)$ first-gap calculation without cardinal arithmetic; the Prikry
fresh sign; the proper-class ACF saturation theorem; the real-form invariant,
agreement with $c_N$ on $\C^N$, several ground real forms and finite forcing
towers; the closed, nowhere-dense and subspace-topology statements. Files
prefixed \texttt{03-fresh-sign-gaps-}.\\
04 & \emph{Surreal Fields Across Set-Theoretic Universes: exact saturation,
forcing, and surcomplex real forms} (25 pages) & The absolute all-ordinal
interval code with closed finite-branch formulas (second route A,
\cref{univ:sub:absolutecode}); the $(0,1)$ clause; the forcing
characterization for preserved uncountable cardinals; $\Add(\kappa,1)$ under
$\kappa^{<\kappa}=\kappa$; the elementarily equivalent nonconjugate
involutions with different spectra and the several-spectra corollary;
external discreteness in the product topology, open prefix cones, and
nowhere-denseness of $\SCM$ in $\SCN$. Files prefixed
\texttt{04-universe-saturation-}.\\
08 & \emph{Forcing and Saturation of Surreal and Surcomplex Fields: exact
preservation criteria and a real-form dichotomy} (24 pages) & The ordinal
block code with its two-bound cone identity and exact separator class (second
route B, \cref{univ:sub:blockcode}); the dense-linear-order clause; the first
new birthday and $\delta\le\cf^N\beta\le\beta$ (\cref{univ:part:beta}); the
Cohen, $\Add(\kappa,1)$ and Prikry values of $\beta$; the lottery-sum caution;
the non-definability of conjugation in the pure field; the saturated
back-and-forth route and the isomorphism over an arbitrary old set. Files
prefixed \texttt{08-forcing-saturation-}.\\
\bottomrule
\end{longtable}}

\textbf{Why one report.} All four manuscripts prove the same spine in the same
setting: the cut criterion, the regular first new length, an omitted
$(\delta,\delta)$ gap, the $\omega$ boundary, the pure/conjugation contrast,
the distributivity corollary and the Cohen and Prikry examples. Their extras
are complementary. \textbf{Source~01 is the base}: its hypotheses are the
weakest (plain ZFC with a definable inner model, Global Choice only for one
negative Boolean result), and its main criterion is the strongest in form,
since one small side suffices.

\textbf{Printed once.} The following results were proved by two or more
sources and appear once, with all proving sources named at the statement:
the absoluteness of old arithmetic and elementarity
(\cref{univ:prop:absolute}); ground-set bounds and envelopes
(\cref{univ:lem:groundbound}); coding of short families
(\cref{univ:lem:shortold}); regularity of $\delta$
(\cref{univ:prop:deltaregular}); cuts versus real-closed saturation
(\cref{univ:lem:RCFsat}); the main spectrum
(\cref{univ:thm:main-summary}); the first gap (\cref{univ:thm:firstgap}); the
$\omega$ criterion (\cref{univ:thm:omega}); pure surcomplex saturation
(\cref{univ:thm:purecomplex}); the conjugation spectrum
(\cref{univ:thm:conjugation}); nonisomorphism (\cref{univ:cor:noniso}); the
distributivity corollary (\cref{univ:cor:distributivity}); the Cohen,
$\Add(\kappa,1)$ and Prikry examples; external discreteness
(\cref{univ:prop:topology}); closedness and nowhere-denseness
(\cref{univ:thm:nowheredense}); the class isomorphism
(\cref{univ:thm:pureiso}); the transported real form
(\cref{univ:thm:transport}); and the several-grounds statements
(\cref{univ:thm:manygrounds}).

\textbf{Marked second routes.} The converse direction of the spine has three
genuinely different constructions. The interval tree of source~01 is printed
as the main route (\cref{univ:sec:decoder}). The absolute interval code of
source~04 (\cref{univ:sub:absolutecode}) and the ordinal block codes of
sources~08 and~03 (\cref{univ:sub:blockcode,univ:prop:blockcodeplus}) are kept
as marked second routes. Two existence proofs for $\delta$ (Mostowski
collapse, sources 01 and 03; least new rank, sources 04 and 08) and two
class back-and-forth routes (algebraic/transcendental extension, sources 03
and 04 as in the collection; saturated elementary maps, source~08) are kept as
well.

\textbf{Where the merge had to choose.} (a) Symbols: every renaming is listed
in \cref{univ:sub:notation}, and none changes a normalization. (b)
Hypotheses: statements are printed at the weakest hypotheses that any source
proved them under. Where sources differ, the difference is recorded at the
statement: regularity of $\delta$ in $M$ (sources 01, 03, 08; source~04 states
$N$ only); cardinal arithmetic for $\Add(\kappa,1)$ (only source~04 assumes
it); the range of $\kappa$ in the cut criterion (every infinite cardinal in
source~01, every infinite ordinal in source~03, uncountable in sources 04 and
08). (c) Open questions: those answered inside the cluster are recorded as
answered, with the answering source (\cref{univ:sec:questions}). (d) Stale
repository statements in the sources are corrected in
\cref{univ:app:audit}. Text written for the merge is confined to this section,
the framing remarks that relate sources and reports, \cref{univ:rem:betafresh},
the question statuses, the non-claims ledger and the appendices. It adds no
mathematical claim beyond the two-line observation of
\cref{univ:rem:betafresh}.

\textbf{Addition of 28 September 2026.} Six further manuscripts of that date,
numbered sources 09--14 in this report's local sequence and pinned to ProveIt
commit \gsfull{} (placed in \texttt{0e53d1034}), answer parts of
\cref{univ:q:cutspectrum,univ:q:completeness,univ:q:families}. They are merged
into one addition, Part~VIII (\cref{univ:gs:part}), placed after
\cref{univ:sec:nonclaims} so that all earlier section, theorem and question
numbers stay fixed; its provenance, bases, notation and non-claims are in
\cref{univ:gs:sec:assembly,univ:gs:sec:nonclaims,univ:gs:sec:audit}. Its labels
carry the sub-prefix \path{univ:gs:}.

\subsection{Relation to the collection}\label{univ:sub:collection}

\textbf{Surcomplex field automorphisms}
(\path{docs/surcomplex/surcomplex-field-automorphisms}, prefix
\path{saut:}). Its Theorem~9.1 (\path{saut:thm:acfhomogeneity}) is the
class back-and-forth used in \cref{univ:part:realforms}; its Section~10
(\path{saut:sec:realforms}) proves that involutions are conjugate exactly
when their fixed real fields are isomorphic, and constructs a real form with a
countable cofinal subset (Theorem~10.2, \path{saut:thm:inequivalentrealform}).
Its Proposition~B.1 (\path{saut:prop:conjugacycriterion}) says that an
involution is conjugate to $c$ if and only if its fixed real field has the
set-cut interpolation property. None of the four sources cites that
proposition. Read in $N$, its easy direction already gives the non-conjugacy
statements of \cref{univ:thm:transport,univ:thm:involutions}, because the
transported fixed field is isomorphic to $\NoM$, which omits an $N$-set cut.
What is new is recorded in \cref{univ:rem:whatisnew}. These results
\emph{continue}, but do not settle, that report's open ``real-form and
continuity classifications'' (its Section~14). Its batch-32 finite-symmetry
sections classify the involutions preserving the valuation, Hahn summation and
$\mathbf{Oz}[i]$ (note after \cref{univ:rem:whatisnew}).

\textbf{First-$\kappa$ coefficients}
(\path{docs/surcomplex/first-kappa-coefficients}, prefix \path{fkc:}; batch 32).
Its Part~II builds, in one universe and without forcing, a proper class of
pairwise nonconjugate pure-field involutions of $\No[i]$ from fields of surreals
with fewer than $\kappa$ terms; this answers one clause of
\cref{univ:q:families} (status note there).

\textbf{Foundations} (\path{docs/foundations-and-computation/foundations},
prefix \path{found:}). \Cref{univ:prop:topology} is an external version of
its Theorem~12.1 (\path{found:thm:discrete}); see \cref{univ:rem:foundations}
for the exact relation, including its hypothesis on bounds and the passage
of its Section~7.7 on external countability.

\textbf{Computable surreals}
(\path{docs/foundations-and-computation/computable-surreals}). Its
\path{prop:topology} is the universe-relative version of set-indexed
convergence, with an index type small in the lower universe. In the present
same-ordinal setting no such smallness is needed, because $\NoM$ is still a
proper class of $N$; the two statements are consistent.

\textbf{Birthday cutoffs and hereditary sets}
(\path{docs/foundations-and-computation/birthday-cutoffs-and-hereditary-sets}),
the sibling report placed in the same batch. It studies bounded birthday
fields \emph{with a birthday symbol}, their interpretation of hereditary sets,
and their behaviour under outer models. It also proves, for every same-ordinal
outer model, that the first new birthday is the least ordinal acquiring a new
subset. This report keeps source~08's statement and proof of that fact,
because source~08 adds cardinality in $N$ and the inequality
$\delta\le\cf^N\beta$ (\cref{univ:rem:sibling}). The two reports study
different structures: here the pure ordered field and its external
saturation, there the birthday-expanded fields and their interpretations.

\textbf{Genetic gaps and primitives}
(\path{docs/surreal/genetic-gaps-and-primitives}) uses ``gap'' for a
proper-class gap of $\No$; the gaps here are set-presented gaps of $\NoM$ seen
from $N$. The same report cites Berenbeim's 2022 dissertation; the 2020 notes
answered in \cref{univ:sec:boolean} are not cited elsewhere in the collection.

\part{Foundations and the common spine}\label{univ:part:spine}

\section{Foundations, notation, and imported facts}\label{univ:sec:foundations}

\subsection{The two universes}

In $N$, $\NoM$ is a proper class: it contains an old all-plus sign of every
ordinal length. The superscript in $\NoM$ means that the sign functions, their
construction and the operations are computed in $M$. It is not an additional
set-sized model assumed to have a cardinality. We never form a set of all
surreals, a proper-class support, or a class-object whose members are all
class maps.

\subsection{Signs, birthdays and cones}

A surreal number is a sign function $s:\alpha\to\{-,+\}$ for an ordinal
$\alpha$; its birthday is $\bd(s)=\alpha$. The empty function represents zero.
At the first place where two signs differ, numerical comparison uses
\[
 -\ <\ \text{termination}\ <\ +.
\]
Write $s\preceq t$ when $s$ is an initial segment (prefix) of $t$, allowing
equality. Simplicity is the proper-prefix relation, not numerical order and not
merely a smaller birthday. The ordinal $\alpha$ is represented by the all-plus
function of length $\alpha$, written $\hatord\alpha$ when confusion with field
arithmetic is possible. Its surreal successor agrees with $\alpha+1$ as an
ordinal; arbitrary ordinal addition and surreal field addition must not be
conflated. Concatenation $\cat$ is ordinal concatenation.

For sets $L,R$ of surreals, $L<R$ means $\ell<r$ for all $\ell\in L$ and
$r\in R$. Empty sides are allowed. A separator means a \emph{strict}
separator: $L<x<R$. The classical cut theorem supplies the unique simplest
separator $\cut LR$, a prefix of every separator
\cite{Conway,Gonshor,vdDE}. We use the birthday bound
\begin{equation}\label{univ:eq:cutbirthday}
 \bd\!\left(\cut LR\right)
 \leq\sup\{\bd(y)+1:y\in L\cup R\},
\end{equation}
with the supremum of the empty set equal to zero. Put
\[
 \Cone(s)=\{t:s\preceq t\}.
\]

\begin{lemma}[Cones are convex; the prefix characterization; sources 03, 04]\label{univ:lem:convex}
\begin{enumerate}[label=\textnormal{(\roman*)}]
\item The cone $\Cone(s)$ is numerically convex. That is, if $s\preceq t$
and $u$ lies numerically between $s$ and $t$, then $s\preceq u$.
\item Let $I$ be a convex class of surreals, let $s\in I$, and suppose no
proper prefix of $s$ belongs to $I$. Then every $x\in I$ extends $s$, and $s$
is the unique simplest element of $I$.
\end{enumerate}
\end{lemma}
\begin{proof}
(i) Two extensions of $s$ have the same signs before $\bd(s)$; anything between
them cannot first disagree with $s$ at an earlier position.

(ii) If $x$ is a proper prefix of $s$, the hypothesis excludes it. If neither
sequence extends the other, their maximal common prefix $u$ is a proper
prefix of $s$. At the first disagreement one sequence has $-$ and the other
$+$, so $u$ lies strictly between them, and convexity gives $u\in I$, a
contradiction. The remaining possibility is that $x$ extends $s$; such an
extension has birthday at least $\bd(s)$, with equality only for $x=s$.
\end{proof}

\begin{lemma}[The cone of a sign; source 01]\label{univ:lem:signcone}
For $s:\alpha\to\{-,+\}$ put
\[
 L_s=\{s\res\gamma:\gamma<\alpha,\ s(\gamma)=+\},\qquad
 R_s=\{s\res\gamma:\gamma<\alpha,\ s(\gamma)=-\}.
\]
A surreal $x$ satisfies $L_s<x<R_s$ if and only if $s\preceq x$.
\end{lemma}
\begin{proof}
If $x$ extends $s$, its comparison with each prefix in the displayed sets is
decided by the next sign of $s$, so the strict inequalities hold. Conversely,
if $x$ terminates at $\gamma<\alpha$ while agreeing up to that point, then
$x=s\res\gamma$ violates the corresponding strict inequality. If it first
disagrees at $\gamma<\alpha$, compare $x$ with the prefix $s\res\gamma$: the
opposite signs place $x$ on the wrong side of it. These exhaust the ways in
which $x$ could fail to extend $s$.
\end{proof}

The numerical bounds
\begin{equation}\label{univ:eq:birthbounds}
 -\hatord{\bd(x)}\leq x\leq\hatord{\bd(x)}
\end{equation}
follow by comparing a sign with the all-minus and all-plus signs of the same
length.

\subsection{Field theory and types}

The full surreal field is real closed, and its extension by a square root of
$-1$ is algebraically closed. Also $\No_{<\lambda}$ is a real closed subfield
whenever $\lambda$ is an epsilon number; in particular for every uncountable
cardinal $\lambda$, since such cardinals satisfy $\omega^\lambda=\lambda$ in
ordinal arithmetic. We use the corrected birthday-fragment results of van den
Dries--Ehrlich and their restatement by Bournez--Guilmant
\cite{vdDE,vdDEerr,BournezGuilmant}. Nothing about exponential functions is
needed.

We use quantifier elimination for real closed fields in
$\Lang_{\mathrm{or}}=\{0,1,+,-,\cdot,<\}$, for algebraically closed fields of
characteristic zero in $\Lang_{\mathrm{fld}}=\{0,1,+,-,\cdot\}$, and for dense
linear orders without endpoints \cite{Marker}. Embeddings between real closed
ordered fields, and between algebraically closed fields of the same
characteristic, are elementary. Over a real closed parameter subfield a
nonalgebraic one-variable type is determined by an order cut; over an
algebraically closed field one-variable types are algebraic or
transcendental.

A class structure is \emph{$\kappa$-saturated} (in $N$) if every complete
one-variable type over an $N$-set of parameters of cardinality less than
$\kappa$ that is finitely satisfiable in the structure is realized there. As
usual this gives the finite-variable version: realize the projection of a
type in $n+1$ variables to the first $n$, then realize the specialized
one-type over the enlarged parameter set. All syntactic data and parameter
sets at every step are sets, and satisfaction is supplied by quantifier
elimination. \emph{$\omega$-saturation} concerns \emph{finite} parameter sets;
countable parameter sets are $\omega_1$-saturation. Even a type over the empty
set can contain countably many formulas.

An ordered class $H$ has the \emph{$\eta_\kappa$ property} if every pair of
$N$-sets $A,B\subseteq H$ with $A<B$ and $\abs{A\cup B}^N<\kappa$ has a strict
separator in $H$. Empty sides are allowed. Following source~03 we allow an
infinite ordinal bound $\kappa$ when needed.

\begin{lemma}[Cuts versus saturation; sources 01, 03, 04, 08]\label{univ:lem:RCFsat}
\begin{enumerate}[label=\textnormal{(\roman*)}]
\item For an uncountable cardinal $\kappa$, a real closed ordered field ---
including the definable class fields used here --- is $\kappa$-saturated if
and only if its order is $\eta_\kappa$.
\item \textup{(source~08)} For a dense linear order without endpoints, the
same equivalence holds for every infinite $\kappa$.
\item \textup{(source~04)} In each universe of set theory, its own full
surreal field is set-saturated as an ordered field.
\end{enumerate}
\end{lemma}
\begin{proof}
(i) Saturation realizes the finitely satisfiable partial type of the
inequalities of a small cut; a complete extension exists by compactness
applied to the set of formulas with those parameters, inside a set-sized real
closed subfield containing them.

Conversely, let $A$ have size less than $\kappa$. The real closure $F$ of
$\Q(A)$ inside the field is a set of cardinality at most
$\max(\abs A,\aleph_0)<\kappa$; there are set many polynomials over $\Q(A)$
and finitely many roots of each. A complete type over $A$ extends to one over
$F$ that is finitely satisfiable in the ambient field. If the extended type
is algebraic, it is realized in $F$. Otherwise it declares, for every $a\in F$,
exactly one of $a<X$ and $X<a$; these inequalities form a separated cut of
size less than $\kappa$. A separator $y$ realizes the entire type: every
nonzero polynomial over the real closed $F$ factors into linear factors and
irreducible quadratics; the quadratics have constant sign, and the signs of the
linear factors are fixed by the cut. Quantifier elimination completes the
argument. Restrict back to $A$.

(ii) A complete one-type of a dense linear order over $A$ either says $X=a$
for some $a\in A$ or specifies a cut of $A$. There is no countable field
closure to add, so the argument works even at $\kappa=\omega$.

(iii) The cut theorem realizes every set-sized separated cut; apply the proof
of (i) with no cardinal bound.
\end{proof}

The uncountability hypothesis in (i) cannot be dropped: the real closure of
even the empty parameter set is countable, not finite. Every dense order
without endpoints has $\eta_\omega$, whereas a real closed field need not be
$\omega$-saturated (\cref{univ:sec:omega}).

\section{Absoluteness and ground-set bounds}\label{univ:sec:bounds}

\begin{proposition}[Old arithmetic and elementary inclusions; sources 01, 03, 04, 08]\label{univ:prop:absolute}
\begin{enumerate}[label=\textnormal{(\roman*)}]
\item $\NoM=\NoN\cap M$ is an initial subclass of the sign tree: every prefix
of an old sign is old. Birthdays and numerical comparison of old signs are
absolute.
\item If $L,R\in M$ are sets of old surreals with $L<R$, then
$\cut LR^M=\cut LR^N$.
\item The inclusion $\NoM\subseteq\NoN$ preserves zero, one, negation,
addition, multiplication and inversion. Every ordinal is old, so $\NoM$ is
cofinal and coinitial in $\NoN$.
\item In $N$, $\NoM$ is real closed and $\SCM$ is algebraically closed, and
\[
 \NoM\preccurlyeq\NoN\quad\text{in }\Lang_{\mathrm{or}},\qquad
 \SCM\preccurlyeq\SCN\quad\text{in }\Lang_{\mathrm{fld}}.
\]
\end{enumerate}
\end{proposition}
\begin{proof}
(i) Being an ordinal-length function into $\{-,+\}$, its domain, and first
disagreement comparison are absolute for transitive models. If $s\in M$ and
$\alpha\leq\bd(s)$, then $\alpha\in M$ and $s\res\alpha\in M$.

(ii) Let $u=\cut LR^M$. It still separates $L,R$ in $N$. The simplest
separator $v$ computed in $N$ is a prefix of $u$ by the simplicity theorem in
$N$, hence old by (i), hence an available separator in $M$, so $u\preceq v$ by
the simplicity theorem in $M$. Thus $u=v$ (source~08). Equivalently, every
proper prefix of $u$ fails in $M$ to separate, witnessed by an old element of
$L$ or $R$, and that witness still works in $N$; \cref{univ:lem:convex}(ii)
applied to the convex class of $N$-separators gives the same conclusion
(source~04). The argument does not minimize birthdays over old candidates,
which would not be absolute since $N$ can contain new short signs; the prefix
witnesses exclude new candidates.

(iii) Conway's addition and multiplication are defined from canonical options
by well-founded, set-like recursion. For old inputs the options are old, the
earlier recursive values agree by induction, the option families belong to
$M$, and (ii) identifies the resulting cuts. Every recursion needed for
particular inputs has a set-sized dependency closure; no class-recursion
principle is assumed. The inverse of an old nonzero number remains its inverse
by uniqueness. A large all-plus or all-minus ordinal lies above or below any
given new surreal.

(iv) A polynomial over $\NoM$ has a finite old coefficient tuple, which
belongs to $M$. In $M$ positive elements have square roots and odd-degree
polynomials have roots, and those witnesses work in $N$. Thus $\NoM$ is
externally real closed, and the same finite-polynomial argument gives external
algebraic closedness of $\SCM$. Quantifier elimination gives both elementary
inclusions.
\end{proof}

This proposition asserts elementarity of the \emph{fields}. It does not
assert $M\preccurlyeq N$ as models of set theory, which is generally false for
a forcing extension. ``Initial'' is a simplicity-tree assertion: $\NoM$ is not
a numerical interval of $\NoN$, and it need not contain a new surreal whose
birthday is below some old birthday.

\noindent[Added 4 October 2026, batch~93: Part~III of
\path{foundations-and-computation/definable-surreals-and-omnific-integers}
(batch-93 manuscript~06; AI-assisted, unrefereed, not formalized) re-proves
(i)--(iii), with Conway's omega map added, as parts (1)--(2) of its
Theorem~34.1 (\path{dsn:op:thm:all-reals-free-absoluteness}), for transitive
models of $\ZFC$, and the ordered-field clause of (iv) as its Corollary~34.5
(\path{dsn:op:cor:inner-rcf-elementarity}). Parts (3)--(5) of that theorem add
normal forms, set-indexed sums and the omnific integers, without assuming that
the two models have the same reals.]

\begin{lemma}[Every external set of old numbers is ground-bounded; sources 01, 04, 08]\label{univ:lem:groundbound}
\begin{enumerate}[label=\textnormal{(\roman*)}]
\item If $A$ is an $N$-set and $A\subseteq\NoM$, there are an ordinal
$\alpha$ and a ground set $B\in M$ with $A\subseteq B=(\No_{<\alpha})^M$.
Every $N$-set of elements of $\SCM$ is likewise contained in a ground set of
old surcomplex numbers, and these ground sets have well-orderings in $M$.
\item Such an $A$ is contained in a set-sized real closed subfield $E\in M$ of
$\NoM$, and an $N$-set of elements of $\SCM$ in a set-sized algebraically
closed subfield $E[i]\in M$; these are elementary substructures. There is an
old element transcendental over any such $E[i]$.
\end{enumerate}
\end{lemma}
\begin{proof}
(i) In $N$ form $\alpha=\sup\{\bd(a)+1:a\in A\}$. It belongs to $M$ because
the universes have the same ordinals. The set of old signs of length below
$\alpha$ exists in $M$ by Power Set and Replacement. For $\SCM$ bound both
coordinates. Choice in $M$ gives the well-orderings.

(ii) Inside $M$, take the field generated by the cover from (i) and its
relative real algebraic closure in $\NoM$; there are set many polynomials and
finitely many roots of each. Quantifier elimination gives elementarity. An
all-plus ordinal above every element of $E$ is not in $E$, hence
transcendental over $E$ (a real closed subfield is relatively algebraically
closed in an ordered extension), hence over $E[i]$.
\end{proof}

The lemma is not a covering property of a forcing notion, and it does not
assert $A\in M$. It holds even when a forcing badly fails familiar small-set
covering properties, and the containing ground set can be vastly larger than
$A$. It works because a set of birthdays is bounded and the two universes
have the same ordinals.

\begin{lemma}[Old bounds for new sets; source 03]\label{univ:lem:oldbounds}
Every $N$-set of elements of $\NoN$ has an upper bound that is an old ordinal.
Consequently every $N$-set of old numbers is bounded above and below in
$\NoM$, and $\NoM$ has no set-sized cofinal or coinitial subset in $N$. Old
positive elements are coinitial in $\NoN_{>0}$.
\end{lemma}
\begin{proof}
For a set $X\subseteq\NoN$, Replacement bounds all birthdays by an ordinal
$\theta$. The all-plus ordinal $\theta+1$ exceeds every member of $X$: at any
earlier minus sign it wins the first comparison, and it exceeds an all-plus
sequence of smaller length. It is old. Apply the same argument to $-X$. For
$r\in\NoN_{>0}$ choose an old positive ordinal $\gamma>1/r$; then
$0<1/\gamma<r$ and $1/\gamma\in\NoM$.
\end{proof}

\begin{lemma}[Coding short families; sources 01, 03, 04]\label{univ:lem:shortold}
Assume $\NS_{<\kappa}(M,N)$. Every $N$-set of old surreal or surcomplex numbers
of cardinality less than $\kappa$ is a member of $M$. More generally, every
sequence of length less than $\kappa$ with values in a fixed ground set is a
ground sequence, and every $N$-subset of cardinality less than $\kappa$ of a
ground set belongs to $M$.
\end{lemma}
\begin{proof}
For a ground set $B$ fix in $M$ an injection $e:B\to\theta$ into an ordinal.
Composing a short sequence in $B$ with $e$ gives an ordinal-valued sequence of
the same length, hence an old one, and decoding with $e$ shows that the
original sequence is old. For an external set $A$ of size less than $\kappa$,
use \cref{univ:lem:groundbound} and an enumeration of $A$ in $N$ of length less
than $\kappa$; the enumeration is old, and so is its range.
\end{proof}

The following asymmetric fact is the efficient direction of the main theorem.

\begin{theorem}[One old side suffices; source 01]\label{univ:thm:oneside}
Let $L,R$ be $N$-sets contained in $\NoM$ with $L<R$. If either $L\in M$ or
$R\in M$, then some $x\in\NoM$ satisfies $L<x<R$.
\end{theorem}
\begin{proof}
Suppose $L\in M$. Choose a ground set $B\supseteq R$ by
\cref{univ:lem:groundbound} and define, in $M$,
$U=\{b\in B:(\forall\ell\in L)\ \ell<b\}$. Order is absolute, so
$R\subseteq U$ and $L<U$. The ground cut $x=\cut LU$ satisfies $L<x<U$, and
therefore $L<x<R$. This also covers empty sides. If $R\in M$, apply the proved
case to $-R<-L$ and negate the separator.
\end{proof}

\begin{corollary}[Source 01]\label{univ:cor:NSoneside}
If $\NS_{<\kappa}(M,N)$ holds, every external set-sized cut with at least one
side of size less than $\kappa$ is filled in $\NoM$.
\end{corollary}
\begin{proof}
The small side belongs to $M$ by \cref{univ:lem:shortold}; apply
\cref{univ:thm:oneside}.
\end{proof}

In particular, every $N$-set of positive elements of $\NoM$ has a positive
lower bound in $\NoM$: take $L=\{0\}$. Every $N$-set of old numbers has an old
strict upper bound. These statements survive every extension under
consideration, including ones that destroy countable saturation.

\section{The first new sequence length}\label{univ:sec:delta}

\begin{definition}\label{univ:def:delta}
When $M\ne N$, the \emph{first new length} is
\begin{equation}\label{univ:eq:delta}
 \delta(M,N)=\min\{\lambda:\exists\theta\
 ({}^\lambda\theta)^N\not\subseteq M\}.
\end{equation}
No value is assigned when $M=N$. A sequence $f:\lambda\to\Ord$ in $N$ is
\emph{fresh over $M$} (source~04) if $f\notin M$ but $f\res\xi\in M$ for every
$\xi<\lambda$. This is the forcing-theoretic notion of a fresh function
\cite{FKW}; the name is not claimed.
\end{definition}

The minimum exists. \emph{First route (sources 01, 03).} Every set of $N$ can
be coded in $N$ by a well-founded extensional relation on an ordinal together
with a marked point, and such a relation is coded by a set of ordinals. If
all sets of ordinals were old, the code and its Mostowski collapse would be
old, so $M=N$. \emph{Second route (sources 04, 08).} Choose a new set $x$ of
least rank $\alpha$. Every set of rank below $\alpha$ is old, so
$x\subseteq V_\alpha^M$, and an $M$-well-ordering of $V_\alpha^M$ identifies
$x$ with a new subset of an ordinal. Either way, a new set of ordinals gives a
new characteristic sequence.

No finite sequence of old ordinals is new, and a sequence of successor length
is determined by its shorter prefix and one old ordinal. Hence
$\delta(M,N)$ is an infinite limit ordinal, and every new sequence of length
$\delta$ is fresh. More generally, a fresh sequence has infinite limit length.

\begin{proposition}[Regularity of the first new length; sources 01, 03, 04, 08]\label{univ:prop:deltaregular}
The ordinal $\delta(M,N)$ is an infinite regular cardinal in both $M$ and $N$.
\end{proposition}
\begin{proof}
Write $\delta=\delta(M,N)$ and choose a new $f:\delta\to\theta$.

If $\delta$ were not an $N$-cardinal, choose in $N$ a bijection
$e:\mu\to\delta$ with $\mu<\delta$. By minimality $e\in M$, and $f\circ e$,
also of length $\mu<\delta$, belongs to $M$. Decoding with $e$ gives $f\in M$,
a contradiction.

If $\cf^N(\delta)=\mu<\delta$, choose a cofinal $e:\mu\to\delta$; it is old by
minimality. Every $f\res e(\xi)$ is old and lies in the ground set
$T=({}^{<\delta}\theta)^M$. Well-order $T$ in $M$. The sequence of codes of
$f\res e(\xi)$, $\xi<\mu$, has length less than $\delta$ and is old; its ground
decoding and union reconstruct $f$, a contradiction. Hence $\delta$ is regular
in $N$. A ground bijection or a shorter ground cofinal sequence would survive
to $N$, so $\delta$ is a regular cardinal of $M$ as well.
\end{proof}

Sources 01, 03 and 08 state regularity in both universes; source~04 states the
$N$ clause, and its proof gives the $M$ clause by the last sentence above.
This is an elementary minimal-length fact about extensions, not a new
characterization of regular cardinals and not a new forcing invariant.

\begin{remark}[Ordinal sequences, not binary ones; source 08]\label{univ:rem:ordinalnotbinary}
Ordinal-valued sequences in \eqref{univ:eq:delta} must not be replaced by
subsets of small ordinals. A new countable cofinal sequence in a large
cardinal may exist without any new real; Prikry forcing is the decisive
example (\cref{univ:sub:prikry}). The codings below preserve the number of
bounds rather than encode large ordinal entries into a short binary string.
\end{remark}

\section{Coding a fresh sequence by an omitted cut}\label{univ:sec:decoder}

The converse direction of the main theorem needs more than the observation
that a new binary sign has a new cut. A short sequence can take values in a
very large ordinal: a new countable cofinal sequence can appear while no real
is added. Encoding it as one binary sign might require a long ordinal domain,
which would destroy the cardinal bound on the cut. Each construction below
keeps the number of inequalities equal to the length of the sequence,
independently of the size of its values. \Cref{univ:sub:tree} is the main
route (source~01); \cref{univ:sub:absolutecode,univ:sub:blockcode} are marked
second routes (sources 04 and 08, with source~03's variant in
\cref{univ:prop:blockcodeplus}).

\subsection{Main route: an arbitrarily branching interval tree (source 01)}\label{univ:sub:tree}

\begin{lemma}[Arbitrarily branching interval tree]\label{univ:lem:tree}
Let $\theta\geq2$ be an ordinal and $\lambda$ a nonzero limit ordinal. In $M$
there is a set-sized table of pairs
\[
 (a_s,b_s)\in(\NoM)^2,
 \qquad s\in T:=\bigl({}^{<\lambda}\theta\bigr)^M,
\]
such that:
\begin{enumerate}[label=\textnormal{(\roman*)}]
\item $a_\varnothing=0<b_\varnothing=1$;
\item if $s$ is a proper initial segment of $t$, then $a_s<a_t<b_t<b_s$;
\item the intervals of distinct immediate successors of one node are pairwise
disjoint and ordered by their successor labels.
\end{enumerate}
\end{lemma}
\begin{proof}
Recurse on the levels below $\lambda$. Given $(a_s,b_s)$ put $w_s=b_s-a_s>0$,
regard $\xi$ and $\theta$ as their all-plus surreal images, and set
\begin{equation}\label{univ:eq:split}
 a_{s^\frown\xi}=a_s+w_s\frac{3\xi+1}{3\theta+3},\qquad
 b_{s^\frown\xi}=a_s+w_s\frac{3\xi+2}{3\theta+3}
 \qquad(\xi<\theta).
\end{equation}
Every operation here is surreal field arithmetic. If $\xi<\eta<\theta$, the
equality of ordinal and surreal successor gives $\xi+1\leq\eta$, hence
$0<3\xi+1<3\xi+2<3\eta+1<3\eta+2<3\theta+3$. This proves the interior and
sibling-separation properties.

At a nonzero limit level $\gamma<\lambda$ and $s\in({}^\gamma\theta)^M$, the
induction hypothesis gives separated ground sets
$A_s=\{a_{s\res\epsilon}:\epsilon<\gamma\}$ and
$B_s=\{b_{s\res\epsilon}:\epsilon<\gamma\}$. Define inside $M$
\begin{equation}\label{univ:eq:limitinterval}
 a_s=\cut{A_s}{B_s},\qquad
 b_s=\cut{A_s\cup\{a_s\}}{B_s}.
\end{equation}
Then $A_s<a_s<b_s<B_s$. Each level and its prefix data are sets in $M$;
Replacement and set-length recursion produce the table in $M$. No recursion
along all ordinals is used.
\end{proof}

\begin{theorem}[Branch decoder and omitted cut]\label{univ:thm:decode}
Suppose $f:\lambda\to\theta$ belongs to $N\setminus M$, where $\lambda$ is a
nonzero limit ordinal and every proper initial segment $f\res\alpha$ belongs
to $M$. Then $\NoM$ omits the finitely satisfiable partial type
\begin{equation}\label{univ:eq:branchtype}
 p_f(X)=\{\ a_{f\res\alpha}<X<b_{f\res\alpha}:\alpha<\lambda\ \}.
\end{equation}
It uses at most $\abs\lambda$ parameters on each side; the left endpoints
strictly increase and the right endpoints strictly decrease.
\end{theorem}
\begin{proof}
All endpoints are old. A finite subset of the inequalities is satisfied by
the midpoint of its deepest interval, which is old. Strict nesting proves
monotonicity.

Suppose an old $x$ realizes all inequalities. Using only the ground table and
$x$, define $\mathrm{dec}_x:\lambda\to\theta$ by recursion in $M$: at a successor step,
with prefix $s$ constructed, choose the unique $\xi<\theta$ with
$a_{s^\frown\xi}<x<b_{s^\frown\xi}$ if it exists, and $0$ otherwise; at limits
take unions. Sibling intervals are disjoint, so the choice is unique, and the
default value makes this a \emph{total ground recursion} that does not
presuppose a branch in $M$.

In $N$, induction gives $\mathrm{dec}_x\res\alpha=f\res\alpha$ for every $\alpha<\lambda$:
at the next step $x$ lies in the interval of $f\res(\alpha+1)$ by the type, and
$\alpha+1<\lambda$ because $\lambda$ is limit. Thus $\mathrm{dec}_x=f$, contradicting
$\mathrm{dec}_x\in M$ and $f\notin M$.
\end{proof}

The proof does not infer that an external branch is old because all its nodes
are old; that inference would be false. An assumed \emph{old realizer}
supplies a total, deterministic ground decoder for the whole branch. There is
also no assertion that the interval widths tend to zero in the fine topology;
\cref{univ:sec:topology} shows they have a common positive old lower bound.

\subsection{Second route A: an absolute all-ordinal interval code (source 04)}\label{univ:sub:absolutecode}

The tree of \cref{univ:lem:tree} fixes an alphabet bound $\theta$ and is built
level by level inside $M$. Source~04's code instead lets the children of a node
be indexed by \emph{all} ordinals and defines the interval of each individual
sequence $s$ by set recursion along $\dom(s)+1$, uniformly in whichever
universe contains $s$. At the root $a_\varnothing=0$, $b_\varnothing=1$. With
$w_s=b_s-a_s$ and $\alpha$ an ordinal,
\begin{equation}\label{univ:eq:childab}
 a_{s\cat\seq\alpha}=a_s+\frac{w_s}{2\hatord\alpha+3},\qquad
 b_{s\cat\seq\alpha}=a_s+\frac{w_s}{2\hatord\alpha+2},
\end{equation}
all operations being surreal field operations. At a nonzero limit length
$\gamma=\dom(s)$,
\begin{equation}\label{univ:eq:limitcenter}
 u_s=\cut{a_{s\res\xi}:\xi<\gamma}{b_{s\res\xi}:\xi<\gamma},\qquad
 e_s=\cut{0}{u_s-a_{s\res\xi},\ b_{s\res\xi}-u_s:\xi<\gamma},
\end{equation}
and $a_s=u_s-e_s/2$, $b_s=u_s+e_s/2$; the right option set of $e_s$ is the
union of the two displayed families. Write $I_s=(a_s,b_s)$.

\begin{theorem}[Absolute ordinal-branching interval code; source 04]\label{univ:thm:intervalcode}
\begin{enumerate}[label=\textnormal{(\roman*)}]
\item For every proper extension $t$ of $s$, $[a_t,b_t]\subset(a_s,b_s)$.
\item If $\alpha<\beta$ then $b_{s\cat\seq\beta}<a_{s\cat\seq\alpha}$; distinct
children have disjoint closed intervals.
\item Every nonroot closed interval lies inside $(0,1)$.
\item If $s\in M$, both endpoints belong to $M$ and have the same values when
computed in $M$ or in $N$.
\end{enumerate}
For a fixed input $s$ the construction needs only set recursion along
$\dom(s)+1$; it forms no set of all nodes at a level and no class-valued
transfinite recursion.
\end{theorem}
\begin{proof}
At a successor node both denominators in \eqref{univ:eq:childab} exceed one,
with $2\hatord\alpha+3>2\hatord\alpha+2>0$, so
$a_s<a_{s\cat\seq\alpha}<b_{s\cat\seq\alpha}<b_s$. If $\alpha<\beta$ then
$\hatord\beta\ge\hatord\alpha+1$, hence
$\frac1{2\hatord\beta+2}\le\frac1{2\hatord\alpha+4}<\frac1{2\hatord\alpha+3}$,
which is (ii). At a limit stage the earlier closed intervals are strictly
nested, so the cut $u_s$ is legitimate and lies strictly inside every earlier
interval; all right options of $e_s$ are positive, so $0<e_s$ is smaller than
each of those distances, and $a_{s\res\xi}<u_s-e_s/2<u_s+e_s/2<b_{s\res\xi}$.
This proves (i)--(iii) by induction. For (iv), run the recursion on the
prefixes of the given $s$: successor steps use absolute arithmetic
(\cref{univ:prop:absolute}(iii)) and limit steps use absoluteness of old cuts
(\cref{univ:prop:absolute}(ii)) twice; all option sets belong to $M$ by
Replacement in $M$. The values on the prefixes of one input form a set-length
sequence of sets, constructed by the ordinary recursion theorem.
\end{proof}

\begin{theorem}[Fresh-branch omitted cut; source 04]\label{univ:thm:freshcut}
Let $f:\mu\to\Ord$ be fresh over $M$, and put
$A_f=\{a_{f\res\xi}:\xi<\mu\}$, $B_f=\{b_{f\res\xi}:\xi<\mu\}$. Then $A_f,B_f$
are $N$-sets of old surreals of cardinality at most $\abs\mu^N$, $A_f<B_f$,
and their strict cut is finitely satisfiable in $\NoM$ but has no separator in
$\NoM$. Removing the root endpoints leaves an omitted cut inside $(0,1)$.
\end{theorem}
\begin{proof}
All endpoints are old by \cref{univ:thm:intervalcode}(iv), and nesting gives
$A_f<B_f$; the midpoint of a deep enough $I_{f\res\eta}$ satisfies finitely
many constraints. Suppose $x\in\NoM$ separates the whole cut, so
$x\in I_{f\res\xi}$ for all $\xi<\mu$. Define $g:\mu\to\Ord$ by set recursion
in $M$: given $g\res\xi$, let $g(\xi)$ be the unique ordinal $\alpha$ with
$x\in I_{(g\res\xi)\cat\seq\alpha}$ if one exists, and $0$ otherwise. The
existential quantifier ranges over ordinals, not over a set of all children,
and uniqueness comes from (ii). By induction in $N$, $g(\xi)=f(\xi)$: since
$\mu$ is a limit, $x\in I_{f\res(\xi+1)}=I_{(g\res\xi)\cat\seq{f(\xi)}}$, an
interval computed absolutely from old data. Hence $g=f\in M$, a contradiction.
For the last clause omit the index $0$ on both sides.
\end{proof}

\begin{proposition}[Closed formulas along a finite branch; source 04]\label{univ:prop:finiteformula}
For an ordinal $\alpha$ put
$r_\alpha=\frac1{2\hatord\alpha+3}$ and
$\ell_\alpha=\frac1{(2\hatord\alpha+2)(2\hatord\alpha+3)}$. If
$s=\seq{\alpha_0,\ldots,\alpha_{n-1}}$, then
\[
 a_s=\sum_{j<n}r_{\alpha_j}\prod_{k<j}\ell_{\alpha_k},\qquad
 b_s=a_s+\prod_{k<n}\ell_{\alpha_k},
\]
with the empty product one and the empty sum zero.
\end{proposition}
\begin{proof}
The child interval is the image of $(0,1)$ under
$t\mapsto a_s+w_s(r_\alpha+\ell_\alpha t)$, since
$\frac1{2\hatord\alpha+2}-\frac1{2\hatord\alpha+3}=\ell_\alpha$. Induct on $n$.
\end{proof}

For a fresh sequence of length $\omega$, such as a Prikry sequence
(\cref{univ:sub:prikry}), all endpoints of the omitted cut are given by these
finite rational expressions in the embedded ordinals; no limit-stage cut is
needed to write them down. The absence of an old separator is still an
infinite assertion, proved by the decoder, not by finite experimentation.

\subsection{Second route B: ordinal block codes of signs (source 08)}\label{univ:sub:blockcode}

Here the fresh sequence is written into a single sign sequence whose prefixes
provide the bounds.

\begin{lemma}[Two bounds force a prefix; source 08]\label{univ:lem:twobounds}
For every sign sequence $u$, the numerical interval $(u\cat-,\ u)$ equals
$\Cone(u\cat-\cat+)$.
\end{lemma}
\begin{proof}
An extension of $u\cat-\cat+$ is below $u$, since its first sign after $u$ is
minus, and above $u\cat-$, since its first sign after that is plus.
Conversely, if $u\cat-<y<u$ and $y$ first differed from or terminated before
the end of $u$, its comparisons with $u$ and with $u\cat-$ would agree,
contradicting the two inequalities. So $u\preceq y$; the upper inequality
forces the next sign to be minus, and the lower inequality, excluding
$y=u\cat-$, forces the following sign to be plus.
\end{proof}

Let $\lambda$ be an infinite limit ordinal and $f:\lambda\to\Ord$. By set
recursion put $p_0=\varnothing$, $p_{\xi+1}=p_\xi\cat+^{f(\xi)}\cat-\cat+$, and
$p_\eta=\bigcup_{\xi<\eta}p_\xi$ at limits, where $+^\alpha$ is a run of plus
signs of ordinal length $\alpha$. Set
\begin{equation}\label{univ:eq:blockendpoints}
 \code(f)=p_\lambda,\qquad
 u_\xi=p_\xi\cat+^{f(\xi)},\qquad
 l_\xi=u_\xi\cat-.
\end{equation}

\begin{lemma}[Decoding and nesting; source 08]\label{univ:lem:blockdecode}
$\code(f)$ determines $f$ and its block boundaries by a set-recursive
decoder. For $\xi<\eta<\lambda$,
$l_\xi<l_\eta<\code(f)<u_\eta<u_\xi$, and
$\bigcap_{\xi<\lambda}(l_\xi,u_\xi)=\Cone(\code(f))$. These assertions are
absolute for transitive universes containing $f$.
\end{lemma}
\begin{proof}
From a block boundary, read the run of plus signs up to the next minus; its
order type is $f(\xi)$. Consume that minus and the following plus; the next
block begins there, and at limits take suprema of boundaries. By
\cref{univ:lem:twobounds}, $(l_\xi,u_\xi)=\Cone(p_{\xi+1})$. Every later
endpoint and $\code(f)$ extend $p_{\xi+1}$; comparing at the final plus marker
of block $\eta$ gives $l_\eta<\code(f)$, and at its minus marker
$\code(f)<u_\eta$. A sequence extends every $p_{\xi+1}$ iff it extends their
union.
\end{proof}

\begin{theorem}[Fresh-sequence cut and its separators; source 08]\label{univ:thm:freshcode}
Let $f:\lambda\to\Ord$ be fresh over $M$, $\lambda$ an infinite limit ordinal.
The endpoints $l_\xi,u_\xi$ are old, but the cut
$L_f=\{l_\xi:\xi<\lambda\}$, $R_f=\{u_\xi:\xi<\lambda\}$ has no separator in
$\NoM$. In $\NoN$ its class of separators is exactly $\Cone(\code(f))$. Every
proper prefix of $\code(f)$ is old, while $\code(f)$ is new. If
$\lambda=\delta(M,N)$, the cut determines a set-presented gap of $\NoM$ of
exact cofinality pair $(\lambda,\lambda)$, and
$\cf^N\bd(\code(f))=\lambda$. \textup{(For an arbitrary fresh sign, the
cofinalities of its gap are computed in \cref{univ:lem:canonicalpresentation}.)}
\end{theorem}
\begin{proof}
Since $\lambda$ is a limit, each $f\res(\xi+1)$ is old, and computing its code
inside $M$ gives $p_{\xi+1}$, $u_\xi$, $l_\xi$. The two endpoint sets need not
belong to $M$. By \cref{univ:lem:blockdecode} the separators in $\NoN$ are the
extensions of $\code(f)$. An old separator $y$ would give
$\code(f)\preceq y$, hence $\code(f)\in M$ by initiality, and decoding inside
$M$ would recover $f$. The same decoding shows $\code(f)\notin M$. Any proper
prefix of $\code(f)$ is a restriction of some old $p_\alpha$.

For the gap, recall that $\lambda=\delta$ is a regular cardinal of $N$
(\cref{univ:prop:deltaregular}). The classes $D^-=\{y\in\NoM:\exists\xi\ y\le l_\xi\}$ and
$D^+=\{y\in\NoM:\exists\xi\ y\ge u_\xi\}$ partition $\NoM$, since a point in
neither would separate the cut; strict nesting shows that neither has an
endpoint toward the other. A cofinal $X\subseteq D^-$ of size below $\lambda$
is impossible: its indices $\xi_x$ with $x\le l_{\xi_x}$ are bounded below
$\lambda$ by regularity, and a later $l_\eta$ exceeds $X$; symmetrically on
the right. Finally $\bd(\code(f))=\sum_{\xi<\lambda}(f(\xi)+2)$; its increasing
partial sums are cofinal, and fewer than $\lambda$ cofinal ordinals would, by
regularity, bound the partial-sum indices below $\lambda$.
\end{proof}

The distinction between a family of old endpoints and an old \emph{set} of
endpoints is the whole mechanism: internal saturation fills cuts whose
bounding sets are in $M$, and the cut constructed here is an external set of
old bounds. Each possibly enormous value $f(\xi)$ is stored in a single old
sign prefix; what is compressed is the \emph{number of parameters}, not the
birthdays of their values. Source~03 uses the variant
$\codeplus(f)=\mathop{\cat}_{\xi}(+^{f(\xi)+1}\cat-)$, with endpoints read off
before the last plus and before the minus of each block; it is printed with
its proof in \cref{univ:prop:blockcodeplus}, where it serves the least-gap
theorem. Both codes use ordinal sums, not the natural sums that describe
addition of ordinal surreals.

\section{The exact spectrum and its first gap}\label{univ:sec:spectrum}

\subsection{Proof of the main equivalence}

\begin{proof}[Proof of \cref{univ:thm:main-summary}]
(i)$\Rightarrow$(ii) is \cref{univ:cor:NSoneside}, and (ii)$\Rightarrow$(iii)
is immediate. Suppose (i) fails. Choose the least ordinal $\lambda<\kappa$ for
which there is a new ordinal-valued sequence of length $\lambda$, choose one
such sequence, and bound its range by an ordinal $\theta\geq2$. As observed in
\cref{univ:sec:delta}, $\lambda$ is an infinite limit ordinal and the
sequence is fresh. \Cref{univ:thm:decode} produces an unfilled cut whose two
sides together have cardinality at most $\abs\lambda<\kappa$ in $N$, so (iii)
fails; \cref{univ:thm:freshcut} gives one inside $(0,1)$. For uncountable
$\kappa$, \cref{univ:prop:absolute,univ:lem:RCFsat}(i) identify (iii) with
ordered-field $\kappa$-saturation, and \cref{univ:lem:RCFsat}(ii) with
$\kappa$-saturation of the dense order.
\end{proof}

$\kappa$ need not be regular. Only the conversion to real closed saturation
requires $\kappa>\omega$; at $\kappa=\omega$ both sides of (i)$\Leftrightarrow$(iii)
hold trivially (source~08), since finite sequences of ordinals and finite sets
of old numbers are old. Source~03 states the cut equivalence for every
infinite \emph{ordinal} bound:
\begin{equation}\label{univ:eq:etaordinal}
 \NoM\text{ has }\eta_\kappa
 \quad\Longleftrightarrow\quad
 \forall\alpha<\kappa\ \bigl({}^{\alpha}\Ord\bigr)^N
 \subseteq M,
\end{equation}
where the right side is shorthand for the class statement that every
ordinal-valued sequence of length $\alpha$ in $N$ belongs to $M$. The proof is
the same: a new sequence of length below $\kappa$ gives $\delta<\kappa$ and an
omitted pair of size $\delta$ (\cref{univ:thm:firstgap}).

\subsection{Exact two-sided cofinalities}

\begin{definition}[Set-presented gaps]\label{univ:def:setgap}
A \emph{set-presented gap} of $\NoM$ is a partition $\NoM=D\sqcup U$ with
$D<U$, neither $D$ having a greatest nor $U$ a least element, together with
$N$-sets $A\subseteq D$ cofinal in $D$ and $B\subseteq U$ coinitial in $U$.
Under \textbf{(H)} a gap is given by such a presentation $(A,B)$, and
$D=\{x\in\NoM:\exists a\in A\ x\le a\}$, $U=\NoM\setminus D$ are defined from
it; statements about all gaps are schemes over presentations. Its
\emph{cofinality pair} $(\cf^N D,\ci^N U)$ is the pair of least sizes, computed
in $N$, of cofinal and coinitial $N$-subsets; these are infinite regular
cardinals. $D$ and $U$ themselves are never required to be sets. This is
source~01's ``set-character gap'', source~03's ``set-presented gap'' and
source~08's ``Dedekind gap'' with its character.
\end{definition}

\begin{theorem}[First gap theorem; sources 01, 03, 08]\label{univ:thm:firstgap}
Let $M\ne N$ and $\delta=\delta(M,N)$. Every set-presented gap of $\NoM$ has
both cofinalities at least $\delta$, and $\NoM$ has one of exact cofinality
pair $(\delta,\delta)$.
\end{theorem}
\begin{proof}
Apply \cref{univ:thm:decode} to a new sequence of minimal length $\delta$ and
write its endpoints as $a_\alpha,b_\alpha$. Define in $N$
$D=\{x\in\NoM:\exists\alpha<\delta\ (x\leq a_\alpha)\}$ and
$U=\{x\in\NoM:\exists\alpha<\delta\ (b_\alpha\leq x)\}$. A point in neither
class would realize the omitted type, so $D\cup U=\NoM$ and $D<U$; strict
monotonicity and the limit length show that neither class has an endpoint
toward the other. The endpoint sequences are witnesses of size $\delta$.

For the lower bound, take any set-presented gap with witnesses $A$ and $B$.
If either could have size less than $\delta$,
\cref{univ:thm:main-summary}(ii) at $\kappa=\delta$ would give an old strict
separator between $A$ and $B$. A point of the lower class cannot lie strictly
above a cofinal subset of it, and a point of the upper class cannot lie
strictly below a coinitial subset of it; since the classes partition $\NoM$,
this is impossible.
\end{proof}

The block code gives a second proof of the upper bound, with the exact
separator class $\Cone(\code(f))$ (\cref{univ:thm:freshcode}, source~08).

\begin{corollary}[Recovering the threshold from the order; sources 01, 03, 08]\label{univ:cor:recover}
For $M\ne N$, $\delta(M,N)$ is the least cardinal $\mu$ for which $\NoM$ has an
unfilled set-sized cut with both sides of size at most $\mu$; equivalently, the
least cofinality on either side of a set-presented gap of $\NoM$.
\end{corollary}
\begin{proof}
The upper bounds follow from \cref{univ:thm:firstgap}. A smaller unfilled cut
contradicts the main theorem at $\delta$, and a smaller gap cofinality
contradicts the first gap theorem.
\end{proof}

\begin{corollary}[The uncountable spectrum; sources 01, 03, 04, 08]\label{univ:cor:spectrum}
For uncountable $N$-cardinals $\kappa$, $\NoM$ is $\kappa$-saturated exactly
when $\kappa\leq\delta(M,N)$. Thus if $\delta>\omega$ it is
$\delta$-saturated but not $(\delta^+)^N$-saturated, and if $\delta=\omega$ it
is not $\omega_1$-saturated. $\NoM$ is externally set-saturated if and only if
$M=N$ (source~04).
\end{corollary}
\begin{proof}
There are no new sequences of length less than $\delta$, and there is one of
length $\delta$; apply the main theorem. If $M=N$, use
\cref{univ:lem:RCFsat}(iii); if $M\ne N$, any uncountable cardinal above
$\delta$ fails.
\end{proof}

That $\NoM$ has no smaller set-presented gap does not exclude cuts with
proper-class cofinality on one side. Those are outside the invariant and
outside the main interpolation theorem. The definition of the invariant uses
only the order of $\NoM$, external sets and their cardinalities, yet its value
is the first length at which $N$ has new ordinal data. Since a real closed
field has a unique field order, it is preserved even by isomorphisms stated
only in the field language (source~03).

\section{The countable boundary: no new reals is weaker}\label{univ:sec:omega}

The absence of new reals and the absence of new countable ordinal-valued
sequences are different preservation properties, and the old surreal field
records the difference exactly.

\begin{theorem}[Finite-parameter saturation; sources 01, 03, 04, 08]\label{univ:thm:omega}
\[
 \NoM\text{ is }\omega\text{-saturated in }N
 \quad\Longleftrightarrow\quad
 \Pow(\omega)^M=\Pow(\omega)^N
 \quad\Longleftrightarrow\quad
 \R^M=\R^N.
\]
In contrast, the pure order $(\NoM,<)$ is always $\omega$-saturated
\textup{(source 08)}.
\end{theorem}
\begin{proof}
Equality of the power sets is equivalent to equality of the reals by ordinary
countable coding in both transitive universes.

Suppose no new subsets of $\omega$ are added. Let $A$ be a finite set of old
parameters and $p$ a complete one-type over $A$ finitely satisfiable in
$\NoM$ according to $N$. The finite set $A$ is in $M$, the formulas over $A$
have a countable enumeration in $M$, and the code of $p$ is a subset of
$\omega$, so $p\in M$. For each finite $p_0\subseteq p$, solvability of its
conjunction is an existential ordered-field formula over old parameters, with
the same truth value in $M$ and $N$ by \cref{univ:prop:absolute} and quantifier
elimination. Hence $M$ recognizes $p$ as finitely satisfiable, and in $M$ the
full surreal field realizes every set-sized type (\cref{univ:lem:RCFsat}(iii)).

Conversely let $r\in\R^N\setminus\R^M$ and consider the rational-cut partial
type
\begin{equation}\label{univ:eq:newrealtype}
 \{q<X:q\in\Q,\ q<r\}
 \ \cup\
 \{X<q:q\in\Q,\ r<q\}.
\end{equation}
It is finitely satisfiable by rationals, and it has no parameters, since
$m/n<X$ is equivalent to $m<nX$ for $n>0$. If an old $x$ realized it, the
ground set $\{q\in\Q:q<x\}$ would be the lower rational cut of $r$, and ground
Dedekind completeness would put $r$ in $M$. For a complete omitted type, take
the type of the image of $r$ in $\NoN$ over the empty set; every finite part is
realized in $\NoM$ because $\NoM\preccurlyeq\NoN$. No class compactness
principle is required. Source~04 gives the same argument with the dyadic
prefixes $q_n$ of a new sign $s:\omega\to\{-,+\}$, using the parameter-free
formulas $q_n<X$ for $s(n)=+$ and $X<q_n$ for $s(n)=-$.

Finally, all finite cuts of the old order are filled, so
\cref{univ:lem:RCFsat}(ii) gives $\omega$-saturation of the pure order.
\end{proof}

\begin{corollary}[Source 01]\label{univ:cor:omegaone}
An extension that adds no reals but adds a countable ordinal-valued sequence
leaves $\NoM$ $\omega$-saturated and destroys $\omega_1^N$-saturation. It
creates an $(\omega,\omega)$ set-presented gap in $\NoM$, but none of its
omitted types can be witnessed over a fixed finite parameter set.
\end{corollary}
\begin{proof}
The first new length is $\omega$. Apply
\cref{univ:thm:omega,univ:thm:firstgap,univ:cor:spectrum}. The last assertion
is the surviving $\omega$-saturation.
\end{proof}

A new countable ordinal-valued sequence need not code a new real. The new
countable parameter family in the interval tree carries information that no
old finite parameter tuple can code; a bare new-real argument would miss this
example. One must not set $\kappa=\omega$ in the field-saturation clause of
\cref{univ:thm:main-summary}: the real closure of a finite parameter set is
usually countable, not finite. This is not a conflict between two definitions
of saturation but an exact distinction between finite and countable parameter
sets in the same language.

\part{Fresh signs classify the set-presented gaps (source 03)}\label{univ:part:gaps}

\section{The gap--fresh-sign correspondence}\label{univ:sec:fresh}

This part is source~03's. It classifies \emph{all} set-presented gaps of
$\NoM$, not only the first. It answers source~01's question about asymmetric
cofinality pairs (\cref{univ:sec:questions}). Its proofs use only the sign
order, prefix-initiality, the old constant-tail closure and the set-cut
principle in $N$; they need no Global Choice.

\begin{definition}[Fresh signs]\label{univ:def:freshsign}
A \emph{fresh sign over $M$ in $N$} is $s\in\NoN\setminus\NoM$ such that
$s\res\alpha\in\NoM$ for every $\alpha<\bd(s)$. Write $\Fresh(M,N)$ for the
class of fresh signs. This is the binary special case of a fresh function
(\cref{univ:def:delta}); the forcing notion of freshness and its spectra are
due to Fischer--Koelbing--Wohofsky \cite{FKW}, and the name is not claimed as a
new invariant. For $s\in\NoN\setminus\NoM$ put
\[
 D_s=\{x\in\NoM:x<s\},\qquad U_s=\{x\in\NoM:s<x\}.
\]
Write $\GapSpec_N(\NoM)$ for the class of cofinality pairs of set-presented
gaps of $\NoM$ (\cref{univ:def:setgap}), computed in $N$.
\end{definition}

\begin{lemma}[Both signs are cofinal]\label{univ:lem:signscofinal}
If $s\in\Fresh(M,N)$ and $\lambda=\bd(s)$, then $\lambda$ is a nonzero limit
ordinal, and both $P_s=\{\alpha<\lambda:s(\alpha)=+\}$ and
$Q_s=\{\alpha<\lambda:s(\alpha)=-\}$ are cofinal in $\lambda$.
\end{lemma}
\begin{proof}
A successor-length sequence is determined by its immediately preceding prefix
and one sign; the empty sequence is old. If, say, $P_s$ were bounded by
$\alpha<\lambda$, then $s$ would be the old prefix $s\res\alpha$ followed by
minus signs up to the old ordinal $\lambda$, an operation in $M$. Interchange
the signs for $Q_s$.
\end{proof}

\begin{lemma}[Canonical cofinal presentations]\label{univ:lem:canonicalpresentation}
For a fresh sign $s$ of length $\lambda$ put
$A_s=\{s\res\alpha:\alpha\in P_s\}$ and $B_s=\{s\res\alpha:\alpha\in Q_s\}$.
Then $A_s$ is cofinal in $D_s$, $B_s$ is coinitial in $U_s$, both are $N$-sets
of old numbers, $A_s<s<B_s$, and
\begin{equation}\label{univ:eq:exactcofs}
 \cf^N D_s=\ci^N U_s=\cf^N\lambda.
\end{equation}
\end{lemma}
\begin{proof}
At a plus position the prefix ends before a plus, so lies below $s$; at a minus
position above. These prefixes are old by freshness, and no old element
extends $s$, since restriction to $\lambda$ would place $s$ in $M$.

Take $x\in D_s$. Either $x$ is a proper prefix of $s$ or it first disagrees
with $s$ at some $\gamma<\lambda$. Choose a plus position beyond $\gamma$ (or
beyond $\bd(x)$); comparison at the earlier position gives
$x<s\res\alpha<s$. A late minus position handles $U_s$. This also shows that
neither side has an endpoint.

Write $\mu=\cf^N\lambda$. A cofinal sequence of positions of length $\mu$ can
be replaced by later plus and later minus positions, giving witnesses of size
at most $\mu$. Conversely, if $X\subseteq D_s$ has size below $\mu$, record for
each $x\in X$ a position where $x$ ends as a prefix or first disagrees with
$s$; these positions are bounded below $\lambda$, and a plus prefix beyond the
bound exceeds all of $X$ and lies in $D_s$. The right side is symmetric.
\end{proof}

\begin{theorem}[Gap--fresh-sign correspondence]\label{univ:thm:gapcorrespondence}
The assignment $s\mapsto(D_s,U_s)$ is a bijection between fresh signs over $M$
in $N$ and the set-presented gaps of $\NoM$. The inverse sends a gap with any
presentation $A\subseteq D$, $B\subseteq U$ to the simplest separator
$\cut AB$ computed in $N$, and the result does not depend on the chosen
presentation. Every set-presented gap is symmetric, with cofinalities given by
\eqref{univ:eq:exactcofs}.
\end{theorem}
\begin{proof}
\Cref{univ:lem:canonicalpresentation} shows the assignment is well defined.
For surjectivity, let $(D,U)$ be presented by $A,B$. No member of $\NoM$
separates $A,B$: a member of $D$ is at most some member of $A$, and a member of
$U$ is at least some member of $B$. Let $s=\cut AB^N$; then $s\notin\NoM$. Every
proper prefix $y=s\res\gamma$ is old: by simplicity $y$ does not separate
$A,B$; if $y<s$, all inequalities $y<b$ already hold, so there is $a\in A$ with
$y\le a<s$, and convexity of $\Cone(y)\ni s$ gives $y\preceq a$, so $y$ is a
prefix of an old number; if $s<y$, argue with some $b\in B$. Thus $s$ is fresh.
Cofinality of $A$ gives $D\subseteq D_s$ and coinitiality of $B$ gives
$U\subseteq U_s$; both are partitions of $\NoM$, so the gap is $(D_s,U_s)$.

For injectivity, take distinct fresh $s,t$. Neither is a proper prefix of the
other, since then a new sign would be a proper prefix of a fresh sign. Their
common prefix $u$ is old and lies strictly between them, on opposite sides of
the two cuts. The same argument shows independence of the presentation.
Symmetry is \cref{univ:lem:canonicalpresentation}.
\end{proof}

\begin{theorem}[The whole filling class]\label{univ:thm:fillingcone}
If $s$ is fresh, then $\{y\in\NoN:D_s<y<U_s\}=\Cone(s)$. Every new
$x\in\NoN\setminus\NoM$ lies in exactly one such cone, that of its shortest new
prefix. Thus $\NoN\setminus\NoM$ is the disjoint union of the cones
$\Cone(s)$, $s\in\Fresh(M,N)$, as class-defined convex pieces.
\end{theorem}
\begin{proof}
An old element cannot extend $s$, so its comparison with any extension of $s$
is decided before $\bd(s)$ and equals its comparison with $s$; this gives
$\supseteq$. If $y$ does not extend $s$, either it is a proper prefix of $s$,
hence old and not a filler of its own ground cut, or it first disagrees before
$\bd(s)$; a later plus prefix in $D_s$ or a later minus prefix in $U_s$ then
lies on the wrong side of $y$. For a new $x$, the least length of a new prefix
gives a fresh prefix, and distinct fresh signs are incomparable, so their cones
are disjoint. The union is shorthand for membership in exactly one piece; it
posits no set of all pieces.
\end{proof}

\begin{corollary}[Exact gap spectrum]\label{univ:cor:gapspectrum}
With all cofinalities taken in $N$,
\begin{equation}\label{univ:eq:gapspectrum}
 \GapSpec_N(\NoM)=
 \bigl\{(\cf^N\bd(s),\cf^N\bd(s)):
                s\in\Fresh(M,N)\bigr\}.
\end{equation}
In particular, an asymmetrically set-presented gap cannot arise in a
same-ordinal ground surreal field.
\end{corollary}

This is stronger than the existence of a missing cut after adding a real. It
classifies all set-presented gaps, including gaps with long sign codes and
small outer cofinality, and it explains why two new surreals can have exactly
the same cut over the whole old field: they lie in one fresh cone, which can
contain a proper class of extensions. The fresh signs themselves, not only
their lengths, parametrize the gaps. The symmetry concerns the two sides of a
bounded gap, not the global cofinality of the field, which is a proper class
(\cref{univ:lem:oldbounds}). Cofinalities must be computed in $N$: the old
universe's cofinality of $\bd(s)$ can be the wrong one (\cref{univ:sub:prikry}).

\section{The least gap}\label{univ:sec:leastgap}

\begin{proposition}[Terminal-plus block code; source 03]\label{univ:prop:blockcodeplus}
Let $\mu$ be an infinite regular cardinal of $N$, and let $f:\mu\to\Ord$ be new
with $f\res\xi\in M$ for every $\xi<\mu$. Define, by ordinal concatenation,
\[
 \codeplus(f)=\mathop{\cat}_{\xi<\mu}\bigl(+^{\,f(\xi)+1}\cat-\bigr),
 \qquad
 \bd(\codeplus(f))=\sum_{\xi<\mu}\bigl(f(\xi)+2\bigr).
\]
Then $\codeplus(f)$ is fresh over $M$ and $\cf^N\bd(\codeplus(f))=\mu$. With
$\sigma_\xi=\sum_{\eta<\xi}(f(\eta)+2)$, the old prefixes
\[
 a_\xi=\codeplus(f)\res(\sigma_\xi+f(\xi)),\qquad
 b_\xi=\codeplus(f)\res(\sigma_\xi+f(\xi)+1)
\]
satisfy $a_\xi<a_\eta<\codeplus(f)<b_\eta<b_\xi$ for $\xi<\eta<\mu$, and
present an omitted $(\mu,\mu)$ gap of $\NoM$.
\end{proposition}
\begin{proof}
Each block is a nonempty plus run of successor length followed by one minus.
A proper prefix of the concatenation lies inside some block and is determined
by the old restriction $f\res(\xi+1)$ and ordinal positions; since $\mu$ is a
limit, $\xi+1<\mu$, so the prefix is old. If $\codeplus(f)$ were in $M$, its
minus positions would delimit the blocks, and the unique predecessors of the
successor run lengths would recover $f$ in $M$. So it is fresh.

The partial sums $\sigma_\xi$ increase and are cofinal, so
$\cf^N\bd(\codeplus(f))\le\mu$; fewer than $\mu$ cofinal positions would, by
regularity, bound the block indices below $\mu$. At length $\sigma_\xi+f(\xi)$ the
next sign is the last plus of block $\xi$, and at $\sigma_\xi+f(\xi)+1$ it is the
minus; first-disagreement comparison gives the nesting. The endpoint lengths
are cofinal, so, as in \cref{univ:lem:canonicalpresentation}, the endpoints are
cofinal and coinitial in the two sides, and
\cref{univ:thm:gapcorrespondence} gives an omitted gap with the asserted
cofinalities.
\end{proof}

The extra plus in each run is not cosmetic: for a limit value $f(\xi)$, a run
of length $f(\xi)$ has no last plus, and the uniformly specified left endpoint
would not exist. Source~08's code in \cref{univ:sub:blockcode} achieves the
same with the two terminal signs $-\cat+$.

\begin{theorem}[Exact least-gap theorem; source 03]\label{univ:thm:leastgap}
For $M\ne N$,
\begin{align}
 \delta(M,N)
 &=\min\{\abs{A\cup B}^N:A<B\text{ is an omitted set pair in }\NoM\}
                                                        \label{univ:eq:leastpair}\\
 &=\min\{\cf^N\bd(s):s\in\Fresh(M,N)\}.                  \label{univ:eq:leastfresh}
\end{align}
There is an omitted gap of exact cofinality pair $(\delta,\delta)$.
\end{theorem}
\begin{proof}
The lower bound in \eqref{univ:eq:leastpair} is the main theorem at
$\kappa=\delta$ (equivalently \cref{univ:lem:shortold}: a pair of size below
$\delta$ is old, and its old cut fills it). Apply
\cref{univ:prop:blockcodeplus} to a new $f:\delta\to\Ord$, using
\cref{univ:prop:deltaregular}; its endpoint families have size $\delta$. Every
fresh sign gives an omitted presentation of size $\cf^N\bd(s)$ by
\cref{univ:lem:canonicalpresentation,univ:thm:gapcorrespondence}, so no such
cofinality is below $\delta$, and the block code attains $\delta$.
\end{proof}

The recovery is deliberately limited. The invariant recovers a threshold, and
\cref{univ:cor:gapspectrum} recovers cofinality data. It does not recover the
inner model, the forcing notion, or the actual sign lengths from a bare
ordered-field presentation. Different grounds can have the same threshold,
and equality of gap spectra is not proved to classify the fields.
\emph{Status (batch 37, 28 September 2026).} It does not classify them: old
fields in one universe can have equal full gap spectra, even equal saturation at
every cardinal, and fail to be isomorphic as ordered or pure fields
(\cref{univ:gs:thm:boolean,univ:gs:thm:amplification,univ:gs:thm:incomplete},
Part~VIII); the full spectrum of a pair is the outer-cofinality image of its
fresh-function spectrum (\cref{univ:gs:thm:transfer}).

\part{The first new birthday (source 08)}\label{univ:part:beta}

\section{First new birthdays and their relation to saturation}\label{univ:sec:beta}

The first new surreal number, measured by birthday, need not appear at the
same ordinal as the first omitted-cut size.

\begin{definition}[First new birthday]\label{univ:def:beta}
For $M\ne N$ put
\begin{equation}\label{univ:eq:beta}
 \beta(M,N)=\min\{\alpha:\Pow(\alpha)^M\ne\Pow(\alpha)^N\},
\end{equation}
the power sets being compared as actual sets of subsets in $N$.
\end{definition}

\begin{proposition}[Birthday threshold; source 08]\label{univ:prop:beta}
$\beta(M,N)=\min\{\bd(s):s\in\NoN\setminus\NoM\}$. It is an infinite cardinal
in both $M$ and $N$. With $\delta=\delta(M,N)$ and $\beta=\beta(M,N)$,
\begin{equation}\label{univ:eq:deltabeta}
 \delta\le\cf^N(\beta)\le\beta.
\end{equation}
In particular, the first new birthday can be much larger than the least size
of an omitted cut.
\end{proposition}
\begin{proof}
Subsets of $\alpha$ and signs of length $\alpha$ are uniformly interdefinable
by characteristic signs. This gives the birthday characterization; existence
follows from \cref{univ:sec:delta}. No finite subset is new.

If $\beta$ were not an $M$-cardinal, an $M$-bijection between $\beta$ and a
smaller ordinal would pull a new subset of $\beta$ back to a new subset of that
smaller ordinal. Suppose there were an $N$-bijection $e:\mu\to\beta$ with
$\mu<\beta$. Transporting the order of $\beta$ along $e$ gives a well-order
$\triangleleft$ of $\mu$ of type $\beta$. Since $\beta$ is an infinite
$M$-cardinal, $M$ codes $\mu\times\mu$ by an ordinal $\gamma<\beta$, so by
minimality $\triangleleft\in M$. Its order type is absolute, so $M$ would see a
well-order of $\mu$ of type $\beta$, contradicting that $\beta$ is an
$M$-cardinal.

For the first inequality, choose a new $X\subseteq\beta$ and let
$\mu=\cf^N(\beta)$. If $\mu<\delta$, a cofinal $c:\mu\to\beta$ belongs to $M$,
and every $X\cap c(\xi)$ is old by minimality of $\beta$. In $M$ well-order
$S=\bigcup_{\alpha<\beta}\Pow(\alpha)^M$ and code these restrictions by
ordinals; the code sequence has length $\mu<\delta$, so it is old, and
decoding and taking the union in $M$ recovers $X$. Hence
$\delta\le\cf^N(\beta)$.
\end{proof}

\begin{remark}[Two resources; source 08]
$\beta$ measures the least \emph{length of a new sign sequence}; $\delta$
measures the least \emph{number of old parameters needed for an omitted cut}.
Long old parameters can carry large ordinal entries, as the block codes show.
Equality of all low birthday fragments therefore does not by itself guarantee
countable saturation of the whole old field (\cref{univ:cor:prikry}).
\end{remark}

\begin{remark}[$\beta$ is the least length of a fresh sign; merge observation]\label{univ:rem:betafresh}
A new sign of least birthday $\beta$ has all its proper prefixes old, since
they are signs of smaller birthday; so it is fresh. Every fresh sign is new, so
$\beta=\min\{\bd(s):s\in\Fresh(M,N)\}$. By \cref{univ:thm:leastgap}, $\delta$
is the least \emph{cofinality} $\cf^N\bd(s)$ of a fresh sign. This two-line
observation, added in the merge, places source~08's $\beta$ in source~03's
fresh-sign picture. It does not answer source~03's question whether an
intrinsic ordered-field invariant recovers such lengths
(\cref{univ:q:length}): $\beta$ is defined through birthdays, which are not
part of the field language.
\end{remark}

\begin{remark}[Relation to the sibling report]\label{univ:rem:sibling}
The sibling report
\path{docs/foundations-and-computation/birthday-cutoffs-and-hereditary-sets}
proves, for every same-ordinal outer model, that the first new birthday is the
least ordinal acquiring a new subset; it also writes that ordinal $\beta$.
\Cref{univ:prop:beta} keeps source~08's statement and
proof because source~08 adds that $\beta$ is a cardinal of $N$ and the
inequality $\delta\le\cf^N\beta$ with the first new length $\delta$ of this
report. The sibling report also studies bounded birthday fields in a language
with a birthday symbol; \cref{univ:sec:fragments} here concerns the
ordered-field saturation of such fragments, without that symbol.
\end{remark}

\noindent[Added 4 October 2026, batch~93: Section~41 of
\path{foundations-and-computation/definable-surreals-and-omnific-integers}
(its Part~III, batch-93 manuscript~06; AI-assisted, unrefereed, not
formalized) re-proves, for $N=V$ and a transitive inner model $M\ne V$ of
$\ZFC$ containing all ordinals, that the first new birthday is the least
ordinal acquiring a new subset and is a cardinal of $N$ and of $M$
(\path{dsn:op:sec:first-omissions}); a dated note there credits
\cref{univ:prop:beta}. The inequality $\delta\le\cf^N\beta$ is not re-proved
there.]

\part{Surcomplex fields, fragments, forcing and topology}\label{univ:part:further}

\section{Surcomplex numbers: a language-dependent dichotomy}\label{univ:sec:complex}

Let $c_M(a+bi)=a-bi$ on $\SCM$ and $c_N$ likewise on $\SCN$. The collection's
automorphism report already emphasizes the difference between a pure
algebraically closed field and a field with a distinguished surreal real axis
\cite{RepoAutomorphisms}. That distinction is standard algebra and is not
claimed as new. The point here is its exact effect on \emph{external
saturation under a change of universe}.

\subsection{The pure field forgets the sequence obstruction}

\begin{lemma}[Algebraic closure stays a set; source 03]\label{univ:lem:algebraicset}
Let $H$ be a class field and $A\subseteq H$ an $N$-set. The field generated by
$A$ is a set, and the elements of $H$ algebraic over it form a set. If $H$ is
algebraically closed, this relative algebraic closure is an algebraically
closed set subfield of $H$.
\end{lemma}
\begin{proof}
Finite rational expressions in $A$ form a set, there are set many polynomials
over that field, and a nonzero polynomial of degree $d$ has at most $d$ roots.
Collect the root sets by Replacement. A root of a polynomial with algebraic
coefficients is algebraic over the original field.
\end{proof}

\begin{theorem}[Set-saturation of the old pure surcomplex field; sources 01, 03, 04, 08]\label{univ:thm:purecomplex}
For every $N$-set $A\subseteq\SCM$, every complete one-variable pure-field type
over $A$ that is finitely satisfiable in $\SCM$ is realized in $\SCM$.
Consequently $\SCM$ is $\kappa$-saturated in $N$ for every $N$-cardinal
$\kappa$, with no assumption on new sequences or new reals. More generally
\textup{(source~03)}, every definable proper-class algebraically closed field is
saturated in the pure field language over every $N$-set of parameters.
\end{theorem}
\begin{proof}
Choose a ground set $B\supseteq A$ by \cref{univ:lem:groundbound}. In $M$ let
$E$ be the set of elements of $\SCM$ algebraic over $\Q(B)$; it is an
algebraically closed ground subfield by \cref{univ:lem:algebraicset}. Choose an
old $t\notin E$ without a class-sized transcendence basis: bound the birthdays
of the coordinates of all elements of $E$ and take a larger all-plus ordinal.
Then $t$ is transcendental over $\Q(B)$, and this persists in $N$, because a
polynomial witnessing algebraicity has a finite old coefficient tuple. In
particular $t$ is transcendental over $\Q(A)$ as formed in $N$.

If the given type $p$ contains no equation $P(X)=0$ with $P\ne0$ over $\Q(A)$,
completeness and quantifier elimination make $p$ the unique transcendental
type over $\Q(A)$, realized by $t$. Otherwise choose such an equation. Its
finitely many roots lie in $\SCM$ by \cref{univ:prop:absolute}. If none realized
$p$, choose for each root a formula of $p$ that it fails; together with
$P(X)=0$ these form an unsatisfiable finite subset of $p$. The argument covers
algebraic parameters and any new organization of $A$.

For a general proper-class algebraically closed field $H$, let $D$ be the
algebraic closure in $H$ of the field generated by $A$
(\cref{univ:lem:algebraicset}); extend the type to $D$ by compactness in that
set model, realize an algebraic extension by a root in $D$, and a
transcendental one by any $y\in H\setminus D$, which exists because $H$ is a
proper class.
\end{proof}

There is no cardinal paradox: $\SCM$ is a proper class of $N$, and the
statement is a scheme over set-sized parameter sets, not a claim that a set of
cardinality $\mu$ realizes types over itself of size $\mu$. The theorem is an
established model-theoretic mechanism rather than a new principle about
algebraically closed fields; what matters is its contrast with the real
field. Nonalgebraic one-types over an algebraically closed field have no order
cut to remember. Nor does the theorem by itself give a class isomorphism
$\SCM\cong\SCN$; that needs a class back-and-forth and Global Choice
(\cref{univ:part:realforms}).

\subsection{Naming conjugation restores the old ordered field}

\begin{lemma}[Real coordinates are interpretable; sources 01, 03, 04, 08]\label{univ:lem:coordinates}
The fixed field of $c_M$ is definable in $(\SCM,c_M)$ by $c_M(x)=x$, and on it
the order is definable by
\begin{equation}\label{univ:eq:orderdef}
 x<y\quad\Longleftrightarrow\quad
 \exists u\,\bigl(c_M(u)=u\ \wedge\ u\ne0\ \wedge\ u^2=y-x\bigr).
\end{equation}
After naming $i$, the map $z\mapsto(\operatorname{re}z,\operatorname{im}z)$ with
\begin{equation}\label{univ:eq:coordinates}
 \operatorname{re}z=\frac{z+c_M(z)}2,\qquad
 \operatorname{im}z=\frac{z-c_M(z)}{2i}
\end{equation}
is a definable bijection $\SCM\to(\NoM)^2$, and field operations and
conjugation become polynomial operations on coordinates:
$(a,b)(u,v)=(au-bv,av+bu)$ and $c_M(a,b)=(a,-b)$.
\end{lemma}
\begin{proof}
The fixed-point equation gives zero imaginary coordinate; in a real closed
field the nonzero squares are exactly the positive elements. The coordinate
formulas follow from $i^2=-1$, with inverse $(a,b)\mapsto a+bi$. Quantifiers
over $\SCM$ become pairs of quantifiers over $\NoM$, and quantifiers over $\NoM$
become quantifiers over the definable fixed field.
\end{proof}

\begin{theorem}[Conjugation spectrum; sources 01, 03, 04, 08]\label{univ:thm:conjugation}
For every infinite $N$-cardinal $\kappa$,
\[
 (\SCM,c_M)\text{ is }\kappa\text{-saturated in }N
 \quad\Longleftrightarrow\quad
 \NoM\text{ is }\kappa\text{-saturated in }N.
\]
For uncountable $\kappa$ this is equivalent to $\NS_{<\kappa}(M,N)$, and for
$\kappa=\omega$ to the absence of new reals. The statements remain equivalent
after naming $i$. Moreover the inclusion $(\SCM,c_M)\preccurlyeq(\SCN,c_N)$ is
elementary \textup{(sources 03, 04, 08)}.
\end{theorem}
\begin{proof}
Naming a fixed finite tuple of constants preserves and reflects
$\kappa$-saturation for infinite $\kappa$: absorb the tuple into the parameter
set, or extend a reduct type by compactness and realize it. So we may name $i$.
If $\NoM$ is $\kappa$-saturated, translate a type in $(\SCM,c_M,i)$ over fewer
than $\kappa$ parameters by \cref{univ:lem:coordinates}: the real and imaginary
coordinates of the parameters are still fewer than $\kappa$, the translated
type has two field variables, and finite-variable saturation of $\NoM$ realizes
it. Conversely translate a type over a small subset of $\NoM$ into
$(\SCM,c_M)$ by \eqref{univ:eq:orderdef}, restricting quantifiers to the fixed
field and adding $c_M(X)=X$; saturation gives a fixed-field realization. The
last assertions follow from \cref{univ:thm:main-summary,univ:thm:omega}.
Elementarity: the same coordinate translation reduces each formula with old
parameters to an ordered-field formula, and $\NoM\preccurlyeq\NoN$.
\end{proof}

In particular the interval type \eqref{univ:eq:branchtype} remains omitted in
the conjugation expansion when written with the definable fixed field and
definable positive squares. In the pure field those order inequalities are
unavailable, and the one-variable transcendental type is realized by
\cref{univ:thm:purecomplex}.

\begin{remark}[Conjugation is not definable in the pure field; source 08]\label{univ:rem:notdefinable}
In a pure algebraically closed field every definable subset of one variable,
with any finite tuple of parameters, is finite or cofinite. The fixed field of
$c_M$ is infinite with infinite complement, so $c_M$ is not definable in the
pure field, even with parameters, and naming a set of constants does not
change this. The contrast in \cref{univ:tab:languages} therefore does not
violate preservation of saturation under definitional expansions.
\end{remark}

\begin{table}[ht]
\centering\small
\caption{Saturation in $N$ by structure and language (source 08). ``No new
reals'' means $\R^M=\R^N$.}\label{univ:tab:languages}
\begin{tabular}{@{}P{0.36\textwidth}P{0.22\textwidth}P{0.3\textwidth}@{}}
\toprule
Structure and language & $\omega$-saturation & Uncountable $\kappa$-saturation\\
\midrule
$(\NoM,<)$, pure order & Always & $\NS_{<\kappa}(M,N)$\\
$\NoM$, ordered field & No new reals & $\NS_{<\kappa}(M,N)$\\
$\SCM$, pure field & Always & Always\\
$(\SCM,c_M)$ & No new reals & $\NS_{<\kappa}(M,N)$\\
\bottomrule
\end{tabular}
\end{table}

\subsection{A nonisomorphism consequence}

\begin{corollary}[Sources 01, 04, 08]\label{univ:cor:noniso}
Assume $M\ne N$. There is no $N$-class order isomorphism $\NoM\to\NoN$ whose
restrictions to sets have set-sized images. Consequently there is no such
field isomorphism, and no such isomorphism $(\SCM,c_M)\cong(\SCN,c_N)$.
\end{corollary}
\begin{proof}
The first gap theorem supplies two $N$-sets of old endpoints with no separator
in $\NoM$. An isomorphism as stated sends them to separated $N$-sets in $\NoN$,
whose cut has a separator; its inverse image would separate the original cut.
A field isomorphism of real closed fields preserves order, since positivity is
being a nonzero square, and an isomorphism preserving conjugation restricts to
an isomorphism of the definable fixed fields.
\end{proof}

The set-image condition is automatic for definable class maps and for class
maps of $\GBC$ with class Replacement; sources 04 and 08, working in a class
theory, state the corollary without it. It is stated here to avoid quantifying
over unrestricted external graphs without an ambient class theory. The
corollary compares an inner universe with its ambient universe only; it says
nothing about whether $\No^{M_1}\cong\No^{M_2}$ forces $M_1=M_2$ for two inner
models of a common universe, whose spectra could agree without settling their
isomorphism problem (source~08).

\section{A sharp theorem for bounded birthday fields (source 01)}\label{univ:sec:fragments}

The previous theorems concern all ground birthdays at once. A set-sized
fragment has a new obstruction: its birthdays can become cofinally approachable
by a short sequence, even if every individual sign of length below the bound
remains old.

Let $\kappa$ be an uncountable regular cardinal in $M$ that remains a cardinal
in $N$, and consider the set $(\No_{<\kappa})^M$. It is a real closed set field
in $M$ by the classical birthday-fragment theorem, and remains real closed in
$N$ by the finite-coefficient argument. Regularity in $N$ is \emph{not}
assumed in advance.

\begin{theorem}[Bounded-birthday preservation]\label{univ:thm:fragments}
The following are equivalent.
\begin{enumerate}[label=\textnormal{(\roman*)}]
\item $(\No_{<\kappa})^M$ is $\kappa$-saturated in $N$ as an ordered field.
\item The order $(\No_{<\kappa})^M$ has the $\eta_\kappa$ property in $N$.
\item $\cf^N(\kappa)=\kappa$ and $({}^\alpha2)^N=({}^\alpha2)^M$ for every
$\alpha<\kappa$.
\item $\kappa$ is regular in $N$ and $(\No_{<\kappa})^M=(\No_{<\kappa})^N$ as
sets of sign sequences.
\end{enumerate}
\end{theorem}
\begin{proof}
(i)$\Leftrightarrow$(ii) is \cref{univ:lem:RCFsat}(i), and
(iii)$\Leftrightarrow$(iv) is the representation of surreals by binary sign
sequences.

Assume (iv), and let $L,R\subseteq(\No_{<\kappa})^M$ be $N$-sets with $L<R$ and
$\abs{L\cup R}<\kappa$. By regularity in $N$,
$\gamma=\sup\{\bd(y)+1:y\in L\cup R\}<\kappa$. The cut constructed in $N$ has
birthday at most $\gamma$ by \eqref{univ:eq:cutbirthday}, so it lies in
$(\No_{<\kappa})^N=(\No_{<\kappa})^M$. This proves (ii).

For necessity, suppose first $\cf^N(\kappa)=\mu<\kappa$. A cofinal sequence of
ordinals below $\kappa$, viewed as all-plus surreals, is cofinal in the
fragment by \eqref{univ:eq:birthbounds}, has size below $\kappa$, and has no
strict upper bound there; with an empty right side it violates (ii). Next
suppose a new binary function of length below $\kappa$ exists, and take one
$s:\alpha\to\{-,+\}$ of least length. Its proper prefixes are old and have
birthday below $\kappa$, so $L_s,R_s$ of \cref{univ:lem:signcone} lie in the
fragment and have size at most $\abs\alpha<\kappa$. Any separator, even in
$\NoN$, extends $s$; an old separator cannot, since restricting it to $\alpha$
would put $s$ in $M$. This violates (ii).
\end{proof}

\begin{remark}[Why the two criteria differ]
In the full-class theorem, $\NS_{<\kappa}$ bans short sequences with
\emph{arbitrarily large range}. The fragment theorem bans new binary signs
below a fixed birthday bound and separately preserves the cofinality of that
bound. A short sequence with unbounded range at a much larger cardinal may be
invisible to one small fragment but visible to the full old field. Conversely,
singularizing the bound can destroy fragment saturation without adding any new
sign below it (\cref{univ:cor:prikry}). In the notation of
\cref{univ:part:beta}, condition (iii) says $\cf^N\kappa=\kappa$ and
$\kappa\le\beta(M,N)$ when $M\ne N$.
\end{remark}

\begin{proposition}[Pure complexification of a set fragment; source 01]\label{univ:prop:setACF}
Let $\mu=\abs{(\No_{<\kappa})^M}^N$. If $\mu$ is uncountable, the algebraically
closed field $(\No_{<\kappa})^M[i]$ is $\mu$-saturated in the pure field
language in $N$, whether or not the conditions of \cref{univ:thm:fragments}
hold.
\end{proposition}
\begin{proof}
An algebraically closed field of characteristic zero and uncountable
cardinality $\mu$ has transcendence degree $\mu$ over $\Q$: a basis of size
$\tau$ would give an algebraic closure of size $\max(\tau,\aleph_0)$. Over fewer
than $\mu$ parameters there is still a transcendental element, and algebraic
types are realized by algebraic closedness; quantifier elimination gives
$\mu$-saturation. This is the usual dimension argument, included as a
consequence rather than a new field theorem.
\end{proof}

\section{Forcing examples and distributivity}\label{univ:sec:examples}

The theorems above apply to every pair of universes in their hypotheses. For
set forcing they translate the usual preservation of short ordinal sequences
into field-theoretic statements. The forcing facts used are standard
\cite{Jech,Prikry,Gitik,Benhamou} and are not reproved.

\subsection{The forcing-theoretic formulation}

A forcing $\mathbb P$ is \emph{$<\kappa$-distributive} if the intersection of
fewer than $\kappa$ dense open subsets is dense, equivalently if its complete
Boolean algebra is $<\kappa$-distributive. With the usual separative
interpretation this is equivalent to adding no new ordinal-valued sequences of
length less than $\kappa$ in any generic extension \cite{Jech}; conventions
about whether $\kappa$ itself is included vary, and the sequence formulation is
the one used here. For the forward direction, take for each coordinate of a
name for a short sequence the dense set of conditions deciding it; a condition
in the intersection decides the whole sequence as a ground sequence.
Conversely, refine the dense open sets by maximal antichains in the complete
Boolean algebra; a generic choice through them is a short ordinal-coded
sequence, and if all such sequences are old, conditions deciding a whole
choice are dense below all the antichains. This is the standard Boolean form of
the equivalence, not a new forcing theorem.

\begin{corollary}[Forcing characterization; sources 01, 03, 04, 08]\label{univ:cor:distributivity}
Let $\kappa$ be an uncountable cardinal of $M$ and $\mathbb P\in M$.
\begin{enumerate}[label=\textnormal{(\alph*)}]
\item \textup{(source 03)} If $\kappa$ is regular in $M$, then $\mathbb P$ adds
no new ordinal sequences of length below $\kappa$ in any generic extension if
and only if, in every generic extension, $\NoM$ has the $\eta_\kappa$
property. The preservation direction preserves $\kappa$ and its regularity.
\item \textup{(sources 04, 08)} If $\mathbb P$ preserves $\kappa$ as an
uncountable cardinal, the following are equivalent: $\mathbb P$ is
$<\kappa$-distributive in $M$; in every generic extension $\NoM$ is externally
$\kappa$-saturated; in every generic extension $(\SCM,c_M)$ is externally
$\kappa$-saturated.
\end{enumerate}
Without conjugation, $\SCM$ is externally set-saturated in every generic
extension, regardless of distributivity.
\end{corollary}
\begin{proof}
Apply \eqref{univ:eq:etaordinal}, \cref{univ:thm:main-summary},
\cref{univ:thm:conjugation} and \cref{univ:thm:purecomplex} in each generic
extension. A new short cofinal map or a collapse to a smaller ordinal would be
a forbidden short sequence, which gives the preservation clause in (a).
\end{proof}

The phrase ``in every generic extension'' is indispensable. A property of one
extension does not identify the behaviour of every condition of a
nonhomogeneous forcing. For example, the lottery sum of trivial forcing and
Cohen forcing has a generic branch that adds nothing and another that adds a
real; the trivial branch preserves all saturation, but the whole forcing is
not countably distributive (source~08). The invariants $\delta(M,N)$ and
$\beta(M,N)$ belong to a particular pair of universes, not to the name of a
poset, and the theorems are formulated pair by pair. One can also apply
\cref{univ:cor:distributivity} below a specified condition.

\subsection{A Cohen real destroys even finite-parameter saturation}

Let $N=M[g]$ add a Cohen real, for instance with
$\mathbb P=\operatorname{Fn}(\omega,2,{<}\omega)$. The generic union is a new
subset of $\omega$: for every ground binary sequence the conditions disagreeing
with it somewhere are dense. No finite ordinal sequence is new, so
\[
 \delta(M,N)=\beta(M,N)=\omega.
\]
$\NoM$ has an $(\omega,\omega)$ gap and is not even $\omega$-saturated; the
rational-cut type \eqref{univ:eq:newrealtype} witnesses this directly. The pure
old order is $\omega$-saturated but not $\omega_1$-saturated. Nevertheless
$\NoM\preccurlyeq\NoN$ and $(\SCM,c_M)\preccurlyeq(\SCN,c_N)$ remain true, and
the pure field $\SCM$ remains saturated over every $N$-set; naming conjugation
makes it fail $\omega$-saturation again. Neither loss of real closedness nor
failure of elementarity explains the omitted type: the change lies in which
countable sets of formulas are available. Nothing here depends on Cohen
combinatorics beyond adding a real.

\subsection{Adding a subset of a regular cardinal}

\begin{proposition}[Any prescribed regular first gap; sources 01, 03, 04, 08]\label{univ:prop:add}
Let $\kappa$ be an infinite regular cardinal in $M$, and let $N=M[G]$ be
obtained by $\Add(\kappa,1)=\operatorname{Fn}_{<\kappa}(\kappa,2)$, the partial
functions $\kappa\to2$ with domain of size below $\kappa$, ordered by reverse
inclusion. Then
\[
 \delta(M,N)=\beta(M,N)=\kappa.
\]
$\NoM$ has an omitted $(\kappa,\kappa)$ gap and no omitted pair of size below
$\kappa$. If $\kappa>\omega$, no real is added, $\NoM$ and $(\SCM,c_M)$ are
$\omega$-saturated and $\kappa$-saturated, and they are not
$(\kappa^+)^N$-saturated. No cardinal arithmetic assumption is needed.
\end{proposition}
\begin{proof}
For uncountable $\kappa$ a descending chain of fewer than $\kappa$ conditions
has its union as a lower bound, by regularity. To decide a name for an
ordinal-valued sequence of length $\lambda<\kappa$, strengthen a condition
coordinate by coordinate, taking unions at limits and at the end; the result
decides the sequence as a ground sequence. So no such sequence is added, and
$\kappa$ stays a regular cardinal. The generic union $g:\kappa\to2$ is total
and new: given an old $h$, it is dense to put the opposite value at an unused
coordinate. Hence $\delta=\kappa$. No subset of an ordinal below $\kappa$ is
added, and $g$ codes a new subset of $\kappa$, so $\beta=\kappa$. Apply
\cref{univ:thm:firstgap,univ:thm:leastgap,univ:cor:spectrum,univ:thm:omega,univ:thm:conjugation}.
For $\kappa=\omega$ this is the Cohen case. The generic sign of length $\kappa$
is itself fresh, so its gap can also be read from
\cref{univ:thm:gapcorrespondence}.
\end{proof}

No assumption such as $2^{<\kappa}=\kappa$ is needed for this first-gap
calculation; the saturation failure is indexed by the successor computed in
the actual extension. The proposition does not compute the full gap spectrum,
does not claim that no larger gaps occur, and asserts nothing about the
preservation of cardinals above $\kappa$. For $\kappa=\omega_1^M$, no reals
are added, $\omega$- and $\omega_1$-saturation survive, and
$\omega_2^N$-saturation is lost: failures at different parameter sizes.
\emph{Status (batch 37, 28 September 2026).} The full spectrum is now computed:
every set-presented gap of $\NoM$ in this extension has character
$(\kappa,\kappa)$, with no cardinal arithmetic and for every ground alphabet
(\cref{univ:gs:thm:cohen}, Part~VIII).

\begin{example}[Under $\kappa^{<\kappa}=\kappa$; source 04]\label{univ:ex:addkappa}
If moreover $M\models\kappa^{<\kappa}=\kappa$ with $\kappa$ uncountable, then
$\abs{\Add(\kappa,1)}=\kappa$, so the forcing has the $\kappa^+$-chain
condition. Together with closure this preserves $\kappa$, $\kappa^+$ and the
regularity of $\kappa$ \cite{Jech}, so the failure of saturation occurs at the
ground model's $\kappa^+$. This hypothesis is stronger than needed for
\cref{univ:prop:add}; it is recorded because source~04 uses it for this extra
preservation clause.
\end{example}

\subsection{Prikry forcing separates the two countable notions}\label{univ:sub:prikry}

Assume $M$ has a measurable cardinal $\rho$, and force with classical Prikry
forcing from a normal measure on $\rho$. Its standard properties are that it
preserves cardinals, adds an increasing cofinal sequence $g:\omega\to\rho$,
changes $\cf(\rho)$ to $\omega$, and adds no bounded subsets of $\rho$; in
particular it adds no reals \cite{Prikry,Gitik,Benhamou}. These are imported
facts, not reproved here. The sequence $g$ is new: $M$ regards $\rho$ as
regular and uncountable. The range of $g$ is a new subset of $\rho$, since in
$M$ its order type $\omega$ and cofinality in $\rho$ would contradict regularity.
Consequently
\begin{equation}\label{univ:eq:prikry}
 \delta(M,N)=\omega,\qquad \beta(M,N)=\rho,\qquad \R^M=\R^N.
\end{equation}

\begin{corollary}[Prikry separation; sources 01, 03, 04, 08]\label{univ:cor:prikry}
In this Prikry extension:
\begin{enumerate}[label=\textnormal{(\roman*)}]
\item $\NoM$ and $(\SCM,c_M)$ are $\omega$-saturated but not
$\omega_1$-saturated, while $\SCM$ is set-saturated.
\item $(\No_{<\rho})^N=(\No_{<\rho})^M$, yet this set field is not
$\rho$-saturated in $N$: its all-plus ordinals acquire a countable cofinal
sequence \textup{(source 01)}.
\item Applying the block code to $g$ gives an omitted $(\omega,\omega)$ gap
whose simplest new separator $\code(g)$ has birthday exactly $\rho$
\textup{(source 08)}.
\item The characteristic sign $s$ of $\ran(g)$, of length $\rho$, is fresh, and
its gap has cofinalities $(\cf^N\rho,\cf^N\rho)=(\omega,\omega)$, not
$(\rho,\rho)$ \textup{(source 03)}.
\item The endpoints of the omitted cut obtained from the absolute interval
code are the finite rational expressions of \cref{univ:prop:finiteformula}
with $\alpha_j=g(j)$ \textup{(source 04)}.
\end{enumerate}
\end{corollary}
\begin{proof}
(i) \Cref{univ:thm:omega} gives $\omega$-saturation, the new countable sequence
and \cref{univ:thm:main-summary} give failure at $\omega_1$, and
\cref{univ:thm:conjugation,univ:thm:purecomplex} give the surcomplex clauses.
(ii) Every binary sign of length below $\rho$ codes a bounded subset of $\rho$;
the regularity clause of \cref{univ:thm:fragments} fails. (iii) Every finite
initial segment of $g$ is old, so \cref{univ:thm:freshcode} applies with
$\lambda=\omega=\delta$; $\bd(\code(g))=\sum_{n<\omega}(g(n)+2)=\rho$, since the
finite partial sums lie below the cardinal $\rho$ and are cofinal in it. (iv)
Every proper initial segment of $\ran(g)$ is finite, so every proper prefix of
$s$ is old, and $\ran(g)$ is new; apply \cref{univ:lem:canonicalpresentation}.
(v) is \cref{univ:prop:finiteformula}.
\end{proof}

\noindent[Added 4 October 2026, batch~93: in a Prikry extension of a model
with a measurable cardinal, Part~II of
\path{foundations-and-computation/definable-surreals-and-omnific-integers}
(batch-93 manuscript~02; AI-assisted, unrefereed, not formalized) shows that
the surreals definable from a real and ordinals are not closed under countable
Hahn sums. A summable sequence of ordinal-definable omnific monomials, one for
each point of the Prikry sequence, has a sum $z$ with
$\omega<z<\omega+\omega^{1/2}$ whose normal form recovers that sequence, so $z$
is definable from no real and ordinals (its Theorem~25.4,
\path{dsn:rp:thm:prikry-sum}); and there is a countable cut with
ordinal-definable endpoints and no interpolant definable from a real and
ordinals (its Corollary~25.5). Clause~(iii) above codes the same kind of new
cofinal sequence into an omitted gap of the old field.]

The parameters of the omitted cut all have birthdays below $\rho$ and were all
present in $M$; what is new is their countable arrangement. The alphabet in the
interval tree must be allowed to reach $\rho$: a test restricted to new reals,
or to new binary signs of countable length, would miss this example. The old
universe's cofinality of the birthday $\rho$ is the wrong cofinality for the
gap computation. \Cref{univ:cor:prikry}(ii) shows directly why regularity is an
independent clause in \cref{univ:thm:fragments}, and (iii) shows how far
$\beta=\rho$ can lie above $\delta=\omega$.

\subsection{Comparison table}

\begin{center}\small
\begin{tabular}{@{}P{0.25\textwidth}ccP{0.11\textwidth}P{0.22\textwidth}P{0.14\textwidth}@{}}
\toprule
Extension & $\delta$ & $\beta$ & New reals? & Old ordered field $\NoM$ & Old pure $\SCM$\\
\midrule
Cohen real & $\omega$ & $\omega$ & Yes & Not $\omega$-saturated & Set-saturated\\[3pt]
$\Add(\kappa,1)$, $\kappa>\omega$ regular & $\kappa$ & $\kappa$ & No & $\kappa$-saturated, not $(\kappa^+)^N$-saturated & Set-saturated\\[3pt]
Prikry at a measurable $\rho$ & $\omega$ & $\rho$ & No & $\omega$-saturated, not $\omega_1$-saturated & Set-saturated\\
\bottomrule
\end{tabular}
\end{center}
Adding conjugation to the last column reproduces the ordered-field column. The
table describes external saturation in the stated outer universe; it does not
alter the internal theorem that each universe's own surreal field realizes all
of that universe's set-sized cuts.

\section{Why order gaps do not become Cauchy limits}\label{univ:sec:topology}

The native fine topology of a surreal field uses all positive radii of that
field, not only real ones; its additive uniformity has the entourages
$\abs{x-y}<\varepsilon$ for $\varepsilon\in\NoM_{>0}$. On the surcomplex field
use either the surreal-valued modulus $\abs{a+bi}=\sqrt{a^2+b^2}$ (source~01)
or $\|a+bi\|_\infty=\max(\abs a,\abs b)$ (source~04); they define the same
product topology and uniformity. All claims are about neighborhoods and
set-indexed families; no set of all open classes is formed.

\begin{proposition}[External discreteness and Cauchy rigidity; sources 01, 03, 04]\label{univ:prop:topology}
For every pair $M\subseteq N$ under \textbf{(H)}:
\begin{enumerate}[label=\textnormal{(\roman*)}]
\item every $N$-set $A\subseteq\NoM$ is closed in the fine topology of $\NoM$ and
uniformly discrete: some old $\varepsilon>0$ is below every nonzero distance in
$A$;
\item every $N$-set-indexed Cauchy net in $\NoM$ is eventually constant, and
every convergent such net is eventually equal to its limit;
\item the same holds for $\SCM$ with either surcomplex modulus.
\end{enumerate}
No short-sequence hypothesis is needed.
\end{proposition}
\begin{proof}
The nonzero pairwise distances in $A$ form an $N$-set of positive old numbers.
\Cref{univ:thm:oneside} with left side $\{0\}$ (source~01), or directly
\cref{univ:lem:oldbounds} applied to the reciprocals (source~04), gives a
positive old lower bound $\varepsilon$. For $x\in\NoM\setminus A$ apply the same
argument to $\{\abs{x-a}:a\in A\}$ to get a neighborhood of $x$ missing $A$.
The range of a set-indexed net is an $N$-set; choosing $\varepsilon$ below all
nonzero pairwise distances in the range, the Cauchy condition for that one
entourage makes late values equal. For a convergent net, apply the separation
argument to its range together with its limit. For $\SCM$, the modulus of an
old nonzero number is old and positive; apply the same argument.
\end{proof}

For the interval tree, the widths $b_{f\res\alpha}-a_{f\res\alpha}$,
$\alpha<\lambda$, are positive old numbers with a common positive old lower
bound; the endpoint sequences are not Cauchy in the fine uniformity, and the
omitted cut is not evidence of a missing fine-Cauchy limit. Even a Prikry cut
with countably many endpoints leaves the proposition intact. A topology that
looks only at selected scales could behave differently, but it is a different
structure and is not used. Elementarity, saturation and fine completeness must
not be treated as interchangeable: the old field remains real closed and
elementary, loses specified types, and keeps this set-indexed rigidity.

\begin{lemma}[Prefix cones are open; source 04]\label{univ:lem:coneopen}
For every $s\in\NoN$, $\Cone(s)$ is open in the fine order topology of $\NoN$.
\end{lemma}
\begin{proof}
By \cref{univ:lem:signcone}, $\Cone(s)$ is the class of separators of
$L_s<R_s$. For $x\in\Cone(s)$, the set
$\{x-l:l\in L_s\}\cup\{r-x:r\in R_s\}$ of positive distances has a positive
lower bound in $\NoN$ by the cut theorem, and a smaller interval about $x$ keeps
all the canonical inequalities. If the set is empty, $\Cone(s)=\NoN$.
\end{proof}

\begin{theorem}[Closed, nowhere dense, and still set-discrete; sources 03, 04]\label{univ:thm:nowheredense}
$\NoM$ is closed in $\NoN$, and its subspace topology is its own order
topology. If $M\ne N$, $\NoM$ is nowhere dense in $\NoN$. The same statements
hold for $\SCM\subseteq\SCN$ in the product topologies.
\end{theorem}
\begin{proof}
If $x\in\NoN\setminus\NoM$, every extension of $x$ is new, because an old sign
has only old prefixes; so the open cone $\Cone(x)$ misses $\NoM$. Source~03's
route uses the open interval $(x\cat-,x\cat+)$, all of whose points extend $x$.

For every positive $r\in\NoN$ there is an old $0<e<r$
(\cref{univ:lem:oldbounds}); this gives equality of the intrinsic and subspace
topologies at every old point. Assume $M\ne N$ and choose a new $z$; replacing
it by $h=\abs z/(1+\abs z)$ gives a new $h\in(0,1)$, since otherwise
$\abs z=h/(1-h)$ would be old. Given $a\in\NoM$ and $r>0$, the point $a+eh$ lies
in $(a-r,a+r)$ and is new, since otherwise $h=((a+eh)-a)/e$ would be old. So no
old point is interior; a closed class with empty interior is nowhere dense.
(Source~03 instead uses $a+z/\gamma$ for an old ordinal $\gamma>\abs z/r$.) For
the complex fields use $\NoM\times\NoM\subseteq\NoN\times\NoN$: products of
closed coordinate classes are closed, the radius comparison gives the subspace
claim, and perturbing one real coordinate by $eh$ shows empty interior.
\end{proof}

The old field is therefore simultaneously cofinal, simplicity-initial,
elementary, closed and nowhere dense in $\NoN$. Order-cofinality must not be
mistaken for topological density, nor elementary inclusion for density. In the
outer field the whole interval of fillers of a fresh gap, the cone of
\cref{univ:thm:fillingcone}, is separated from every old point. No convergence
of an ordinary sequence to a missing point is asserted.

\begin{remark}[Relation to the foundations report]\label{univ:rem:foundations}
The internal statement is Theorem~12.1 (\path{found:thm:discrete}) of
\path{docs/foundations-and-computation/foundations}, which proves (i)--(iii) of
\cref{univ:prop:topology} for sets of the ambient universe, assuming its
Lemma~3.3 (\path{found:lem:bounds}, bounds and small positive lower bounds)
``for the permitted small families''. \Cref{univ:prop:topology} extends it to
\emph{new} $N$-sets of old numbers. Its only input is the existence of old
positive lower bounds, which holds in every same-ordinal extension even where
cut filling fails. In the Lean ledger (\path{docs/FORMALIZATION.md}), the
Cauchy-net clause of \path{found:thm:discrete} is stated under the hypothesis
\texttt{HasSmallCutFillers}, that is, under small cut \emph{filling}. The
external version shows that positive lower bounds suffice for that clause.
This is recorded as information for formalization; nothing here is formalized.

The foundations report's Section~7.7 (\path{found:sub:countable}) warns:
``\,`every set family has a bound' means every family represented by a set
\emph{in the relevant model}; it is not an absoluteness claim about every family
visible from a stronger metatheory.'' That warning concerns countable models
and different ordinals, and it stands. For same-ordinal outer universes this
report sharpens it: bounds of outer sets of old numbers \emph{are} absolute
(\cref{univ:lem:groundbound,univ:lem:oldbounds}), while cut filling is not
(\cref{univ:thm:main-summary}). The different-height analogue, the embedding
$\No_U\to\No_V$ for Grothendieck universes, is treated in that report and is
not used here. The catalogue's remark that universe-relative formulations
``require smaller-universe subsets and index types'' also concerns carriers of
different heights; in the present setting the index sets need not be
ground-small.
\end{remark}

\part{Class isomorphisms and real forms (Global Choice)}\label{univ:part:realforms}

\section{Pure-field class isomorphisms}\label{univ:sec:classiso}

\emph{Throughout this part assume \textbf{(GC)}}: the ambient theory is
$\GBC$ with Global Choice, and $M$ is a class of the class universe. Source~04
works with a definable $M$ and definable fields, which is a special case. All
assertions concern individual class maps; $\Aut(\cdot)$ is shorthand for
quantification over class field automorphisms, and no class-object of all class
maps is formed. At every ordinal stage of the extension arguments, the partial
fields and maps are sets; their histories below a given ordinal are sets by
class Replacement, and deterministic choices from a fixed set-like global
well-order permit set-valued transfinite recursion with class parameters and
its class union. No general recursion with class-valued stages, no class Zorn
lemma and no unrestricted class truth predicate is assumed.

\begin{theorem}[Class back-and-forth; prior result]\label{univ:thm:classbf}
Let $H_1,H_2$ be proper-class algebraically closed fields of the same
characteristic. Every isomorphism between set-sized subfields of $H_1$ and
$H_2$ extends to a class field isomorphism $H_1\to H_2$.
\end{theorem}

This is Theorem~9.1 (\path{saut:thm:acfhomogeneity}) of
\path{docs/surcomplex/surcomplex-field-automorphisms}, with the same
hypothesis (Global Choice). Sources~03 and~04 reprove it by the same route:
enumerate both fields along $\Ord$ by set-like global well-orders, extend at a
forth step by a root of the transported irreducible polynomial or by an element
transcendental over the current target field --- which exists by
\cref{univ:lem:algebraicset} and proper-class size --- do the symmetric back
step, and take unions at limits. Source~03 included the proof because applying
a theorem stated for set-sized fields to proper-class fields without such a
check would leave a foundational gap; the collection's statement is already
for proper-class fields.

\emph{Second route (source 08).} Let $H_1,H_2$ be proper-class algebraically
closed fields of characteristic zero that are saturated over every set of
parameters. Every set-sized partial elementary map $H_1\to H_2$ extends to a
class isomorphism. At a forth step, transport the complete type of the next
uncovered element over the current domain, a set by quantifier elimination;
it is finitely satisfiable in $H_2$ by partial elementarity, and set-saturation
of $H_2$ realizes it, outside the current range because the type records
inequality with every element of the domain. Back steps are symmetric, and the
least realization in a fixed global well-order makes each step specified. The
method is classical \cite{Ehrlich1989,vdDE2018}.

\begin{theorem}[Old and new pure surcomplex fields; sources 03, 04, 08]\label{univ:thm:pureiso}
For every $N$-set $A\subseteq\SCM$ there is a class field isomorphism
$\Phi:\SCM\to\SCN$ fixing $A$ pointwise. If $M\ne N$, no such isomorphism maps
$\NoM$ onto $\NoN$, and none satisfies $\Phi\circ c_M=c_N\circ\Phi$.
\end{theorem}
\begin{proof}
Both fields are proper classes, since each contains every surreal ordinal in its
real coordinate, and are algebraically closed in $N$
(\cref{univ:prop:absolute}). The identity on the field generated by $A$ is an
isomorphism of set-sized subfields; apply \cref{univ:thm:classbf}. (Along the
second route: the identity on $A$ is partial elementary by
\cref{univ:prop:absolute}, both fields are set-saturated by
\cref{univ:thm:purecomplex}.) If $\Phi[\NoM]=\NoN$, its restriction would be
an isomorphism of real closed fields, contrary to \cref{univ:cor:noniso}; class
Replacement supplies the set-image condition. An isomorphism intertwining the
conjugations would carry one fixed field onto the other.
\end{proof}

The inclusion $\SCM\hookrightarrow\SCN$ and the isomorphism $\Phi$ are different
maps: for $M\subsetneq N$ the inclusion is a proper embedding, and the existence
of $\Phi$ does not make it surjective. $\Phi$ is noncanonical. It need not
preserve old coordinates, the natural valuation, the real axis, simplicity, an
exponential, or a derivation, and no such compatibility is claimed.
\Cref{univ:thm:pureiso} settles, under Global Choice, source~01's question on
class isomorphisms $\SCM\cong\SCN$ (\cref{univ:sec:questions}).

\section{Gap invariants of surcomplex real forms}\label{univ:sec:realforms}

A \emph{real form} of an algebraically closed field $H$ of characteristic zero
is a real closed subfield $F\subseteq H$ with $H=F(i)$; its conjugation fixes
$F$ and sends $i$ to $-i$. This is the usage of the automorphism report
(\path{saut:sec:realforms}). For a real closed class field $F$, let
$\GapSpec_N(F)$ be the class of cofinality pairs of its set-presented gaps, and
let $g_N(F)$ be the least size of an omitted separated $N$-set pair, with the
bookkeeping value $g_N(F)=\infty$ (not a cardinal) when every set pair is
filled.

\begin{lemma}[Invariance; source 03]\label{univ:lem:invariant}
A field isomorphism between real closed class fields preserves their gap
spectra and least omitted-pair sizes. Two involutions of $\SCN$ are conjugate
by a class field automorphism if and only if their fixed real fields are
isomorphic.
\end{lemma}
\begin{proof}
A field isomorphism preserves nonzero squares and hence the unique order; its
restriction to an $N$-set is a set bijection by class Replacement, so it
preserves cardinalities, cofinalities and the existence of separators. The
second statement is the conjugacy criterion of \path{saut:sec:realforms}:
restrict a conjugating automorphism, or extend $f:F_1\to F_2$ by
$x+iy\mapsto f(x)+if(y)$. It is the transport mechanism, not a new invariant.
\end{proof}

\begin{theorem}[A forcing-sensitive real form without set-sized cofinal subsets; sources 03, 04, 08]\label{univ:thm:transport}
Let $M\subsetneq N$, choose $\Phi:\SCM\to\SCN$ by \cref{univ:thm:pureiso}, and put
\[
 j=\Phi c_M\Phi^{-1},\qquad H=\Phi(\NoM).
\]
Then $j$ is a field involution of $\SCN$ with fixed field $H$ and
$\SCN=H(i)$, and
\begin{align}
 \GapSpec_N(H)
 &=\{(\cf^N\bd(s),\cf^N\bd(s)):s\in\Fresh(M,N)\},\label{univ:eq:realformspectrum}\\
 g_N(H)&=\delta(M,N).\label{univ:eq:realformmin}
\end{align}
$H$ has no $N$-set cofinal or coinitial subset, and every $N$-set of elements
of $H$ is bounded above and below in $H$. Nevertheless $j$ is not conjugate to
$c_N$ by any class field automorphism of $\SCN$.
\end{theorem}
\begin{proof}
Transport of structure gives $\Fix(j)=\Phi(\Fix(c_M))=H$ and $\SCN=H(i)$; even
if $\Phi(i)=-i$ the quadratic extension is unchanged. The displayed equations
follow from \cref{univ:cor:gapspectrum,univ:thm:leastgap,univ:lem:invariant}.
A cofinal $N$-set in $H$ would pull back to a cofinal set in $\NoM$, contrary
to \cref{univ:lem:oldbounds}; the same holds for coinitial sets and for
bounds. Finally $\NoN=\Fix(c_N)$ fills every $N$-set cut, so
$g_N(\NoN)=\infty\ne\delta(M,N)=g_N(H)$; the fixed fields are not isomorphic,
and \cref{univ:lem:invariant} gives non-conjugacy.
\end{proof}

\begin{theorem}[Elementarily equivalent, not conjugate; source 04]\label{univ:thm:involutions}
In the situation of \cref{univ:thm:transport}:
\begin{enumerate}[label=\textnormal{(\roman*)}]
\item $\Fix(c_N)=\NoN$ and $\Fix(j)$ is isomorphic as an ordered field to
$\NoM$; both fixed fields are real closed, and $\SCN=\Fix(c_N)(i)=\Fix(j)(i)$;
\item $c_N$ and $j$ are not conjugate by any class field automorphism;
\item $(\SCN,c_N)$ and $(\SCN,j)$ are elementarily equivalent, but their
external saturation spectra differ: for uncountable $\kappa$ the second is
$\kappa$-saturated exactly when $\kappa\le\delta(M,N)$, whereas the first is
set-saturated; at $\omega$ the second is saturated exactly when $\R^M=\R^N$.
\end{enumerate}
\end{theorem}
\begin{proof}
(i) and (ii) are \cref{univ:thm:transport}; source~04 derives (ii) from
\cref{univ:cor:noniso}. Both expansions are interpreted, with the same
coordinate rules, in real closed fields, and $\RCF$ is complete, so they are
elementarily equivalent. Their spectra are those of their fixed fields by
\cref{univ:thm:conjugation} transported along $\Phi$, and
\cref{univ:cor:spectrum,univ:thm:omega,univ:lem:RCFsat}(iii) compute them.
\end{proof}

Three statements must not be confused: the embedded real axes are different;
the real axes are not isomorphic as fields; and the involutions are not
conjugate. Applying a phase automorphism to one real axis proves the first but
produces a \emph{conjugate} involution and an isomorphic fixed field. The
nonisomorphism is detected by external types, not by a first-order sentence.

\begin{remark}[What is new relative to the automorphism report]\label{univ:rem:whatisnew}
Proposition~B.1 (\path{saut:prop:conjugacycriterion}) of
\path{docs/surcomplex/surcomplex-field-automorphisms} states that an involution
of $\No[i]$ is conjugate to $c$ if and only if its fixed real field has the
set-cut interpolation property. It was present at the sources' pin as
\path{prop:conjugacycriterion}, and none of the four sources cites it. Read
in $N$, its easy direction (the property holds in $\NoN$ by Conway cuts and is
preserved by ordered isomorphisms) already shows that $j$ is not conjugate to
$c_N$: $\Fix(j)\cong\NoM$ omits an $N$-set cut whenever $M\ne N$
(\cref{univ:thm:main-summary}). The non-conjugacy clauses of
\cref{univ:thm:transport,univ:thm:involutions} are therefore consequences of
that proposition together with the spectrum theorem. What is genuinely new in
this part is:
\begin{enumerate}[label=(\alph*)]
\item \emph{no set-sized cofinal subset}: the transported real forms, like
$\NoN$ itself, have no set-sized cofinal subset, so the cofinality obstruction
of the automorphism report's Theorem~10.2
(\path{saut:thm:inequivalentrealform}, a real form with a countable cofinal
subset) cannot separate them; the obstruction lies inside a bounded cut;
\item \emph{several forms separated from one another} by the invariant $g_N$
and the gap spectrum (\cref{univ:thm:manygrounds}), which the interpolation
criterion alone does not do, since it only separates a form from $c$;
\item \emph{agreement with $c_N$ on ordinary complex numbers}
(\cref{univ:cor:agreement});
\item \emph{elementary equivalence with different external spectra}
(\cref{univ:thm:involutions}), and the exact values of the spectra.
\end{enumerate}
These continue, but do not settle, the automorphism report's open
``real-form and continuity classifications'' (its Section~14): no
classification of real forms or involutions is obtained.
\end{remark}

\emph{Related (batch 32).} The automorphism report's finite-symmetry sections
classify a different class of involutions: those of $\No[i]$ that preserve the
valuation, Hahn summation and $\mathbf{Oz}[i]$. Two of them are conjugate by such an
automorphism if and only if their restrictions to $\C$ are conjugate in $\Aut(\C)$,
equivalently if and only if the coefficient fixed fields are isomorphic, and
the fixed pair of each is a Hahn field over a real closed coefficient field of
transcendence degree $2^{\aleph_0}$ with its integer part
(\path{saut:fs:thm:main-involutions}, \path{saut:fs:prop:conjugacy}). That
classification assumes the valuation data; it says nothing about the
pure-field real forms transported here, whose compatibility with a valuation
structure is \cref{univ:q:compatible}.

\begin{corollary}[Agreement on prescribed ordinary data; source 03]\label{univ:cor:agreement}
Let $A\subseteq\SCM$ be a set-sized subfield stable under $c_M$ (as a subfield
of $\SCN$ it is then stable under $c_N$). The involution $j$ can be chosen to
agree with $c_N$ on every element of $A$. If $\R^M=\R^N$, it can be chosen to
agree with $c_N$ on all of $\C^N$.
\end{corollary}
\begin{proof}
Choose $\Phi$ fixing $A$ pointwise. For $a\in A$, $c_M(a)\in A$, so
$j(a)=\Phi(c_M(\Phi^{-1}(a)))=\Phi(c_M(a))=c_M(a)=c_N(a)$. When no reals are
added, $\C^M=\C^N$ is a set-sized common subfield stable under conjugation.
\end{proof}

Thus agreeing with ordinary conjugation on every ordinary complex constant does
not identify the surreal real axis even up to conjugacy. The obstruction is
not an unusual choice of sign for $i$ or a moved ordinary coefficient.

\subsection{Several ground real forms in one field}

\begin{theorem}[Different thresholds separate real forms; sources 03, 04]\label{univ:thm:manygrounds}
\begin{enumerate}[label=\textnormal{(\roman*)}]
\item \textup{(source 03)} Let $r\ge1$, and suppose that the models $M_a$,
$a<r$, are proper transitive inner models of $N$ with the same ordinals, each
available as a class and satisfying the standing hypotheses, and that the
cardinals $\delta(M_a,N)$ are pairwise distinct. Then $\SCN$ has $r+1$ pairwise nonconjugate field
involutions, including $c_N$, none of whose fixed fields has a set-sized
cofinal or coinitial subset. If each $M_a$ has the same reals as $N$, all of
them can be chosen to agree with $c_N$ on $\C^N$.
\item \textup{(source 04)} Suppose $M_1,\ldots,M_n$ satisfy the standing
hypotheses and their old surreal fields have pairwise different external
saturation spectra in $N$. Then $\SCN$ has $n$ pairwise nonconjugate
real-closed-field involutions realizing those spectra, and $c_N$ supplies one
more whenever none of the listed spectra is the full set-saturation spectrum.
\end{enumerate}
\end{theorem}
\begin{proof}
Apply \cref{univ:thm:transport} to each $M_a$ with its own isomorphism
$\Phi_a:\No^{M_a}[i]\to\SCN$. The fixed fields have
least omitted-pair sizes $\delta(M_a,N)$, and that of $c_N$ is $\infty$; all
values differ, so \cref{univ:lem:invariant} separates the conjugacy classes.
Cofinality and agreement come from
\cref{univ:thm:transport,univ:cor:agreement}. For (ii), a conjugacy between two
transported involutions would give an isomorphism of fixed fields, which
preserves external saturation; $\NoN$ is set-saturated.
\end{proof}

Part (ii) realizes a \emph{given} finite family of inner models; it does not
assert that arbitrary families of saturation spectra can be produced by
forcing without further hypotheses.

\begin{corollary}[Arbitrarily long finite lists; source 03]\label{univ:cor:finitetower}
For every positive integer $r$ there are finite forcing-tower constructions
whose final surcomplex field admits $r+1$ involutions as in
\cref{univ:thm:manygrounds}(i), all agreeing on every ordinary complex number.
No large-cardinal assumption is needed. This is a relative forcing
construction over a suitable ambient model, not an assertion that every
universe contains the requisite grounds.
\end{corollary}
\begin{proof}
Starting from $M_0$, build $M_0\subset M_1\subset\cdots\subset M_r=N$: at stage
$a$ choose in $M_a$ an uncountable regular cardinal $\kappa_a$ above all
previously chosen ones, and force with $\Add(\kappa_a,1)$. Each later stage is
closed under decreasing sequences of length below its own larger regular
cardinal, so it adds no ordinal sequences of length below $\kappa_a$ and
preserves $\kappa_a$ as a regular cardinal: a sequence in the final model of
length below $\kappa_a$ was present before the last forcing by closure, and
repeating through the finitely many later stages puts it in $M_a$. The generic
subset of $\kappa_a$ added at stage $a$ remains new over $M_a$ in $N$. Thus
$\delta(M_a,N)=\kappa_a$, pairwise distinct in $N$. No stage adds reals, so all
the models have the same ordinary complex field. Apply
\cref{univ:thm:manygrounds}(i).
\end{proof}

The finite tower is intentional: an infinite iteration needs support and
preservation arguments not supplied by the finite proof, and no claim is made
that an arbitrary proper class of gap invariants can be realized
simultaneously. \emph{Status (batch 37, 28 September 2026).} Infinite
set-indexed families are obtained in Part~VIII from product forcing rather than
iteration, with one common threshold and distinct higher spectra
(\cref{univ:gs:sec:realforms}); infinite iterations remain unanalysed
(\cref{univ:gs:q:iteration}).

\part{A Boolean completion of reverse simplicity (source 01)}\label{univ:part:boolean}

\section{A Boolean completion of reverse simplicity}\label{univ:sec:boolean}

Berenbeim's Question~2 asks for a complete class Boolean algebra in which the
reverse-simplicity order of surreal numbers embeds densely
\cite[p.~32]{Berenbeim}. The answer depends on the meaning of ``complete.'' We
give an explicit positive answer for joins of sets, and a negative answer for
joins of all subclasses under Global Choice. This section works in a single
universe and does not use the pair $M\subseteq N$.

This is a different use of the sign tree from numerical cuts: \emph{stronger}
means extending a sign, regardless of its numerical value. Identifying $-$ and
$+$ with $0$ and $1$, the partial order is
\[
 P={}^{<\Ord}2,\qquad p\leq_P q\ \Longleftrightarrow\ q\subseteq p.
\]
Each condition is a set, and $P$ is a proper class; it is the reverse of the
non-strict initial-segment simplicity order.

\subsection{A direct limit with set-coded elements}

For each ordinal $\alpha$ let $B_\alpha=\Pow({}^\alpha2)$, and for
$\alpha\leq\gamma$ define the cylinder embedding
\begin{equation}\label{univ:eq:cylinder}
 j_{\alpha\gamma}(A)=\{s\in{}^\gamma2:s\res\alpha\in A\}.
\end{equation}
These are injective Boolean homomorphisms preserving arbitrary set-indexed
unions and intersections, with $j_{\gamma\epsilon}j_{\alpha\gamma}=j_{\alpha\epsilon}$.
Identify pairs $(\alpha,A)$ with $A\subseteq{}^\alpha2$ when their images agree
at a common later level. A proper-class equivalence class must not be made an
element; instead represent each class by its unique \emph{least-level pair}.
Among the levels at most a representative's level there is a least one from
which it is a cylinder, and injectivity makes its subset unique. Taking the
least-level pair is a definable operation on set-coded representatives.

Let $\mathbb B$ be the class of canonical pairs. To perform a finite Boolean
operation, lift the arguments to a common level, operate in the power set
there, and take the canonical pair; the result is independent of the common
level by \eqref{univ:eq:cylinder}. Thus $\mathbb B$ is a class Boolean algebra
whose elements are sets.

\begin{theorem}[Set-complete dense completion]\label{univ:thm:booleanyes}
$\mathbb B$ is complete for every set-sized family of its elements. The map
\[
 e:P\longrightarrow\mathbb B\setminus\{0\},\qquad
 e(p)=[(\bd(p),\{p\})],
\]
is an injective dense order embedding; in particular
$e(p)\leq_{\mathbb B}e(q)\Leftrightarrow q\subseteq p$.
\end{theorem}
\begin{proof}
Let $S$ be a set of canonical pairs and choose an ordinal $\gamma$ at least
every level occurring in $S$. Lift all of them to $B_\gamma$ and take the union;
its canonical pair is their join, since any upper bound can be moved to a
common later level where the power-set union lies below it, and cylinder
embeddings preserve that union. Complements give meets; the empty join and
meet are $0$ and $1$.

A condition $p$ determines its nonempty cylinder of extensions. If $p$ extends
$q$, its cylinder lies in that of $q$; if they disagree, the cylinders are
disjoint; if $p$ is a proper prefix of $q$, the cylinder of $p$ is not contained
in that of $q$ (extend $p$ with the opposite sign at the next coordinate of
$q$). This gives the order equivalence and injectivity. Finally, if $b\ne0$ is
represented by $(\alpha,A)$ with $A\ne\varnothing$ and $p\in A$, then
$e(p)\leq b$, so the embedding is dense.
\end{proof}

The construction answers the existence question when ``complete'' means
set-complete, the size convention in which a class Boolean algebra has joins of
its set-sized subsets, as in Holy--Krapf--L\"ucke--Njegomir--Schlicht
\cite{HolyEtAl}. General existence and uniqueness questions for class forcing
are subtle; the construction works directly for this binary tree. It is not
claimed that every class partial order admits the same construction, nor that
this completion is the unique class completion.

\subsection{Why joins of all subclasses are impossible}

For the negative statement only, work in $\GBC$ with Global Choice. A class
Boolean algebra is \emph{class-complete} if every subclass of its carrier,
including proper subclasses, has a supremum in the carrier. This is strictly
stronger than set-completeness.

\begin{lemma}[Diagonal obstruction]\label{univ:lem:classdiag}
In $\GBC$ with Global Choice, a class Boolean algebra with an ordinal-indexed
proper-class antichain of nonzero elements is not class-complete.
\end{lemma}
\begin{proof}
Let $(a_\alpha)_{\alpha\in\Ord}$ be such an antichain. The carrier is a proper
class, and Global Choice gives a class surjection $h:\Ord\to\mathbb B$ (restrict
a set-like global well-order). Let $A=\{\alpha:a_\alpha\not\leq h(\alpha)\}$. If
$\{a_\alpha:\alpha\in A\}$ had a supremum $b$, then
\begin{equation}\label{univ:eq:diagiff}
 a_\alpha\leq b\quad\Longleftrightarrow\quad\alpha\in A.
\end{equation}
Indeed, for $\alpha\notin A$, $a_\alpha$ is disjoint from every member of the
subclass, so $\neg a_\alpha$ is an upper bound and $b\leq\neg a_\alpha$; since
$a_\alpha\ne0$, $a_\alpha\not\leq b$. The other direction is the definition of
supremum. Choose $\gamma$ with $h(\gamma)=b$; then
$\gamma\in A\Leftrightarrow a_\gamma\not\leq b\Leftrightarrow\gamma\notin A$, a
contradiction. All classes used are legitimate comprehension instances; no
collection of all classes is formed.
\end{proof}

\begin{theorem}[No all-subclass completion]\label{univ:thm:booleanno}
In $\GBC$ with Global Choice there is no class-complete Boolean algebra with a
dense order embedding of $P={}^{<\Ord}2$ into its nonzero part.
\end{theorem}
\begin{proof}
The signs $p_\alpha=+^\alpha\cat-$, $\alpha\in\Ord$, are pairwise incompatible.
Under a dense embedding their images are pairwise disjoint, since a nonzero
meet would contain the image of a condition, a common extension of
incompatible signs. They form an ordinal-indexed proper-class antichain; apply
\cref{univ:lem:classdiag}.
\end{proof}

The two answers are consistent: a set-sized family of cylinders has a common
level at which its join is computed, and a proper-class family need not.
Set-completeness neither implies nor purports to imply all-subclass
completeness. This also prevents the misleading claim that Question~2 requires
an inconsistency of ordinary set or class theory.

\part{Scope, questions and limits}\label{univ:part:scope}

\section{Scope and significance}\label{univ:sec:scope}

\subsection{What the new information consists of}

The usual statement that the full surreal class realizes every set-sized cut
is internal to a universe. The main theorem explains exactly how it changes
when the elements are held fixed while the allowed families of parameters
change. Its most distinctive features are these:
\begin{enumerate}
\item a cut with one actual ground side is filled even if its other side is a
new and arbitrarily large set;
\item a first new sequence can be encoded by a cut of the same cardinal length,
however large its ordinal range;
\item every set-presented failure of the cut principle comes from a unique
fresh sign, and is symmetric;
\item removing the distinguished real structure by passing to the pure
surcomplex language erases the obstruction entirely.
\end{enumerate}
Together they identify a concrete set-theoretic invariant inside the external
order of the old surreal field, and show that ``the model theory of surcomplex
numbers'' has no language-independent saturation answer.

The interval method (source~01) uses only set-sized interpolation inside the
ground order, enough room to split an interval into any prescribed set of
pairwise disjoint subintervals, and a ground-set containment property for
external parameter sets. Surreal numbers supply all three explicitly. Other
definable universal orders with the same properties may be worth studying; no
classification of such orders is claimed.

\subsection{Limits of the assertions}

The argument does not address exponential or differential expansions of $\No$,
the Berarducci--Mantova derivation, or expansions naming simplicity. Their
types are not classified by real closed cuts, so the sufficiency proof cannot
simply be copied after adding those operations; even absoluteness of the
individual old operations would not give a type-realization theorem in the
expanded language. An elementary reduct can be well behaved while an expansion
has additional omitted types \cite{EhrlichKaplan,vdDE2018}.

The results do not assert that every class function or arbitrary proper-class
cut is available in ZFC. Set-sized data, the ground predicate, and the stated
uses of class principles are hypotheses, not conventions to be supplied later.
The report is not a consistency proof for ZFC, $\GBC$, measurable cardinals or
the model configurations; the Prikry examples are conditional on a measurable
cardinal. No countable set model whose ordinals stop at an external height is
used; the same-ordinals hypothesis is essential to the birthday cover.

No theorem here is Lean-checked. The collection's existing formal proofs and
the source manuscripts are separate evidence. Successful typesetting is not
mathematical verification, and the finite Python checks
(\cref{univ:app:checks}) verify no transfinite claim.

\subsection{A formalization route without class-sized type objects}

These are proposed proof tasks, not existing declarations. A first target is
the branch-decoder theorem: a set ordinal, a set alphabet, a table of
intervals, and an abstract predicate for old sequences; the obligations are
sibling uniqueness, limit nesting, totality of the default decoder and its
agreement with an externally coded branch (source~01). Source~03 proposes
starting instead with a generic ordered sign tree and an initial subclass
predicate, proving
\cref{univ:lem:signscofinal,univ:lem:canonicalpresentation,univ:thm:gapcorrespondence}
as ordinal-combinatorial statements. The one-old-side theorem needs a
bounded-stage cover and an internal cut operation, and should quantify over
actual set-sized sides. For a universe-polymorphic proof assistant, the
fragment theorem (\cref{univ:thm:fragments}) is the cleanest end-to-end target:
all structures can live in a large universe while the two set models remain
distinct; full-class formulations become metatheorems or universe-relative
schemes, not an illicit same-universe type of all surreals. The class
back-and-forth should be formalized with set-stage maps and a class-theory
interface. The surcomplex results separate into two reusable lemmas: external
type realization for a proper-class algebraically closed field with ground-set
parameter covers, and interpretation of the ordered base field through fixed
squares in the quadratic complexification, which avoids repeating the
collection's quadratic-algebra development. The relation of
\cref{univ:prop:topology} to the Lean form of \path{found:thm:discrete} is
recorded in \cref{univ:rem:foundations}.

\section{Questions, and which of them the merge answered}\label{univ:sec:questions}

The sources' questions are recorded here with their status. They are
directions suggested by the results, not claims that a particular published or
catalogued question remains open.

\begin{question}[External spectra of expansions; sources 01, 08; open]\label{univ:q:exponential}
Determine the external saturation spectrum of old surreal exponential fields,
and more generally of canonical expansions (surreal exponentiation, a specified
derivation). What extra obstruction appears after naming them?
\end{question}
The main theorem gives a necessary condition, since the ordered-field reduct
must be saturated; no sufficiency theorem for an expanded language is proved.

\begin{question}[The cut-cofinality spectrum above $\delta$; source 01; asymmetric part answered by source 03]\label{univ:q:cutspectrum}
Source~01 asked to classify more of the external cut-cofinality spectrum than
its first threshold, and to decide which asymmetrical pairs can occur.
\end{question}
\emph{Status.} The asymmetric part is answered by source~03: every
set-presented gap is symmetric, and the spectrum is exactly
$\{(\cf^N\bd(s),\cf^N\bd(s)):s\in\Fresh(M,N)\}$
(\cref{univ:cor:gapspectrum}); source~01's ``set-character gaps'' are source~03's
``set-presented gaps''. What stays open is which symmetric pairs above $\delta$
occur for a specified forcing, that is, computing the outer cofinalities of the
lengths of fresh signs; \cref{univ:prop:add} computes only the first gap. This
is related to, and must not be conflated with, the fresh-function spectra of
Fischer--Koelbing--Wohofsky \cite{FKW}, whose spectrum-realization questions
are not resolved here.
\emph{Status (batch 37, 28 September 2026): answered for the following
forcings, open in general.} Part~VIII reduces every case to the fresh-function
spectrum of the pair: the characters are the outer cofinalities of the ground
regular cardinals that carry fresh functions (\cref{univ:gs:thm:transfer}). It
computes the full spectrum for one Cohen subset or one collapse of a regular
$\kappa$, with any ground alphabet and no cardinal arithmetic
($\{(\kappa,\kappa)\}$, \cref{univ:gs:thm:cohen}); for any nonzero number of
Cohen reals (\cref{univ:gs:prop:cohenreals}); for GCH Easton products (the Easton
closure, \cref{univ:gs:thm:easton}); for Prikry forcing ($\{(\omega,\omega)\}$)
and for Namba forcing under CH ($\{(\omega,\omega),(\omega_1,\omega_1)\}$)
(\cref{univ:gs:cor:prikry,univ:gs:cor:namba}); and for the intermediate grounds
of sparse products and of countable GCH products
(\cref{univ:gs:thm:finitecomplement,univ:gs:thm:coordinate}). What stays open is
the spectrum of an arbitrary forcing, of intermediate grounds of products with
uncountable index sets at their accumulation points, and of these constructions
without GCH (\cref{univ:gs:q:beyondeaston,univ:gs:q:accumulation,univ:gs:q:gch});
the Fischer--Koelbing--Wohofsky realization questions remain unresolved.

\begin{question}[Class isomorphisms of the pure surcomplex fields; source 01; settled under Global Choice]\label{univ:q:classiso}
Under which weak class-theoretic assumptions is there a class isomorphism
$\SCM\cong\SCN$?
\end{question}
\emph{Status.} Under $\GBC$ with Global Choice there is one, even fixing any
$N$-set of old parameters: \cref{univ:thm:pureiso} (sources 03, 04, 08), by the
collection's class back-and-forth (\path{saut:thm:acfhomogeneity}). Only the
question whether an assumption \emph{weaker} than Global Choice suffices
remains open. Set-saturation (\cref{univ:thm:purecomplex}) gives local
realizations but is not itself a global construction.

\begin{question}[Completeness of the invariant; sources 03, 08; open, one part answered]\label{univ:q:completeness}
For same-ordinal inner models $M_0,M_1$ of a fixed $N$, which invariants beyond
the least omitted-cut size distinguish $\No^{M_0}$ and $\No^{M_1}$ as class
orders or real closed fields in $N$? Under what hypotheses does equality of
$\GapSpec_N(\No^{M_0})$ and $\GapSpec_N(\No^{M_1})$ imply isomorphism?
\end{question}
\emph{Status.} Source~08's remark that its construction ``provides one family of
gap witnesses, but not a complete spectrum of all omitted cut characters'' is
answered for a fixed pair by source~03's \cref{univ:cor:gapspectrum}. The
classification question stays open: the spectrum records which cofinalities
occur, not how many gaps of each kind occur or how they are arranged, and a
classification may need an enriched incidence structure of cuts.
\emph{Status (batch 37, 28 September 2026): answered negatively for real closed
fields; open for pure orders and under extra hypotheses.} Equality of the full
gap spectra does not imply isomorphism as real closed fields, nor as pure fields:
\cref{univ:gs:thm:boolean} (source~12; ZFC and Cohen forcing only, with equal
saturation at every cardinal and field embeddability exactly inclusion of index
sets), \cref{univ:gs:thm:amplification} (source~09; continuum many fields in one
universe, CH sufficing for an instance) and \cref{univ:gs:thm:incomplete}
(source~11; a measurable cardinal). Invariants beyond the least omitted-cut size
that do distinguish old fields are the higher characters (sources 10, 11, 13, 14),
the real trace $\Tr(\NoM)=\R^M$ (sources 12, 09) and $\omega$-saturation
(source~11); counting how often a character occurs does not help, since every
occurring character occurs in every interval a proper class of times
(\cref{univ:gs:prop:ubiquity}). Within the countable coordinate family the full
spectrum does classify (\cref{univ:gs:prop:coordorder}). Still open: whether equal
full gap spectra \emph{and} equal real traces imply isomorphism
(\cref{univ:gs:q:trace}); the pure-order version, where no construction above
applies (\cref{univ:gs:q:order}); and positive classification theorems under
further hypotheses.

\begin{question}[Recovering more than outer cofinality; source 03; open]\label{univ:q:length}
Can an intrinsic invariant of the real closed field recover the least
\emph{length} of a fresh sign, rather than only outer cofinalities of such
lengths?
\end{question}
Birthday length is not part of the field language. The least length of a fresh
sign is $\beta(M,N)$ (\cref{univ:rem:betafresh}), but $\beta$ is defined through
the ambient sign tree, so this does not answer the question.

\begin{question}[Additional compatible structure; source 03; open]\label{univ:q:compatible}
Which transported real forms can also be required to preserve a prescribed
valuation structure, a surreal exponential, or a specified derivation on an old
set-sized substructure?
\end{question}
The pure-field extension lemma fixes arbitrary set-sized field data but does not
extend isomorphisms in these stronger languages. Homogeneous-universal results
for suitable surreal $H$-fields \cite{vdDE2018} have additional hypotheses and
should not be substituted without checking them.
\emph{Status (batch 32): related, not answered.} Among involutions of $\No[i]$
required to preserve the valuation, Hahn summation and $\mathbf{Oz}[i]$, the
automorphism report classifies the conjugacy classes by the coefficient
involution (\path{saut:fs:thm:main-involutions}, \path{saut:fs:prop:conjugacy};
see the note after \cref{univ:rem:whatisnew}). It does not decide which
transported real forms preserve such a structure.

\begin{question}[Infinite families and relative real forms; sources 03, 08; open, finite version answered]\label{univ:q:families}
Which external saturation spectra, and which infinite collections of regular
thresholds or entire gap spectra, can be realized simultaneously by fixed fields
of involutions of one $\SCN$ that agree on $\C^N$ and have no set-sized cofinal
subset?
\end{question}
\emph{Status.} The finite-threshold version is answered by
\cref{univ:cor:finitetower} (source~03), and spectra arising from given inner
models are realized by \cref{univ:thm:transport,univ:thm:manygrounds}. A
classification of arbitrary real forms, or of infinite families, needs genuine
iteration and spectrum analysis. This is also where the results continue the
automorphism report's open real-form classification (\cref{univ:rem:whatisnew}).
\emph{Status (batch 32).} Read in a single universe, with the threshold of a
real form taken to be the least cofinality of a set-presented gap of its fixed
field (the invariant of \cref{univ:cor:recover}), the clause on infinite
collections of regular thresholds is answered without forcing by the report
\path{docs/surcomplex/first-kappa-coefficients} (\path{fkc:sb:thm:involutions},
\path{fkc:sb:rem:univ}). Under Global Choice, for every class $\mathcal C$ of
infinite regular cardinals it gives pairwise nonconjugate involutions of
$\No[i]$ that agree with conjugation on $\C$, whose fixed fields have no
set-sized cofinal subset and have thresholds exactly the members of
$\mathcal C$; their saturation spectra are initial segments of the cardinals
and each gap spectrum is a single pair. Arbitrary saturation spectra, gap
spectra with more than one character, and families of old fields $\NoM$ remain
open, and no real form is classified.
\emph{Status (batch 37, 28 September 2026): answered for multicharacter
spectra and for families of old fields; open for arbitrary prescribed spectra.}
Part~VIII realizes, under Global Choice, arbitrarily large set-indexed families of
pairwise nonconjugate involutions of one $\SCN$ whose fixed fields are
isomorphic to actual old fields $\No^{M_A}$, share one saturation threshold, have
distinct multicharacter gap spectra and no set-sized cofinal subset, and agree with
$c_N$ on $\C^N$, or on any prescribed set-sized conjugation-stable core
(\cref{univ:gs:thm:boolinvolutions0,univ:gs:thm:samesatinvolutions,univ:gs:thm:coordinvolutions,univ:gs:thm:classinvolutions});
also continuum many with one \emph{identical} full spectrum, agreeing only on
$\C^M$, since agreement on $\C^N$ is impossible there
(\cref{univ:gs:thm:ampinvolutions,univ:gs:prop:coreobstruction}), and a $\ZFC$
version for one set field (\cref{univ:gs:thm:setforms}). Still open: which
arbitrary saturation spectra or gap spectra, prescribed in advance, can be
realized simultaneously; infinitely many real forms with one identical full
spectrum agreeing on $\C^N$ (\cref{univ:gs:q:finalconstants}); and any
classification of real forms.

\emph{Berenbeim's Question~2} (2020 notes) is answered in the precise
set-complete interpretation of \cref{univ:thm:booleanyes}, with the
class-complete interpretation refuted under Global Choice
(\cref{univ:thm:booleanno}). It is not an open question of the collection.

\section{What this report does not claim}\label{univ:sec:nonclaims}

Every limitation, disclaimer and priority caveat of the four sources is kept.
They are listed per source; items written for the merge are marked M. The
non-claims of sources 09--14 are listed in \cref{univ:gs:sec:nonclaims}.

\subsection*{Source 01}
\begin{enumerate}[label=01.\arabic*,leftmargin=3.2em]
\item The principal spectrum and decoding results are proposed contributions
with proofs, not certified historical priority; no Lean verification and no
independent referee review.
\item Classical surreal facts, quantifier elimination and standard forcing facts
are not claimed; several consequences are deliberately labelled consequences.
\item Berenbeim's Question~2 comes from 2020 research notes, not a peer-reviewed
conjecture. The construction may be folklore in class forcing; it is not a
claim that Boolean completions were previously unavailable, and general
class-forcing completion theory predates the notes. The notes' broader
suggestions about embeddings of universes are not settled.
\item External means computed in $N$, not in a larger metatheory; countable
transitive models must be read internally. No unrestricted class truth
predicate is assumed.
\item Field elementarity does not assert $M\preccurlyeq N$. The ground-bound
lemma is not a covering property of a forcing, and the containing set may be
much larger than the external set.
\item The decoder does not infer that an external branch is old because its
nodes are old. Interval widths are not claimed to tend to zero.
\item Regularity of $\delta$ is an elementary minimal-length fact, not a new
characterization of regular cardinals. The gap invariant excludes cuts with
proper-class cofinality on one side.
\item Set-saturation of the pure field does not produce a class isomorphism;
the nonisomorphism corollary requires set-sized images;
\cref{univ:prop:setACF} is a consequence, not a new field theorem.
\item The distributivity equivalence is standard; statements are pair by pair,
not about every generic extension of a poset unless so quantified; no
preservation of larger cardinals is asserted.
\item A topology looking only at selected scales is a different structure and
is not used.
\item The Boolean completion is not claimed for every class partial order nor
claimed unique; set-completeness does not imply class-completeness; the
negative result assumes Global Choice, with no optimality claim for that
assumption.
\item No classification of other orders with the interval-method properties is
claimed.
\item Exponential and differential expansions are not addressed. Arbitrary class
functions and proper-class cuts are not claimed available in ZFC. Not a
consistency proof for ZFC, $\GBC$, measurables or the model configurations;
Prikry is conditional on a measurable. Nothing is Lean-checked; typesetting is
not verification.
\item Its three questions are recorded with their current status in
\cref{univ:sec:questions}.
\item Its repository review was limited: the foundations report was read to
line~180, the automorphism report only in its opening text, and code searches
used a default-branch index. The repository was not modified or built.
Literature searches were not exhaustive. Its novelty claim was: candidate
contributions not located in the consulted literature or the stated repository
review, subject to independent checking and further priority research.
\end{enumerate}

\subsection*{Source 03}
\begin{enumerate}[label=03.\arabic*,leftmargin=3.2em]
\item Historical novelty is not certified by an exhaustive literature review,
peer review or formal verification; no named published conjecture is claimed
solved.
\item The forcing theory of fresh functions is not renamed as a new invariant;
quantifier elimination is standard; the collection's class back-and-forth,
conjugacy/fixed-field criterion and countably cofinal real form are not
claimed.
\item No set of all surreals, no class-object of all automorphisms, no class
Zorn lemma and no general class-valued recursion; sizes are relative to $N$;
proper-class size in $N$ is not proper-class size in every outside universe.
\item Principal cuts and cuts needing a proper-class presentation are not
considered.
\item The invariant recovers a threshold and cofinality data only, not the
inner model, the forcing or the sign lengths; equality of spectra is not proved
to classify the fields.
\item $\kappa=\omega$ must not be substituted into the field-saturation clause.
The invariant belongs to a pair of universes, not to a poset.
\item The $\Add(\kappa,1)$ computation does not compute the full gap spectrum or
exclude larger gaps, and asserts nothing about preservation of higher cardinals.
No Prikry property is reproved.
\item Set-saturation of a proper-class algebraically closed field is an
established mechanism, not a novelty claim; class back-and-forth already exists
in the collection; the inclusion is not the isomorphism $\Phi$; the conjugacy
criterion is the transport mechanism, not the new invariant.
\item The conjugation structures are elementarily equivalent; $\Phi$ need not
preserve simplicity, the natural valuation, the real exponential or a
derivation, and no such compatibility is claimed.
\item Finite towers only: no infinite iteration, and no arbitrary proper class of
invariants realized simultaneously.
\item Set-discreteness is classical and already in the collection; no
convergence of an ordinary sequence to a missing point is asserted.
\item Tempting shortcuts are false or unproved: replacing a new ordinal-valued
sequence by a binary one of the same length without a coding proof;
interchanging ground and outer cofinalities; letting global cofinality control
bounded gaps; inferring type preservation over new parameter sets from
elementary equivalence; assuming a pure-field isomorphism preserves the real
axis; inferring a full-class property from one set-sized workspace.
\item Not a Lean formalization; the repository was not rebuilt or
axiom-audited; the Python program cannot test freshness, forcing, uncountable
regularity or the class arguments; its proposed formalization order names
proof tasks, not existing declarations.
\item Its questions are directions, not assertions that a published question
remains open.
\item The repository comparison was focused, not exhaustive. Fischer--Koelbing--Wohofsky's
spectrum-realization questions are not resolved, and their convention for
spectra of a forcing must not be conflated with cofinalities in a particular
outer universe. The novelty statement is narrow; a specialist may identify an
earlier theorem subsuming part or all of the arguments.
\item Its shipped \path{03-fresh-sign-gaps-PROOF_STATUS.md} lists what was
established by argument, the established ingredients and seven boundaries, and
states that no Lean proof, theorem-prover verification, forcing simulation or
repository build was performed and that independent specialist review is still
appropriate.
\end{enumerate}

\subsection*{Source 04}
\begin{enumerate}[label=04.\arabic*,leftmargin=3.2em]
\item Proposed contributions; no exhaustive priority claim, no solution of a
named published problem, no Lean certification, no peer review.
\item Ehrlich's absolutely saturated models and the universal $H$-fields of van
den Dries--Ehrlich are precedents. External means internal to the specified
outer universe.
\item Definability of $M$ is a convenience. No claims about arbitrary externally
countable models; class isomorphisms add Global Choice.
\item Not new: cut realization in $\No$, quantifier elimination, saturation of
algebraically closed fields with enough independent elements, conjugation naming
a real closed fixed field, and many surcomplex automorphisms and real axes in the
collection.
\item Birthday minimization over old candidates is not claimed absolute;
``initial'' means under simplicity, not numerically; the cover lemma is not the
$\kappa$-cover property; no equivalence between $\eta_{\aleph_0}$ and
$\aleph_0$-saturation of a real closed field; no class-valued recursion.
\item No regularity assumption on $\kappa$; cardinals are computed in the outer
universe; its name for the first new length was local notation, not a claim to
a new forcing invariant; regularity is a set-theoretic fact.
\item In the pure-field section the contrast is the contribution, not a
classification of types.
\item Forcing statements quantify over every generic extension; one extension
is not enough. The comparison table does not alter internal saturation.
\item The class back-and-forth is noncanonical and preserves no coordinates,
valuation, real axis, simplicity or exponential. The several-spectra corollary
concerns a given finite family; arbitrary spectra are not claimed producible.
\item Full-class fine-topology results are already in the collection.
\item The full text of Ehrlich's 1989 paper was not accessible to it. No
extension to exponential or differential languages; the Rangel--Mariano and
Bertram questions are not resolved. Its novelty claim: proposed original
results, with priority remaining to be independently established.
\item A code-search response marked incomplete was not treated as an exhaustive
negative; the nonduplication statement is not a claim about later commits.
\item Its sensitive proof points (old elements are not an old set; a minimal
fresh length is a limit; ordinal labels cannot be replaced by binary digits;
class recursion is not silently assumed; external cardinal bounds belong to the
outer universe; finite parameters need not give a finite type; forcing
preservation quantifies over all generics; a pure field does not remember a
chosen real form) and its finite-check limits are kept; no verified priority or
solved named problem.
\end{enumerate}

\subsection*{Source 08}
\begin{enumerate}[label=08.\arabic*,leftmargin=3.2em]
\item Proposed contributions; classical theory credited; no Lean, no peer
review, no named problem solved, no exhaustive priority claim.
\item The algebraically closed field statement is not a new principle; class
back-and-forth is classical and the isomorphism is a corollary; exponentials,
derivations and unrestricted automorphisms are not classified.
\item Global Choice is used only for the class isomorphisms. No countable set
model and no set-valued field object is used.
\item The cover lemma is not a covering or approximation property; $\delta$ is
undefined when $M=N$; ordinal sequences must not be replaced by binary ones.
\item The algebraically closed field argument is a classical instance;
conjugation is not definable in the pure field.
\item Distributivity needs the uniform quantifier (lottery sum). Closure facts
are standard; no GCH; no claim about preserving higher cardinals; Prikry
properties are imported, not reproved.
\item Its real-form corollary does not classify involutions; the collection's
back-and-forth, conjugacy criterion and countably cofinal real form are prior
work.
\item An omitted cut is not a Cauchy sequence, and no topological completion
theorem is claimed. The ordered exponential field, the Berarducci--Mantova
derivation and simplicity expansions are not treated. Nonisomorphism compares
inner and ambient universes only; nothing is claimed about
$\No^{M_1}\cong\No^{M_2}$.
\item Its questions are not claims about previously catalogued open questions.
\item No repository file was modified and no independent build was done; the
finite checks do not establish transfinite claims.
\item It inspected a scan of Ehrlich's 1989 paper; its repository and literature
searches were limited, and absence of a search hit is weak evidence. Its
warranted conclusion: a specific relative-universe theorem package not found in
the inspected literature or repository material, with historical novelty
subject to broader expert checking.
\end{enumerate}

\subsection*{Added by the merge}
\begin{enumerate}[label=M.\arabic*,leftmargin=3.2em]
\item The merge adds no mathematical claim beyond the observation of
\cref{univ:rem:betafresh}; its other additions relate statements and reports,
or restate a source's statement in another source's notation.
\item The real-form results continue, and do not settle, the real-form
classification question of the automorphism report.
\item The relation to the Lean form of \path{found:thm:discrete}
(\cref{univ:rem:foundations}) is information for formalization; nothing in this
report is formalized.
\item The sibling report is cited by directory and by what it is described to
prove; this report makes no claim about its numbering or its other statements.
\item The source manuscripts were not rebuilt for the merge; their page counts
are as delivered.
\end{enumerate}
This gives 15, 16, 13 and 11 items from sources 01, 03, 04 and 08, and 5 added
by the merge: 60 in all.

\part{Full gap spectra, relative grounds and classification (sources 09--14)}\label{univ:gs:part}

\section{How Part VIII was assembled}\label{univ:gs:sec:assembly}

This part was added on 28 September 2026. It merges six manuscripts of that
date, called \emph{sources 09--14} here, into one addition to the report. The
numbers continue the report's own local sequence (its sources 01, 03, 04 and 08
are the manuscripts of Parts I--VII), and they are the file prefixes of the
delivered code and data. Parts I--VII, their numbering and their labels are
unchanged; the only edits outside this part are dated status pointers (on the
title page, in \cref{univ:sub:provenance}, after
\cref{univ:thm:leastgap,univ:prop:add,univ:cor:finitetower}, in the statuses of
\cref{univ:q:cutspectrum,univ:q:completeness,univ:q:families}, and at the head
of \cref{univ:sec:nonclaims,univ:app:checks}), five new bibliography entries,
and the notation macros of the preamble.

\subsection{Six manuscripts, one addition}\label{univ:gs:sub:provenance}

All six manuscripts were written against commit \gsfull{} of
\texttt{VladimirReshetnikov/ProveIt} (short \gspin, 25 September 2026), at which
this report was the four-source merge of Parts I--VII, unchanged since. They
arrived in commit \texttt{37e61c1fd} and were placed, as additions to this
report, in commit \texttt{0e53d1034} (batch~37 of the collection's intake,
manuscripts 01--06). They do not know one another. No source manuscript, source
README or source PDF is shipped; their code, data and audit files are, with the
prefixes listed in the README.

{\small
\begin{longtable}{@{}P{0.07\textwidth}P{0.34\textwidth}P{0.51\textwidth}@{}}
\toprule
Source & Manuscript (batch-37 number, archive, pages) & Contribution here\\
\midrule
\endhead
09 & \emph{Full Gap Spectra of Old Surreal Fields: forcing transfer, Easton
realization, and continuum many nonisomorphic fields with identical gap
spectra} (01, \path{proveit_gap_spectra}, 23 pages) & Transfer with
fixed-width blocks; the weak-approximation reading; Easton realization (base);
pointwise versus union spectra; common-small-forcing invariance by canonical
names (\cref{univ:gs:sec:amplification}); the rational-cut invariant (second
route to the real trace); Cohen amplification: continuum many old fields with
one full spectrum, pairwise nonisomorphic as pure fields; exact Prikry and
Namba spectra; same-spectrum involutions agreeing on $\C^M$ and the obstruction
to agreement on $\C^N$. Files prefixed \texttt{09-cohen-amplification-}.\\
10 & \emph{Exact Gap Spectra of Ground-Model Surreal Fields: fresh functions,
forcing, intermediate grounds, and many nonconjugate surcomplex real forms}
(02, \path{ProveIt_Exact_Surreal_Gap_Spectra}, 22 pages) & The delimiter
code; the filter-base bound (base of the density bound); the initial-segment
route to the Cohen and collapse singletons; Easton examples; isolated-coordinate
localization; the $2^\chi$ Boolean family with one threshold
(\cref{univ:gs:thm:booleanfamily}); involutions agreeing on any set-sized
conjugation-stable core, elementarily equivalent over it. Files prefixed
\texttt{10-isolated-coordinates-}.\\
11 & \emph{Beyond the First Omitted Cut: full gap spectra of old surreal fields,
Cohen collapse, and surcomplex real forms} (03,
\path{ProveIt_Full_Gap_Spectra}, 25 pages) & Base of the transfer; the
lower barrier; every omitted set cut is a genuine gap; the external density
bound; the generalized Cohen theorem for every ground alphabet (base); relative
finite-complement spectra of sparse products; same-saturation families of any
length; the Cohen--Prikry tower (\cref{univ:gs:sec:tower}); the tower bound.
Files prefixed \texttt{11-finite-complements-}.\\
12 & \emph{Beyond Gap Spectra: real traces and embeddability of old surreal
fields} (04, \path{proveit_beyond_gap_spectra}, 22 pages) & The real trace
and its functoriality (base); the Boolean family with constant spectrum
$\{(\omega,\omega)\}$ and embeddability exactly inclusion, with no large
cardinal, CH or GCH (\cref{univ:gs:sec:trace}); realization of every ground
set-sized partial order; the square-chain-condition bound; the generalized
Cohen theorem by explicit collapse; equivariant embeddings of involutions.
Files prefixed \texttt{12-real-traces-}.\\
13 & \emph{Cofinality Compression and Coordinate Spectra of Old Surreal Fields}
(05, \path{Surreal_Gap_Spectra_Research}, 22 pages) & The transfer for
arbitrary same-ordinal pairs; the Cohen and single-collapse singletons by
substring enumeration; name localization and intersections of mutually generic
extensions (base); the exact spectra of \emph{all} coordinate grounds of a
countable GCH product (base, \cref{univ:gs:thm:coordinate}); binary
intersections; local ubiquity (base); the continuum of equal-threshold real
forms. Files prefixed \texttt{13-coordinate-spectra-}.\\
14 & \emph{Beyond the First Gap: Exact Spectra and Boolean Families of Old
Surreal Fields} (06, \path{Surreal_Gap_Spectra_Research (1)}, 26 pages) &
Localization under distributivity only; finite Boolean cubes without
sparseness; the countable theorem at $\kappa_n=\aleph_{n+1}$ with its own
preservation proof (second route); exact gap pairs of set-sized birthday
cutoffs, including asymmetric boundary pairs (\cref{univ:gs:sec:cutoff}); a
ZFC theorem on real forms of one set field; the class version. Files prefixed
\texttt{14-relative-cutoffs-}.\\
\bottomrule
\end{longtable}}

\textbf{Why one addition.} Each manuscript answers one or more of this report's
named questions \cref{univ:q:cutspectrum,univ:q:completeness,univ:q:families},
and all six prove the same spine: the outer-cofinality pushforward from fresh
functions to gap characters. Printed side by side, they would repeat the
fresh-sign correspondence (already \cref{univ:thm:gapcorrespondence}) six times,
the transfer six times, the Cohen singleton five times and the class
back-and-forth (already \cref{univ:thm:classbf}) six times. Otherwise they are
complementary, and several answer each other's questions
(\cref{univ:gs:sec:questions}).

\textbf{Bases, and what is printed once.} Where two or more sources prove the
same statement it is printed once, with every proving source named at the
statement, from the \emph{base} named below; genuinely different proofs are
kept as marked second routes.
\begin{itemize}
\item Fresh-sign correspondence: all six reprove
\cref{univ:thm:gapcorrespondence,univ:thm:fillingcone}; it is not reprinted
(\cref{univ:gs:rem:signroutes} records their variants).
\item Transfer: base source~11, whose proof is explicitly collapse-sensitive;
the three block codes of the sources are all kept.
\item Density bound: base source~10 (an arbitrary cofinal base of the actual
generic filter, not required to lie in $M$); sources 11, 13, 14 (a ground dense
set) and 09 (the whole poset) are special cases, and source~12's square chain
condition is a second route.
\item Generalized Cohen spectrum: base source~11 (every ground alphabet with at
least two letters); second routes from sources 10 (initial-segment bases) and
12 (explicit collapse); sources 13 and 14 use source~11's route.
\item Easton realization: base source~09; sources 10 and 11 prove the same.
\item Name localization: statement from source~14 (distributivity suffices),
proof also by closure from sources 10, 11, 13 as a second route;
intersections: base source~13.
\item Countable coordinate spectra: base source~13 (any strictly increasing
uncountable regular $\kappa_n$); source~14 proves the case
$\kappa_n=\aleph_{n+1}$, with its own preservation proof as a second route.
\item Real trace: base source~12; source~09's rational-cut invariant is the same
set, with a second proof.
\item Equal full spectra without isomorphism: three constructions, all kept
(sources 12, 09, 11); source~12 has the weakest hypotheses.
\item Local ubiquity: base source~13; source~14's affine route kept.
\item Class back-and-forth: all six reprove \cref{univ:thm:classbf}; not
reprinted (\cref{univ:gs:rem:bfroutes}).
\end{itemize}

\textbf{Where the merge had to choose.} (a) Notation, fixed in
\cref{univ:gs:sub:notation}; no normalization changes. (b) Index direction:
subscripts on intermediate grounds always list \emph{retained} coordinates;
three sources indexed by omitted ones, and their statements are rewritten
(\cref{univ:gs:sub:notation}). (c) Questions answered inside the cluster are
recorded as answered (\cref{univ:gs:sec:questions}). (d) Two sentences of the
sources overstate what this report said; they are corrected where they are
printed (\cref{univ:gs:rem:overstatement1,univ:gs:rem:overstatement2}). Text
written for the merge is confined to this section, the framing remarks that
relate sources, the question statuses, the non-claims marked M and
\cref{univ:gs:sec:audit}; it adds no mathematical claim beyond the
comparisons stated as merge observations.

\subsection{What the addition answers}\label{univ:gs:sub:answers}

\begin{itemize}
\item \textbf{Transfer} (\cref{univ:gs:thm:transfer}; all six). For every
same-ordinal pair $M\subseteq N$, the characters of the set-presented gaps of
$\NoM$ are the outer cofinalities of the ground regular cardinals that carry
fresh ordinal-valued functions:
$\Char_N(\NoM)=\{\cf^N\lambda:\lambda\in\FF(M,N)\}$.
\item \textbf{Question~19.2, computed families.} One Cohen subset, or one
collapse, of an infinite regular $\kappa$ gives exactly $\{(\kappa,\kappa)\}$,
with no cardinal arithmetic (\cref{univ:gs:thm:cohen}); GCH Easton products give
exactly the Easton closure of their coordinates (\cref{univ:gs:thm:easton});
Prikry gives $\{(\omega,\omega)\}$ and Namba under CH gives
$\{(\omega,\omega),(\omega_1,\omega_1)\}$ (\cref{univ:gs:cor:prikry,univ:gs:cor:namba});
every coordinate ground of a countable GCH product has an explicit spectrum
(\cref{univ:gs:thm:coordinate}).
\item \textbf{Question~19.4, negative for fields.} Equality of the entire gap
spectrum does not imply isomorphism of old real closed fields, even with equal
saturation at every cardinal: \cref{univ:gs:thm:boolean} (source~12, ZFC),
\cref{univ:gs:thm:amplification} (source~09) and \cref{univ:gs:thm:incomplete}
(source~11, a measurable cardinal). The least threshold and even the whole
saturation spectrum do not classify (\cref{univ:gs:thm:booleanfamily},
source~10, and others).
\item \textbf{Question~19.7, families of old fields.} Arbitrarily large
set-indexed families of pairwise nonconjugate involutions of one $\SCN$, with
one saturation threshold and distinct multicharacter spectra, agreeing with
$c_N$ on $\C^N$ or on any prescribed set-sized conjugation-stable core
(\cref{univ:gs:sec:realforms}).
\end{itemize}
What stays open is recorded in the dated statuses of
\cref{univ:q:cutspectrum,univ:q:completeness,univ:q:families} and in
\cref{univ:gs:sec:questions}. In particular it is open whether equal real
traces \emph{and} equal full gap spectra imply isomorphism
(\cref{univ:gs:q:trace}), and whether equal full spectra imply isomorphism of
the pure orders (\cref{univ:gs:q:order}).

\subsection{Conventions for Part VIII}\label{univ:gs:sub:notation}

The standing hypothesis \textbf{(H)} and the class discipline of
\cref{univ:sub:levels} apply throughout. Every forcing is a set forcing; forcing
statements are relative, in the usual ground/extension sense, and none asserts
that $\ZFC$ proves the existence of a transitive model of itself or of a generic
filter over the universe. The real-form theorems use \textbf{(GC)}, except
\cref{univ:gs:thm:setforms}, which is a $\ZFC$ theorem about a set field.
``Set-presented gap'' has the meaning of \cref{univ:def:setgap}; all six sources
count exactly these, identify different presentations of one partition, and
exclude cuts with a proper-class side. For an ordered class $F$ of $N$ write
\[
 \Char_N(F)=\{\mu:(\mu,\mu)\in\GapSpec_N(F)\},
\]
the \emph{scalar} (diagonal) spectrum. For an old field $\NoM$ no information is
lost, since every set-presented gap is symmetric (\cref{univ:cor:gapspectrum}).
For the set-sized cutoff fields of \cref{univ:gs:sec:cutoff}, where asymmetric
pairs occur, the full pair set $\GapSpec_N$ is used.

The sources used different letters, and several letters also carry a meaning
elsewhere in this report. One symbol per notion is fixed here.

\begin{center}\small
\begin{longtable}{@{}P{0.3\textwidth}P{0.19\textwidth}P{0.43\textwidth}@{}}
\toprule
Notion & Symbol here & In sources 09--14\\
\midrule
\endhead
Scalar gap spectrum & $\Char_N(\NoM)$ & $\mathsf G(M,N)$ (09); $\mathcal G=\mathrm{Gap}_N(M)$ (10); $g(M,N)$ (11); $\mathrm{Char}_N$ (13); $\Gamma^N_M$, $\mathrm{DiagGap}$ (14)\\
Pair spectrum & $\GapSpec_N$ & $\GapSpec$ (09); $\mathrm{GapSpec}$ written $\mathrm{Gap}$ (12); $\mathrm{Gap}_N$ (13); $\mathrm{GapPairs}$ (14)\\
Fresh-function spectrum of the pair & $\FF(M,N)$ & $\mathsf F$ (09); $\mathcal F_{M,N}$ (10); $F(M,N)$ (11); $\mathfrak F$ (12); $\mathrm{Fresh}(M,N)$ (13); $\mathcal F$ (14)\\
``Some $f:\gamma\to\Ord$ is fresh'' & $\Fr_\gamma(M,N)$ & $\mathsf{Fresh}_\delta(M,N)$ (09)\\
Union-over-generics spectrum of a poset & $\FRESH(\mathbb P)$ & same (09, 12, 14; after \cite{FKW})\\
Easton closure & $\Ecl(A)$ & $\mathrm{EC}$ (09); $\mathrm{Ecl}$ (10, 13); $E(A)$ (11)\\
Real trace & $\Tr(F)$ & $\mathrm{RCut}(K)$ (09); $\mathrm{Tr}_{\R}$ (12)\\
Fixed-width one-plus block code & $\bcode_\theta(f)$ & unnamed (09, 11, 12, 14)\\
Delimiter code $\mathop{\cat}_\xi(+^{f(\xi)}\cat-)$ & $\dcode(f)$ & $\mathrm{Code}(f)$ (10)\\
Terminal-plus block code & $\codeplus(f)$ & $s_f$ (13); as in \cref{univ:prop:blockcodeplus}\\
Retained-coordinate ground & $M_A$ ($A$ retained) & $M_S$ (09), $M_A$ (12, 14): retained; $M_S$, $M_B$, $M_F$ (10, 11, 13): omitted\\
Old class fields & $\No^{M_A}$ written out & $K_S$ (09); $F_S$ (10); $K$, $K_\xi$ (11); $F_A$ (12); $K_F$ (13); $K_A$ (14)\\
Old set-sized cutoff fields & $(\No_{<\theta})^{M_A}$ & $F_A$ (14)\\
Surcomplex class & $\SC$ (roman) & $\mathbf{SC}$ (10--14)\\
Final universe & $N$ & $N$ (09--13); $W$ (14)\\
Intermediate universe & $W$ & $W$ (09, 10); $B$, $U$ (13, 14 in proofs)\\
Coordinate Cohen reals & $y_n$, $y_i$ & $c_n$ (09); $c_i$ (12); $c$ (11)\\
Transported involutions & $j_A$; isomorphisms $\Phi_A$ & $\jmath_S,\Phi_S$ (09); $j_S,\phi_S$ (10); $j_\xi,\varphi_\xi$ (11); $j_A,\phi_A$ (12, 14); $j_F,\phi_F$ (13); $J_A$ (14)\\
Generic filter of a small or Cohen factor & $G_{\mathbb Q}$, $G_{\mathbb C}$ & $H$ (09, 13)\\
Index cardinal of a family & $\chi$ & $\theta$ (10, 11)\\
$(2^{<\kappa})^M$ & $\sigma_\kappa$ & $\theta$ (12)\\
Bounding cardinals of sparse recursions & $\varsigma_i$, $\varsigma_\xi$ & $\sigma_i$ (10); $\delta_\xi$ (11)\\
Chain-condition cardinal & $\mu$ & $\chi$ (12)\\
Least omitted coordinate & $n_*$ & $b$ (14)\\
Range bound of a fresh function & $\theta$ & $\rho$ (09 Lemma~3.3, 14 Theorem~4.1); $\chi$ (11); $\theta$ (12)\\
\bottomrule
\end{longtable}
\end{center}

\emph{Index direction.} A subscript on an intermediate ground always lists the
coordinates it \emph{retains}: $M_A=V[G\res A]$. Sources 09, 12 and 14 used this
convention. Sources 10, 11 and 13 indexed by the \emph{omitted} coordinates; their
statements are rewritten. Thus source~10's $M_S$ is $M_{K_S}$ with
$K_S=\{\kappa_i:i\in\chi\setminus S\}$, source~11's $M_B$ is $M_{A\setminus B}$,
and source~13's $M_F$ is $M_{\omega\setminus F}$. Omitted sets are written $B$ and
never appear as a subscript. Source~13's subfamily ``$0\in F$'' and source~14's
``$0\notin A$'' are the same subfamily.

\emph{Letters reserved elsewhere in this report} keep their meaning here:
$\delta=\delta(M,N)$ is the first new length and $\beta=\beta(M,N)$ the first new
birthday (never a code length); $\theta$ is an ordinal bound on a range or a set
of birthdays (so source~14's cutoff $\lambda^{++}$ keeps the letter $\theta$,
and source~13's collapse alphabet $\theta$ is a range bound); $\rho$ is the
measurable cardinal of Prikry forcing; $K$ is not used for proper-class fields;
$c_M,c_N$ are conjugations, so Cohen reals are renamed; $H=\Phi(\NoM)$ keeps its
meaning in \cref{univ:thm:transport}, so generic filters are renamed; $g_N(F)$ is
still the least omitted-pair size (for an old field with $M\ne N$,
$g_N(\NoM)=\delta=\min\Char_N(\NoM)$), while $g$ also names generic functions, as in
\cref{univ:sub:prikry}; $\Fresh(M,N)$ is still the class of fresh \emph{signs},
and fresh-function spectra are written $\FF(M,N)$. The letter $\nu$ is not used.
Source-local letters for generic ordinals were renamed silently where they
clashed with these reservations.

\emph{Tempting false readings.} $\FF(M,N)$ is a set of \emph{ground} regular
cardinals, while $\Char_N(\NoM)$ is a set of regular cardinals \emph{of $N$}; they
are compared only through $\lambda\mapsto\cf^N\lambda$. $\FF(M,N)$ is the spectrum
of one actual pair of universes, not $\FRESH(\mathbb P)$, which quantifies over
possible generic extensions. An equality ``$\Char_N=\FF$'' is an identification of
the common regular cardinals of cofinality-preserving extensions only.
``Isomorphic'' for class fields always means a class isomorphism taking sets to
sets. A statement about the \emph{inclusion} order of old fields is not a
statement about their intersections or composita.

\section{The cofinality transfer}\label{univ:gs:sec:transfer}

\subsection{Fresh-function spectra of a pair}

\begin{definition}[Sources 09--14]\label{univ:gs:def:ff}
For a same-ordinal pair $M\subseteq N$ and an ordinal $\gamma$, write
$\Fr_\gamma(M,N)$ if some $f:\gamma\to\Ord$ in $N$ is fresh over $M$
(\cref{univ:def:delta}). The \emph{fresh-function spectrum} of the pair is
\[
 \FF(M,N)=\{\lambda\in\Reg^M:\Fr_\lambda(M,N)\},
\]
where $\Reg^M$ is the class of infinite regular cardinals of $M$.
\end{definition}

The range of a fresh function is bounded by an ordinal, even though the notation
permits arbitrary ordinal values; this makes all codes and name collections below
sets (source~09). The notion is Fischer--Koelbing--Wohofsky's \cite{FKW},
specialized to one actual pair. For a ground poset $\mathbb P$, their
$\FRESH(\mathbb P)$ collects the ground regular cardinals at which \emph{some}
condition forces a fresh function; the pointwise spectrum of one generic $G$ is
$\FF(M,M[G])$, and for homogeneous forcing the two agree
\cite[Proposition~4.3]{FKW}, homogeneity meaning that for any two conditions an
automorphism moves one to be compatible with the other. The distinction prevents
a false realization argument: a lottery sum can have a union spectrum with many
cardinals, although one generic chooses one summand, and the old surreal field of
one universe sees the pointwise spectrum (sources 09, 11, 12, 14). For example
(source~13), under GCH a lottery choice between adding a Cohen real and adding a
Cohen subset of $\omega_2$ has both thresholds across its generics, while a
particular generic chooses one summand; the lottery-sum caution of
\cref{univ:sec:examples} is the same point.

\begin{remark}[The fresh-sign correspondence, six times; merge note]\label{univ:gs:rem:signroutes}
All six sources reprove \cref{univ:thm:gapcorrespondence,univ:thm:fillingcone}
and credit them to Part~II: source~09 Theorem~2.5, source~10 Theorem~3.3,
source~11 Theorem~3.2 and Corollary~3.3, source~12 Theorem~3.4 (citing it as
``Theorems 8.4--8.5''), source~13 Theorem~3.4 (``now established without a
forcing hypothesis'') and source~14 Theorem~3.3. It is not reprinted. Their
surjectivity arguments take a separator $x$ of a presentation in $N$ and use its
\emph{least new prefix} (sources 09, 11, 12), or the simplest separator and the
minimality of its birthday (sources 10, 13, 14), where Part~II uses convexity of
$\Cone(y)$; these are one argument in three wordings. Source~09 adds the
endpoint argument of \cref{univ:gs:lem:omittedgap}, which is specific to old
surreal fields: for arbitrary dense orders an omitted set pair need not define an
endpoint-free gap. Source~11 records that the root of a fresh cone is the
simplest separator of its gap, the unique member of least birthday, so birthday
information is available in the ambient sign tree although it is not
automatically preserved by an ordered-field isomorphism. Source~10 records that
the correspondence classifies cuts, not Cauchy limits: a fresh sign may have very
large birthday even when its gap has countable character, a whole cone rather
than one point fills the cut, and nothing says the old field is dense in the new
one. Source~14 warns that the correspondence cannot be used blindly at a birthday
cutoff (\cref{univ:gs:sec:cutoff}).
\end{remark}

The old ordinal bounds of \cref{univ:lem:oldbounds} are likewise reproved by all
six. \emph{Second route to the positive lower bound (source~11).} If $X$ is an
$N$-set of positive old numbers, choose an old ordinal $\gamma_0$ above all their
birthdays and let $\varepsilon$ have a plus at position~$0$ and minus signs at
every later position through $\gamma_0$. Every positive $x$ begins with a plus,
and its first difference from $\varepsilon$ is a further plus or the end of $x$,
both greater than a minus; so $0<\varepsilon<x$. The argument does not require
$X\in M$, and it is not a covering property asserting that $X$ lies in an old set
of its own size.

\begin{lemma}[Every omitted set cut is a genuine gap; sources 11, 09]\label{univ:gs:lem:omittedgap}
Let $A,B$ be $N$-sets of elements of $\NoM$ with $A<B$ and no $x\in\NoM$ satisfying
$A<x<B$. Then both sides are nonempty, $A$ has no largest and $B$ no least
element, and their downward and upward closures partition $\NoM$ into a
set-presented gap.
\end{lemma}
\begin{proof}
If one side is empty, the ordinal bounds of \cref{univ:lem:oldbounds} give a
separator. If $A$ has a largest element $a$, apply the positive lower bound to
$\{b-a:b\in B\}$ and choose an old $0<\varepsilon<b-a$ for all $b\in B$; then
$a+\varepsilon$ separates the cut. A least element of $B$ is analogous. Let
$L=\{x:(\exists a\in A)\ x\le a\}$ and $R=\{x:(\exists b\in B)\ b\le x\}$. They are
disjoint, and a point in neither would be a strict separator, so they partition
$\NoM$; the endpoint observations and the cofinality of $A,B$ finish the proof.
\end{proof}

Source~09 gives the same argument from the reciprocal form of the bounds, and
stresses that different presentations of one partition are one gap. The lemma is
why an omitted pair is not counted merely because a chosen point of $\NoN$ is
absent from $\NoM$ (source~11). Source~12 records that a countable rational cut
around a new real can describe a larger interval in $\NoN$, but describes a
genuine partition with no separator in $\NoM$ when that real is new over $M$.

\subsection{Compression and three block codes}

\begin{lemma}[Compression to a ground regular domain; sources 09--14]\label{univ:gs:lem:compress}
Let $f:\alpha\to\Ord$ be fresh over $M$ in $N$, and put $\lambda=\cf^M\alpha$.
Then there is a fresh $h:\lambda\to\Ord$, so $\lambda\in\FF(M,N)$, and
$\cf^N\lambda=\cf^N\alpha$.
\end{lemma}
\begin{proof}
Freshness makes $\alpha$ a nonzero limit ordinal. Choose in $M$ a strictly
increasing cofinal $e:\lambda\to\alpha$, and an ordinal $\theta$ bounding the range
of $f$; it is an ordinal of $M$. The ground set
$T=\bigcup_{\xi<\alpha}({}^{\xi}\theta)^M$ has in $M$ an injective ordinal code
$c:T\to\eta$. Put $h(i)=c(f\res e(i))$ for $i<\lambda$. For $j<\lambda$, the
positions $e(i)$, $i<j$, are bounded below $\alpha$, since $e$ is increasing of
regular order type $\lambda$; so $h\res j$ is definable in $M$ from one old
proper restriction of $f$ and the old maps $e,c$, and is old. If $h$ were old,
decoding its values and taking their union would put $f$ in $M$. So $h$ is fresh.

The map $e$ remains strictly increasing and cofinal in $N$. A cofinal subset of
$\lambda$ maps to a cofinal subset of $\alpha$, and the least indices of $e$ above
a cofinal subset of $\alpha$ are cofinal in $\lambda$; so the two cofinalities
agree in $N$. No preservation of $\lambda$ as a regular cardinal of $N$ is assumed.
\end{proof}

Sources 09 and 10 state the lemma for ordinal-valued $f$, sources 11--14 for a
fresh sign; the existence part is a version of \cite[Proposition~5.2]{FKW}.
Recording \emph{both} the ground and the outer cofinality is what makes it useful
for old surreal fields (source~09).

A fresh ordinal-valued function need not give a fresh \emph{binary} function of
the same length: ordinal-valued and binary functions have different first
appearances in Prikry extensions (\cref{univ:rem:ordinalnotbinary}, source~11).
The sources use three codes, each of which lets the birthday grow while
preserving the outer cofinality. The terminal-plus code $\codeplus(f)$ of
\cref{univ:prop:blockcodeplus} is source~13's $s_f$; its freshness and cofinality
proof is the one printed there.

\begin{lemma}[Fixed-width block code; sources 09, 11, 12, 14]\label{univ:gs:lem:bcode}
Let $f:\lambda\to\Ord$ be fresh over $M$, $\lambda$ an infinite limit ordinal, and
let $\theta\ge2$ be an ordinal strictly above every value of $f$. Put
\[
 \bcode_\theta(f)(\theta\cdot\xi+\zeta)=
 \begin{cases}+,&\zeta=f(\xi),\\ -,&\zeta\ne f(\xi),\end{cases}
 \qquad(\xi<\lambda,\ \zeta<\theta).
\]
Then $\bcode_\theta(f)$ is a fresh sign of length $\theta\cdot\lambda$, and
$\cf^N(\theta\cdot\lambda)=\cf^N\lambda$.
\end{lemma}
\begin{proof}
The blocks $[\theta\cdot\xi,\theta\cdot(\xi+1))$ are disjoint and consecutive, so
this is one sign. Every proper prefix ends within, or at the start of, some block
$\xi<\lambda$ and is determined by $\theta$, its length and $f\res(\xi+1)$, which
is old since $\xi+1<\lambda$. If the whole sign were old, reading the unique plus
of each block would recover $f$ in $M$. The map $\xi\mapsto\theta\cdot\xi$ is
strictly increasing, continuous and cofinal, so the cofinalities agree.
\end{proof}

The four sources choose $\theta$ differently: $\theta=\theta_0+2$ with
$\theta_0>\sup\ran f$ (source~09, so that every block contains a minus and both signs
occur cofinally), an infinite cardinal of $M$ above the range (source~11), an
infinite ordinal above the range (source~12), or any $\theta\ge2$ above the range
(source~14). The code may increase the birthday greatly, and no optimal-birthday
assertion is made; an unbounded ordinal alphabet may need uncountable blocks,
which is why fresh ordinal functions need not be fresh reals (source~09).

\begin{lemma}[Delimiter code; source 10]\label{univ:gs:lem:dcode}
If $f:\alpha\to\Ord$ is fresh over $M$, the sign
\[
 \dcode(f)=\mathop{\cat}_{\xi<\alpha}\bigl(+^{f(\xi)}\cat-\bigr),
 \qquad \bd(\dcode(f))=\eta=\sum_{\xi<\alpha}(f(\xi)+1),
\]
is fresh, and $\cf^N\eta=\cf^N\alpha$. The sums are ordinal sums.
\end{lemma}
\begin{proof}
Let $b_0=0$, $b_{\xi+1}=b_\xi+(f(\xi)+1)$ and $b_\zeta=\sup_{\xi<\zeta}b_\xi$ at
limits; put plus signs on $[b_\xi,b_\xi+f(\xi))$ and a minus at its end. The
sequence $(b_\xi)_{\xi\le\alpha}$ is strictly increasing and continuous with
supremum $\eta=b_\alpha$, so $\cf^N\eta=\cf^N\alpha$. A proper prefix ends before
some $b_\zeta$, $\zeta<\alpha$, so it is a restriction of the code of the old
$f\res\zeta$. To decode, enumerate the minus positions in increasing order; the
order type of the plus interval between consecutive boundaries is the value of
$f$, at a limit boundary one takes the supremum of the earlier ones, and
consecutive minus signs encode zero values. This set recursion works in either
universe, so an old code would put $f$ in $M$.
\end{proof}

The assertion concerns cofinality, not equality of lengths: large ordinal entries
are not squeezed into a bounded number of binary coordinates, and the delimiters
prevent the loss of information that ordinal sums alone would cause
(source~10). Unlike $\code(f)$ and $\codeplus(f)$ of Part~I, $\dcode(f)$ has no
uniformly specified left endpoint in a block with a limit value; it is used here
only for freshness and cofinality.

\subsection{The transfer theorem}

\begin{theorem}[Cofinality transfer; sources 11 (base), 09, 10, 12, 13, 14]\label{univ:gs:thm:transfer}
For every same-ordinal pair $M\subseteq N$ under \textbf{(H)},
\begin{align}
 \Char_N(\NoM)&=\{\cf^N\lambda:\lambda\in\FF(M,N)\},\label{univ:gs:eq:transfer}\\
 \GapSpec_N(\NoM)&=\{(\cf^N\lambda,\cf^N\lambda):\lambda\in\FF(M,N)\}.\label{univ:gs:eq:pairtransfer}
\end{align}
If $M$ and $N$ have the same cofinalities, then $\Char_N(\NoM)=\FF(M,N)$,
identifying the common regular cardinals. In general the right side is a
many-to-one image of the ground spectrum.
\end{theorem}
\begin{proof}
A set-presented gap comes from a fresh sign $s$ (\cref{univ:thm:gapcorrespondence}),
with character $(\cf^N\bd(s),\cf^N\bd(s))$ (\cref{univ:lem:canonicalpresentation}).
Applying \cref{univ:gs:lem:compress} to $s$ gives $\lambda=\cf^M\bd(s)\in\FF(M,N)$
with $\cf^N\lambda=\cf^N\bd(s)$. Conversely, a fresh function on
$\lambda\in\FF(M,N)$ gives, by any of
\cref{univ:gs:lem:bcode,univ:gs:lem:dcode} or \cref{univ:prop:blockcodeplus}, a
fresh sign whose length has outer cofinality $\cf^N\lambda$, and hence a gap of
that character. Symmetry gives the pair form. If cofinalities are preserved, each
ground regular $\lambda$ remains regular with the same value.
\end{proof}

Source~13 stresses that the transfer holds for every same-ordinal inner-model
pair, whether or not $N$ is a set-forcing extension of $M$. Source~09 notes that a
spectrum may be a proper class for a general pair; in the set-forcing examples of
this part it is a set (\cref{univ:gs:thm:filterbase}). The coding lemmas are
closely related to \cite[Propositions~5.2--5.3]{FKW}; the transfer makes their
interaction with the fresh-sign correspondence explicit, and it is not claimed as
a new forcing-spectrum theorem (sources 09, 12).

\begin{corollary}[What a forcing calculation must compute; source 10]\label{univ:gs:cor:fibres}
In a cofinality-preserving extension, excluding a gap at a regular $\mu$ is
equivalent to excluding fresh ordinal-valued functions on $\mu$. Without
preservation, all ground regular domains of outer cofinality $\mu$ must be
considered together.
\end{corollary}

\begin{remark}[Two sources of information, three measurements; sources 09, 11]\label{univ:gs:rem:measurements}
Equation \eqref{univ:gs:eq:transfer} separates two operations: determine which
ground regulars carry fresh functions, then compute their cofinalities in the
actual outer universe. The second operation can identify distinct members of
the first spectrum, and forgetting it can invent uncountable gaps that do not
exist; computing only the least new length discards possibly noninterval
behaviour higher in the spectrum (source~09). The length of a fresh ordinal
function, the birthday of a fresh sign encoding it, and the character of its gap
can all differ. A ground regular cardinal can become singular, and then the last
measurement is $\cf^N\lambda$, not $\lambda$. Neither a ground cardinality bound
on a forcing nor its ground fresh-function spectrum can be relabelled as an
external gap spectrum without this step (source~11). An apparently large
ground-domain spectrum can give a singleton external spectrum (source~10;
\cref{univ:gs:thm:cohen}).
\end{remark}

\emph{A weak-approximation reading (source~09).} Following the terminology
discussed in \cite{FKW}, the weak $\lambda$-approximation property for $(M,N)$ is
the absence of a fresh ordinal-valued function on $\lambda$. When ground regular
cofinalities are preserved, the gap characters of $\NoM$ are exactly the failures
of these weak approximation properties at ground regular cardinals; in general,
they are the outer-cofinality image of those failures. This concerns \emph{weak}
approximation, not the stronger approximation property involving every small
intersection of an arbitrary new set.

\begin{corollary}[Lower barrier; sources 11, 10, 13, 14]\label{univ:gs:cor:lowerbarrier}
If $N$ adds no ordinal-valued sequences over $M$ of length less than an infinite
$N$-cardinal $\kappa$, then $\Char_N(\NoM)\cap\kappa=\varnothing$.
Conversely, if $\Char_N(\NoM)$ has no member below an uncountable $\kappa$, then no
ordinal-valued sequence of length below $\kappa$ is new (source~13).
\end{corollary}
\begin{proof}
The hypothesis says $\kappa\le\delta(M,N)$ when $M\ne N$, and
\cref{univ:thm:firstgap} bounds every character below by $\delta$. The sources'
direct proof: if a fresh sign $s$ had $\mu=\cf^N\bd(s)<\kappa$, an increasing
cofinal $e:\mu\to\bd(s)$ in $N$ would be old by hypothesis, and the old ordinal
codes of the prefixes $s\res e(i)$, $i<\mu$, would form an old sequence whose
decoded union is $s$. For the converse, the shortest new restriction of a new
sequence of length $\alpha<\kappa$ is fresh of limit length $\gamma\le\alpha$, and a
block code gives a gap of character $\cf^N\gamma\le\abs\alpha^N<\kappa$; this does
not need cofinality preservation (source~13). Both directions are the content of
\cref{univ:cor:recover}.
\end{proof}

\begin{corollary}[Uncountable saturation is read off the least character; sources 11, 13, 14]\label{univ:gs:cor:saturation}
For every uncountable cardinal $\lambda$ of $N$, $\NoM$ is $\lambda$-saturated as
an ordered field if and only if $\Char_N(\NoM)\cap\lambda=\varnothing$. If
$\Char_N(\NoM)$ has least element $\mu$, the uncountable saturation spectrum is
exactly $\{\lambda:\lambda\le\mu\}$.
\end{corollary}
\begin{proof}
This is \cref{univ:cor:spectrum} with $\mu=\delta(M,N)$ (\cref{univ:thm:leastgap}).
Source~11's direct proof: a gap of character $(\mu,\mu)$ with $\mu<\lambda$ gives a
finitely satisfiable omitted type over fewer than $\lambda$ parameters; conversely
every omitted cut presentation of size below $\lambda$ would, by
\cref{univ:gs:lem:omittedgap} and symmetry, give such a gap, and the real closure
argument of \cref{univ:lem:RCFsat}(i) converts filled small cuts into saturation.
\end{proof}

No regularity of $\lambda$ is required; the size estimate uses one parameter set
and its countable algebraic closure, not an unbounded union of $\lambda$ smaller
sets (source~11). The finite-parameter boundary is different: $\NoM$ is
$\omega$-saturated exactly when $\R^M=\R^N$ (\cref{univ:thm:omega}; reproved as
source~09 Proposition~6.3 and source~11 Theorem~9.3). A new rational cut can fail
to have an old realization although its gap character alone does not reveal
whether the cut came from a new real or from a much later fresh sign
(source~11). Pure dense orders are $\omega$-saturated in this finite-parameter
sense, since a finite parameter set specifies only an interval, an endpoint or an
exterior ray; the criterion is a \emph{field-language} phenomenon and does not
prove that the underlying pure orders of the examples below are nonisomorphic
(source~09). Source~09 recalls the criterion because it identifies a particularly
simple pair in the family of \cref{univ:gs:thm:amplification}; quantifier
elimination and the internal saturation argument are classical
\cite{EhrlichBSL,Marker}, and no truth predicate for all formulas over a class
field is used.

\section{External density and single-coordinate spectra}\label{univ:gs:sec:density}

\subsection{A ceiling measured in the extension}

\begin{definition}[Source 10]\label{univ:gs:def:filterbase}
For a generic filter $G\subseteq\mathbb P\in M$, a \emph{cofinal filter base} is a
set $B\subseteq G$ such that every $p\in G$ has some $b\in B$ with $b\le p$
(smaller conditions are stronger). Let $\mathfrak b_N(G)$ be the least
$N$-cardinality of such a base. This is notation local to this part, not the
bounding number of cardinal-invariant theory.
\end{definition}

\begin{theorem}[Filter-base ceiling; source 10 (base), sources 09, 11, 13, 14]\label{univ:gs:thm:filterbase}
Let $N=M[G]$ for a set forcing $\mathbb P\in M$, and let $B\subseteq G$ be a cofinal
filter base. If $f:\alpha\to\Ord$ is fresh over $M$, then $\cf^N\alpha\le\abs B^N$.
Consequently every $\mu\in\Char_N(\NoM)$ satisfies
\[
 \mu\le\mathfrak b_N(G)\le\abs{\mathbb P}^N,
\]
and $\mu\le\abs D^N$ for every dense $D\subseteq\mathbb P$ with $D\in M$. In
particular $\Char_N(\NoM)$ is a set.
\end{theorem}
\begin{proof}
Fix a ground name $\dot f$ and strengthen inside $G$ to a condition forcing that
$\dot f$ is a function on $\alpha$ with range in a specified ordinal. For each
$\xi<\alpha$ the old restriction $f\res\xi$ has a canonical ground name, and the
truth lemma gives $p_\xi\in G$ forcing $\dot f\res\xi=\check{(f\res\xi)}$. By
directedness of $G$ and cofinality of $B$ choose $b_\xi\in B$ below $p_\xi$ and
the fixed condition. If $\cf^N\alpha>\abs B^N$, some $b\in B$ equals $b_\xi$ for an
unbounded set of $\xi<\alpha$, since otherwise $\alpha$ would be a union of fewer
than $\cf^N\alpha$ bounded sets. Then for every $\zeta<\alpha$ the single ground
condition $b$ decides the value of $\dot f(\zeta)$, and in $M$ the forcing relation
and Replacement construct the function sending $\zeta$ to that value. In $N$ it is
$f$, contradicting freshness. Apply this to the fresh sign of each gap
(\cref{univ:thm:gapcorrespondence}). The generic $G$ is a base inside the ground
set $\mathbb P$, and $G\cap D$ is a base for every dense $D\in M$.
\end{proof}

The base need not belong to $M$: a poset may be large in $M$ while its actual
generic filter has a much smaller cofinal chain in $N$, and the argument
reconstructs $f$ from one ground condition, not from an allegedly ground sequence
of conditions (source~10). The use of $\abs D^N$ rather than $\abs D^M$ is the
point: a forcing with a large ground presentation may enumerate it by a much
shorter generic sequence, and the sharper bound is invisible if ground
cardinalities are substituted throughout (source~13). If a forcing collapses its
own old family of conditions to size $\kappa$, the ceiling rules out every
character above $\kappa$, even when the ground forcing had fresh functions on
several larger ground cardinals; their outer cofinalities must fall below the
ceiling (source~11). Sources 11, 13 and 14 state the dense-set form (source~11
Theorem~5.1, source~13 Theorem~5.1, source~14 Lemma~5.1), and source~09 the case
$D=\mathbb P$ with outer cofinality (source~09 Lemma~5.3, where $\mathbb P$ is a
small forcing $\mathbb Q$ over a universe $V$ and $U=V[G_{\mathbb Q}]$): if
$\cf^U\gamma>\abs{\mathbb Q}^U$ then $\Fr_\gamma(V,U)$ fails. For ground regular
cardinals this is the familiar size bound recalled in \cite[Corollary~2.3]{FKW};
the outer-cofinality form is convenient when another forcing has already changed
cofinalities. Source~14 calls it the usual repeated-condition argument. For every
nontrivial set-forcing extension, the character spectrum is therefore a set,
although the field and the class of gaps are proper classes; this separates three
issues, namely large birthdays, many fresh ground domains and large externally
visible characters, of which only the last is the cut invariant (source~13).

\begin{corollary}[Two-sided singleton criterion; sources 10, 11, 13]\label{univ:gs:cor:singleton}
Suppose $N=M[G]$ adds no ordinal sequence shorter than $\kappa$, has a cofinal
filter base (for instance $G\cap D$ for a ground dense $D$) of $N$-size at most
$\kappa$, and contains a fresh function whose domain has $N$-cofinality $\kappa$.
Then $\Char_N(\NoM)=\{\kappa\}$.
\end{corollary}
\begin{proof}
\Cref{univ:gs:cor:lowerbarrier,univ:gs:thm:filterbase} and the given fresh function.
\end{proof}

\subsection{Generalized Cohen forcing}

Fix an infinite regular cardinal $\kappa$ of $M$ and a set $A\in M$ with at least
two elements, and let $\mathbb P=\Fn_{<\kappa}(\kappa,A)^M$ be the partial functions
$\kappa\to A$ of ground domain size below $\kappa$, ordered by extension. The
superscript matters: this is one fixed ground poset. Let $g=\bigcup G:\kappa\to A$.

\begin{lemma}[Universal occurrence; source 11]\label{univ:gs:lem:word}
For every ground word $w:\eta\to A$ with $\eta<\kappa$ there is $\gamma<\kappa$ with
$\gamma+\eta<\kappa$ and $g(\gamma+\xi)=w(\xi)$ for all $\xi<\eta$. Consequently the
ground set $T=(\bigcup_{\eta<\kappa}{}^\eta A)^M$ has cardinality at most $\kappa$
in $N$.
\end{lemma}
\begin{proof}
Conditions placing $w$ on some interval are dense: regularity bounds $\dom(p)$
below $\kappa$, choose $\gamma$ beyond that bound, note $\gamma+\eta<\kappa$ since
$\kappa$ is an infinite initial ordinal, and add the prescribed values; the domain
still has size below $\kappa$. Each such dense set belongs to $M$, so genericity
gives the first assertion for every ground word. In $N$ there are only $\kappa$
pairs $(\gamma,\eta)$ of positions and lengths, each determining at most one word
by reading $g$, and every member of $T$ appears; so $\abs T^N\le\kappa$.
\end{proof}

\begin{lemma}[The old poset becomes small; source 11]\label{univ:gs:lem:posetsize}
In $N$, $\abs{\mathbb P}=\kappa$, and $\mathbb P$ adds no ordinal-valued sequence of
length below $\kappa$; in particular $\kappa$ stays regular.
\end{lemma}
\begin{proof}
Choose a ground finite $m\ge1$ and an injection $\iota:A\sqcup\{\square\}\to A^m$
(for finite $A$ take $m$ large; for infinite $A$, Choice gives $m=1$), the extra
symbol recording an undefined coordinate. For $p\in\mathbb P$ let
$\eta=\sup\{\xi+1:\xi\in\dom p\}<\kappa$, write the word of length $\eta$ with
$\xi$-th symbol $p(\xi)$ or $\square$, and replace each symbol by its $m$-letter
code. The resulting old word of length $m\cdot\eta<\kappa$ determines $p$ by
fixed-size block decoding, so $\mathbb P$ injects in $M$ into $T$ and
$\abs{\mathbb P}^N\le\kappa$ by \cref{univ:gs:lem:word}. The single-coordinate
conditions with a fixed letter give $\kappa$ distinct conditions. The forcing is
$<\kappa$-closed in $M$ (a union of fewer than $\kappa$ descending conditions is a
condition, by regularity), so the usual recursive decision argument adds no
ordinal-valued sequence of length below $\kappa$ and preserves the regularity of
$\kappa$; these are therefore inequalities of $N$-cardinals.
\end{proof}

The conclusion is not that the generic has listed all \emph{new} short words over
an arbitrary alphabet; only the specified ground set is counted, which is exactly
the set the ceiling needs (source~11).

\begin{theorem}[Generalized Cohen spectrum; source 11 (base), sources 10, 12, 13, 14]\label{univ:gs:thm:cohen}
For every infinite regular $\kappa$ of $M$ and every $A\in M$ with at least two
elements, if $N=M[G]$ for $G$ generic over $\Fn_{<\kappa}(\kappa,A)^M$, then
$\kappa$ remains regular,
\[
 \Char_N(\NoM)=\{\kappa\},\qquad \GapSpec_N(\NoM)=\{(\kappa,\kappa)\},
\]
and the least birthday of a new surreal number is $\kappa$. No assumption on
$2^{<\kappa}$ or $\kappa^{<\kappa}$ is made, and no assertion that cardinals above
$\kappa$ are preserved.
\end{theorem}
\begin{proof}
By \cref{univ:gs:lem:posetsize} no ordinal-valued sequence of length below
$\kappa$ is added, so all shorter binary signs are old and
\cref{univ:gs:cor:lowerbarrier} excludes characters below $\kappa$. Fix a ground
surjection $\pi:A\to\{-,+\}$ and put $s=\pi\circ g$. Every proper prefix of $g$ is
old (conditions whose domain contains a given $\eta<\kappa$ are dense), so every
proper prefix of $s$ is old. For each old $r:\kappa\to\{-,+\}$ it is dense to make
$\pi\circ g$ differ from $r$ at a coordinate outside the condition's domain, using
a letter that projects to the opposite sign. So $s$ is fresh of length $\kappa$,
giving a gap of character $(\kappa,\kappa)$ and a new surreal of birthday $\kappa$.
Finally \cref{univ:gs:lem:posetsize} and \cref{univ:gs:thm:filterbase} with
$D=\mathbb P$ exclude every character above $\kappa$.
\end{proof}

This answers the single-coordinate case of \cref{univ:q:cutspectrum} at the level
of the complete external spectrum (source~10); \cref{univ:prop:add} computed only
the first gap and $\delta=\beta=\kappa$, and its remark that larger gaps were not
excluded is now settled for $\Add(\kappa,1)$. The upper bound concerns \emph{all}
gaps, not only a chosen generic cut, and applies when the alphabet is so large that
the forcing collapses many ground cardinals (source~11). The theorem does not say
that $\kappa$ is the only ground regular domain supporting a fresh function; the
transfer, not an unqualified equality of ground and outer indices, controls that
issue (sources 10, 12, 14). Indeed the ground-indexed spectrum of $\Add(\kappa,1)$
is $\FRESH(\Add(\kappa,1))=[\kappa,(2^{<\kappa})^M]\cap\Reg^M$
\cite[Proposition~6.1]{FKW}; the theorem says that all these ground regulars have
outer cofinality $\kappa$ (\cref{univ:gs:lem:collapse}), which resolves the
single-Cohen case left open beyond the first gap (source~14). Fresh signs of larger
birthdays may exist, but their outer cofinalities cannot give larger characters
(source~13).

\begin{example}[Binary Cohen forcing and large alphabets; source 11]\label{univ:gs:ex:cohen}
Taking $A=2$ gives the complete spectrum of $\Add(\kappa,1)=\Fn_{<\kappa}(\kappa,2)$
at every infinite regular $\kappa$; an assumption such as $2^{<\kappa}=\kappa$ is not
needed for the external gap calculation, and the proof does not assert that the
ground fresh-function spectrum of this forcing is a singleton. If instead $A$ is a
ground infinite cardinal $\tau\ge\kappa$, the generic function maps onto $A$
(requiring each $a\in A$ in its range is dense), so the forcing collapses $\tau$ to
cardinality at most $\kappa$; its old surreal gap spectrum is still exactly
$\{\kappa\}$. This covers the partial-function presentation of the corresponding
collapse forcing. Different ground regulars that contribute fresh functions are
merged by their new cofinalities, as \eqref{univ:gs:eq:transfer} requires.
\end{example}

\emph{A constructive certificate (source~11).} The lower bound supplies an explicit
sign representative, not merely an existential type: its cofinal and coinitial
witnesses are the prefixes ending immediately before a plus or a minus, of external
size $\kappa$, and no smaller subfamily suffices by the first-disagreement argument.
At the other end, universal occurrence compresses all ground conditions into
$\kappa$ actual intervals of the generic word; a hypothetical larger fresh sign would
force one deciding condition to repeat cofinally, turning the supposedly new sign
into a ground set.

\emph{Second route: initial-segment bases (source~10).} For $A=2$ (source~10
Theorem~5.5), the set $B=\{g\res\alpha:\alpha<\kappa\}$ is a cofinal filter base of
size $\kappa$: every condition has bounded domain, and a sufficiently long prefix of
$g$ extends it; each $g\res\alpha$ is a ground condition in $G$, since conditions
whose domain contains $\alpha$ are dense. \Cref{univ:gs:cor:singleton} applies
directly, with no counting of the ground tree. \Cref{univ:gs:thm:filterbase} applied
to $g$ also shows that no cofinal base has size below $\kappa$. For $\kappa=\omega$,
the closure assertions mean finite closure and the unchanged regularity of~$\omega$.

\begin{corollary}[The same singleton after an ordinal collapse; source 10]\label{univ:gs:cor:collapse}
Let $\tau\ge2$ be an ordinal and $N$ a generic extension by
$\Coll(\kappa,\tau)=\Fn_{<\kappa}(\kappa,\tau)^M$. Then $\Char_N(\NoM)=\{\kappa\}$. If
$\tau\ge\kappa$, the generic map is onto $\tau$, so arbitrarily large ground
cardinals can be collapsed while the external spectrum stays a singleton.
\end{corollary}
\begin{proof}
Union gives $<\kappa$-closure; the generic $g$ has all proper prefixes in $M$, is new
by the disagreement dense sets, and its initial-segment conditions form a cofinal
base of size $\kappa$. Apply \cref{univ:gs:cor:singleton}. For $\zeta<\tau$, putting
$\zeta$ into the range is dense.
\end{proof}

The comparison proves that the external gap spectrum does not remember which
cardinals the forcing collapsed. It does \emph{not} prove that the resulting old
fields are isomorphic, nor does it give two nonisomorphic fields with equal full
spectrum; that classification issue is separate (source~10, and
\cref{univ:gs:sec:trace,univ:gs:sec:amplification,univ:gs:sec:tower}).

\emph{Third wording: substring enumeration (sources 13, 14).} Source~13 (alphabet
$A=2$ or an ordinal $A\ge\kappa$, its Theorem~1.2) and source~14 ($A=2$, its
Theorem~5.3) prove \cref{univ:gs:thm:cohen} by the argument of
\cref{univ:gs:lem:word}: initial-segment conditions are dense, so the old tree
$T=({}^{<\kappa}A)^M$ is a dense presentation of the forcing; every old $t\in T$
occurs in $g$ as a contiguous block $\seq{g(\gamma+\xi):\xi<\eta}$; there are at
most $\kappa$ pairs $(\gamma,\eta)$ in the extension, all blocks are old by
closure, and $\kappa$ distinct constant strings give $\abs T^N=\kappa$. For
$\kappa=\omega$ the argument uses only finite closure and Cohen-style genericity;
for $\kappa>\omega$ it does not assume $\kappa^{<\kappa}=\kappa$; for $A\ge\kappa$
the union takes every ground value, which explains the single-collapse
interpretation; and the result concerns the entire spectrum, not only gaps whose
separator has birthday $\kappa$ (source~13).

\subsection{The square chain condition}\label{univ:gs:sub:square}

\begin{lemma}[Mitchell-type square-chain-condition lemma; source 12]\label{univ:gs:lem:square}
Work in $M$. Let $\mu\le\lambda$ be infinite regular cardinals, and suppose
$\mathbb P\times\mathbb P$ has the $\mu$-chain condition. Then no $\mathbb P$-generic
extension contains a fresh function on $\lambda$ over $M$.
\end{lemma}
\begin{proof}
Suppose $p$ forces that $\dot f:\lambda\to\Ord$ is fresh. Below $(p,p)$ in the square,
let $D$ be the set of pairs whose coordinates decide different values of $\dot f$ at
some common position. $D$ is dense: otherwise below some $(r,s)$ no disagreement is
possible; for each $\alpha<\lambda$ choose an extension of $s$ deciding
$\dot f(\alpha)=a_\alpha$; every extension of $r$ deciding that coordinate gives the
same value, so $r$ forces $\dot f$ to be the ground function $\alpha\mapsto a_\alpha$,
contradicting freshness. Choose a maximal antichain $\mathcal A\subseteq D$ and for
each member a disagreement position $\alpha_a$. The chain condition gives
$\abs{\mathcal A}<\mu\le\lambda$, and regularity gives
$\beta=\sup\{\alpha_a+1:a\in\mathcal A\}<\lambda$. Every member of $\mathcal A$ forces
that the two copies of $\dot f\res\beta$ differ, so by maximality $(p,p)$ does. But
freshness gives $q\le p$ deciding $\dot f\res\beta$ as one old sequence, and $(q,q)$
forces the copies to agree.
\end{proof}

This is \cite[Proposition~2.2]{FKW}, which credits Mitchell and Unger antecedents;
the proof is reproduced for transparency. The \emph{square} is essential: a c.c.c.\
Suslin-tree forcing can add a fresh function on $\omega_1$, so the hypothesis is not
replaced by the c.c.c.\ of $\mathbb P$ alone (source~12).

\begin{corollary}[Source 12]\label{univ:gs:cor:size}
If a forcing has a dense subset of infinite cardinality $\tau$ in $M$, it adds no
fresh function on a ground regular cardinal greater than $\tau$.
\end{corollary}
\begin{proof}
Pass to the dense subset; its square has size at most $\tau$, hence the
$\tau^+$-chain condition. Apply \cref{univ:gs:lem:square}.
\end{proof}

This is a second route to the ground-size case of \cref{univ:gs:thm:filterbase},
through the chain condition instead of the generic filter; it bounds ground regular
indices, not outer cofinalities.

\begin{proposition}[Real-adding square-c.c.c.\ extensions; source 12]\label{univ:gs:prop:squareccc}
If $\mathbb P\times\mathbb P$ is c.c.c.\ in $M$ and $N=M[G]$ contains a new real, then
$\GapSpec_N(\NoM)=\{(\omega,\omega)\}$.
\end{proposition}
\begin{proof}
A new binary real is fresh on $\omega$. \Cref{univ:gs:lem:square} excludes fresh
functions on uncountable ground regulars. The c.c.c.\ preserves cardinals and
cofinalities: each coordinate of a name for an ordinal sequence has only countably
many possible values on a maximal antichain, so a name for a sequence of length
below an uncountable regular $\lambda$ has range in a ground set of size below
$\lambda$ and is not cofinal in it. Apply \cref{univ:gs:thm:transfer}.
\end{proof}

\begin{proposition}[Any nonzero number of Cohen reals; source 12]\label{univ:gs:prop:cohenreals}
For every nonempty set $J\in M$, if $N=M[G]$ for $\Add(\omega,J)^M$, then
$\FF(M,N)=\{\omega\}$ and $\GapSpec_N(\NoM)=\{(\omega,\omega)\}$.
\end{proposition}
\begin{proof}
Conditions are finite partial binary functions on $J\times\omega$. An uncountable
family of their domains has an uncountable $\Delta$-system subfamily; after fixing
the finitely many values on its root, any two members are compatible. This is the
Knaster property. The square is itself a Cohen forcing on two disjoint copies of the
coordinates, hence c.c.c.\ by the same argument. Each coordinate adds a fresh binary
real, so \cref{univ:gs:prop:squareccc} applies.
\end{proof}

The fresh-spectrum equality is noted after \cite[Corollary~2.4]{FKW}. What
\cref{univ:gs:sec:trace} uses is that it remains true in \emph{every intermediate
coordinate model}, where the quotient is again a Cohen forcing; no generic filter is
required to meet dense sets from an inappropriate universe (source~12).

\emph{Fourth route to \cref{univ:gs:thm:cohen} for $A=2$: explicit collapse
(source~12, its Theorem~10.1).} It needs only a size bound, closure and an explicit
collapse, and applies to each actual generic.

\begin{lemma}[Dense tree and closure; source 12]\label{univ:gs:lem:densetree}
Write $\sigma_\kappa=(2^{<\kappa})^M$. The forcing $\Fn(\kappa,2,{<}\kappa)^M$ has a
dense subset of cardinality $\sigma_\kappa$, namely the binary sequences of length
below $\kappa$. It is $<\kappa$-closed and adds no ordinal-valued sequence of length
below $\kappa$.
\end{lemma}
\begin{proof}
A condition has bounded domain by regularity, so it extends to a binary function on
an initial segment; there are $2^{<\kappa}$ of these. A union of a decreasing sequence
of fewer than $\kappa$ conditions is a condition. To decide a name for a sequence of
length $\mu<\kappa$, decide each coordinate recursively, taking lower bounds at limits
and at the end. This also prevents $\kappa$ from acquiring a shorter cofinal sequence
or being collapsed.
\end{proof}

\begin{lemma}[Explicit collapse of the relevant ground cardinals; source 12]\label{univ:gs:lem:collapse}
If $\lambda$ is regular in $M$ and $\kappa\le\lambda\le\sigma_\kappa$, then
$\cf^N\lambda=\kappa$.
\end{lemma}
\begin{proof}
For each nonzero cardinal $\mu<\kappa$ of $M$, partition $\kappa$ by a ground bijection
into $\kappa$ labelled blocks of size $\mu$. For a fixed $a\in(2^\mu)^M$ it is dense
that the generic function equal $a$ on some block: a condition meets fewer than
$\kappa$ blocks, and filling an unused block still gives a condition. So every old
binary $\mu$-sequence appears on some block, each block readout is old by
\cref{univ:gs:lem:densetree}, and $\abs{(2^\mu)^M}^N\le\kappa$. The claim is immediate
for $\lambda=\kappa$. If $\kappa<\lambda<\sigma_\kappa$, some $\mu<\kappa$ has
$\lambda\le(2^\mu)^M$, so $\abs\lambda^N=\kappa$. If $\lambda=\sigma_\kappa>\kappa$ is
regular in $M$, the supremum defining $\sigma_\kappa$ is attained, since otherwise its
cofinality in $M$ would be at most $\kappa<\lambda$; the same conclusion follows. So
$\cf^N\lambda\le\kappa$, and it cannot be smaller, since a shorter cofinal sequence
would be a new ordinal sequence of length below $\kappa$.
\end{proof}

\begin{proof}[Proof of \cref{univ:gs:thm:cohen} for $A=2$ along this route]
Closure adds no fresh function on a ground regular below $\kappa$;
\cref{univ:gs:cor:size} and the dense tree exclude ground regulars above
$\sigma_\kappa$. So $\FF(M,N)\subseteq[\kappa,\sigma_\kappa]\cap\Reg^M$, and
\cref{univ:gs:lem:collapse} sends every member to outer cofinality $\kappa$. The
generic binary sequence on $\kappa$ is fresh, so $\kappa\in\FF(M,N)$; apply
\cref{univ:gs:thm:transfer}.
\end{proof}

This route illustrates a second kind of information loss: the field spectrum
forgets not only which rational cuts are present (\cref{univ:gs:sec:trace}) but also
which ground regular lengths merged to the same cofinality after forcing. An analysis
in which ``regular cardinal'' and ``cofinality'' are silently read in different
universes would miss this collapse (source~12).

\section{Easton products}\label{univ:gs:sec:easton}

Work in $V\models\GCH$. For a regular $\alpha$ write $\mathbb C_\alpha=\Add(\alpha,1)^V$,
and for a set $A\subseteq\Reg^V$ let $\mathbb P_A=\prod^{\mathrm E}_{\alpha\in A}\mathbb
C_\alpha$ be the Easton product: coordinatewise order, supports bounded below every
inaccessible cardinal and unrestricted at other limits \cite[Definition~4.7]{FKW}. Under
GCH the regular limit cardinals relevant to the product are inaccessible. All posets
and supports are ground objects.

\begin{definition}[Fischer--Koelbing--Wohofsky; sources 09, 10, 11]\label{univ:gs:def:easton}
A set $A$ of infinite regular cardinals is \emph{Easton closed} if every limit point
$\mu$ of $A$ satisfies: if $\mu$ is regular and not Mahlo, then $\mu\in A$; if $\mu$ is
singular, then $\mu^+\in A$. The least Easton-closed superset of $A$ is $\Ecl(A)$.
\end{definition}

The Mahlo exception is part of the definition and cannot be dropped; Mahlo has its
usual meaning (inaccessible, with stationarily many inaccessibles below). The final
published formulation is used, not a simplified rule that closes at every regular
limit, nor the more restrictive wording of earlier talk notes (sources 09, 10, 11).

\begin{theorem}[Published input \cite{FKW}; imported]\label{univ:gs:thm:fkw}
Assume $V\models\GCH$. Then $\mathbb P_A$ preserves cardinals and cofinalities, is
homogeneous, and for every $V$-generic $G$ has $\FF(V,V[G])=\Ecl^V(A)$.
\end{theorem}

Preservation is \cite[Proposition~4.10]{FKW}; the spectrum is their Theorem~4.24 with
the realization statement in Corollary~4.25; the passage to the actual generic is
their homogeneity discussion, Proposition~4.3. The proof includes substantive forcing
and pcf arguments, which are imported and not claimed (sources 09, 10, 11).

\begin{theorem}[Exact Easton realization; source 09 (base), sources 10, 11]\label{univ:gs:thm:easton}
If $G$ is $\mathbb P_A$-generic over $V\models\GCH$ and $N=V[G]$, then
\[
 \Char_N(\No^V)=\Ecl^V(A),\qquad \GapSpec_N(\No^V)=\{(\mu,\mu):\mu\in\Ecl^V(A)\}.
\]
Every ground Easton-closed set of regular cardinals is the exact full set-presented
gap-character set of an old surreal field in a cofinality-preserving set-forcing
extension.
\end{theorem}
\begin{proof}
Homogeneity identifies the published spectrum with the pointwise spectrum of this $G$,
and cofinality preservation makes $\lambda\mapsto\cf^N\lambda$ the identity on
$\Reg^V$. Apply \cref{univ:gs:thm:transfer,univ:gs:thm:fkw}; for the last assertion
take $A$ Easton closed.
\end{proof}

This computes the entire spectrum for the specified family of forcings, a broad family
within \cref{univ:q:cutspectrum}. It is not a characterization of the spectra of all
cofinality-preserving forcings, does not assert that Easton closure is necessary for
arbitrary forcing spectra, and is not a proof of a more general spectrum-realization
conjecture (sources 09, 10, 11).

\begin{example}[A genuine hole; sources 09, 10]\label{univ:gs:ex:hole}
For $A=\{\aleph_1,\aleph_3\}$ there are no limit points, so $\Ecl(A)=A$. The old field
has exactly the characters $\aleph_1$ and $\aleph_3$ and no set-presented gap of
character $\aleph_2$. The smallest character does not determine the full spectrum even
in very elementary product examples; the example is relative to a GCH ground, and all
three cardinals are preserved.
\end{example}

\begin{example}[A character with no coordinate; sources 09, 10, 11]\label{univ:gs:ex:singular}
For $A=\{\aleph_{2n+1}:n<\omega\}$ the only limit point is the singular $\aleph_\omega$,
so $\Ecl(A)=A\cup\{\aleph_{\omega+1}\}$ (adding the successor creates no new limit
point), and
\[
 \Char_{V[G]}(\No^V)=\{\aleph_{2n+1}:n<\omega\}\cup\{\aleph_{\omega+1}\}.
\]
The character $\aleph_{\omega+1}$ occurs although no coordinate sits there, and there
are no others; the rule ``one character per coordinate'' is false for the full
spectrum. The extra character is the surreal image of the published singular-limit
phenomenon, not the result of a finite-support heuristic. The product is
$<\aleph_1$-closed and adds no reals, so multicharacter behaviour here is not a
new-real phenomenon.
\end{example}

\begin{example}[Same first character, different full spectra; source 09]\label{univ:gs:ex:samefirst}
The products indexed by $\{\aleph_1\}$ and $\{\aleph_1,\aleph_3\}$ have the same first
omitted character but different entire spectra. This compares the two extensions; by
itself it is not a simultaneous construction of two old fields in one universe, which
needs the further arguments of \cref{univ:gs:sec:relative}.
\end{example}

\begin{example}[Why the Mahlo clause matters; source 10]\label{univ:gs:ex:mahlo}
At a regular non-Mahlo limit point the closure adds that cardinal; at a Mahlo limit
point it imposes no such requirement by that clause alone. Whether a cardinal belongs
to the closure must be computed from the stated operator, including the original set,
not from the slogan ``close under regular limits''.
\end{example}

At the saturation threshold the higher characters carry strictly more information:
for uncountable $\kappa$, external saturation is controlled by the least new length
(\cref{univ:thm:main-summary}), which in the nonempty Easton examples is $\min A$,
while \cref{univ:gs:sec:relative,univ:gs:sec:trace} show that even \emph{all}
characters together do not classify the fields (source~09).

\emph{Why quotients need a new argument (sources 10, 11).} Suppose some coordinates
have been added, producing an intermediate $W$. A remaining factor is the old set
$\Add^V(\kappa,1)$, not automatically $\Add^W(\kappa,1)$: if $W$ has new short binary
sequences, it has new partial functions allowed by the latter poset, and the closure
of the old set can fail for newly available descending sequences. So
\cref{univ:gs:thm:fkw} cannot simply be applied in every intermediate model. The next
section proves exactly the distributivity needed, placing the closed high forcing
before the small forcing, instead of assuming that closure survives every change of
universe.

\section{Relative spectra of intermediate grounds}\label{univ:gs:sec:relative}

This section keeps the final universe fixed and varies the intermediate ground. The
object computed is not the spectrum of a forcing over its original ground but the
spectrum of each inclusion $M_A\subseteq N$ inside one product extension. A ground
Cohen poset, viewed after another coordinate has been added, need not be the Cohen
poset defined afresh there, and its ground closure proof cannot simply be reused
(sources 10, 11, 13, 14).

\subsection{Forcing tools}

\begin{lemma}[Short-name localization; source 14 (statement), sources 10, 11, 13]\label{univ:gs:lem:localization}
In a ground $V$ let $\mathbb P,\mathbb Q$ be set forcings and $\mu$ an infinite
cardinal with $\abs{\mathbb P}\le\mu$, and suppose that $\mathbb Q$ adds no
ordinal-valued sequences of length $\mu$ (for instance, $\mathbb Q$ is
$\le\mu$-closed in $V$). For mutually generic $G_{\mathbb P},G_{\mathbb Q}$,
\[
 {}^\mu\Ord\cap V[G_{\mathbb P},G_{\mathbb Q}]={}^\mu\Ord\cap V[G_{\mathbb P}].
\]
More generally, if $\mathbb Q$ adds no sequences of length below a cardinal
$\kappa$, the same holds for each length $\eta$ with
$\max(\abs{\mathbb P},\abs\eta,\aleph_0)<\kappa$.
\end{lemma}
\begin{proof}
Reorder the product as $\mathbb Q*\check{\mathbb P}$. In $V[G_{\mathbb Q}]$, a
$\mathbb P$-name for a function on $\mu$ is represented below a condition in
$G_{\mathbb P}$ by a maximal deciding antichain for each coordinate together with the
ordinal assigned to each member. A ground enumeration of $\mathbb P$ has length at most
$\mu$, so the whole array of deciding conditions and values is coded by an ordinal
sequence of length at most $\mu$. By the hypothesis on $\mathbb Q$ the array belongs
to $V$; predensity and incompatibility of its antichains are absolute, since the
elements and order of $\mathbb P$ are unchanged. So the array is a name in $V$ whose
interpretation by $G_{\mathbb P}$ is the given function. Replacing $\mu$ by a cardinal
below $\kappa$ that is large enough to code the array proves the last assertion.
\end{proof}

It is enough that $\mathbb Q$ be $<\kappa$-closed \emph{in the original ground}: its
original distributivity shows that a \emph{name array} is old, and no claim is made
that $\check{\mathbb Q}$ remains closed after forcing with $\mathbb P$. This is the key
safeguard when intermediate grounds contain only some coordinates (source~14).

\emph{Second route, by closure (sources 13, 10, 11).} Assume $\mathbb Q$ is
$\le\mu$-closed in $V$. Fix a product name $\dot f$ and $(p_0,q_0)$ forcing
$\dot f:\mu\to\Ord$. For fixed $\alpha<\mu$, enumerate $\mathbb P\res p_0$ and build an
antichain $A_\alpha$ below $p_0$ while strengthening the $\mathbb Q$-coordinate: at an
enumerated $p$ compatible with a chosen member do nothing; otherwise choose $a\le p$
and a stronger $\mathbb Q$-condition so that the pair decides $\dot f(\alpha)$, add
$a$ and record the value. Closure gives lower bounds at limits and after the
enumeration, and $A_\alpha$ is maximal below $p_0$. Repeating for all $\alpha<\mu$ with
decreasing $\mathbb Q$-conditions gives at most $\mu$ choices at each of at most $\mu$
stages; a decreasing chain indexed by an ordinal of cardinality at most $\mu$ has a
cofinal subchain of length at most $\mu$, so the stated closure suffices even if the
recursion has length $\mu\cdot\mu$. A common lower bound $q_*$ and the recorded
antichains define a $\mathbb P$-name $\dot h$ with $(p_0,q_*)\forces\dot f=\dot h$,
and such representations are dense (source~13, its Lemma~6.1). Source~10 (its
Lemma~7.1) states the case $\abs{\mathbb P}\le\kappa$, $\mathbb Q$ $\le\kappa$-closed,
lengths at most $\kappa$; source~11 (its Lemma~8.2) the case $\abs{\mathbb
P}<\mu$, $\mathbb Q$ $<\mu$-closed, concluding that $\check{\mathbb Q}$ is
$<\mu$-distributive over the $\mathbb P$-extension, building one maximal antichain per
coordinate in at most $\abs{\mathbb P}$ steps. Both are small-poset instances of
Easton's lemma: if $\mathbb P$ is $\kappa^+$-c.c.\ and $\mathbb Q$ is $\le\kappa$-closed,
then after $\mathbb P$ the ground poset $\mathbb Q$ is $\le\kappa$-distributive
\cite[Lemma~3.8]{CodyGitman}. The assertion is distributivity, not closure in the
intermediate model, and that distinction is why it is useful (source~11). A poset
closed in $V$ need not retain the same closure in $V[G_{\mathbb P}]$ (source~13).

\begin{lemma}[Product restrictions and high tails; sources 10, 11]\label{univ:gs:lem:high}
An Easton product over a set of coordinates factors as a product over any finite
partition of the coordinates. If every coordinate of a factor exceeds an infinite
cardinal $\mu$, that factor is closed under descending sequences of length at most
$\mu$ in $V$, and adds no ordinal-valued sequences of length at most $\mu$.
\end{lemma}
\begin{proof}
Restriction to a subset of coordinates preserves the support condition, and a finite
union of supports bounded below an inaccessible is bounded there; restriction and
union are inverse product maps. For closure, take the coordinatewise union of at most
$\mu$ descending conditions. At a coordinate $\alpha>\mu$, regularity keeps the domain
below $\alpha$. Below an inaccessible $\tau'\le\mu$ the high factor has empty support;
below an inaccessible $\tau'>\mu$ the union of at most $\mu$ bounded subsets is bounded.
The recursive decision proof, with a lower bound after $\mu$ steps, gives the last
assertion.
\end{proof}

\begin{lemma}[Intersections of mutually generic extensions; source 13 (base), source 14]\label{univ:gs:lem:intersection}
If $G_{\mathbb P}\times G_{\mathbb Q}$ is generic for a product of two set forcings over
$V$, every set of ordinals, and every ordinal-valued function, belonging to both
$V[G_{\mathbb P}]$ and $V[G_{\mathbb Q}]$ belongs to $V$. Consequently, for coordinate
subextensions of a product and ground sets $A,A'$ of coordinates,
$V[G\res A]\cap V[G\res A']=V[G\res(A\cap A')]$ when restricted to sets of ordinals or
ordinal-valued functions.
\end{lemma}
\begin{proof}
Let $x\subseteq\gamma$ lie in both, with a $\mathbb P$-name $\dot x$, a $\mathbb Q$-name
$\dot y$ and $(p,q)$ in the generic forcing $\dot x=\dot y\subseteq\gamma$. For each
$\xi<\gamma$, $p$ decides $\xi\in\dot x$: otherwise $p_1,p_2\le p$ force opposite
answers, a strengthening of $q$ decides $\xi\in\dot y$, and one of $(p_1,q'),(p_2,q')$
forces a contradiction. Separation in $V$ collects the positive decisions, which form
$x$. Functions are coded by sets of ordinals using a ground range bound and pairing.
For the coordinate form, factor out the common coordinates first; the two remaining
groups are mutually generic over that subextension, and unused coordinates do not
affect the intersection.
\end{proof}

Source~14 gives the same proof strengthening $q$ first. Source~09 proves a different
intersection statement for a common small forcing over two nested universes
(\cref{univ:gs:lem:nestedintersection}).

\begin{corollary}[Freshness survives an independent base extension; sources 13, 14]\label{univ:gs:cor:independent}
If $f\in V[G_{\mathbb P}]$ is fresh over $V$ and $G_{\mathbb P},G_{\mathbb Q}$ are mutually
generic, then $f$ is fresh over $V[G_{\mathbb Q}]$ in $V[G_{\mathbb P},G_{\mathbb Q}]$.
\end{corollary}
\begin{proof}
Its proper prefixes lie in $V\subseteq V[G_{\mathbb Q}]$, and
\cref{univ:gs:lem:intersection} prevents $f\in V[G_{\mathbb Q}]$.
\end{proof}

Ordinary persistence of freshness to a larger ambient universe keeps the \emph{ground}
fixed; this corollary changes the ground, which is the extra ingredient needed here
(source~14).

\subsection{Isolated coordinates (source 10)}

Fix $A\subseteq\Reg^V$, assume $V\models\GCH$, let $G$ be $\mathbb P_A$-generic and
$N=V[G]$. For $K\subseteq A$ in $V$ put $M_K=V[G\res K]$: the coordinates in $K$ are
retained, those in $A\setminus K$ are missing.

\begin{theorem}[Isolated-coordinate localization; source 10]\label{univ:gs:thm:isolated}
Suppose $\kappa\in A$ and $\abs{\mathbb P_{A\cap\kappa}}^V<\kappa$. Then
\[
 \kappa\in\Char_N(\No^{M_K})\iff\kappa\notin K.
\]
\end{theorem}
\begin{proof}
The full product preserves cardinals and cofinalities (\cref{univ:gs:thm:fkw}), so
every intermediate model has the same cardinals and cofinalities as $V$ and $N$: a
collapsing bijection or shorter cofinal map in an intermediate model would persist to
$N$.

\emph{A missing coordinate gives a gap.} If $\kappa\notin K$, the coordinate generic
$g_\kappa$ has all proper initial segments in $V\subseteq M_K$. By product factorization
it is generic over $M_K$ for the \emph{old} Cohen poset, and for any $h:\kappa\to2$ in
$M_K$ the conditions disagreeing with $h$ are dense there: extend a ground condition at
one unused ordinal by the opposite bit, which is still a ground condition even when $h$
is new over $V$. So $g_\kappa$ is fresh over $M_K$; apply \cref{univ:gs:thm:transfer}.

\emph{A kept coordinate gives no gap.} If $\kappa\in K$, put
\[
 W=V[G\res(K\setminus(\kappa+1))],\quad
 \mathbb R=\mathbb P_{A\cap(\kappa+1)},\quad
 \mathbb Q=\mathbb P_{(A\setminus K)\setminus(\kappa+1)},\quad
 \mathbb L=\mathbb P_{(A\setminus K)\cap\kappa},
\]
all ground posets, with corresponding restrictions $G_{\mathbb R},G_{\mathbb Q},G_{\mathbb L}$ of
$G$. Factorization gives $N=W[G_{\mathbb R}][G_{\mathbb Q}]$ and $W[G_{\mathbb R}]=M_K[G_{\mathbb L}]$,
the latter because $\kappa\in K$ puts all missing low coordinates strictly below
$\kappa$. The forcing from $V$ to $W$ is $\le\kappa$-closed (\cref{univ:gs:lem:high}),
so it adds no sequences of length at most $\kappa$, and the old set $\mathbb Q$ remains
$\le\kappa$-closed in $W$: a short descending sequence of its old conditions is coded,
through a ground well-order of $\mathbb Q$, by an old ordinal sequence, and has a ground
lower bound. Under GCH, $\abs{\Add^V(\kappa,1)}=\kappa$, so $\abs{\mathbb R}^V\le\kappa$, also
in $W$. By \cref{univ:gs:lem:localization} the final high forcing adds no functions on
$\kappa$ over $W[G_{\mathbb R}]$. A fresh $f:\kappa\to\Ord$ over $M_K$ in $N$ would therefore
lie in $W[G_{\mathbb R}]=M_K[G_{\mathbb L}]$; but $\abs{\mathbb L}^{M_K}<\kappa$ and $\kappa$ is
regular there, so \cref{univ:gs:thm:filterbase} excludes it. Cofinalities are preserved
between $M_K$ and $N$, so \cref{univ:gs:thm:transfer} excludes the gap.
\end{proof}

The theorem computes one isolated character of a quotient extension, not its whole
spectrum. It uses the smallness of the \emph{full} low product, not merely of one
coordinate; at accumulation points such as the singular limit of
\cref{univ:gs:ex:singular} the hypothesis fails and extra characters can occur. No step
identifies $\Add^V(\lambda,1)$ with a Cohen poset recomputed over $M_K$ (source~10).
The low product used with the localization lemma \emph{includes} the kept coordinate at
$\kappa$, so its size may be $\kappa$; only the subsequent missing-low factor has size
strictly below $\kappa$, and confusing the two bounds would invalidate the argument
(source~10).

\subsection{Sparse products and finite complements (source 11)}

\begin{definition}[Source 11]\label{univ:gs:def:sparse}
Under $V\models\GCH$, a set $A$ of infinite regular cardinals is \emph{sparse} if
$\abs{\prod_{\alpha\in A\cap\mu}\mathbb C_\alpha}^V<\mu$ for every $\mu\in A$, the
product being the full-support product (so the bound holds for every Easton
subproduct below $\mu$).
\end{definition}

``Sparse'' is local terminology for this hypothesis, not a claim that every Easton set
or increasing sequence has it; the strong form makes preservation at intermediate
grounds transparent.

\begin{lemma}[Missing upper coordinates remain distributive; source 11]\label{univ:gs:lem:reldist}
Let $A$ be sparse, $B\subseteq A$ finite, $N=V[G]$ the full Easton extension and
$M_{A\setminus B}=V[G\res(A\setminus B)]$. If $\mu\in B$ and
$T=\prod_{\alpha\in B,\ \alpha\ge\mu}\mathbb C_\alpha$ is the finite ground product of the
missing coordinates at or above $\mu$, then $T$ is $<\mu$-distributive over
$M_{A\setminus B}$.
\end{lemma}
\begin{proof}
Factor the retained coordinates $A\setminus B$ into the high part $H'$ above $\mu$ and
the low part $L$ below $\mu$ (there is no retained coordinate at $\mu$), so
$M_{A\setminus B}=V[G_{H'}][G_L]$. In $V$ the finite product $T$ is $<\mu$-closed; by
\cref{univ:gs:lem:high}, $H'$ adds no sequences of ground objects of length below $\mu$,
so every such descending sequence in the old set $T$ in $V[G_{H'}]$ was already in $V$
and has its ground lower bound, and $T$ is still $<\mu$-closed in $V[G_{H'}]$.
Sparseness gives $\abs L^V<\mu$, which persists; all subproducts preserve cardinals and
cofinalities (\cref{univ:gs:thm:fkw}), so $\mu$ is regular at this stage. Apply
\cref{univ:gs:lem:localization} in $V[G_{H'}]$.
\end{proof}

\begin{theorem}[Exact finite-complement spectrum; source 11]\label{univ:gs:thm:finitecomplement}
Assume $V\models\ZFC+\GCH$ and $A$ sparse, and let $N=V[G]$ be generic for $\mathbb P_A$.
For each finite nonempty $B\subseteq A$,
\[
 \FF(M_{A\setminus B},N)=\Char_N(\No^{M_{A\setminus B}})=B,
\]
identifying cardinals preserved between the models.
\end{theorem}
\begin{proof}
Removing finitely many coordinates does not affect the support restriction, so
$\mathbb P_A\cong\mathbb P_{A\setminus B}\times\prod_{\alpha\in B}\mathbb C_\alpha$, and the
missing-coordinate generic is generic over $M_{A\setminus B}$ for the finite product of
the old posets. All intermediate models have the cardinals and cofinalities of $V$ and
$N$.

\emph{Lower inclusion.} For $\alpha\in B$, the coordinate function $g_\alpha:\alpha\to2$
is new over $M_{A\setminus B}$ (disagreement with any prescribed $h:\alpha\to2$ is dense
over it, a condition leaving an unused coordinate), and each bounded prefix lies in $V$,
because in the original coordinate it is dense to decide the whole prefix by one ground
condition. So $g_\alpha$ is fresh and $\alpha\in\FF(M_{A\setminus B},N)$.

\emph{Upper inclusion.} Let $\lambda\notin B$ be infinite regular, and split the finite
missing product into $S=\prod_{\alpha\in B,\alpha<\lambda}\mathbb C_\alpha$ and
$T=\prod_{\alpha\in B,\alpha>\lambda}\mathbb C_\alpha$. If $S$ is nontrivial, its ground
size is below $\lambda$, since under GCH $\abs{\mathbb C_\alpha}=\alpha$ and a finite
product has the maximum of its factors' sizes; the bound persists. Suppose a fresh
$f:\lambda\to\Ord$ over $M_{A\setminus B}$ lies in $N$, and force first with $T$, writing
$U=M_{A\setminus B}[G_T]$. If $f\notin U$, it is fresh over $U$ and $N=U[G_S]$ with
$\abs S^N<\lambda=\cf^N\lambda$, contradicting \cref{univ:gs:thm:filterbase}; so
$f\in U$. If $T$ is empty, $f\in M_{A\setminus B}$, a contradiction. Otherwise let
$\mu=\min(B\setminus(\lambda+1))$; then $\lambda<\mu$ and \cref{univ:gs:lem:reldist}
says $T$ adds no sequences of length $\lambda$ over $M_{A\setminus B}$, again a
contradiction. So $\FF(M_{A\setminus B},N)=B$, and cofinality preservation and
\cref{univ:gs:thm:transfer} give the gap spectrum.
\end{proof}

The ambient spectrum $\Char_N(\No^V)$ can have additional limit-generated characters;
they do not automatically remain in the relative spectrum over $M_{A\setminus B}$,
because the generic information responsible for them may already be in
$M_{A\setminus B}$. The upper-inclusion proof, not a subtraction of ground spectra,
establishes their absence (source~11). The central point is that the missing
coordinates are \emph{ground posets}, not Cohen posets recomputed in the intermediate
model, so distributivity must be proved there.

\begin{proposition}[Sparse sets of any set length; source 11]\label{univ:gs:prop:sparse}
Assume $V\models\GCH$. Given an uncountable regular $\mu$ and an infinite cardinal
$\chi$, there are increasing regular cardinals $(\kappa_\xi)_{\xi<\chi}$ above $\mu$ such
that $\{\mu\}\cup\{\kappa_\xi:\xi<\chi\}$ is sparse.
\end{proposition}
\begin{proof}
At stage $\xi$ choose an infinite cardinal $\varsigma_\xi$ at least $\mu$, $\abs\xi$ and
every earlier $\kappa_\eta$, and set $\kappa_\xi=(2^{\varsigma_\xi})^+$, a regular
successor. Each earlier coordinate is at most $\varsigma_\xi$ and its Cohen poset has
that size under GCH; there are at most $\varsigma_\xi$ of them, so the product below
$\kappa_\xi$ has size at most $\varsigma_\xi^{\varsigma_\xi}=2^{\varsigma_\xi}<\kappa_\xi$.
At $\mu$ the lower product is trivial. The recursion has set length $\chi$.
\end{proof}

\subsection{Finite Boolean cubes (source 14)}

\begin{theorem}[Finite relative spectrum; source 14]\label{univ:gs:thm:cube}
Let $M\models\ZFC+\GCH$, let $\kappa_0<\cdots<\kappa_d$ be uncountable regular cardinals
of $M$, force with $\mathbb P=\prod_{i\le d}\Add(\kappa_i,1)^M$, and put $N=M[G]$ and
$M_A=M[G_i:i\in A]$ for $A\subseteq\{0,\ldots,d\}$. Then
\[
 \Char_N(\No^{M_A})=\{\kappa_i:i\notin A\}.
\]
If $n_*=\min(\{0,\ldots,d\}\setminus A)$ exists, the first new ordinal sequence and the
first new sign over $M_A$ both have length $\kappa_{n_*}$; that is,
$\delta(M_A,N)=\beta(M_A,N)=\kappa_{n_*}$.
\end{theorem}
\begin{proof}
Under GCH each coordinate has size $\kappa_i$. For preservation, factor at any regular
$\mu$ into the coordinates at most $\mu$ and those above: the lower factor has size at
most $\mu$, the upper adds no sequences of length $\mu$, and
\cref{univ:gs:lem:localization} shows the upper factor cannot add short cofinal maps
after the lower. If a coordinate sits at $\mu$ it is $<\mu$-closed and the still lower
finite product has size below $\mu$; if not, the lower factor is small. Either way the
regularity of $\mu$ is preserved, hence all cofinalities and cardinals, also for every
subproduct.

For $i\notin A$, $g_i$ is fresh over $M$ and stays fresh over $M_A$
(\cref{univ:gs:cor:independent}). Conversely, let $\mu$ be regular and not a missing
coordinate, and suppose $f:\mu\to\Ord$ is fresh over $M_A$. Put
$U=M[G_i:\kappa_i>\mu]$ and $V_*=U[G_i:i\in A,\ \kappa_i\le\mu]$. The remaining factor
$\mathbb R$ consists of the missing lower coordinates; it has size strictly below $\mu$,
since no missing coordinate equals $\mu$ and finitely many lie below it; and
$M_A\subseteq V_*$, $N=V_*[G_{\mathbb R}]$. If $f\notin V_*$, it is fresh over $V_*$ and
added by a forcing of size below $\mu$, contradicting \cref{univ:gs:thm:filterbase}; so
$f\in V_*$. The high product adds no $\mu$-sequences over $M$ and the retained lower
product has size at most $\mu$, so \cref{univ:gs:lem:localization} gives
$f\in M[G_i:i\in A,\ \kappa_i\le\mu]\subseteq M_A$, a contradiction. Cofinality
preservation and \cref{univ:gs:thm:transfer} finish. For the last assertion, all factors
below $\kappa_{n_*}$ are retained; factor into these and the $<\kappa_{n_*}$-closed tail,
and the general form of \cref{univ:gs:lem:localization} puts every sequence of length
below $\kappa_{n_*}$ into $M_A$, while $g_{n_*}$ is a fresh sign of length
$\kappa_{n_*}$.
\end{proof}

\begin{example}[Four fields, one threshold; source 14]\label{univ:gs:ex:four}
With three coordinates $\kappa_0<\kappa_1<\kappa_2$ and $0\notin A$, the four spectra
are $\{\kappa_0\}$, $\{\kappa_0,\kappa_1\}$, $\{\kappa_0,\kappa_2\}$ and
$\{\kappa_0,\kappa_1,\kappa_2\}$: four old fields in one universe with the same first
gap and the same uncountable saturation spectrum, pairwise nonisomorphic. With $d+1$
coordinates one obtains $2^d$ such fields. The countable construction does not simply
replace $d$ by $\omega$: an additional singular-successor character appears.
\end{example}

\subsection{All coordinate grounds of a countable product}

Fix in $M\models\GCH$ a strictly increasing sequence $\seq{\kappa_n:n<\omega}$ of
uncountable regular cardinals, put $\lambda=\sup_n\kappa_n$, and let
\begin{equation}\label{univ:gs:eq:productsetup}
 \mathbb P=\prod^{\mathrm{full}}_{n<\omega}\Add(\kappa_n,1)^M,\qquad N=M[G],\qquad
 M_A=M[G\res A]\quad(A\subseteq\omega),
\end{equation}
with full support and ground conditions throughout. For $A\subseteq\omega$ write
$B=\omega\setminus A$ for the omitted coordinates. Source~13 treats this generality;
source~14 treats $\kappa_n=\aleph_{n+1}^M$, $\lambda=\aleph_\omega^M$ (its final universe
is called $W$ there).

\begin{proposition}[Published product input; imported; sources 13, 14]\label{univ:gs:prop:published}
Assume GCH and let $\seq{\kappa_n}$ be strictly increasing infinite regular cardinals
with supremum $\lambda$. The full product of the Cohen forcings at the $\kappa_n$
preserves all cardinals and cofinalities, and every infinite coordinate subproduct adds a
fresh function on $\lambda^+$; more precisely, the full product over an unbounded set of
regulars below the singular $\lambda$ adds a $\lambda^+$-Cohen subset over the ground.
\end{proposition}

\emph{Source, not a new proof.} Preservation is \cite[Proposition~4.10]{FKW}; these
full products are Easton products, since only finitely many coordinates lie below any
ordinal below $\lambda$. The witness is \cite[Corollary~4.21]{FKW}, obtained there from
Shelah's Cohen-algebra embedding theorem \cite{Shelah2000}. Such a $\lambda^+$-Cohen
subset is fresh over the original ground: its own Cohen extension is
$<\lambda^+$-closed and adds no shorter sequences, even though the full product does.
The full support and the singular limit are indispensable. Neither assertion is
reproved, and every novelty claim excludes them (sources 13, 14). The finite theorem,
the single-Cohen theorem and every exclusion argument below are independent of this
input; it is used only for the presence of $\lambda^+$ when $B$ is infinite
(source~14).

\begin{lemma}[Sizes and unchanged reals; source 13]\label{univ:gs:lem:sizes}
Let $\mathbb P_F$ be the subproduct on $F\subseteq\omega$, formed in $M$. Then
$\abs{\mathbb P_F}^M$ is at most $1$ if $F=\varnothing$, at most $\max_{n\in F}\kappa_n$ if
$F$ is finite and nonempty, and at most $\lambda^+$ if $F$ is infinite; the same bounds
hold in every intermediate extension and in $N$. Moreover $\Pow(\omega)^N=\Pow(\omega)^M$.
\end{lemma}
\begin{proof}
Each coordinate has size $\kappa_n$ under GCH, and for infinite $F$,
$\abs{\mathbb P_F}\le\lambda^{\aleph_0}\le(2^\lambda)^{\aleph_0}=2^\lambda=\lambda^+$. Ground
enumerations persist, and by \cref{univ:gs:prop:published} the cardinal bounds are
preserved; preservation in the final model gives it in intermediate models, since a
collapse or a shortened cofinality cannot be undone. Since $\kappa_0$ is uncountable, a
descending countable sequence of product conditions has a coordinatewise union, so the
product is countably closed and adds no reals.
\end{proof}

\begin{lemma}[Preservation and localization for the full product; source 14, second route]\label{univ:gs:lem:preservation}
For $\kappa_n=\aleph_{n+1}^M$: (i) $\mathbb P$ adds no countable ordinal sequences; (ii) every
$\kappa_n$, and $\lambda$ and $\lambda^+$, remain cardinals with their original
cofinalities; (iii) all cardinals and cofinalities are preserved; (iv) for every
infinite cardinal $\mu<\lambda$, some finite head extension $M[G_i:i\le n]$ contains every
$\mu$-sequence of ordinals of $N$. All intermediate subproduct models have the cardinals
and cofinalities of $M$ and $N$.
\end{lemma}
\begin{proof}
(i) Full support permits coordinatewise unions of countable descending chains, every
$\kappa_n$ being uncountable regular. (iv) The head through $n$ has size $\kappa_n$ and
the tail after $n$ is $<\kappa_{n+1}$-closed in $M$, so \cref{univ:gs:lem:localization}
gives ${}^\mu\Ord\cap N={}^\mu\Ord\cap M[G_i:i\le n]$ for $\mu\le\kappa_n$. To preserve
$\kappa_n$, split immediately \emph{before} coordinate $n$: the head has size below
$\kappa_n$ and the tail from $n$ on is $<\kappa_n$-closed, so every $\eta$-sequence,
$\eta<\kappa_n$, lies in the head extension, and a forcing of size below a ground regular
cannot change its cofinality. The limit $\lambda$ remains a cardinal (collapsing it
would collapse a preserved cardinal below it) and its ground countable cofinal sequence
stays cofinal. GCH gives $\abs{\mathbb P}^M\le\lambda^{\aleph_0}\le(2^\lambda)^{\aleph_0}
=\lambda^+$. If $\lambda^+$ acquired cofinality $\mu<\lambda^+$, then $\mu<\lambda$ (an
infinite cofinality is regular, $\lambda$ is singular, and forcing creates no cardinals),
so choose $n$ with $\mu\le\kappa_n$; by (iv) the cofinal sequence lies in a head extension
by a forcing of size $\kappa_n<\lambda^+$, which preserves the regularity of $\lambda^+$.
This also rules out its collapse. The $\lambda^{++}$-chain condition preserves the
cofinalities of ground regulars at least $\lambda^{++}$. So all ground regulars keep their
cofinalities; a change in the cofinality of another ordinal would change that of a ground
regular, and no cardinal collapses while all ground regulars are preserved. Intermediate
models follow by the same factorization, or because a collapse there would persist.
\end{proof}

This is a self-contained second route to the preservation half of
\cref{univ:gs:prop:published} in the case $\kappa_n=\aleph_{n+1}$. Every $A\subseteq\omega$
of $N$ is in $M$, so all the $M_A$ are defined by ground subproducts;
$\R^{M_A}=\R^M=\R^N$; and $N\models2^{\aleph_0}=\aleph_1$ (source~14, its
Corollary~8.2).

\begin{theorem}[Exact spectra of all coordinate grounds; source 13 (base), source 14]\label{univ:gs:thm:coordinate}
In \eqref{univ:gs:eq:productsetup}, for every $A\subseteq\omega$ with $B=\omega\setminus A$,
\[
 \Char_N(\No^{M_A})=\{\kappa_n:n\in B\}\cup
 \begin{cases}\varnothing,&B\text{ finite},\\ \{\lambda^+\},&B\text{ infinite}.\end{cases}
\]
All cardinals in this formula, including $\lambda^+$, have the same meaning in $M$, every
$M_A$ and $N$. There are no asymmetric set-presented gaps. Every character on the right
has a fresh \emph{binary} witness whose length is that character (source~14).
\end{theorem}

\begin{proof}
All relevant cofinalities are preserved, so by \cref{univ:gs:thm:transfer} it suffices to
compute $\FF(M_A,N)$.

\emph{Lower inclusion.} For $n\in B$, the coordinate union $g_n:\kappa_n\to2$ is fresh over
$M$ in the one-coordinate extension (closure makes proper prefixes old, genericity makes it
new), and the retained coordinates are independent of it, so it is fresh over $M_A$ by
\cref{univ:gs:cor:independent}. If $B$ is infinite, it is unbounded in $\omega$, so its
coordinate cardinals still have supremum $\lambda$; by \cref{univ:gs:prop:published} the
subextension $M[G\res B]$ contains a $\lambda^+$-Cohen subset $r$, fresh over $M$, and
\cref{univ:gs:cor:independent} makes it fresh over $M_A$.

\emph{Nothing unexpected below $\lambda$ (source~13).} Let $\mu<\lambda$ be infinite regular
with $\mu\notin\{\kappa_n:n\in B\}$, and suppose $f:\mu\to\Ord$ is fresh over $M_A$. Let
$J=\{n:\kappa_n\le\mu\}$, a finite head of size at most $\mu$ in $M$; the tail is
$\le\mu$-closed in $M$, so \cref{univ:gs:lem:localization}, applied to the whole original
product before changing the base model, gives $f\in L:=M[G\res J]$. Let
$W_0=M[G\res(J\cap A)]$. For $\alpha<\mu$, $f\res\alpha$ lies in $M_A$ and in $L$, so in
$W_0$ by \cref{univ:gs:lem:intersection}; and $f\notin W_0\subseteq M_A$. So $f$ is fresh
over $W_0$ in $L$. But $L$ is obtained from $W_0$ by the finite old product on $J\cap B$,
whose coordinates are all \emph{strictly} below $\mu$ (equality would put
$\mu\in\{\kappa_n:n\in B\}$); it has size below $\mu$ in $W_0$ and $L$. If it is empty
there is no extension; otherwise \cref{univ:gs:thm:filterbase} excludes a fresh function
on the still-regular $\mu$.

\emph{The ceiling.} If $B=\varnothing$ then $M_A=N$ and the internal cut theorem leaves no
set-presented gap. Otherwise $N/M_A$ is a forcing extension by the old poset $\mathbb P_B$,
whether or not it keeps its former closure. If $B$ is finite,
\cref{univ:gs:lem:sizes,univ:gs:thm:filterbase} exclude every character above
$\max_{n\in B}\kappa_n$, and the rest is below $\lambda$. If $B$ is infinite, the ceiling
excludes every character above $\lambda^+$; $\lambda$ remains singular and no cardinal lies
strictly between $\lambda$ and $\lambda^+$, so $\lambda^+$ is the only possible character
not below $\lambda$. Below $\lambda$ use the previous paragraph.
\end{proof}

\emph{Second route to the exclusion of retained coordinates (source~14).} Fix $n\in A$,
$\mu=\kappa_n$, and a hypothetical fresh $f:\mu\to\Ord$ over $M_A$. Put
$U=M[G_i:i>n]$ and $V_*=U[G_i:i\in A,\ i\le n]$; then $M_A\subseteq V_*$, since $V_*$
contains every retained coordinate and all missing high coordinates, and $N/V_*$ is the
finite ground product $\mathbb R=\prod_{i\in B,\,i\le n}\Add(\kappa_i,1)^M$, all of whose
factors have index below $n$, so $\abs{\mathbb R}<\mu$. If $f\notin V_*$ it is fresh over
$V_*$, contradicting \cref{univ:gs:thm:filterbase}; so $f\in V_*$. Reordering, the high
factor is $<\kappa_{n+1}$-closed in $M$ and adds no $\mu$-sequences, and the retained low
factor has size at most $\mu$, so \cref{univ:gs:lem:localization} gives
$f\in M[G_i:i\in A,\ i\le n]\subseteq M_A$. This works even when $A$ and $B$ are both
infinite. The remaining characters are excluded as follows: no countable sequence is new,
so $\omega$ is out; every regular below $\lambda$ is $\omega$ or some $\kappa_n$, and
$\lambda$ is singular; above $\lambda^+$ the quotient $\mathbb P_B$ has size at most
$\lambda^+$, and if $B$ is finite at most its largest coordinate, which excludes
$\lambda^+$ too. The argument deliberately allows $M_A$ to contain new sequences of old
conditions: closure is invoked only in $M$, at the head--tail factorization where it is
known, and fresh prefixes are pulled back to the common base by mutual genericity. This is
the relative step that a direct appeal to the ground product's spectrum would not supply
(sources 13, 14).

\begin{example}[A concrete spectrum table; sources 13, 14]\label{univ:gs:ex:table}
For $\kappa_n=\aleph_{n+1}$, $\lambda=\aleph_\omega$, the following are full character
spectra in one extension, not lower bounds.
\begin{center}\small
\begin{tabular}{@{}P{0.3\textwidth}P{0.55\textwidth}@{}}
\toprule
Omitted coordinates $B$ & $\Char_N(\No^{M_{\omega\setminus B}})$\\
\midrule
$\varnothing$ & $\varnothing$\\
$\{0\}$ & $\{\aleph_1\}$\\
$\{0,2\}$ & $\{\aleph_1,\aleph_3\}$\\
$\{2m:m<\omega\}$ & $\{\aleph_{2m+1}:m<\omega\}\cup\{\aleph_{\omega+1}\}$\\
$\omega$ & $\{\aleph_{n+1}:n<\omega\}\cup\{\aleph_{\omega+1}\}$\\
\bottomrule
\end{tabular}
\end{center}
The two infinite rows show the successor-of-singular contribution, which listing the
generic coordinate lengths would miss. All nonempty rows have least omitted coordinate
$0$ and hence the same saturation spectrum; the isolated character $\aleph_{\omega+1}$
detects an infinite complement.
\end{example}

\begin{proposition}[An exact coordinate invariant; source 13, with source 14]\label{univ:gs:prop:coordorder}
For $A,A'\subseteq\omega$ with omitted sets $B,B'$:
\begin{align*}
 \No^{M_A}\subseteq\No^{M_{A'}}&\iff A\subseteq A',\\
 \No^{M_A}\cap\No^{M_{A'}}&=\No^{M_{A\cap A'}},\\
 \Char_N(\No^{M_A})\subseteq\Char_N(\No^{M_{A'}})&\iff B\subseteq B',
\end{align*}
and $A=\{n<\omega:\kappa_n\notin\Char_N(\No^{M_A})\}$. Consequently $\No^{M_A}\cong\No^{M_{A'}}$
as class real closed fields (or as orders with a set-respecting class isomorphism) if and
only if $A=A'$.
\end{proposition}
\begin{proof}
If $A\subseteq A'$ then $M_A\subseteq M_{A'}$. If $n\in A\setminus A'$, the generic sign
$g_n$ lies in $\No^{M_A}$ but not in $\No^{M_{A'}}$, by mutual genericity. A sign in both
fields has an ordinal code in both models, and \cref{univ:gs:lem:intersection} puts it in
$M_{A\cap A'}$. The coordinate cardinals recover $A$, since $\lambda^+$ differs from all
of them, and the $\lambda^+$ clause respects inclusion. An isomorphism of real closed
fields preserves the order (positivity is being a nonzero square), and as a class map
preserves set presentations and their sizes, hence the spectrum.
\end{proof}

The fields form a poset anti-isomorphic to $\Pow(\omega)$ under omitted sets. The
intersection assertion is a \emph{binary} one; field composita are not identified, and it
is not asserted that arbitrary intersections commute with unions of omitted sets
(source~13). Within this displayed family the gap spectrum is complete, which is a
completeness result for the indexed family only, not a theorem that equal spectra
classify arbitrary old fields, arbitrary real closed class fields or arbitrary
involutions (sources 13, 14). The models do not recover the first new birthday from the
pure field language in every forcing extension; cofinality compression is specifically a
loss of ordinal-length information (source~13).

\begin{proposition}[Exact saturation in the coordinate family; sources 13, 14]\label{univ:gs:prop:threshold}
Let $B\ne\varnothing$ and $n_*=\min B$. Then
$\delta(M_A,N)=\beta(M_A,N)=\kappa_{n_*}$; $\No^{M_A}$ is
$\omega$-saturated, and for uncountable $\tau$ it is $\tau$-saturated in $N$ exactly when
$\tau\le\kappa_{n_*}$. If $B=\varnothing$, $\No^{M_A}=\No^N$ realizes all $N$-set-sized types.
\end{proposition}
\begin{proof}
By \cref{univ:gs:thm:coordinate} there is no character below $\kappa_{n_*}$ and one at
$\kappa_{n_*}$, so by \cref{univ:gs:cor:lowerbarrier} no ordinal sequence shorter than
$\kappa_{n_*}$ is new, and \cref{univ:thm:main-summary,univ:gs:cor:saturation} give the
uncountable clause; no new reals and \cref{univ:thm:omega} give $\omega$-saturation.
The fresh sign $g_{n_*}$ has length $\kappa_{n_*}$, so $\beta(M_A,N)\le\kappa_{n_*}$, and
$\delta\le\beta$ (\cref{univ:prop:beta}) gives equality. \emph{Source~14's route} proves the short-sequence clause directly: all coordinates below
$n_*$ are retained, the head below $n_*$ has size below $\kappa_{n_*}$, the tail from
$n_*$ on is $<\kappa_{n_*}$-closed, so every sequence of length below $\kappa_{n_*}$ lies
in the retained head extension, while $g_{n_*}$ is a fresh binary sequence of length
$\kappa_{n_*}$ (its Lemma~10.1). \emph{Source~13's route} proves saturation directly: an
$N$-set $A_0\subseteq\No^{M_A}$ of size below $\tau$ lies in a bounded old set of signs,
so an ordinal-coded enumeration of it is old; a complete one-type over $A_0$ has at most
$\max(\aleph_0,\abs{A_0})$ formulas and is old by the same coding; finite satisfiability
is absolute by quantifier elimination; and inside $M_A$ the real closure of the field
generated by $A_0$ and the surreal cut theorem realize it, including end cuts and
infinitesimal sides (its Lemma~8.2). Either way, a gap of character $\kappa_{n_*}$ gives a
finitely satisfiable omitted type over $\kappa_{n_*}$ parameters, so every
$\tau>\kappa_{n_*}$ fails.
\end{proof}

The threshold is proved directly in the setting at hand and does not depend on formal
verification of the general external saturation theorem (source~13). The fields with one
least omitted coordinate are indistinguishable by this threshold but distinguishable by
their spectra.

\begin{corollary}[One threshold, many nonisomorphic fields; sources 14, 13]\label{univ:gs:cor:onethreshold}
The fields $\No^{M_A}$ with $0\notin A$ are $2^{\aleph_0}$ in number in $N$, pairwise
nonisomorphic as class ordered fields under set-respecting class isomorphisms, contain the
same real field, agree on the common subfield $(\No_{<\omega_1})^M=(\No_{<\omega_1})^N$, and
are all $\kappa_0$-saturated but not $\kappa_0^+$-saturated, with the same entire saturation
spectrum and $\omega$-saturated. If $A\subseteq A'$ then
$\No^{M_A}\subseteq\No^{M_{A'}}\subseteq\No^N$ are elementary ordered-field inclusions, and all
these fields have the same first-order theory. For $\kappa_n=\aleph_{n+1}$ they are
$\omega_1$-saturated but not $\omega_2$-saturated.
\end{corollary}

No reals are added, so here $2^{\aleph_0}=\aleph_1$ in $N$; ``continuum many'' does not
assert an arbitrarily large continuum (sources 13, 14). One may restrict further to infinite
$B$ without reducing the size of the family, since only countably many $B$ are finite.

\subsection{Families of any size with one threshold}

\emph{Source~10's Boolean family.} Fix an infinite cardinal $\chi$ in a GCH ground $V$ and an
uncountable regular $\kappa_*>\chi$. For $i<\chi$ let
$\varsigma_i=\sup(\{\kappa_*\}\cup\{\kappa_j:j<i\})$ and $\kappa_i=(2^{\varsigma_i})^+$
(cardinal arithmetic in $V$), so $\kappa_i>\varsigma_i$ and the sequence increases. Let
$I=\{\kappa_i:i<\chi\}$, $A=\{\kappa_*\}\cup I$ and $N=V[G]$ for $G$ generic over
$\mathbb P_A$.

\begin{lemma}[Every designated coordinate is isolated enough; source 10]\label{univ:gs:lem:sparsesize}
For every $i<\chi$, $\abs{\mathbb P_{A\cap\kappa_i}}^V<\kappa_i$. The full product is
$<\kappa_*$-closed, preserves all cardinals and cofinalities, and adds no subsets of $\chi$
and no reals.
\end{lemma}
\begin{proof}
Every factor below $\kappa_i$ has size at most $\varsigma_i$, and there are at most
$\chi+1\le\varsigma_i$ of them, so even the full-support product has size at most
$\varsigma_i^{\varsigma_i}=2^{\varsigma_i}<\kappa_i$. For a descending chain of length
below $\kappa_*$, coordinatewise union preserves condition sizes and the bounds below
inaccessibles, as in \cref{univ:gs:lem:high}. Preservation is \cref{univ:gs:thm:fkw}.
Closure adds no ordinal sequences of length below $\kappa_*>\chi,\omega$.
\end{proof}

For $S\subseteq\chi$ let $K_S=\{\kappa_i:i\in\chi\setminus S\}$: the fields
$\No^{M_{K_S}}$ retain the coordinates outside $S$, omit the test coordinates in $S$, and
omit the common coordinate $\kappa_*$ always. Since no subset of $\chi$ is added, all
$S\in\Pow(\chi)^N$ are ground sets.

\begin{theorem}[Boolean-indexed nonisomorphic old fields with one threshold; source 10]\label{univ:gs:thm:booleanfamily}
In the single universe $N$, the fields $\No^{M_{K_S}}$, $S\subseteq\chi$, satisfy:
\begin{enumerate}[label=\textnormal{(\roman*)}]
\item their isolated-character profiles are exact:
$\Char_N(\No^{M_{K_S}})\cap A=\{\kappa_*\}\cup\{\kappa_i:i\in S\}$;
\item $\No^{M_{K_S}}\subseteq\No^{M_{K_T}}\iff T\subseteq S$, and every such inclusion is
elementary in the ordered-field language;
\item distinct indices give nonisomorphic ordered fields and nonisomorphic pure fields;
there are $2^\chi$ of them, computed in $N$;
\item all have the same real field $\R^V=\R^N$, no set-sized cofinal or coinitial subset,
least gap size $\kappa_*$ and first new birthday $\kappa_*$ relative to $N$;
\item for every infinite cardinal $\tau$ of $N$, each is $\tau$-saturated exactly when
$\tau\le\kappa_*$.
\end{enumerate}
\end{theorem}
\begin{proof}
(i) Apply \cref{univ:gs:thm:isolated} with \cref{univ:gs:lem:sparsesize} at each $\kappa_i$:
$\kappa_i\notin K_S$ iff $i\in S$. The generic $g_{\kappa_*}$ is fresh over every
$M_{K_S}$ by the missing-coordinate argument. (ii) If $T\subseteq S$, then
$K_S\subseteq K_T$. If $i\in T\setminus S$, the sign obtained from $g_{\kappa_i}$ lies in
$\No^{M_{K_S}}$ but is fresh over $M_{K_T}$. Quantifier elimination makes inclusions of real
closed fields elementary. (iii) Distinct indices give different spectra by (i); an order
isomorphism preserves characters, and a pure-field isomorphism of real closed fields is an
order isomorphism. There are $\abs{\Pow(\chi)^N}=2^\chi$ indices. (iv) No ordinal sequence
of length below $\kappa_*$ is added from $V$, hence from any $M_{K_S}$, so
\cref{univ:gs:cor:lowerbarrier} excludes smaller characters; $g_{\kappa_*}$ has birthday
$\kappa_*$ while every shorter sign is in $V$; no reals are added; old ordinals bound every
$N$-set (\cref{univ:lem:oldbounds}). (v) Source~10 proves saturation without importing a
preservation theorem: a parameter set $A_0$ of size below $\kappa_*$ has bounded birthdays,
an $N$-enumeration of it coded through a well-order of $(\No_{<\gamma})^{M_{K_S}}$ is a short
ordinal sequence, hence in $V$ and in $M_{K_S}$; a complete finite-variable type over
$A_0$ is coded by a binary sequence of length at most $\max(\abs{A_0},\aleph_0)<\kappa_*$,
old for the same reason, and its finite satisfiability is absolute; internally the
surreal field realizes it (real closure of the parameters, cut filling, successive
variables). The gap of $g_{\kappa_*}$ gives a finitely satisfiable omitted type over
$\kappa_*$ parameters, so $\kappa_*^+$-saturation fails; monotonicity gives the rest,
including $\omega$.
\end{proof}

The least gap has sides of size $\kappa_*$, and saturation at $\kappa_*$ only tests
parameter sets of size \emph{strictly} less than $\kappa_*$, so it is not contradicted by
that gap: the first failure is at $\kappa_*^+$ (source~10).

\begin{corollary}[The endpoint members; source 10]\label{univ:gs:cor:endpoints}
$\Char_N(\No^{M_I})=\{\kappa_*\}$ and $\Char_N(\No^{M_\varnothing})=\Char_N(\No^V)=\Ecl^V(A)$.
\end{corollary}
\begin{proof}
With no coordinate retained, $M_\varnothing=V$ and \cref{univ:gs:thm:easton} applies. With
all test coordinates retained, the only missing factor is the Cohen forcing at $\kappa_*$;
the retained product is $\le\kappa_*$-closed in $V$, so adds no sequences of length at most
$\kappa_*$ and no new Cohen conditions at $\kappa_*$, and the remaining old poset is
exactly $\Add^{M_I}(\kappa_*,1)$. Apply \cref{univ:gs:thm:cohen}. This equality of posets
is justified here and not assumed for arbitrary intermediate members.
\end{proof}

\begin{example}[Continuum many fields with the same first gap; source 10]\label{univ:gs:ex:omegafamily}
Take $\chi=\omega$, $\kappa_*=\aleph_1$ in a GCH ground; the recursion gives
$\kappa_i=\aleph_{2i+3}$. The final universe has $2^{\aleph_0}$ pairwise nonisomorphic old
fields, all $\aleph_1$-saturated but not $\aleph_2$-saturated, with
$\aleph_{2i+3}\in\Char_N(\No^{M_{K_S}})\iff i\in S$, the same reals and least new
birthday $\aleph_1$. The largest in inclusion has spectrum $\{\aleph_1\}$; the smallest is
$\No^V$ with spectrum $\{\aleph_{2n+1}:n<\omega\}\cup\{\aleph_{\omega+1}\}$. Source~10 does
not infer the remaining full spectra from the isolated profiles and leaves room for
accumulation characters. \emph{Merge note:} $\kappa_*<\kappa_0<\kappa_1<\cdots$ is a
strictly increasing $\omega$-sequence of uncountable regulars and, having no inaccessible
limit point, its Easton product is the full product; so
\cref{univ:gs:thm:coordinate} (source~13) computes every member:
$\Char_N(\No^{M_{K_S}})=\{\aleph_1\}\cup\{\aleph_{2i+3}:i\in S\}\cup\{\aleph_{\omega+1}:S\text{ infinite}\}$.
This answers source~10's own question for $\chi=\omega$ (\cref{univ:gs:q:accumulation}).
\end{example}

\begin{corollary}[The least threshold is not a classification invariant; source 10]\label{univ:gs:cor:notclassified}
Even after fixing the real field, the first new birthday, the least omitted-cut character
and the entire saturation spectrum, there can be arbitrarily large set-indexed families of
pairwise nonisomorphic old surreal fields in one universe, and their elementary-subfield
inclusion order can contain the opposite of $\Pow(\chi)$.
\end{corollary}

The last assertion is about inclusion order only; it does not identify intersections or
field composita with Boolean operations (source~10).

\begin{remark}[A corrected overstatement of source 10]\label{univ:gs:rem:overstatement2}
Source~10 presents this corollary as ``a concrete negative answer to the suggestion that the
least gap size, or equivalently the saturation spectrum in this family, could classify the
fields''. This report made no such suggestion: \cref{univ:q:completeness} asks which
invariants \emph{beyond} the least omitted-cut size distinguish old fields, and
\cref{univ:sec:leastgap} states that different grounds can have the same threshold. The
corollary is correctly read as showing that the least threshold, and even the whole saturation
spectrum, is not a complete invariant, which that question presupposes; source~10 itself states
that it does not answer the other half of the question, whether the \emph{full} gap spectrum is
complete (answered negatively by sources 12, 09, 11: \cref{univ:gs:sec:trace,univ:gs:sec:amplification,univ:gs:sec:tower}).
\end{remark}

\emph{Source~11's family of any length.} Here every family member omits exactly two
coordinates.

\begin{theorem}[A large same-saturation family in one universe; source 11]\label{univ:gs:thm:samesat}
Assume $V\models\ZFC+\GCH$, $\mu$ uncountable regular and $\chi$ an infinite cardinal. There
are a cardinal- and cofinality-preserving set-forcing extension $N$ and intermediate grounds
$M_\xi$, $\xi<\chi$, with $\Char_N(\No^{M_\xi})=\{\mu,\kappa_\xi\}$ for pairwise distinct
regular $\kappa_\xi>\mu$. All the $\No^{M_\xi}$ are $\omega$-saturated and, for uncountable
$\lambda$, $\lambda$-saturated exactly when $\lambda\le\mu$; they are pairwise nonisomorphic as
ordered fields, have no set-sized cofinal subsets, and all contain every real of $N$.
\end{theorem}
\begin{proof}
Take $A=\{\mu\}\cup\{\kappa_\xi:\xi<\chi\}$ sparse by \cref{univ:gs:prop:sparse}, force with
$\mathbb P_A$, and let $M_\xi=M_{A\setminus\{\mu,\kappa_\xi\}}$. By
\cref{univ:gs:thm:finitecomplement} the spectrum is exactly $\{\mu,\kappa_\xi\}$. The full
product is $<\mu$-closed (no coordinate below $\mu$; bounded supports at inaccessibles above
$\mu$ stay bounded under unions of fewer than $\mu$), so no reals are added. Apply
\cref{univ:thm:omega,univ:gs:cor:saturation,univ:lem:oldbounds}. A class order isomorphism
with Replacement for its graph preserves the spectrum.
\end{proof}

\begin{corollary}[Saturation is strictly coarser than the full spectrum; source 11]\label{univ:gs:cor:coarser}
Even among old fields containing all reals of the ambient universe and having no set-sized
cofinal subset, the complete saturation spectrum determines neither the ordered-field
isomorphism type nor the complete gap spectrum.
\end{corollary}

This cannot be obtained by assigning distinct least thresholds: the threshold is held fixed,
and the second omitted character carries the distinguishing information (source~11). An
infinite family obtained by changing only the least character would not establish
\cref{univ:gs:thm:samesat}.

\section{Set-sized cutoffs and local ubiquity}\label{univ:gs:sec:cutoff}

\subsection{The exact birthday boundary (source 14)}

For $\NoM$ every proper prefix of an old sign is in the field. For a birthday cutoff
$(\No_{<\theta})^M$ a missing separator can instead be an \emph{old} sign of length
$\theta$, and such signs can have asymmetric characters; so the correspondence of
\cref{univ:thm:gapcorrespondence} cannot be used blindly at a cutoff (source~14). The
nonendpoint clause of \cref{univ:def:setgap} is essential here: strict inequalities $x>a$
with a coinitial family above $a$ can fail to have a realization without describing a
genuine gap, since the cut may have the existing endpoint $a$; only genuine gaps are
counted, not all omitted systems of strict inequalities.

\begin{lemma}[Cutoff gap decomposition; source 14]\label{univ:gs:lem:boundary}
Suppose $M\subseteq N$ have the same cofinalities, $\theta$ is a common uncountable regular
cardinal and $F_\theta=(\No_{<\theta})^M$. Put
$H_\theta(M,N)=\{\cf^N\alpha:\alpha<\theta,\text{ some fresh sign has length }\alpha\}$. Then
\[
 \GapSpec_N(F_\theta)=\{(\mu,\mu):\mu\in H_\theta(M,N)\}\cup
 \{(\mu,\theta),(\theta,\mu):\mu\le\theta\text{ infinite regular in }N\}.
\]
\end{lemma}
\begin{proof}
Take a genuine gap of the set $F_\theta$; its two sides form a set presentation, and their
simplest separator $s$ in $\NoN$ has $\bd(s)\le\theta$ by \eqref{univ:eq:cutbirthday}. If
$\bd(s)<\theta$, $s$ is not old (it would lie in $F_\theta$); its shortest new prefix has
the same comparisons with old elements and separates the cut, so minimality makes $s$ fresh;
its proper prefixes lie in $F_\theta$, and the canonical-prefix argument gives the character
$(\cf^N\bd(s),\cf^N\bd(s))$. If $\bd(s)=\theta$ and $s$ is new, it is fresh by the same
argument, both sign sets are cofinal in $\theta$, and the character is $(\theta,\theta)$. If
$s$ is old of length $\theta$, its plus and minus positions $P,Q$ index a strictly increasing
cofinal chain on the left and a strictly decreasing coinitial chain on the right; neither is
empty or has a last index, as the gap is genuine; $P\cup Q=\theta$ and $\theta$ is regular, so
one of them has cofinality $\theta$ and the other some infinite regular cofinality at most
$\theta$. Conversely, a fresh sign of length below $\theta$ gives its usual gap in $F_\theta$.
For regular $\mu<\theta$ the old sign $+^\mu\cat-^\theta$ has length $\theta$ and a genuine
gap of character $(\mu,\theta)$: its left canonical chain has cofinality $\mu$, its right one
$\theta$, and no sign shorter than $\theta$ extends it. Negation gives $(\theta,\mu)$, and an
old sign of length $\theta$ with plus at limit and minus at successor positions gives
$(\theta,\theta)$.
\end{proof}

With $\mu=\omega$ the lemma gives an $(\omega,\theta)$ gap in the full internal cutoff
$\No_{<\theta}$, which is not a fresh-sign gap across universes: ``all genuine gaps of a
birthday cutoff are symmetric'' is false even without forcing. The full-class theorem avoids
this because the old sign of length $\theta$ is then an available separator (source~14). The
symmetry of \cref{univ:cor:gapspectrum} is a statement about the class $\NoM$ only. The lemma
is the gap-pair counterpart, across universes, of \cref{univ:thm:fragments}, which gives the
saturation of such fragments.

\noindent[Added 3 October 2026, batch~81: for $M=N$ the lemma is extended to every infinite
cardinal $\theta$, singular $\theta$ and $\theta=\omega$ included, by source~17 of the report
\nolinkurl{Algebra/SurrealNumbers/docs/foundations-and-computation/surreal-well-orders/}
(its Theorem~29.4, \nolinkurl{swo:gs:thm:alphabet-gaps}): with $r=\cf\theta$ the gap pairs of
$\No_{<\theta}$ are $(r,r)$ and $(\rho,r)$, $(r,\rho)$ for the infinite regular $\rho<\theta$.]

\begin{theorem}[Exact gap pairs at a common cutoff; source 14]\label{univ:gs:thm:cutoff}
In the setting of \eqref{univ:gs:eq:productsetup} with $\kappa_n=\aleph_{n+1}$, put
$\theta=\lambda^{++}=\aleph_{\omega+2}$ and $S_A=\Char_N(\No^{M_A})$. Then
\[
 \GapSpec_N\bigl((\No_{<\theta})^{M_A}\bigr)=\{(\mu,\mu):\mu\in S_A\}\cup
 \{(\mu,\theta),(\theta,\mu):\mu\le\theta\text{ infinite regular}\},
\]
so $\Char_N((\No_{<\theta})^{M_A})=S_A\cup\{\theta\}$. Every $(\No_{<\theta})^{M_A}$ is a real
closed field of cardinality exactly $\theta$, and distinct indices give nonisomorphic fields
of this same cardinality.
\end{theorem}
\begin{proof}
The fields are real closed since $\theta$ is uncountable regular in every model. Each
character in $S_A$ has a fresh binary witness of length equal to it
(\cref{univ:gs:thm:coordinate}), at most $\lambda^+<\theta$; conversely a fresh sign of
length below $\theta$ is a fresh sign for the full old field. So $H_\theta(M_A,N)=S_A$ and
\cref{univ:gs:lem:boundary} applies. Each field contains every ordinal below $\theta$. For the
upper bound count names in $M$: $\abs{\mathbb P}\le\lambda^+$; for $\eta<\theta$ a binary
$\eta$-sequence has a name coded by a subset of $\mathbb P\times\eta$, and
$\abs\eta^M\le\lambda^+$ with $M\models2^{\lambda^+}=\theta$, so there are at most $\theta$
such names, and $\theta$ choices of $\eta$; so $\abs{(\No_{<\theta})^N}^N\le\theta$. Distinct
$A$ give different $S_A$ at coordinate cardinals strictly below $\theta$, which the common
boundary character does not erase, and field isomorphisms preserve order and gap pairs.
\end{proof}

\begin{corollary}[Equal-sized real closed counterexamples to classification by saturation; source 14]\label{univ:gs:cor:setfields}
The fields $(\No_{<\theta})^{M_A}$, $0\notin A\subseteq\omega$, are $2^{\aleph_0}$ pairwise
nonisomorphic real closed fields of cardinality $\theta$. They are elementary subfields of
$(\No_{<\theta})^N$, contain the same real field and the same birthday fragment below
$\omega_1$, and have exactly the same saturation spectrum: $\omega$- and
$\omega_1$-saturated, not $\tau$-saturated for any uncountable $\tau>\omega_1$.
\end{corollary}
\begin{proof}
For fewer than $\kappa_{n_*}$ endpoints, the coding argument of
\cref{univ:gs:prop:threshold} makes the endpoint sets old; their birthday supremum is below
$\theta$ by regularity, so the internal simplest separator stays in the cutoff. The fresh
gap at $\kappa_{n_*}$ lies below the cutoff. A countable type uses countably many elements
of bounded birthday, so $\omega$-saturation stays within the cutoff. Take $n_*=0$ and apply
\cref{univ:gs:thm:cutoff}.
\end{proof}

\subsection{Local ubiquity}

\begin{proposition}[Every character in every interval; source 13 (base), source 14]\label{univ:gs:prop:ubiquity}
Let $F$ be a real closed class field and $a<b$ in $F$. There is a parameter-definable order
isomorphism from $F$ onto $(a,b)_F$. Consequently every set-presented gap character of $F$
is represented inside every such interval, and if $F$ contains all ordinals as an ordered
subclass, every occurring character has a proper class of distinct gap realizations,
in pairwise disjoint subintervals.
\end{proposition}
\begin{proof}
With positive square roots put $u(x)=x/\sqrt{1+x^2}$ and $v(y)=y/\sqrt{1-y^2}$ for
$-1<y<1$; direct substitution gives $u(v(y))=y$ and $v(u(x))=x$. The map $u$ preserves signs,
and for $0\le x<z$, $x^2/(1+x^2)<z^2/(1+z^2)$ by cross-multiplication; oddness gives the
negative half, so $u$ is strictly increasing and $x\mapsto a+\frac{b-a}2(1+u(x))$ is the
required isomorphism. Transport a set presentation of a gap through it and add the parts of
$F$ below and above the interval to the lower and upper classes; the transported cofinal and
coinitial sets keep their exact sizes. Doing this for one fixed gap into the disjoint
intervals $(\alpha,\alpha+1)$, $\alpha$ an ordinal, gives proper-class many distinct gaps.
\end{proof}

\emph{Second route (source~14, its Proposition~10.4), for old fields.} Choose old endpoints
$a,b$ bounding a fixed gap. A positive affine map with old coefficients sends $a,b$ to any two
old interior points of a target interval; it is an order automorphism of the old field and
moves the gap into the interval with unchanged character. Inside $(0,1)$ the old intervals
$I_\alpha=(\omega^{-(\alpha+1)},2\omega^{-(\alpha+1)})$, $\alpha\in\Ord$, are pairwise disjoint,
a later exponent giving a scale at least a factor $\omega$ smaller; an affine change of
variables places the family inside any interval, and transporting the same presented gap into
each $I_\alpha$ by explicitly given affine maps needs no selection of a proper class of
witnesses.

Recording merely whether a character has finitely many, set many or proper-class many
occurrences, or whether it occurs in a chosen interval, therefore adds no information: every
occurring character has proper-class multiplicity, even locally. This does \emph{not} rule
out finer incidence, dependence or algebraic information about families of cuts (sources 13,
14). In the coordinate family the exact \emph{list} of characters already distinguishes the
fields; when two other fields have the same list, a subtler invariant is needed.

\section{The real trace and a Boolean family with one spectrum}\label{univ:gs:sec:trace}

A gap spectrum records the sizes of cofinal witnesses but not the rational cuts already
realized. The latter is an intrinsic invariant of the field, even when no copy of the real
numbers is named in its language. This section is source~12's, with source~09's
rational-cut invariant as a second route.

\subsection{The real trace}

For an ordered subfield $F\subseteq\NoN$ put
\[
 \mathcal O_F=\{x\in F:\abs x<n\text{ for some integer }n\ge1\},\qquad
 \mathfrak m_F=\{x\in F:\abs x<1/n\text{ for every integer }n\ge1\},
\]
the finite valuation ring and its infinitesimal ideal. The quantifiers over positive
integers are external, set-indexed quantifiers; no first-order definability of these classes
in a real closed field is claimed. Every finite surreal has a unique standard part in
$\R^N$, and one useful cut formula, avoiding a possible greatest rational in a lower cut, is
\[
 C_x=\{q\in\Q:q+1/n<x\text{ for some }n\ge1\}.
\]
If $r=\st_N(x)$ then $C_x=\{q\in\Q:q<r\}$; the extra $1/n$ matters when $x$ lies
infinitesimally above a rational, since the unmodified set $\{q:q<x\}$ would contain that
rational.

\begin{definition}[Real trace; source 12, source 09]\label{univ:gs:def:trace}
The \emph{real trace} of $F$ in $N$ is
$\Tr(F)=\{\st_N(x):x\in\mathcal O_F\}\subseteq\R^N$, equivalently the set of real Dedekind
cuts represented by $C_x$ for finite $x\in F$. Source~09's \emph{rational-cut invariant}
$\mathrm{RCut}(F)$ is the set of $r\in\R^N$ such that some $x\in F$ satisfies $a<x<b$
whenever $a,b\in\Q$ and $a<r<b$; it is the same set.
\end{definition}

``Intrinsic'' refers to invariance under field embeddings; the canonical realization inside
$\R^N$ uses the rational prime field. It is stronger than remembering the cardinality of an
abstract residue field (source~12). At a rational boundary, the definition deliberately
forgets whether that rational lies in the strict lower cut and records its real supremum;
the two comparisons are made in their own fields, and no subtraction of an element of $F$
and an external real is needed. It is an isomorphism invariant, not an assertion that
``Archimedean-finite'' is first-order definable by one formula (source~09).

\begin{proposition}[Residue realization and functoriality; source 12]\label{univ:gs:prop:tracefunctor}
$\mathcal O_F$ is a ring with maximal ideal $\mathfrak m_F$, and standard part is a surjective
ring homomorphism $\st_N:\mathcal O_F\to\Tr(F)$ with kernel $\mathfrak m_F$, whose image is an
Archimedean ordered subfield of $\R^N$. If $j:F\hookrightarrow E$ is an order-preserving field
embedding, then $\st_N(j(x))=\st_N(x)$ for $x\in\mathcal O_F$, and $\Tr(F)\subseteq\Tr(E)$.
\end{proposition}
\begin{proof}
Sums and products of finite elements are finite, and a finite multiple of an infinitesimal
is infinitesimal; a finite non-infinitesimal $x$ has $\abs x\ge1/n$ for some $n$, so $x^{-1}$
is finite. So $\mathfrak m_F$ is the nonunit ideal. Writing $x=r+\epsilon$, $y=s+\epsilon'$
with $r,s\in\R^N$ and $\epsilon,\epsilon'$ infinitesimal gives $\st_N(x+y)=r+s$ and
$\st_N(xy)=rs$; inversion of a finite element with nonzero standard part gives inverses. The
embedding $j$ fixes every rational, preserves every rational inequality, and preserves and
reflects finiteness and infinitesimality; so $C_{j(x)}=C_x$, and two reals with the same lower
rational cut are equal.
\end{proof}

The notation $\mathcal O_F/\mathfrak m_F$ is harmless only if its meaning is specified: for a
proper-class field an equivalence class modulo infinitesimals can itself be a proper class
and cannot be an element of a set or of a class of sets. The residue field is
\emph{represented by the set of real cuts} $\Tr(F)$, with the displayed standard-part map;
it has the usual quotient universal property without treating proper classes as elements,
and ``residue field'' always has this meaning here (source~12).

\begin{theorem}[Exact trace of an old surreal field; source 12 (base), source 09]\label{univ:gs:thm:oldtrace}
For every same-ordinal pair $M\subseteq N$, $\Tr(\NoM)=\R^M\subseteq\R^N$. If
$M_0,M_1\subseteq N$ and there is any unital field embedding $\No^{M_0}\hookrightarrow\No^{M_1}$,
then $\R^{M_0}\subseteq\R^{M_1}$. Consequently an isomorphism forces equality of the real
traces as actual subfields of the ambient real line, and $\R^{M_0}\ne\R^{M_1}$ rules out an
isomorphism of ordered fields and of pure fields.
\end{theorem}
\begin{proof}
Let $x\in\NoM$ be finite. Its cut $C_x$ is a set in $M$, defined from the old parameter $x$
by rational inequalities whose truth is absolute; completeness of the reals \emph{inside $M$}
gives $r\in\R^M$ representing it, and in $N$ it is the same cut, so $r=\st_N(x)$. Every old
real is an old finite surreal with itself as standard part. A unital embedding between real
closed fields preserves their unique orders ($a<b$ iff $b-a$ is a nonzero square), so
\cref{univ:gs:prop:tracefunctor} applies; an isomorphism and its inverse give equality.
\end{proof}

\emph{Second route (source~09, its Lemma~6.2).} Every $r\in\R^M$ realizes its own rational
cut. Conversely, an old surreal $x$ bounded by an integer has, inside $M$, a real supremum
$r=\st(x)\in\R^M$ of its rational cut; equivalently $x-r$ is smaller in absolute value than
every positive rational, and these individual rational inequalities are absolute, so the same
$r$ is determined by $x$ in $N$. An ordered-field isomorphism fixes $\Q$ and hence the real
label of every rational cut it realizes, and for real closed fields a pure-field isomorphism
is order preserving.

An embedding is \emph{not} asserted to fix every real element pointwise as a surreal; it
could move a transcendental real by an infinitesimal amount. The rigid conclusion is that its
\emph{residue} has the same rational cut, which is exactly what the nonembedding proof needs
(source~12).

\emph{The trace as a spectrum of countable types (source~12).} For $r\in\R^N$ let
\[
 p_r(X)=\{q<X:q\in\Q,\ q<r\}\cup\{X<q:q\in\Q,\ r<q\},
\]
rational constants abbreviating terms, or inequalities with cleared integer denominators.
Every finite subset is satisfiable in $\Q$, and $p_r$ is realized in an ordered subfield
$F\subseteq\NoN$ exactly when $r\in\Tr(F)$: a realization is finite with standard part $r$,
and conversely. So $p_r$ is realized in $\NoM$ iff $r\in\R^M$ (compare \eqref{univ:eq:newrealtype}).
The sentence $\Phi_r=\exists X\bigwedge p_r(X)$ is an $L_{\omega_1,\omega}$ sentence of the
ordered-field language and, by replacing order with nonzero-square formulas, of the pure field
language; it is not a finitary sentence. The real $r$ occurs only in the \emph{external code of
its rational inequalities}, not as a named parameter, so there is no conflict with the
elementary equivalence of all real closed fields.

\subsection{The Boolean family}

\begin{theorem}[Boolean embeddability with a constant full gap spectrum; source 12]\label{univ:gs:thm:boolean}
Let $V$ be a ground universe, $I\in V$ a set and $*\notin I$ a reserved coordinate. Let
$G\subseteq\Add(\omega,I\cup\{*\})^V=\Fn((I\cup\{*\})\times\omega,2,{<}\omega)^V$ be generic and
$N=V[G]$. For $A\in\Pow(I)^V$ let $G_A$ be the restriction to $A$ and $M_A=V[G_A]$. Then, in
$N$:
\begin{enumerate}[label=\textnormal{(\roman*)}]
\item $\GapSpec_N(\No^{M_A})=\{(\omega,\omega)\}$ for every $A$;
\item for arbitrary unital class field embeddings,
$\No^{M_A}\hookrightarrow\No^{M_{A'}}\iff A\subseteq A'$;
\item $\No^{M_A}\cong\No^{M_{A'}}$ if and only if $A=A'$;
\item all $\No^{M_A}$ are elementarily equivalent real closed fields, and none is
$\kappa$-saturated for any infinite $N$-cardinal $\kappa$;
\item as pure orders, all are $\omega$-saturated and fail $\kappa$-saturation for every
uncountable $\kappa$;
\item every $N$-set of elements of $\No^{M_A}$ is bounded above and below in $\No^{M_A}$.
\end{enumerate}
\end{theorem}

The reserved coordinate is not cosmetic: it makes $N$ add a real over \emph{every} $M_A$,
including $M_I$, so no endpoint of the family has a different spectrum or saturation status.
The family is uniformly class-defined over the set of indices $\Pow(I)^V$; it is not a set
whose elements are proper classes. No large cardinal, CH or GCH is used (source~12).

Let $J=I\cup\{*\}$. For $A\in\Pow(I)^V$ restriction gives
$\Add(\omega,J)^V\cong\Add(\omega,A)^V\times\Add(\omega,J\setminus A)^V$. Finite partial functions
on a ground set are the same in an intermediate extension, so in $M_A$ the second factor is
precisely $\Add(\omega,J\setminus A)^{M_A}$, not an obsolete ground version of a forcing with
new finite conditions missing, and $N=M_A[G_{J\setminus A}]$; the quotient is nontrivial
because it contains the reserved coordinate. For $i\in J$ let $y_i:\omega\to2$ be the $i$-th
generic sequence and $r_i=\sum_{n<\omega}2y_i(n)/3^{n+1}\in\R^N$. This ternary coding has a
canonical inverse on its image and, unlike the binary expansion map, is injective even at
endpoints: at the first different digit the leading difference exceeds the largest possible
opposing tail.

\begin{lemma}[Coordinate detection; source 12]\label{univ:gs:lem:detection}
For $i\in I$ and $A\in\Pow(I)^V$, $r_i\in\R^{M_A}\iff i\in A$; and $r_*\notin\R^{M_A}$ for
every such $A$.
\end{lemma}
\begin{proof}
If $i\in A$, $y_i\in M_A$ and so is its code. If $i\notin A$, $y_i$ is Cohen generic over
$M_A$, and for every binary $d\in M_A$ disagreement on the $i$-coordinate is dense (a finite
condition leaves a position free); so $y_i\notin M_A$, and the canonical ternary decoder in
$M_A$ would recover $y_i$ from an old $r_i$. The same holds for $*$.
\end{proof}

\begin{proof}[Proof of \cref{univ:gs:thm:boolean}]
(i) Apply \cref{univ:gs:prop:cohenreals} inside each $M_A$ to its quotient; this excludes
\emph{every} uncountable character, not only the first. (ii) If $A\subseteq A'$, inclusion of
old signs is an ordered-field embedding. If $j$ is any unital embedding,
\cref{univ:gs:thm:oldtrace} gives $\R^{M_A}\subseteq\R^{M_{A'}}$, and for $i\in A$,
$r_i\in\R^{M_A}$ forces $i\in A'$ by \cref{univ:gs:lem:detection}. (iii) Apply (ii) both ways.
(iv) Real closed fields share one complete theory. The type $p_{r_*}$ is finitely satisfiable
and omitted in every $\No^{M_A}$; for a complete omitted type take the complete type of $r_*$ in
$\NoN$ over the empty set, every finite part of which is realized in $\No^{M_A}\preccurlyeq\NoN$.
So $\omega$-saturation fails, hence every larger saturation. (v) A dense linear order without
endpoints is $\omega$-saturated (types over a finite set are realized by named points, the
finitely many intervening intervals or the two unbounded ones), and a countably presented
omitted gap from (i) gives an omitted type over a countable set. (vi) is
\cref{univ:lem:oldbounds}.
\end{proof}

No field-theoretic preservation requirement was imposed beyond a unital embedding: the
nonembedding direction allows maps that move $\omega$, change birthdays, fail to preserve
simplicity or normal-form monomials; rational-cut rigidity still blocks them (source~12).

\noindent[Added 4 October 2026, batch~90: a countable contrast. The report
\path{foundations-and-computation/cantor-families-of-surreal-subfields}
(batch-90 manuscript~04; AI-assisted, unrefereed, not formalized) also
indexes real closed fields by the subsets $A$ of a set: the countable
relative real closures $K_A$ inside one countable real closed subfield of
$\No$. There embeddability is not inclusion: on its colored-order subfamily
it is embeddability of colored rank orders (its Theorem~7.3,
\path{csf:thm:colored-coding}), and on the whole family it is a complete
analytic quasiorder (its Theorem~1.1, \path{csf:thm:main}). No proof is
shared.]

\begin{example}[A four-field diagram; source 12]\label{univ:gs:ex:fourfield}
With $I=\{0,1\}$ and a reserved coordinate, the four fields embed exactly as
$\varnothing\subset\{0\},\{1\}\subset\{0,1\}$. The middle two do not embed into one another;
each embeds properly into $\No^{M_{\{0,1\}}}$, and $\No^{M_\varnothing}$ into each. All four
have only countable-character set-presented gaps and none is $\omega$-saturated as a field.
A diagram this small already refutes the implication from equal full spectra to field
isomorphism.
\end{example}

\begin{corollary}[One Cohen real suffices; source 12]\label{univ:gs:cor:onecohen}
In a one-Cohen-real extension of $V$ there is a family indexed by $\Pow(\omega)^V$ of pairwise
nonisomorphic old surreal fields with all the properties of \cref{univ:gs:thm:boolean}. Its
size is the continuum of the extension, and each member's real trace has that cardinality.
\end{corollary}
\begin{proof}
Use $I=\omega$; a ground bijection of $(\omega\cup\{*\})\times\omega$ with $\omega$ identifies
the forcing with one-real Cohen forcing. With $\mathfrak c^V=(2^{\aleph_0})^V$, the index set has
size $\mathfrak c^V$, preserved by the c.c.c. A countable forcing has at most $\mathfrak c^V$
nice names for reals (each is coded by a subset of a countable set), every real has one, and
all ground reals remain; so $\mathfrak c^N=\mathfrak c^V$. Each $M_A$ contains $\R^V$ and lies
in $N$, so $\abs{\R^{M_A}}^N=\mathfrak c^N$.
\end{proof}

The family is indexed by \emph{ground} subsets of $\omega$, not by all subsets in $N$; the
equality of cardinalities just proved is what licenses ``continuum''. So the examples do not
exploit different sizes of residue fields: a two-element inclusion chain and a two-element
incomparability pattern both occur with the spectrum $\{(\omega,\omega)\}$ (source~12).

\begin{corollary}[Universality for ground set-sized partial orders; source 12]\label{univ:gs:cor:poset}
For every set-sized partial order $(P,\le)\in V$, a Cohen extension contains old surreal fields
$E_p$, $p\in P$, with $E_p\hookrightarrow E_q\iff p\le q$ and
$\GapSpec_N(E_p)=\{(\omega,\omega)\}$; distinct indices give nonisomorphic fields with the
saturation profiles of \cref{univ:gs:thm:boolean}.
\end{corollary}
\begin{proof}
Take $I=P$ and $A_p=\{q:q\le p\}\in V$. Then $p\le q\iff A_p\subseteq A_q$, and
$E_p=\No^{M_{A_p}}$; antisymmetry and \cref{univ:gs:thm:boolean}(iii) give nonisomorphism.
\end{proof}

``Universality'' means exactly this realization of ground set-sized partial orders; it is not
a statement about Borel or effective reducibility, nor about a proper-class collection of all
partial orders. The order need not be well founded, countable or a lattice, so prescribed
antichains, descending chains and mixed diagrams occur inside one gap-spectrum class. The
trace does not classify arbitrary old fields; it classifies embeddability \emph{within this
family} (source~12).

\begin{theorem}[Failure of the unrestricted field-classification implication; source 12]\label{univ:gs:thm:incompletetrace}
There are same-ordinal $M_0,M_1\subseteq N$ with
$\GapSpec_N(\No^{M_0})=\GapSpec_N(\No^{M_1})$ and $\No^{M_0}\not\cong\No^{M_1}$ as fields. They
can be chosen with the same saturation profile at every infinite cardinal, the same least new
birthday $\omega$, and equally large Archimedean residue fields; in one example neither field
embeds into the other.
\end{theorem}
\begin{proof}
Use two incomparable coordinate sets in the countable-coordinate construction; apply
\cref{univ:gs:thm:boolean,univ:gs:cor:onecohen}. Finite signs are absolute and the reserved
coordinate gives a new sign of length $\omega$, so the first new birthday is $\omega$ over both.
\end{proof}

This is a negative answer to the unrestricted \emph{real-closed-field} implication asked about
in \cref{univ:q:completeness}. It does not answer the pure linear-order question: an arbitrary
order isomorphism need not preserve $0$, $1$, addition or the rational copy, so the trace
obstruction does not apply to it. Nor does it rule out positive classification under
additional hypotheses, such as equality of real traces together with suitable higher
compatibility data (source~12).

\begin{proposition}[The trace is exact on the family; source 12]\label{univ:gs:prop:traceexact}
For the fields of \cref{univ:gs:thm:boolean} the following are equivalent: (i)
$\No^{M_A}\hookrightarrow\No^{M_{A'}}$; (ii) $\Tr(\No^{M_A})\subseteq\Tr(\No^{M_{A'}})$; (iii) every
type $p_r$ realized in $\No^{M_A}$ is realized in $\No^{M_{A'}}$; (iv) $A\subseteq A'$. Equality
of traces is equivalent to $A=A'$, hence to field isomorphism within the family.
\end{proposition}
\begin{proof}
\Cref{univ:gs:thm:oldtrace} gives (i)$\Rightarrow$(ii), the type characterization gives
(ii)$\Leftrightarrow$(iii), \cref{univ:gs:lem:detection} gives (ii)$\Rightarrow$(iv), and
inclusion gives (iv)$\Rightarrow$(i).
\end{proof}

The fields do not have different first-order theories; they realize different \emph{sets of
countable types} while agreeing on saturation at every cardinal, so appending the entire
saturation profile to the gap spectrum does not repair the proposed classification invariant.
There is no claim that $\Tr(F)$ is efficiently computable or finitary definable; its
significance is structural (source~12).

\begin{center}\small
\begin{tabular}{@{}P{0.44\textwidth}P{0.46\textwidth}@{}}
\toprule
Invariant or property (source~12) & Value throughout the countable-coordinate family\\
\midrule
First-order ordered-field theory & $\RCF$\\
Full set-presented gap spectrum & $\{(\omega,\omega)\}$\\
Field saturation at infinite cardinals & Fails at every infinite cardinal\\
Pure-order saturation at infinite cardinals & Holds exactly at $\omega$\\
Least new birthday over the inner universe & $\omega$\\
Set-sized cofinal or coinitial subset & None\\
Cardinality of the real trace & $\mathfrak c^N$\\
The actual real trace & $\R^{V[G_A]}$, different for different $A$\\
Field embeddability & Exactly inclusion of the coordinate sets\\
\bottomrule
\end{tabular}
\end{center}

\section{Common small forcing and Cohen amplification}\label{univ:gs:sec:amplification}

This section is source~09's. Its technical ingredient compares different old fields in one
ambient universe. It concerns a square, not a two-step iteration with an arbitrarily
reinterpreted second forcing:
\[
 \begin{array}{ccc}
 M&\subseteq&W\\[-1pt]
 \cap&&\cap\\[-1pt]
 M[G_{\mathbb Q}]&\subseteq&W[G_{\mathbb Q}],
 \end{array}
\]
where $G_{\mathbb Q}$ is generic over $W$ for the \emph{same set partial order}
$\mathbb Q\in M$; the outer version of $\mathbb Q$ has no additional conditions.

\subsection{Two name facts and the base-change theorem}

\begin{lemma}[Atomic forcing is absolute for old names; source 09]\label{univ:gs:lem:atomic}
Let $M\subseteq W$ be transitive, $\mathbb Q\in M$ the identical partial order in both, and
$\sigma,\tau\in M$ $\mathbb Q$-names. For every $q\in\mathbb Q$, $q\forces\sigma\in\tau$ and
$q\forces\sigma=\tau$ are absolute between $M$ and $W$, and so are their forced negations.
\end{lemma}
\begin{proof}
Induct on the ranks of the names through the recursive definitions of atomic forcing. The
quantifiers range over the same conditions, extensions and members of old names, with
unchanged compatibility and order. Forced negation is the absence of an extension forcing the
positive assertion, again over the same set of conditions.
\end{proof}

\begin{lemma}[Intersection for ordinal-valued functions; source 09]\label{univ:gs:lem:nestedintersection}
Let $G_{\mathbb Q}$ be $\mathbb Q$-generic over $W$ as above. An ordinal-valued function
belonging to both $W$ and $M[G_{\mathbb Q}]$ belongs to $M$.
\end{lemma}
\begin{proof}
Let $f:\gamma\to\Gamma$ be in both, $\Gamma$ an old ordinal bound. There is an old name $\tau$
with $\tau^{G_{\mathbb Q}}=f$, and in $W$ the truth lemma gives $q\in G_{\mathbb Q}$ forcing
$\tau=\check f$. For every $(\alpha,\zeta)\in\gamma\times\Gamma$, membership in $f$ is
equivalent to $q\forces\check{(\alpha,\zeta)}\in\tau$, absolute by
\cref{univ:gs:lem:atomic}; Separation in $M$ recovers the graph of $f$ from
$q,\tau,\gamma,\Gamma$.
\end{proof}

Only this function version of the usual intersection argument is needed. A fresh function
already in $W$ remains new over $M[G_{\mathbb Q}]$, and its prefixes remain old because they
were in $M$ to start with.

\begin{theorem}[Common-small-forcing invariance; source 09]\label{univ:gs:thm:basechange}
Let $M\subseteq W$ have the same ordinals, $\mathbb Q\in M$ the same set partial order in both,
$G_{\mathbb Q}$ generic over $W$ and $U=W[G_{\mathbb Q}]$. For every limit ordinal $\gamma$ with
\begin{equation}\label{univ:gs:eq:smallbound}
 \cf^U\gamma>\abs{\mathbb Q}^U
\end{equation}
we have $\Fr_\gamma(M,W)\iff\Fr_\gamma(M[G_{\mathbb Q}],U)$. The forward implication does not
need \eqref{univ:gs:eq:smallbound}.
\end{theorem}
\begin{proof}
Forward: a fresh $f\in W$ over $M$ has proper prefixes in $M\subseteq M[G_{\mathbb Q}]$, and
$f\notin M[G_{\mathbb Q}]$ by \cref{univ:gs:lem:nestedintersection}.

Converse: take $f:\gamma\to\Gamma$ in $U$ fresh over $M[G_{\mathbb Q}]$, a $W$-name $\dot f$ and
$q_0\in G_{\mathbb Q}$ forcing that it is a function with this domain and range bound. A
\emph{ground set} of names is needed, not an unspecified collection of all names: identify a
graph name with a subset of $(\gamma\times\Gamma)\times\mathbb Q$, using check names for graph
pairs; let $T$ be the set of such names in $M$ and $T_\alpha$ those supported on
$(\alpha\times\Gamma)\times\mathbb Q$. Restricting a graph name to $\alpha$ is a ground operation,
and every subset of $\alpha\times\Gamma$ in $M[G_{\mathbb Q}]$ has a name in $T_\alpha$ by the
nice-name construction. Fix in $M$ a bijection $e:\tau_0\to T$ from an ordinal. For $q\le q_0$
define in $W$
\[
 B_q=\{\alpha<\gamma:\exists\sigma\in T_\alpha\ \ q\forces_W\dot f\res\alpha=\sigma\}.
\]
Every $\alpha<\gamma$ lies in some $B_q$ with $q\in G_{\mathbb Q}$, since $f\res\alpha\in
M[G_{\mathbb Q}]$ and the truth lemma supplies a condition. By \eqref{univ:gs:eq:smallbound}
one $q\in G_{\mathbb Q}$, $q\le q_0$, has $B_q$ unbounded in $\gamma$, an absolute property of
this old set $B_q\in W$. Restricting a witness at a larger member of $B_q$ shows that for
\emph{every} $\alpha<\gamma$ some $\sigma\in T_\alpha$ has $q\forces_W\dot f\res\alpha=\sigma$.
Define in $W$
\[
 h(\alpha)=\min\{\xi<\tau_0:e(\xi)\in T_\alpha,\ q\forces_W\dot f\res\alpha=e(\xi)\}.
\]
For $\zeta<\gamma$ fix $\sigma_\zeta\in T_\zeta$ with $q\forces_W\dot f\res\zeta=\sigma_\zeta$;
for $\alpha<\zeta$ the minimum is the least index of a $\sigma\in T_\alpha$ with
$q\forces_W\sigma=\sigma_\zeta\res\alpha$, an equality-forcing relation between old names,
computed in $M$ by \cref{univ:gs:lem:atomic}. So $h\res\zeta$ is constructed in $M$ from
$q,\sigma_\zeta,e,\zeta$. If $h\in M$, the union $\upsilon=\bigcup_{\alpha<\gamma}e(h(\alpha))$ is
in $M$; each selected name evaluates to $f\res\alpha$ since $q\in G_{\mathbb Q}$, graph-name union
commutes with evaluation, and $\upsilon^{G_{\mathbb Q}}=f$ as $\gamma$ is a limit, putting $f$
in $M[G_{\mathbb Q}]$. So $h$ is fresh over $M$.
\end{proof}

\begin{corollary}[Common dense suborder; source 09]\label{univ:gs:cor:densesub}
The theorem holds with $\abs{\mathbb Q}^U$ replaced by $\abs D^U$ for a dense suborder $D\in M$
of $\mathbb Q$ that is the identical dense suborder in $M$ and $W$.
\end{corollary}
\begin{proof}
Use the induced $D$-generic filter; the extensions are unchanged.
\end{proof}

The bound is strict; at the size of the forcing the pigeonhole step can fail. The argument is
not replaced by a chain-condition assertion, and the countable boundary is not claimed to
behave like uncountable cofinalities; the counterexamples below exploit exactly the difference
between countable ordinal sequences and reals. The theorem keeps the underlying poset fixed: a
ground forcing that is closed need not remain closed after a small generic is added, since the
new universe has new descending sequences of old conditions, and recomputing a forcing from a
definition may supply additional conditions; neither operation is the fixed-poset square. So
there is no conflict with the small-then-strategically-closed gap-forcing phenomena of Hamkins
\cite{Hamkins}, also discussed in \cite[Section~2]{FKW} (source~09). No assertion of first
discovery of the general forcing technique is made.

\subsection{Cohen amplification}

\begin{theorem}[Cohen-coordinate amplification; source 09]\label{univ:gs:thm:amplification}
Let $M\subseteq W$ be same-ordinal transitive universes with (i) $\R^M=\R^W$ and (ii) some
$g:\omega\to\Ord$ in $W\setminus M$. Let $G_{\mathbb C}$ be generic over $W$ for the ground
countable forcing $\mathbb C=\Fn(\omega\times\omega,2,{<}\omega)$ and $N=W[G_{\mathbb C}]$. For
$S\in\Pow(\omega)^M$ let $G_S$ be the restriction to $S\times\omega$ and
$M_S=M[G_S]$, $W_S=W[G_S]$. Then, in $N$:
\begin{enumerate}[label=\textnormal{(\alph*)}]
\item $\Char_N(\No^{M_S})=\Char_W(\NoM)$ for every $S$: all these fields have the same full gap
spectrum, and the cofinality labels on the right are unchanged by the final Cohen forcing;
\item for $S\ne T$, $\No^{M_S}\not\cong\No^{M_T}$, even as pure fields;
\item the index set $\Pow(\omega)^M$ has cardinality $2^{\aleph_0}$ in $N$;
\item $\No^{M_\omega}$ is $\omega$-saturated as a real closed field, and each $\No^{M_S}$ with
$S\subsetneq\omega$ is not;
\item every $\No^{M_S}$ is a real closed elementary subfield of $\NoN$ with no $N$-set-sized
cofinal subset.
\end{enumerate}
\end{theorem}
\begin{proof}
\emph{Step 1: coordinate factorization.} For old $S$ the ground poset is
$\mathbb C_S\times\mathbb C_{\omega\setminus S}$ with the expected restriction maps, all
countable and identical in $M$ and $W$. So $G_S$ is generic over $W$, the complementary filter
over $W_S$, and $M_S\subseteq N$ is an inner universe. Cohen forcing preserves cardinals and
cofinalities, so $\Reg^{M_S}=\Reg^M$, and the cofinalities of a fixed ordinal in $W$, $W_S$ and
$N$ agree.

\emph{Step 2: agreement above countable cofinality.} For $\lambda\in\Reg^M$ with
$\cf^N\lambda>\omega$, \cref{univ:gs:thm:basechange} for $\mathbb C_S$ gives
$\Fr_\lambda(M,W)\iff\Fr_\lambda(M_S,W_S)$ (empty $S$ is trivial). Passing from $W_S$ to $N$ by
the countable complementary forcing: a function fresh over $M_S$ in $W_S$ stays fresh in $N$;
conversely, a function fresh over $M_S$ in $N$ but not in $W_S$ would have its proper prefixes
in $M_S\subseteq W_S$ and be fresh over $W_S$, contrary to \cref{univ:gs:thm:filterbase} since
$\cf^N\lambda>\omega$. So $\Fr_\lambda(M_S,W_S)\iff\Fr_\lambda(M_S,N)$, and with
\cref{univ:gs:thm:transfer} the character sets agree above $\omega$.

\emph{Step 3: the countable character occurs everywhere.} Finite restrictions of $g$ are old,
so $g$ is fresh over $M$ in $W$; by \cref{univ:gs:lem:nestedintersection}, $g\notin
M[G_{\mathbb C}]$, so $g\notin M_S$, and the same prefixes make $g$ fresh over $M_S$ in $N$. The
transfer gives character $\omega$ in every field, including $\NoM$ in $W$. This proves (a).

\emph{Step 4: coordinate reals distinguish the fields.} Let $y_n:\omega\to2$ be the $n$-th
coordinate real, identified through a fixed ground injection with a real of the Cantor set (a
coding that preserves membership in each universe). If $n\in S\setminus T$, then $y_n\in M_S$;
and $y_n\notin M_T$, since by product factorization over $M$ the $n$-th coordinate filter is
Cohen generic over the extension by all coordinates in $T$. So $\R^{M_S}\ne\R^{M_T}$, and
\cref{univ:gs:thm:oldtrace} gives (b), without distinguishing saturation levels of two proper
subsets.

\emph{Step 5: all final reals already appear over $M$.} Every real of $W[G_{\mathbb C}]$ has a
nice $\mathbb C$-name in $W$, coded by a subset of the countable ground set
$\omega\times\mathbb C$, hence by a real of $W$, which lies in $M$ by (i); the decoded name is an
$M$-name, so $\R^N=\R^{M[G_{\mathbb C}]}$. This uses the countability of the \emph{poset}, not
merely its chain condition. The old real codes of valid names map onto all final reals (and
all old codes do, by a default value), so in $N$,
$2^{\aleph_0}=\abs{\R^N}\le\abs{\Pow(\omega)^M}\le\abs{\Pow(\omega)^N}=2^{\aleph_0}$, which is
(c), even if $W$ changed the ground cardinal label of the continuum. By
\cref{univ:thm:omega}, $\No^{M_\omega}$ is $\omega$-saturated; for proper $S$ choose $n\notin S$,
and $y_n\in N\setminus M_S$ shows that $\No^{M_S}$ is not. This is (d).

\emph{Step 6.} Real closedness and elementarity are \cref{univ:prop:absolute} for
$M_S\subseteq N$, and \cref{univ:lem:oldbounds} excludes set-sized cofinal subsets.
\end{proof}

The theorem holds for \emph{every} character, not only the least, and produces the examples in
one fixed universe, where a class-field isomorphism question between them is meaningful. The
same-spectrum conclusion rests on \cref{univ:gs:thm:basechange}; it does not follow from the
mere fact that a countable forcing adds no fresh function over its immediate ground at
uncountable cofinality, because the old base has changed from $M$ to $M_S$ and that change must
be controlled. Nor is the distinction between the fields a difference of names of embedded
reals: a field isomorphism fixes every rational and preserves positivity, hence every rational
cut, so the fields are abstractly nonisomorphic, not merely unequal subclasses of $\NoN$
(source~09).

\begin{corollary}[An immediate obstruction to completeness; source 09]\label{univ:gs:cor:noncomplete}
Whenever the hypotheses of \cref{univ:gs:thm:amplification} are realized, equality of the entire
set-presented gap spectrum is not a complete invariant for old surreal real closed fields in one
universe; the failure has continuum many pairwise nonisomorphic witnesses with exactly the same
invariant.
\end{corollary}

\subsection{Two exact spectra: Prikry and Namba}

Let $M$ have a measurable cardinal $\rho$ with a normal measure $\mathcal U$, and let $W=M[G]$
for the usual Prikry forcing $\mathbb P_{\mathcal U}$ (finite increasing stems and measure-one
reservoirs). The standard Prikry properties give no new bounded subsets of $\rho$, preservation
of cardinals, and $\cf^W\rho=\omega$; so no reals are added and a countable ordinal sequence is.

\begin{theorem}[Published Prikry spectrum; imported]\label{univ:gs:thm:prikryinput}
$\FRESH(\mathbb P_{\mathcal U})=\{\omega,\rho\}$; both members occur in every Prikry generic
extension, and no others occur \cite[Theorem~7.21(1)]{FKW}.
\end{theorem}

The lower witnesses are the Prikry sequence on $\omega$ and its range, a fresh subset of $\rho$
all of whose proper initial segments are finite; the upper exclusions are part of the cited
theorem, and no GCH is required (source~09). Source~11 recalls the same facts: the actual
spectrum bound follows from the cited theorem, and both entries occur in each generic extension
by the displayed witnesses, so nothing is inferred from a possible-generic union.

\begin{corollary}[Prikry's exact surreal gap spectrum; source 09, source 13]\label{univ:gs:cor:prikry}
For $W=M[G]$ as above, $\Char_W(\NoM)=\{\omega\}$. After adjoining countably many Cohen reals
over $W$, the final universe contains continuum many pairwise nonisomorphic old surreal fields,
all with $\GapSpec=\{(\omega,\omega)\}$; one of them is $\omega$-saturated as a real closed
field, and those indexed by proper coordinate subsets are not.
\end{corollary}
\begin{proof}
Both ground regulars of \cref{univ:gs:thm:prikryinput} have cofinality $\omega$ in $W$; apply
\cref{univ:gs:thm:transfer} and then \cref{univ:gs:thm:amplification}.
\end{proof}

This refines \cref{univ:cor:prikry}, which exhibited $(\omega,\omega)$ gaps: every gap is of that
character. Reporting a $(\rho,\rho)$ gap would use the wrong universe's cofinality. A measurable
cardinal supplies a particularly small common spectrum, but the failure of classification itself
needs no large cardinal (source~09).

Let $M\models\CH$ and let $\mathbb P_{\mathrm{Nb}}$ be Namba forcing in the form of
\cite[Section~7.6]{FKW}: trees on finite sequences from $\omega_2$ in which every node extends to
a node with $\omega_2$ immediate successors, ordered by inclusion; let $W=M[G]$. Under CH it adds
no reals but adds a countable sequence cofinal in $\omega_2^M$.

\begin{theorem}[Published Namba spectrum; imported]\label{univ:gs:thm:nambainput}
Under CH, $\FRESH(\mathbb P_{\mathrm{Nb}})=\{\omega,\omega_1^M,\omega_2^M\}$; all three lower
witnesses occur in each generic extension, the upper exclusion holds for all generics, $\omega_1$
is preserved and $\cf^W(\omega_2^M)=\omega$ \cite[Theorem~7.23(1)]{FKW}.
\end{theorem}

The no-new-reals property and cofinality facts are recalled in the same section of \cite{FKW};
CH is essential to the no-new-reals assertion. No claim is made that the forcing preserves all
cardinals or all other cofinalities; $\omega_2^M$ is an ordinal label, not the second uncountable
cardinal of $W$ (source~09).

\begin{corollary}[Namba's exact surreal gap spectrum; source 09, source 13]\label{univ:gs:cor:namba}
For this Namba extension $\Char_W(\NoM)=\{\omega,\omega_1^W\}$. After adjoining countably many
Cohen reals over $W$, the final universe contains continuum many pairwise nonisomorphic old
surreal fields inside $\NoN$, all with $\GapSpec_N=\{(\omega,\omega),(\omega_1^N,\omega_1^N)\}$.
This needs CH in the ground but no large-cardinal hypothesis.
\end{corollary}
\begin{proof}
The outer cofinalities of the three ground regulars are $\omega$, $\omega_1^W$ and $\omega$.
Apply the transfer; the hypotheses of \cref{univ:gs:thm:amplification} hold, and the final
countable Cohen forcing preserves $\omega_1^W$.
\end{proof}

\begin{remark}[A corrected overstatement of source 09]\label{univ:gs:rem:overstatement1}
Source~09 concludes here that ``the unconditional implication proposed as a possible
classification principle in the repository'' cannot be a theorem. This report proposed no such
principle: \cref{univ:q:completeness} only \emph{asks} under what hypotheses equality of gap
spectra implies isomorphism, and \cref{univ:sec:leastgap} states that equality of gap spectra is
not proved to classify the fields. Read correctly, the corollary shows that the unconditional
implication about which that question asks fails for real closed fields, even over the usual
CH-based forcing setting. It is a relative construction in the ordinary sense of forcing, not a
proof of the consistency of $\ZFC$.
\end{remark}

\noindent[Added 4 October 2026, batch~93: Part~II of
\path{foundations-and-computation/definable-surreals-and-omnific-integers}
(batch-93 manuscript~02; AI-assisted, unrefereed, not formalized) uses Namba
forcing over a model of $V=L$, in the Laver-style presentation in which every
node at or above the stem has $\omega_2$ immediate successors, to make
countable assembly of definitions from a real and ordinals fail while
$\mathrm{HOD}_{\R}=\mathrm{HOD}=L$ (its Theorem~25.6, \path{dsn:rp:thm:namba}).
It imports the no-new-reals theorem under CH for that presentation from
expository sources and records the missing textbook reference as its
Question~28.17 (\path{dsn:rp:q:nambainput}). \Cref{univ:gs:thm:nambainput} is
imported for the classical presentation of \cite[Section~7.6]{FKW}, in which
every node extends to a node with $\omega_2$ immediate successors; the two
presentations are not forcing equivalent.]

\begin{center}\small
\begin{longtable}{@{}P{0.25\textwidth}P{0.25\textwidth}P{0.19\textwidth}P{0.22\textwidth}@{}}
\toprule
Extension (sources) & Ground-regular fresh spectrum & Outer characters $\Char$ & Hypothesis\\
\midrule
\endhead
One Cohen subset of regular $\kappa$ (13; \cref{univ:gs:thm:cohen}) & $[\kappa,(2^{<\kappa})^M]\cap\Reg^M$ \cite[Proposition~6.1]{FKW} & $\{\kappa\}$ & none\\[3pt]
Easton Cohen product (09) & $\Ecl^M(A)$ & $\Ecl^M(A)$ & GCH in $M$\\[3pt]
Prikry at $\rho$ (09, 13) & $\{\omega,\rho\}$ & $\{\omega\}$ & normal measure on $\rho$\\[3pt]
Namba (09, 13) & $\{\omega,\omega_1^M,\omega_2^M\}$ & $\{\omega,\omega_1^W\}$ & CH in $M$\\[3pt]
Cohen-amplified old bases $M_S$ (09) & not asserted equal as ground-regular sets at countable outer cofinality & exactly $\Char_W(\NoM)$ & no new reals in $W$, a new countable ordinal sequence\\
\bottomrule
\end{longtable}
\end{center}

For the Prikry and Namba rows, $\cf^N\rho=\omega$, $\cf^N\omega_2^M=\omega$ and
$\omega_1^M=\omega_1^N$; the forcing facts are standard \cite{Jech} and the fresh-domain lists are
credited to \cite{FKW}. The Cohen row is proved independently of its middle entry
(\cref{univ:gs:thm:cohen}). The table shows why the transfer must use the map
$\lambda\mapsto\cf^N\lambda$, not merely intersect a ground spectrum with the new cardinals
(source~13). The last row is intentionally not strengthened to equality of the
\emph{ground-regular} fresh spectra: the proof controls their images under outer cofinality, with
a separate argument at $\omega$, which is exactly what surreal gap spectra need (source~09).
Countable-character gaps do not by themselves say that a new real appeared: Prikry adds no reals,
nor does Namba under CH; the gap sees a new arrangement of old parameters, not necessarily a new
binary sequence on $\omega$ (source~13).

\section{A Cohen--Prikry tower}\label{univ:gs:sec:tower}

This section is source~11's. It uses two standard inputs: small forcing preserves a measurable
cardinal \cite{LevySolovay}, and for Prikry forcing at a normal measure on a measurable $\rho$ no
bounded subsets of $\rho$ are added, $\cf(\rho)$ becomes $\omega$, the actual ground-regular
spectrum is $\{\omega,\rho\}$ (\cref{univ:gs:thm:prikryinput}), and cardinals and all other
regular cofinalities are preserved. These are imported.

\begin{lemma}[An elementary tower bound; source 11]\label{univ:gs:lem:tower}
Let $M_0\subseteq M_1\subseteq N$ have the same ordinals, with $M_0\subseteq M_1$ preserving
regular cardinals. If $f:\lambda\to\Ord$ is fresh over $M_0$ and $\lambda$ is regular in $M_0$,
then $\lambda\in\FF(M_0,M_1)$ or $\lambda\in\FF(M_1,N)$. Consequently
$\FF(M_0,N)\subseteq\FF(M_0,M_1)\cup\FF(M_1,N)$.
\end{lemma}
\begin{proof}
If $f\in M_1$, its proper prefixes are in $M_0$ and it witnesses the first alternative;
otherwise the same prefixes lie in $M_1$ and $f$ is fresh over $M_1$. Regularity preservation
lets $\lambda$ index both spectra.
\end{proof}

This is an upper inclusion only: freshness over the intermediate model does not imply that all
proper prefixes were already in the earlier model.

\begin{theorem}[Failure of completeness of the full gap spectrum; source 11]\label{univ:gs:thm:incomplete}
Suppose $V\models\ZFC$ has a measurable cardinal $\rho$. There is a set-forcing tower
$V\subset V[y]\subset N$ with
\[
 \Char_N(\No^V)=\Char_N(\No^{V[y]})=\{\omega\},
\]
but $\No^V$ and $\No^{V[y]}$ are not isomorphic as ordered fields in $N$: only the second is
$\omega$-saturated in $N$.
\end{theorem}
\begin{proof}
Add one Cohen real $y$. The forcing is countable, so $\rho$ stays measurable in $V[y]$; choose
a normal measure there and force with \emph{its} Prikry forcing (genuinely the one of $V[y]$, not
an unchanged ground poset), obtaining $N$. For $V[y]\subseteq N$ the imported spectrum and
\cref{univ:gs:thm:transfer} give $\Char_N(\No^{V[y]})=\{\cf^N\omega,\cf^N\rho\}=\{\omega\}$. For
the whole tower: countable Cohen forcing has no fresh functions at uncountable ground regulars
(\cref{univ:gs:thm:filterbase}) and adds a fresh $\omega$-sequence, so $\FF(V,V[y])=\{\omega\}$;
the first step preserves cofinalities, so \cref{univ:gs:lem:tower} gives
$\FF(V,N)\subseteq\{\omega,\rho\}$. Both indices have cofinality $\omega$ in $N$, and $y$ stays new
over $V$ with finite prefixes in $V$, so $\omega\in\FF(V,N)$ and $\Char_N(\No^V)=\{\omega\}$.
Prikry forcing adds no reals, so $\R^N=\R^{V[y]}\ne\R^V$; by \cref{univ:thm:omega}, $\No^{V[y]}$ is
$\omega$-saturated in $N$ and $\No^V$ is not. Saturation is invariant under a class isomorphism with
Replacement for set images.
\end{proof}

\begin{corollary}[A strict hierarchy of information; source 11]\label{univ:gs:cor:hierarchy}
The full gap spectrum does not determine the ordered-field isomorphism type of an old surreal field:
after all its characters are known, the finite-parameter saturation boundary can supply independent
information. Conversely, fixing every saturation cardinal does not determine the full gap spectrum
(\cref{univ:gs:thm:samesat}).
\end{corollary}

There is no contradiction between the two directions. At uncountable cardinals the gap criterion
controls saturation, but at $\omega$ types may contain infinitely many formulas over a \emph{finite}
parameter set, and their real codes are information that the cardinality of a full gap forgets
(source~11).

\begin{proposition}[The simplest missing separators have different birthdays; source 11]\label{univ:gs:prop:birthdays}
In this tower $\beta(V,N)=\omega$ and $\beta(V[y],N)=\rho$, while both full gap spectra are
$\{\omega\}$.
\end{proposition}
\begin{proof}
No finite sign is new, and $y$ is a new sign of length $\omega$ over $V$. Over $V[y]$ no sign of
length below $\rho$ is added, since such a sign codes a bounded subset of $\rho$; the Prikry
generic subset of $\rho$ is new with every bounded prefix old, a fresh sign of birthday $\rho$ and
external cofinality $\omega$.
\end{proof}

The full gap spectrum remembers the external cofinality of a root but loses its birthday. The
birthday profile is richer in the ambient sign tree, but its invariance under pure ordered-field
isomorphism would need a separate argument, so it is not used as an invariant; the nonisomorphism
follows from saturation, which is intrinsic (source~11).

\subsection{Three constructions of equal spectra without isomorphism (merge comparison)}

{\small
\begin{longtable}{@{}P{0.2\textwidth}P{0.22\textwidth}P{0.22\textwidth}P{0.26\textwidth}@{}}
\toprule
& Source~12 (\cref{univ:gs:thm:boolean}) & Source~09 (\cref{univ:gs:thm:amplification}) & Source~11 (\cref{univ:gs:thm:incomplete})\\
\midrule
\endhead
Hypotheses & ZFC; Cohen forcing only & a no-new-reals extension adding a countable sequence; instances need a measurable (Prikry) or CH (Namba) & a measurable cardinal\\[3pt]
Common spectrum & $\{(\omega,\omega)\}$ & any $\Char_W(\NoM)$; $\{(\omega,\omega)\}$ or $\{(\omega,\omega),(\omega_1,\omega_1)\}$ & $\{(\omega,\omega)\}$\\[3pt]
Number of fields & $\abs{\Pow(I)^V}$ for any ground set $I$ & $2^{\aleph_0}$ & two\\[3pt]
Saturation profiles & all equal (no field is $\omega$-saturated) & $\No^{M_\omega}$ is $\omega$-saturated, the others not & one $\omega$-saturated, one not\\[3pt]
Distinguishing invariant & real trace; embeddability is exactly inclusion & real trace (rational cuts) & $\omega$-saturation\\[3pt]
Equal reals? & no & no & no\\
\bottomrule
\end{longtable}}

\noindent In all three the fields have different real traces, so none decides whether equal real
traces together with equal full spectra imply isomorphism (\cref{univ:gs:q:trace}), and all
three distinguish fields, not pure orders (\cref{univ:gs:q:order}). Source~11's example has
equal gaps but unequal saturation; source~12's has equal gaps \emph{and} equal saturation at every
cardinal, which answers source~11's question to that effect (\cref{univ:gs:q:equalspectra}).
Source~09's Theorem~7.1(d) and source~12's Theorem~1.1(iv) are consistent: they concern different
constructions, and source~12's reserved coordinate is what prevents an $\omega$-saturated member.

\section{Surcomplex real forms from families of old fields}\label{univ:gs:sec:realforms}

\emph{Throughout this section assume \textbf{(GC)}, except in \cref{univ:gs:thm:setforms}.}
Every family of class maps below is one class relation whose sections are indexed by a set;
no set whose members are proper classes, and no choice from a collection of classes of
possible isomorphisms, is used. Each theorem is an instance of the transport mechanism of
\cref{univ:thm:transport,univ:cor:agreement,univ:lem:invariant}: an isomorphism
$\Phi_A:\SC^{M_A}\to\SCN$ fixing a common set-sized subfield turns $c_{M_A}$ into an involution
$j_A=\Phi_Ac_{M_A}\Phi_A^{-1}$ of $\SCN$ with $\Fix(j_A)=\Phi_A(\No^{M_A})$, and two such
involutions are conjugate exactly when their fixed fields are isomorphic. The proofs below print
what each instance adds. A field isomorphism of real closed fields preserves the order, and as a
class map with Replacement it preserves set presentations, their sizes, and the absence of
set-sized cofinal subsets.

\begin{remark}[The class back-and-forth, six times; merge note]\label{univ:gs:rem:bfroutes}
All six sources reprove \cref{univ:thm:classbf} (source~09 Lemma~9.1, source~10 Lemma~9.1,
source~11 Lemma~11.1, source~12 Theorem~9.2 with Lemma~9.1, source~13 Lemma~8.4, source~14
Lemma~12.2), by set-valued recursion along $\Ord$ with class parameters and a set-like global
well-order: at each stage an isomorphism between set-sized (algebraically closed) subfields,
extended at a forth step by a root of a transported minimal polynomial or by an element
transcendental over the current range, which exists because the algebraic closure of a set-sized
subfield is a set (\cref{univ:lem:algebraicset}; source~12 Lemma~9.1) while the field is a proper
class; symmetric back steps; unions at limits, which stay algebraically closed because a
polynomial's coefficients occur at an earlier stage. They add these points. Source~10 fixes a
common set-sized algebraically closed subfield $E$ pointwise, and notes that a set-sized
conjugation-stable $D\subseteq\SC^V$ can be enlarged, inside the common algebraically closed class
$\SC^V$, to a set-sized conjugation-stable algebraically closed field containing $D\cup\{i\}$,
since conjugation carries roots of a polynomial to roots of the conjugate polynomial. Source~11
uses no preselected class transcendence basis and orders elements first by ambient rank. Source~13
extends any prescribed isomorphism of set-sized subfields and shows that every element is
eventually included, since the process cannot keep choosing distinct predecessors past the
Hartogs ordinal of the element's set of predecessors; its remark stresses that the lemma builds in
no order preservation and does not make old and new real closed surreal fields isomorphic.
Source~14 assumes that algebraic closures of set-generated subfields are sets and, for a
set-indexed family, retains at each stage the set of stage maps for all indices. Source~12 proves
that any two proper-class algebraically closed fields of characteristic zero are isomorphic, and
records that the theorem fails for ``real closed'': the analogous real-closure argument still
gives ample transcendence, but a real transcendental has a cut type that cannot be matched
arbitrarily, and its whole main construction exploits this difference. No source uses a class
Zorn lemma, a class-valued recursion, or an unrestricted class truth predicate.
\end{remark}

\begin{proposition}[Forgetting the involution loses the obstruction; source 12, source 11]\label{univ:gs:prop:pureacf}
In the setting of \cref{univ:gs:thm:boolean}, the pure class field $\SCN$ realizes every
complete type in finitely many variables over every $N$-set of parameters, whereas every
$(\SCN,j_A)$ of \cref{univ:gs:thm:boolinvolutions} omits a complete type over the empty set and is
not $\omega$-saturated.
\end{proposition}
\begin{proof}
The first assertion is \cref{univ:thm:purecomplex} (source~11 reproves it for every proper-class
algebraically closed field, as its Proposition~11.4). For the expansion, impose $j_A(X)=X$ and
translate $p_{r_*}$ into the fixed field, positive fixed elements being exactly the nonzero squares
of fixed elements; the parameter-free type is finitely satisfiable and omitted, and one extends it
to a complete type or transports the complete fixed-field cut type through the two-coordinate
interpretation of \cref{univ:lem:coordinates}.
\end{proof}

The obstruction is not the algebraically closed carrier but the real structure selected by the
involution; the same class field looks maximally saturated in the pure language while naming one
involution exposes specific omitted cuts, and nonconjugacy concerns the added structure, not
different isomorphism types of the pure carrier (sources 11, 12).

\subsection{Same spectrum, agreement on old constants (source 09)}

\begin{theorem}[Same-spectrum nonconjugate real forms; source 09]\label{univ:gs:thm:ampinvolutions}
In the setting of \cref{univ:gs:thm:amplification}, under Global Choice, there are involutions
$j_S$, $S\in\Pow(\omega)^M$, of the single field $\SCN$ such that (i) $\Fix(j_S)\cong\No^{M_S}$;
(ii) all fixed fields have the same full gap spectrum and none has a set-sized cofinal subset;
(iii) $j_S$ and $j_T$ are not conjugate by any class field automorphism when $S\ne T$; (iv) all
$j_S$ restrict to ordinary conjugation on the common old complex field $\C^M$. For the Prikry
instance the common spectrum is $\{(\omega,\omega)\}$; for the Namba instance it is
$\{(\omega,\omega),(\omega_1,\omega_1)\}$.
\end{theorem}
\begin{proof}
Each $\SC^{M_S}$ is algebraically closed and a proper class; extend the identity on the common
set-sized subfield $\C^M$ to $\Phi_S:\SC^{M_S}\to\SCN$ and put $j_S=\Phi_Sc_{M_S}\Phi_S^{-1}$. Then
(i), (ii) follow from \cref{univ:gs:thm:amplification}(a),(e); a conjugacy would give
$\No^{M_S}\cong\No^{M_T}$, contrary to \cref{univ:gs:thm:amplification}(b); $\Phi_S$ fixes $\C^M$
and $c_{M_S}$ acts there as usual. The spectra are \cref{univ:gs:cor:prikry,univ:gs:cor:namba}.
\end{proof}

A uniform use of the global well-order makes the construction one class relation indexed by the
set $\Pow(\omega)^M$ (source~09).

\begin{proposition}[Obstruction to agreement on all final constants; source 09]\label{univ:gs:prop:coreobstruction}
If $\R^{M_S}\ne\R^N$, no real form isomorphic to $\No^{M_S}$ contains $\R^N$; so no involution
whose fixed field is isomorphic to $\No^{M_S}$ agrees with standard conjugation on all of $\C^N$.
\end{proposition}
\begin{proof}
An inclusion of $\R^N$ into such a real form, followed by an isomorphism to $\No^{M_S}$, is an
ordered-field embedding fixing $\Q$, so each real's image realizes that real's rational cut, and
\cref{univ:gs:thm:oldtrace} would put every real of $N$ into $\R^{M_S}$. Agreement with $c_N$ on
$\C^N$ fixes $\R^N$ pointwise.
\end{proof}

The common core of \cref{univ:gs:thm:ampinvolutions} is $\C^M$, not $\C^N$, and this boundary is
necessary for this construction, not an omitted choice of isomorphism: the family answers an
infinite-family problem with a prescribed \emph{old} core, not the all-final-constants version that
\cref{univ:q:families} asks for (source~09). The families of
\cref{univ:gs:thm:boolinvolutions0,univ:gs:thm:samesatinvolutions,univ:gs:thm:coordinvolutions,univ:gs:thm:classinvolutions}
do agree on $\C^N$, because their grounds share the reals of $N$; they have different spectra.

\subsection{Agreement on any set-sized core (source 10)}

\begin{theorem}[Many equally saturated nonconjugate involutions; source 10]\label{univ:gs:thm:boolinvolutions0}
In a \textbf{(GC)} version of the extension of \cref{univ:gs:thm:booleanfamily}, let
$D\subseteq\SC^V$ be any prescribed set-sized subfield stable under ordinary conjugation. There is a
uniform family $(j_S:S\subseteq\chi)$ of involutions of $\SCN$ such that:
\begin{enumerate}[label=\textnormal{(\roman*)}]
\item $j_S$ agrees with $c_N$ on $D$; in particular all can agree on every element of $\C^N$;
\item $\Fix(j_S)\cong\No^{M_{K_S}}$, with isolated characters as in
\cref{univ:gs:thm:booleanfamily}(i), least character $\kappa_*$ and no set-sized cofinal subset;
\item the $j_S$ are pairwise nonconjugate, and none is conjugate to $c_N$;
\item all fixed fields, and all structures $(\SCN,j_S)$ in the ring language with the involution
named, have the same saturation spectrum: $\tau$-saturation exactly for infinite
$\tau\le\kappa_*$;
\item all $(\SCN,j_S)$ have the same first-order theory, and the maps can be chosen so that they
have the same theory over $D$, each element of $D$ named by a constant.
\end{enumerate}
\end{theorem}
\begin{proof}
Choose a common set-sized conjugation-stable algebraically closed $E\subseteq\SC^V$ containing
$D\cup\{i\}$ (\cref{univ:gs:rem:bfroutes}), and $\Phi_S:\SC^{M_{K_S}}\to\SCN$ with
$\Phi_S\res E=\mathrm{id}_E$, uniformly in the set parameter $S$. For $d\in D$, $d$ and $c_N(d)$ lie
in $E$, so $j_S(d)=\Phi_S(c_{M_{K_S}}(d))=c_N(d)$; no reals were added, so $\C^N=\C^V$ is an
allowed core. (ii) and the first part of (iv) transfer along $\Phi_S$. A conjugacy $j_S\sim j_T$
would give $\No^{M_{K_S}}\cong\No^{M_{K_T}}$, contrary to \cref{univ:gs:thm:booleanfamily}; $c_N$
has fixed field $\NoN$, which fills every set-sized cut, while every displayed fixed field has a
$\kappa_*$-gap. For the expansions, $F(i)$ with conjugation and $i$ named is interpreted in $F^2$
with $(a,b)(u,v)=(au-bv,av+bu)$ and $c(a,b)=(a,-b)$, and conversely the fixed field and its order
are definable (\cref{univ:lem:coordinates}); these finite interpretations preserve saturation at
infinite cardinals, a parameter set giving at most twice as many coordinates, and naming or
removing $i$ changes nothing (\cref{univ:thm:conjugation}). For (v), $E_0=E\cap\No^V$ is real
closed with $E=E_0(i)$: real and imaginary parts of elements of $E$ lie in $E$, the conjugation's
fixed field is $E_0$, and positive square roots and real roots of odd-degree polynomials over
$E_0$, which exist in $\No^V$ and are algebraic over $E$, lie in $E_0$. Each transported fixed
field contains $E_0$ with its original embedding, so quantifier elimination makes them
elementarily equivalent over $E_0$, and the coordinate translation gives equivalence of the
conjugation structures over $E$, hence over $D$.
\end{proof}

Elementary equivalence over $D$ is distinct from saturation in a language with arbitrarily many
named constants; the saturation claim in (iv) is for the fixed finite language (source~10).

\begin{corollary}[A simultaneous multicharacter answer; source 10]\label{univ:gs:cor:multichar}
The multicharacter and old-ground aspects of \cref{univ:q:families} have positive simultaneous
realizations of arbitrary set-indexed size: $2^\chi$ pairwise nonconjugate real forms share one
saturation threshold and agree on all ordinary complex numbers, while their higher gap profiles
range over every subset of $\chi$ designated coordinates.
\end{corollary}

The conjugacies and field isomorphisms are pure-field maps: nothing is asserted about the ambient
simplicity order, the natural valuation, Hahn summation, an exponential, an omnific integer part
or a derivation, and requiring any of these would define a different and potentially far more
rigid classification problem (source~10).

\subsection{\texorpdfstring{Agreement on $\C^N$ with one saturation spectrum (source 11)}{Agreement on C\^{}N with one saturation spectrum (source 11)}}

\begin{theorem}[Many nonconjugate involutions, one field; source 11]\label{univ:gs:thm:samesatinvolutions}
For the family of \cref{univ:gs:thm:samesat}, under Global Choice, there is a uniformly indexed
family $(j_\xi)_{\xi<\chi}$ of involutive automorphisms of $\SCN$ with $\Fix(j_\xi)\cong\No^{M_\xi}$
as ordered fields, $j_\xi$ ordinary conjugation on $\C^N$, $j_\xi$ and $j_\eta$ not conjugate for
$\xi\ne\eta$, and all $(\SCN,j_\xi)$ elementarily equivalent with the same saturation spectrum.
\end{theorem}
\begin{proof}
$\C^N$ lies in every $\SC^{M_\xi}$, all reals being shared; take $\Phi_\xi$ fixing it. A conjugacy
maps fixed fields isomorphically, and positivity is field-definable by
$x>0\iff x\ne0\wedge\exists y\ x=y^2$, so it would identify $\No^{M_\xi}$ with $\No^{M_\eta}$
contrary to \cref{univ:gs:thm:samesat}. With $i$ named, every $z$ is uniquely
$x+iy$ with $x=(z+j_\xi(z))/2$ and $y=(z-j_\xi(z))/(2i)$ in $\Fix(j_\xi)$, a finite-coordinate
bi-interpretation with real closed fields; completeness of $\RCF$ gives one theory, and type
translation gives saturation at each infinite cardinal; then forget $i$.
\end{proof}

These are field isomorphisms: the proof does not make them preserve the natural valuation,
birthdays, simplicity, Conway normal forms, exponentiation, a derivation or the omnific integer
ring, so the involutions are real forms of the pure surcomplex field, not classified symmetries of
an enriched structure (source~11).

\subsection{Conjugacy and equivariant embeddings in the Boolean family (source 12)}

In the setting of \cref{univ:gs:thm:boolean}, for each $A$ choose, uniformly using the global
choices, $\Phi_A:\No^{M_A}(i)\to\SCN$ and put $j_A=\Phi_Ac_{M_A}\Phi_A^{-1}$; it has order exactly two
and $\Fix(j_A)=\Phi_A(\No^{M_A})$ is real closed.

\begin{theorem}[Boolean patterns of surcomplex involutions; source 12]\label{univ:gs:thm:boolinvolutions}
(i) Every fixed field has full gap spectrum $\{(\omega,\omega)\}$ with its unique real-closed-field
order. (ii) $j_A$ and $j_{A'}$ are conjugate in $\Aut(\SCN)$ exactly when $A=A'$. (iii) There is a
unital field embedding $h:\SCN\hookrightarrow\SCN$ with $h\circ j_A=j_{A'}\circ h$ exactly when
$A\subseteq A'$. (iv) All $(\SCN,j_A)$ are elementarily equivalent.
\end{theorem}
\begin{proof}
(i) Isomorphisms of real closed fields preserve order and gap characters. (ii) A conjugacy
restricts to an isomorphism of fixed fields, hence $\No^{M_A}\cong\No^{M_{A'}}$ and $A=A'$. (iii) An
intertwining embedding restricts to a field embedding of fixed fields, forcing $A\subseteq A'$;
conversely, when $A\subseteq A'$ the inclusion $\iota:\No^{M_A}(i)\hookrightarrow\No^{M_{A'}}(i)$
commutes with conjugation and $h=\Phi_{A'}\iota\Phi_A^{-1}$ works. (iv) $u+iv\mapsto(u,v)$ represents
$F(i)$ with conjugation in two coordinates, so every sentence translates into one about a real
closed field; choosing $i$ or $-i$ changes coordinates but not truth.
\end{proof}

An equivariant embedding need not be onto, which is why a Boolean \emph{embedding} order appears
although conjugacy is equality of indices. In the one-Cohen-real case there are continuum many
nonconjugate involutions with one fixed-field spectrum; for arbitrary $I$ the number of indexed
conjugacy classes is $\abs{\Pow(I)^V}^N=(2^{\abs I})^V$, since the forcing preserves cardinals. All
statements concern set-indexed families of class maps (source~12). No preservation of the canonical
conjugation, a valuation or real constants is asserted for the maps $\Phi_A$; the real-field
counterexamples of \cref{univ:gs:sec:trace} do not need this transport.

\subsection{A continuum of equal-threshold real forms (source 13)}

\begin{theorem}[Simultaneous involutions of the coordinate family; source 13]\label{univ:gs:thm:coordinvolutions}
In \eqref{univ:gs:eq:productsetup}, under Global Choice in $N$, there are involutions $j_A$,
$A\subseteq\omega$, of $\SCN$ with $\Fix(j_A)\cong\No^{M_A}$, $j_A\res\C^N$ ordinary conjugation,
$\Char_N(\Fix(j_A))$ given by \cref{univ:gs:thm:coordinate}, and $j_A$, $j_{A'}$ conjugate by an
ambient field automorphism if and only if $A=A'$.
\end{theorem}
\begin{proof}
All intermediate grounds share the reals of $N$ (\cref{univ:gs:lem:sizes}), so each $\SC^{M_A}$
contains $\C^N=\R^N[i]$; take $\Phi_A$ fixing it, uniformly. A conjugacy would equate spectra and
\cref{univ:gs:prop:coordorder} gives $A=A'$. \Cref{univ:lem:oldbounds} excludes set-sized cofinal
subsets, a property that transfers.
\end{proof}

\begin{corollary}[A continuum of equal-threshold real forms; source 13]\label{univ:gs:cor:continuumforms}
The subfamily with $0\notin A$ has cardinality $2^{\aleph_0}$ in $N$; its members are pairwise
nonconjugate, agree on $\C^N$, and their fixed fields have saturation threshold $\kappa_0$, are
$\omega$-saturated and lack set-sized cofinal subsets, while their full gap spectra differ. One may
restrict further to infinite omitted sets without reducing the size. In this GCH construction the
continuum is $\aleph_1$.
\end{corollary}
\begin{proof}
Fixing the bit at $0$ does not change the number of subsets; distinct parameters give distinct
spectra; apply \cref{univ:gs:prop:threshold,univ:gs:thm:coordinvolutions}. Removing the countably
many finite sets leaves continuum many.
\end{proof}

For $\kappa_n=\aleph_{n+1}$ these fixed fields are all $\aleph_1$-saturated and none is
$\aleph_2$-saturated; on the infinite-omitted-set subfamily they all also have an
$\aleph_{\omega+1}$ gap, while the intervening characters encode an arbitrary subset of $\omega$
with its first bit fixed, which makes precise the information lost by recording only the least
omitted type size (source~13). All $\kappa_n$ are uncountable in this construction; otherwise the
common complex-field clause would need separate analysis (source~13). The isomorphisms $\Phi_A$
are noncanonical: they preserve field operations and the named ordinary complex field, but no
ambient simplicity, monomial map, natural valuation, summability, exponentiation or derivation,
and a real form of an algebraically closed field is not an embedding theorem in any richer
language (source~13).

\subsection{Real forms of one set field, and the class version (source 14)}

\begin{theorem}[Continuum many set-sized real structures, one saturation spectrum; source 14; $\ZFC$]\label{univ:gs:thm:setforms}
In the model $N$ of \cref{univ:gs:thm:cutoff} there are one algebraically closed field
$\Omega=(\No_{<\theta})^N[i]$ of cardinality $\theta=\aleph_{\omega+2}$, containing $\C^N$, and
involutions $j_A$, $0\notin A\subseteq\omega$, such that: (i) $\Fix(j_A)$ is isomorphic over $\R^N$
to $(\No_{<\theta})^{M_A}$; (ii) all $j_A$ restrict to ordinary conjugation on $\C^N$; (iii) no two are
conjugate by an automorphism of $\Omega$; (iv) all $(\Omega,j_A)$ are elementarily equivalent with
identical saturation spectra, although their fixed fields have different complete gap-pair spectra.
No class choice principle is needed.
\end{theorem}
\begin{proof}
Each $(\No_{<\theta})^{M_A}[i]$ is algebraically closed of characteristic zero and cardinality
$\theta$. The common algebraically closed subfield $\C^N$ has size $\aleph_1<\theta$, and an algebraic
closure of a field generated by $\C^N$ and $\tau$ transcendentals has size at most
$\max(\aleph_1,\tau)$; so both fields have transcendence degree $\theta$ over $\C^N$, and a bijection
of transcendence bases extends to $\Phi_A:(\No_{<\theta})^{M_A}[i]\to\Omega$ fixing $\C^N$. There are
set many $A$ and the isomorphisms are sets, so ordinary Choice picks them. A conjugacy would map
fixed fields isomorphically, which \cref{univ:gs:thm:cutoff} forbids. $(E[i],c)$ is interpretable in
$E$ on pairs, $E$ is definable in it by $c(x)=x$ and positivity by nonzero squares, so all expansions
have one theory; naming $i$ and translating types preserves the saturation properties, which are then
those of \cref{univ:gs:cor:setfields}.
\end{proof}

\begin{theorem}[Full-class real forms with a common threshold; source 14]\label{univ:gs:thm:classinvolutions}
Assume \textbf{(GC)} in the product extension of \cref{univ:gs:thm:coordinate} with
$\kappa_n=\aleph_{n+1}$, and the ground-class predicates of \textbf{(H)}. There are $2^{\aleph_0}$
pairwise nonconjugate involutions $J_A$ of $\SCN$, indexed by $0\notin A\subseteq\omega$, such that
$\Fix(J_A)\cong\No^{M_A}$ with the spectrum of \cref{univ:gs:thm:coordinate}; $J_A$ agrees with
standard conjugation on $\C^N$; all fixed fields are $\omega_1$-saturated but not
$\omega_2$-saturated, with the same entire saturation spectrum; none has a set-sized cofinal or
coinitial subset; and the expansions $(\SCN,J_A)$ are elementarily equivalent but pairwise
nonisomorphic as class structures.
\end{theorem}
\begin{proof}
Apply the class back-and-forth to $(\SC^{M_A},\SCN)$ over $\C^N$: the fields are proper classes,
containing every old ordinal, algebraic closures of set-generated subfields are sets, and real
closedness makes each complexification algebraically closed. The properties transfer from
\cref{univ:gs:thm:coordinate,univ:gs:prop:threshold,univ:gs:cor:onethreshold,univ:lem:oldbounds};
conjugacy would give a set-respecting isomorphism of fixed fields, excluded by
\cref{univ:gs:prop:coordorder}; first-order equivalence is the interpretation argument of
\cref{univ:gs:thm:setforms}.
\end{proof}

The class version needs an actual class construction, not a class-sized transcendence basis selected
by ordinary Choice; the set-sized theorem isolates the spectrum result from that foundational issue
(source~14). Unstructured does not mean compatible with every expansion: the isomorphisms need not
preserve the natural valuation, Conway normal forms, simplicity, the surreal exponential, a
derivation or omnific integers; the involutions are field automorphisms, not strong Hahn
automorphisms or automorphisms of a differential or exponential field, and their agreement on
$\C^N$ is the one compatibility condition proved (source~14).

\subsection{Relation to Part V and to the collection}

\Cref{univ:thm:manygrounds,univ:cor:finitetower} separated finitely many real forms by distinct
\emph{first} thresholds $\delta(M_a,N)$. The families above hold the threshold fixed and separate by
higher characters (sources 10, 11, 13, 14), or hold the \emph{whole spectrum} fixed and separate by
the real trace (sources 09, 12). The report
\path{docs/surcomplex/first-kappa-coefficients} already gives, in one universe and without forcing,
families of real forms with individually prescribed singleton thresholds
(\path{fkc:sb:thm:involutions}); the families here are of actual old fields with multicharacter
spectra and a common threshold (sources 13, 14). None of these theorems classifies real forms or
involutions; they continue, but do not settle, the automorphism report's open real-form
classification (\cref{univ:rem:whatisnew}), and none addresses \cref{univ:q:compatible}.

\section{Questions of sources 09--14, and which of them the cluster answered}\label{univ:gs:sec:questions}

The six sources pose 58 questions: 9 from source~09 (its Questions 11.1--11.9), 12 from
source~10 (Questions 1--12), 10 from source~11 (Research questions 14.1--14.10), 9 from
source~12 (Questions 12.1--12.9), 8 from source~13 (Questions 10.1--10.8) and 10 from source~14
(Questions 14.1--14.10). Questions asking the same thing are merged below, with every asking
source named; those answered by a sibling source are recorded as answered. They are directions
suggested by the results, not claims that each has been catalogued as an open problem in the
published literature.

\begin{question}[Equal full spectra, nonisomorphic fields; sources 10, 11, 13, 14; answered for fields]\label{univ:gs:q:equalspectra}
Construct, or rule out under precise hypotheses, two old surreal fields in one universe with
identical full set-presented gap spectra that are not isomorphic as ordered fields (source~10
Question~2, source~13 Question~10.5, source~14 Question~14.4, first clause); can this happen
with equal gap \emph{and} equal saturation spectra (source~11 Question~14.4); and without a
measurable cardinal (source~11 Question~14.3)? Which invariants of the arrangement or
multiplicity of cuts detect the difference (sources 10, 13)?
\end{question}
\emph{Status: answered for fields.} \Cref{univ:gs:thm:boolean} (source~12) gives such fields in
ZFC with Cohen forcing only, with equal saturation at every infinite cardinal, and
\cref{univ:gs:thm:amplification} (source~09) gives continuum many, without large cardinals under
CH (\cref{univ:gs:cor:namba}); \cref{univ:gs:thm:incomplete} (source~11) gives a pair with unequal
$\omega$-saturation from a measurable. Source~09's Prikry family also answers the
equal-saturation clause for distinct proper index sets, none of which is $\omega$- or
$\omega_1$-saturated. The detecting invariants found are the real trace (sources 12, 09) and
$\omega$-saturation (source~11); coarse multiplicities cannot help
(\cref{univ:gs:prop:ubiquity}). Incidence of set-indexed cuts, their joint realizability in
algebraic extensions, algebraic relations among parameters defining gaps and invariants of
finitely generated subfields remain candidates (sources 13, 14). The pure-order clause is
\cref{univ:gs:q:order}.

\begin{question}[Equal real traces and equal full spectra; sources 09, 12; open]\label{univ:gs:q:trace}
Suppose $M_0,M_1\subseteq N$ have the same ordinals, $\R^{M_0}=\R^{M_1}$ and
$\GapSpec_N(\No^{M_0})=\GapSpec_N(\No^{M_1})$. Under which additional hypotheses are the old
fields isomorphic? Is there a same-universe counterexample with equal real sets, and one using
only set forcing with no large cardinal (source~09 Question~11.1, source~12 Question~12.2)?
\end{question}
All three constructions of \cref{univ:gs:sec:tower}'s comparison vary the trace and do not decide
this. Cuts over a common non-Archimedean countable subfield, where real agreement no longer
controls parameterized types, are a promising next step (source~09). Forcing that adds no reals
is a natural testing ground, but adding different subsets of $\omega_1$ is not by itself a proof
of field nonisomorphism: an arbitrary field map can move the chosen infinite scale (source~12).

\begin{question}[Parameterized cut profiles and a trace hierarchy; sources 09, 12; open]\label{univ:gs:q:profile}
(a) For every finitely generated ordered subfield $F$ of an old field $\NoM$, record which cuts of
$\rc(F)$ are realized in $\NoM$, compatibly with embeddings between such subfields; with the gap
spectrum, does this profile admit a usable classification of old surreal fields (source~09
Question~11.2)? (b) Can residues of canonically specified convex valuations, or realizable cut
types over finitely generated subfields, give an intrinsic hierarchy beyond $\Tr(F)$; which
levels are invariant under all field embeddings, and which depend on naming a noncanonical scale
(source~12 Question~12.3)?
\end{question}
The trace supplies only the $F=\Q$ part of the profile; a serious formulation must account for
compatibility between parameter systems, not merely collect cardinalities (source~09). An
adequate definition should not name $\omega$ merely because one presentation has it: the
Boolean theorem succeeds because its rational reference scale is fixed by every unital
embedding (source~12).

\begin{question}[The pure-order version; sources 09, 12, 14; open]\label{univ:gs:q:order}
Can two old surreal fields with equal full gap spectra have nonisomorphic pure order reducts in
one universe, and which arrangement or multiplicity invariant of fresh cones would detect it
(source~09 Question~11.3, source~14 Question~14.4, second clause)? For the Boolean Cohen family,
which fields are isomorphic as pure linear orders; can the whole family collapse to one
pure-order type although its field embeddability order is Boolean (source~12 Question~12.1)?
\end{question}
Field arithmetic anchors the rational cuts in every construction above; an arbitrary order
isomorphism has no obligation to preserve them, so a solution needs an order-intrinsic
obstruction or an order back-and-forth that does not rely on field saturation. This is the most
direct unresolved part of \cref{univ:q:completeness} (sources 09, 12).

\begin{question}[Embeddings rather than isomorphisms; sources 12, 14; open]\label{univ:gs:q:embeddings}
(a) Outside the Cohen family, characterize the pairs with $\R^{M_0}\subseteq\R^{M_1}$ for which
$\No^{M_0}\hookrightarrow\No^{M_1}$, and the extra compatibility of omitted cuts that is needed
(source~12 Question~12.4). (b) Can every ground set-sized partial order be represented by
embeddability of old fields with one common real trace and one common full spectrum; even a
two-element antichain (source~12 Question~12.5)? (c) For the coordinate family of
\cref{univ:gs:thm:coordinate}, what is the exact condition for a class ordered-field embedding
$\No^{M_A}\to\No^{M_{A'}}$; does an initial-image condition or a fixed common Hahn core recover
inclusion information (source~14 Question~14.5)?
\end{question}
\Cref{univ:gs:thm:oldtrace} gives necessity, not a converse; a positive theorem would have to
extend embeddings through larger old sign domains without meeting an unmatchable cut, and (b)
asks for a new mechanism rather than another arrangement of independent Cohen reals
(source~12). In (c) the traces coincide, and an embedding need not preserve omitted gaps as
omitted gaps, since its image can acquire separators; so the isomorphism argument does not settle
the embedding order (source~14).

\begin{question}[Small-forcing base change beyond a common dense suborder; source 09; open]\label{univ:gs:q:basechange}
Can the cardinality hypothesis of \cref{univ:gs:thm:basechange} be replaced by a chain-condition,
approximation or Boolean-algebra invariant while keeping both directions, and what is the sharp
failure mechanism at $\cf\gamma=\abs{\mathbb Q}$ (source~09 Question~11.4)?
\end{question}
\Cref{univ:gs:cor:densesub} already allows the size of a common dense suborder; the canonical-name
proof shows where a stronger theorem would need a replacement for the cofinal pigeonhole step.

\begin{question}[Equal spectra and agreement on $\C^N$; source 09; open]\label{univ:gs:q:finalconstants}
Can a no-new-reals family of intermediate bases give infinitely many nonconjugate real forms with
identical entire gap spectra, all agreeing with standard conjugation on $\C^N$ (source~09
Question~11.5)?
\end{question}
\Cref{univ:gs:prop:coreobstruction} rules out rechoosing the isomorphisms in the Cohen family; new
invariants with the same real set are needed. The families of sources 10, 11, 13 and 14 agree on
$\C^N$ but have different spectra.

\begin{question}[Family sizes; sources 11, 14 answered; source 09 open]\label{univ:gs:q:size}
(a) Can one force $2^\chi$ pairwise nonisomorphic old fields in one universe with a fixed
uncountable saturation threshold, all containing the ambient reals (source~11 Question~14.2)?
(b) For fixed regular $\kappa$ and set cardinal $\chi'$, when are there $\chi'$ pairwise
nonisomorphic old real closed fields in one universe with first gap $\kappa$, the same real field
and explicitly different higher spectra (source~14 Question~14.3)? (c) Can arbitrarily large
set-indexed families of pairwise nonisomorphic old fields share one fixed \emph{full} spectrum in
one universe while retaining a common prescribed real field (source~09 Question~11.6)?
\end{question}
\emph{Status.} (a) is answered by \cref{univ:gs:thm:booleanfamily} (source~10): $2^\chi$ fields,
threshold $\kappa_*$, all containing $\R^V=\R^N$. (b) is answered for every uncountable regular
$\kappa$ and every infinite $\chi'$ by \cref{univ:gs:thm:samesat} (source~11), and with $2^\chi$
fields when $\kappa>\chi$ by source~10; the case $\kappa=\omega$ is not covered by these
constructions.
(c) stays open: a larger Cohen poset changes the range
where base change is automatic, and names for reals are no longer coded by one old real, so more
coordinates alone do not suffice (source~09).

\begin{question}[Birthday information; sources 09, 10, 11, 12, 13, 14; open]\label{univ:gs:q:birthday}
For each gap character, which birthdays do fresh roots of that cofinality have, which profiles
occur when several ground regulars collapse to one external cofinality, and which parts survive an
abstract field isomorphism rather than depending on the ambient sign tree (source~09
Question~11.7, source~11 Question~14.5)? Can a pure-field invariant recover the least fresh
length or first new birthday, or would naming simplicity or a birthday comparison be needed, and
how much (source~10 Question~3, source~13 Question~10.6)? Can two old fields have the same
pushed-forward spectrum but different least fresh sign lengths, and is there a useful subclass
where lengths are recovered (source~14 Question~14.6)? Which information about fresh-sign
\emph{lengths} can be recovered intrinsically (source~12 Question~12.7, second half)?
\end{question}
This continues \cref{univ:q:length}. \Cref{univ:gs:prop:birthdays} gives equal spectra with
different least new birthdays, but for nonisomorphic fields; a decisive negative answer would need
two pairs with different least new birthdays and \emph{isomorphic} old fields, and a positive one
needs more than the sign representation, which is not in the field language (source~13). The
coincidence of first birthday and $\kappa_*$ in \cref{univ:gs:thm:booleanfamily} is a feature of
that construction, not a reconstruction theorem (source~10). Prikry and Namba give explicit tests
(source~09); \cref{univ:gs:thm:cohen} shows a complicated ground spectrum becoming a singleton
externally (sources 12, 14).

\begin{question}[Spectra beyond Easton products; sources 09, 10, 11, 12; open]\label{univ:gs:q:beyondeaston}
For a specific forcing with known fresh-function spectrum and changed cofinalities, compute the
complete pushforward \eqref{univ:gs:eq:transfer}, and determine which distinct fresh-function
spectra give the same surreal spectrum (source~09 Question~11.8). Which pairs
$(\FF(M,N),\lambda\mapsto\cf^N\lambda)$ occur, and can a nontrivial collapse pattern and a
multicharacter spectrum be prescribed together while keeping the reals (source~10 Question~5)?
Which sets occur as $\Char_N(\NoM)$ for cofinality-preserving set forcing (source~11
Question~14.6)? Given a ground spectrum, which identifications among its members can the
pushforward make (source~12 Question~12.7, first half)?
\end{question}
The transfer reduces the surreal question exactly to these realization problems but supplies no
general realization theorem; collapse forcings and further Prikry-type forcings are natural
candidates, with the generic and the cofinality map specified separately (sources 09, 10, 11).

\begin{question}[Relative spectra at accumulation points; sources 10, 11, 13, 14; answered for countable sequences]\label{univ:gs:q:accumulation}
For the fields of \cref{univ:gs:thm:booleanfamily}, determine $\Char_N(\No^{M_{K_S}})$ at every
regular cardinal: is it $\Ecl^V(\{\kappa_*\}\cup\{\kappa_i:i\in S\})$, or can kept coordinates
change the answer (source~10 Question~1)? For infinite omitted sets $B$ in the sparse
construction, determine $\Char_N(\No^{M_{A\setminus B}})$ (source~11 Question~14.1). Extend the
coordinate calculation to longer Easton products with several regular or singular accumulation
points (source~13 Question~10.3), and characterize $\FF(M[G_A],M[G_I])$, $A\subseteq I$, for a
general index set $I$ (source~14 Question~14.1).
\end{question}
\emph{Status (merge).} For index length $\omega$ this is answered by \cref{univ:gs:thm:coordinate}
(sources 13, 14): source~10's family with $\chi=\omega$ is an instance, and its spectra are
$\{\kappa_*\}\cup\{\kappa_i:i\in S\}\cup\{\lambda^+:S\text{ infinite}\}$, which is exactly
$\Ecl^V(\{\kappa_*\}\cup\{\kappa_i:i\in S\})$, so the answer to source~10's question is yes for
$\chi=\omega$ (\cref{univ:gs:ex:omegafamily}); likewise infinite omitted sets forming an
$\omega$-sequence are covered. Open for uncountable index sets, where heads may be large
products, several accumulation points interact, regular-limit and Mahlo phenomena need the
chain-condition analysis of the published theory, and the support convention at inaccessible
limits must be kept (sources 13, 14). Any proof must treat the quotient as a product of old posets,
not a product recomputed in the intermediate model (source~10), and writing $\Ecl(B)$ proves
nothing (source~11).

\begin{question}[Weaker hypotheses than GCH; sources 10, 13, 14; open]\label{univ:gs:q:gch}
Find the weakest cardinal-arithmetic or forcing hypotheses for the simultaneous families
(source~10 Question~4); hypotheses under which \cref{univ:gs:thm:coordinate} still holds, and its
replacement when products collapse cardinals or $\lambda^{\aleph_0}>\lambda^+$ (source~13
Question~10.4); in particular whether local assumptions such as $2^{<\kappa_n}=\kappa_n$ and a
true-cofinality hypothesis at $\lambda$ suffice (source~14 Question~14.2).
\end{question}
The single-coordinate theorem needs no GCH; the obstacles are preservation and localization
(small low products, a chain condition for the low part including the test coordinate, a closed
high tail), the $\lambda^+$ ceiling and the singular-product witness, while transfer and density
are unaffected (sources 10, 13, 14). A changed forcing spectrum need not change the gap spectrum,
since outer cofinalities can merge its entries (source~13).

\begin{question}[Approximation beyond initial segments; source 10; open]\label{univ:gs:q:approximation}
Absence of fresh functions is an initial-segment approximation condition. Which intrinsic
expansions or configurations of surreal cuts detect the stronger approximation properties defined
by all small ground subsets, and is there a converse reconstruction (source~10 Question~6)?
\end{question}

\begin{question}[Accumulation points and large-cardinal structure; source 10; open]\label{univ:gs:q:mahlo}
Give a direct surreal interpretation of the distinction between non-Mahlo regular accumulation
points, Mahlo accumulation points and singular-successor characters in the Easton closure; can the
relevant organization of gap families be recognized internally (source~10 Question~7)?
\end{question}
No large cardinal is needed for the Boolean family, whose test coordinates are successors; the
question concerns what becomes visible when regular-limit structure is included (source~10).

\begin{question}[Compatible and enriched structure; sources 10, 11, 12, 13, 14; open]\label{univ:gs:q:enriched}
Which of the involutions constructed here can also preserve prescribed valuation data, an omnific
integer structure, a surreal exponential or a derivation, perhaps on more than the ordinary
complex subfield or on an independently chosen set-sized differential subfield; is there an
obstruction already for a small differential or exponential core (source~10 Question~8, source~11
Question~14.8, source~13 Question~10.7, source~14 Question~14.8)? Which parts of the
constructions survive naming exponentiation, a derivation, a valuation, a monomial cross-section
or an omnific subring (source~13 Question~10.7, source~14 Question~14.7)? What extra assumptions
become available for involutions preserving a valuation, integer part or exponential (source~12
Question~12.6, second half)?
\end{question}
This continues \cref{univ:q:compatible,univ:q:exponential}. The pure-field back-and-forth uses
only algebraic and transcendental one-element extensions and supplies no extension property in
these categories; merely transporting the field involution does not transport a preassigned
differential or exponential structure. The ordered-field lower bounds remain necessary for
saturation of an expansion, but sufficiency used real-closed-field quantifier elimination
(sources 10, 11, 13, 14). A tractable first problem is to fix one set-sized expanded
substructure and prove the extension lemma it requires (source~13).

\begin{question}[Class-theoretic strength; sources 10, 11, 12, 13, 14; open]\label{univ:gs:q:classstrength}
Over a theory weaker than $\GBC$ with Global Choice, what is the exact strength of the relative
class-field isomorphism principle and of its uniform set-indexed version; can the particular
surreal fields here, or the Cohen-coded family, admit the needed isomorphisms or definable maps
without it (source~10 Question~9, source~11 Question~14.7, source~12 Question~12.8, source~13
Question~10.8, first half, source~14 Question~14.9)?
\end{question}
This continues \cref{univ:q:classiso}. The set-forcing and sign arguments need no Global Choice;
the set-valued recursion avoids class truth recursion but not the global choices of enumerations
and extensions; a positive answer should give an actual uniform class construction, not infer one
from local set-saturation (sources 10, 11, 14). \Cref{univ:gs:thm:setforms} isolates the spectrum
result from this issue.

\begin{question}[Proper-class spectra and class forcing; source 10; open]\label{univ:gs:q:classforcing}
Under which tameness and preservation assumptions can class forcing produce a prescribed
proper-class pattern of gap characters, and can a class-indexed isolated-coordinate argument be
carried out while preserving Replacement (source~10 Question~10)?
\end{question}
The forcings here are sets, and \cref{univ:gs:thm:filterbase} gives a set bound; a formal union of
set-forcing constructions would not establish the class version.

\begin{question}[Richer inclusion patterns and symmetry; source 10; partly answered]\label{univ:gs:q:inclusion}
Which partial orders, incidence structures or group actions are realized by inclusions and field
automorphisms of such families with one first threshold, and can a refined invariant encode more
than isolated membership bits (source~10 Question~11)?
\end{question}
\emph{Status.} \Cref{univ:gs:cor:poset} (source~12) realizes every ground set-sized partial order by
inclusion, and exactly by field embeddability, among old fields with one full spectrum and
threshold $\omega$; \cref{univ:gs:prop:coordorder} (source~13) adds binary intersections for the
coordinate family. Group actions, incidence structures and uncountable thresholds remain open, and
the inclusion order must still be distinguished from intersections and composita (source~10).

\begin{question}[Intersections and composita; source 13; open]\label{univ:gs:q:intersections}
For $\seq{A_m:m<\omega}$, determine $\bigcap_m\No^{M_{A_m}}$ and when it equals
$\No^{M_{\bigcap_mA_m}}$, in particular for tails of coordinates (source~13 Question~10.1). Inside
$\NoN$, compare the compositum $\No^{M_A}\No^{M_{A'}}$ and its real closure with
$\No^{M_{A\cup A'}}$: which elements are algebraic over the compositum, and when is the inclusion
an equality (source~13 Question~10.2)?
\end{question}
The binary intersection proof does not iterate through a countable limit, where a reduced product
or a common singular-limit quotient may obstruct; a first target is the intersection of the models
restricted to sets of ordinals. A finite rational expression in elements of two old fields need not
capture every set-theoretic combination of their coordinates; a criterion via names with
restricted coordinate dependence is an intermediate goal (source~13).

\begin{question}[Fixed-field invariants for involutions; source 12; open]\label{univ:gs:q:involutions}
Does the pair of the real trace and the full gap spectrum classify any broad, naturally defined
family of involutions of $\SC$ up to conjugacy (source~12 Question~12.6, first half)?
\end{question}
The involutions here are deliberately unconstrained, and the transport maps preserve no analytic
or normal-form data; the questions about enriched automorphisms are separate
(\cref{univ:gs:q:enriched}).

\begin{question}[Product versus iteration calculus; source 11; open]\label{univ:gs:q:iteration}
Develop a relative spectrum calculus distinguishing fixed ground products from iterations whose
next factor is recomputed in the intermediate universe (source~11 Question~14.9).
\end{question}
\Cref{univ:gs:lem:tower} gives one inclusion; a converse needs a criterion for whether intermediate
fresh prefixes were already present earlier. Finite towers of this report
(\cref{univ:cor:finitetower}) are iterations; the families of \cref{univ:gs:sec:relative} are
products.

\begin{question}[Formalization; sources 09--14; open]\label{univ:gs:q:formal}
Formalize the fresh-sign correspondence, block codes, transfer and cutoff decomposition first for an
abstract initial predicate on a set-sized sign tree or two set-sized transitive models; then the
canonical-name lemma, the filter-base and density bounds, name localization and intersections as
abstract forcing lemmas with the truth lemma as an explicit interface; then the Boolean
embeddability theorem from real-trace preservation and coordinate detection with its forcing
hypotheses kept visibly external; and the class-isomorphism layer as a separate theorem with an
explicit Global Choice interface (source~09 Question~11.9, source~10 Question~12, source~11
Question~14.10, source~12 Question~12.9, source~13 Question~10.8, second half, source~14
Question~14.10). Which existing ProveIt and Mathlib interfaces support this?
\end{question}
The sources propose staged routes: ordinal prefix operations, cofinal subsets and a cut-filling
axiom for the sign layer; atomic absoluteness for a fixed poset, graph-name restriction, nice-name
coding and the truth lemma for the forcing layer (source~09); a sign tree with the three values
$-<0<+$, then ordinal block combinatorics with the two cofinality computations kept separate, then
forcing, then the ground family and class maps, with set-sized models in a larger universe first
(source~10); published forcing imports kept separate from the independently checkable sign and
cofinality lemmas, and the final spectrum equalities not introduced as axioms (source~11); residues
represented by Dedekind cuts, not proper-class equivalence classes, and any change from genuine old
fields to an abstract family of real closed fields exposed (source~12); the advanced product facts as
named hypotheses rather than unadvertised axioms (sources 13, 14). None of them is a same-universe
Lean type of all surreals, no Lean declaration or build is included or claimed, and importing one
checked surreal theorem would not certify the forcing above it (sources 10, 13).

\section{What sources 09--14 do not claim}\label{univ:gs:sec:nonclaims}

Every limitation, disclaimer and priority caveat of the six sources is kept, per source; items
written for the merge are marked M. The shipped audit files of each source
(\cref{univ:gs:sec:audit}) state their boundaries in their own words.

\subsection*{Source 09}
\begin{enumerate}[label=09.\arabic*,leftmargin=3.2em]
\item AI-assisted manuscript with written proofs; not independently refereed or Lean-verified;
priority in the wider literature not certified; open to independent proof review, with the
canonical-name lemma and the simultaneous field construction the priority targets.
\item The fresh-sign correspondence and symmetry are the repository's, reproved and credited, not
claimed; the transfer is not claimed as a new forcing-spectrum theorem; FKW Propositions~4.3, 4.10,
5.2, Definitions~4.7, 4.8, Theorems~4.24, 7.21, 7.23 and Corollary~4.25 are imported; the Easton
result is an application, not a solution of FKW's realization problems, which remain untouched.
\item No assertion of first discovery of the common-base-change technique; the lemma keeps the poset
fixed and is not a chain-condition assertion; the bound is strict; the countable boundary is not
claimed to behave like uncountable cofinalities.
\item No classification of all old surreal fields, of pure orders or of all forcing spectra; no claim
that all gap spectra are Easton closed; no assertion that fields with the same reals and gap spectra
must be isomorphic; nonisomorphism is a field claim, not one about pure order reducts.
\item The Namba instance needs CH and the Prikry instance a measurable; not a consistency proof of
$\ZFC$ or of a large-cardinal configuration.
\item The involutions agree on $\C^M$; agreement on $\C^N$ is impossible for most members and is not
claimed; no exponential, differential, valuation-preserving or simplicity-preserving isomorphism is
claimed, and the pure surcomplex isomorphisms generally preserve none of these.
\item Global Choice for the simultaneous class isomorphisms, as one indexed class relation; for
general pairs the spectrum can be a proper class.
\item The last row of its comparison table is deliberately not strengthened to equality of
ground-regular spectra.
\item The finite checks test coding only: no generic extension, freshness of an infinite sequence or
class recursion is modelled; successful compilation is not a proof check.
\item Its questions are research directions, not certified novelty statements.
\end{enumerate}

\subsection*{Source 10}
\begin{enumerate}[label=10.\arabic*,leftmargin=3.2em]
\item Proposed contributions with written proofs; not refereed, not Lean-verified; no priority
certification; the content may overlap unlocated consequences of existing literature.
\item The gap correspondence, class back-and-forth and freshness predate it; the Easton spectrum is
an imported published theorem (FKW Definition~4.8, Propositions~4.3, 4.10, Theorem~4.24), and
Easton's lemma is used in Cody--Gitman's form.
\item No proof that equal full spectra imply isomorphism, and no explicit pair with equal full spectra
but nonisomorphic fields is asserted (supplied by sources 09, 11, 12; M).
\item No full-spectrum formula at accumulation points of intermediate quotients; only isolated
characters are computed; ground Cohen posets are identified with recomputed ones only at the singleton
endpoint where proved.
\item GCH for the Easton and simultaneous constructions; $V=L$ as a concrete relative setting; no large
cardinals; set forcing only; GBC with Global Choice for class maps, as a uniformly indexed relation.
\item Saturation is claimed for the fixed finite language; elementary equivalence over a core is a
separate assertion; the inclusion order does not identify intersections or composita.
\item No preservation of valuation, simplicity, exponentiation, derivation, Hahn summation or omnific
integer parts by the class-field maps; no solution of a general forcing-spectrum conjecture; not a
characterization of all cofinality-preserving spectra.
\item The finite tests check coding and bookkeeping only; no Lean declarations or build.
\end{enumerate}

\subsection*{Source 11}
\begin{enumerate}[label=11.\arabic*,leftmargin=3.2em]
\item AI-assisted manuscript with conventional proofs; not refereed; no Lean; not a certified priority
claim; independent review remains appropriate; novelty not certified by the search.
\item Its source audit was targeted: the recursive tree, collection and report guides and the
\emph{opening} of this report's LaTeX source; not a proof audit of the repository.
\item Imported: surreal, quantifier-elimination and forcing background; FKW Definition~4.8,
Proposition~4.10, Theorem~4.24 and Theorem~7.21; Lévy--Solovay; the Prikry facts.
\item The finite-complement theorem does not extend to infinite complements without proof; the sparse
hypothesis is local terminology.
\item The Prikry forcing in the tower is computed in the Cohen extension; the tower needs a measurable;
nonisomorphism is via $\omega$-saturation, not birthdays, whose field invariance is not claimed.
\item No claim that equal gaps \emph{and} equal saturation suffice for, or obstruct, isomorphism
(answered by source~12; M).
\item The involutions preserve no valuation, birthdays, simplicity, normal forms, exponential,
derivation or omnific ring; they are real forms of the pure field.
\item No named conjecture is solved; its review priorities (its §§8 and 11) are kept; the finite checks
do not verify forcing, freshness, cofinality, saturation or class arguments.
\end{enumerate}

\subsection*{Source 12}
\begin{enumerate}[label=12.\arabic*,leftmargin=3.2em]
\item A proposed negative answer to the \emph{real-closed-field} direction of
\cref{univ:q:completeness} only, not to the pure-order question; not refereed or Lean-formalized;
priority not certified; not a solution of a named published conjecture.
\item Credited ingredients: Conway, Gonshor, the repository's Part~II, FKW Proposition~2.2,
Corollary~2.4, Propositions~5.2--5.3 and 6.1, and the Mitchell and Unger antecedents via FKW; the
Archimedean residue is standard.
\item Forcing statements are relative, not claims that ZFC proves the existence of transitive models or
generics; no large cardinal, CH or GCH for the main theorem; Global Choice only for the transports.
\item ``Universality'' is only realization of ground set-sized posets, not Borel or effective
reducibility; the trace classifies embeddability only within the family, says nothing about equal
traces, and is not claimed finitary definable or efficiently computable.
\item The maps $\Phi_A$ preserve no conjugation, valuation or real constants; equivariant embeddings are
not surjective.
\item No repository file was modified and no formalization compiled; the literature check is not
exhaustive; the finite checks cover local identities only, and there are no finite fresh signs.
\end{enumerate}

\subsection*{Source 13}
\begin{enumerate}[label=13.\arabic*,leftmargin=3.2em]
\item Written proofs with an internal audit; not refereed, not Lean-verified; to be reviewed by a set
theorist familiar with product forcing and class model theory.
\item FKW Proposition~4.10 and Corollary~4.21 (Shelah) are imported and not reproved; Proposition~6.1 and
Theorems~7.21, 7.23 are used only in the comparison table.
\item Contributions may be folklore consequences of known results; no first-in-literature breakthrough,
no classification of all real forms or all forcings, no solution of the general spectrum-realization
problem.
\item Classification only within the displayed family; binary intersections only, no composita, no
arbitrary intersections; the first new birthday is not recovered from the pure field language.
\item The continuum is $\aleph_1$ in the construction; uncountable $\kappa_n$ are needed for the
complex-core clause; the isomorphisms are noncanonical and preserve no richer structure.
\item ``The user is not represented as having reviewed or endorsed these proofs.''
\item The finite checks do not verify freshness, genericity, cofinality, preservation, the
successor-of-singular input, saturation or class arguments; in particular not the $\lambda^+$ term.
\end{enumerate}

\subsection*{Source 14}
\begin{enumerate}[label=14.\arabic*,leftmargin=3.2em]
\item Proposed contribution with written proofs; not refereed, not Lean-verified; priority of the
relative formulations not certified; not to be represented as an independently certified
breakthrough.
\item The singular-limit witness (FKW Corollary~4.21, Shelah 2000) is imported and not verified; no new
pcf theorem; the fresh-sign correspondence and the saturation criterion are credited.
\item GBC with Global Choice and ground predicates for the class version; the set version uses ZFC
only.
\item No preservation of valuation, normal forms, simplicity, exponential, derivation or omnific
integers; the involutions are field automorphisms.
\item No general equal-spectrum classification, no arbitrary spectrum realization, no elimination of
the singular-successor character, no arbitrarily large continuum, no consistency proof.
\item Its review order (the localization rectangle, the retention of the singular witness, the
fresh-versus-boundary distinction) is kept; the finite checks exclude forcing, pcf, saturation and
class recursion; a TeX build is not proof review.
\end{enumerate}

\subsection*{Added by the merge}
\begin{enumerate}[label=M.\arabic*,leftmargin=3.2em,start=6]
\item Part VIII adds no mathematical claim beyond comparisons stated as merge observations (the
instance of \cref{univ:gs:thm:coordinate} in \cref{univ:gs:ex:omegafamily}, and question statuses
derived from sibling theorems); proofs printed once from a base are the base's.
\item The corrections of \cref{univ:gs:rem:overstatement1,univ:gs:rem:overstatement2} concern how two
sources describe this report, not their mathematics.
\item The source manuscripts were not rebuilt for the merge; their page counts are as delivered and
match their build records.
\item Nothing in Part VIII is formalized, and placement beside the collection's Lean development
confers no formal status.
\end{enumerate}
This gives 10, 8, 8, 6, 7 and 6 items from sources 09--14 and 4 added by the merge: 49 in all.
(The merge items continue the numbering M.1--M.5 of \cref{univ:sec:nonclaims}.)

\section{Provenance, repository comparison and finite checks of sources 09--14}\label{univ:gs:sec:audit}

\subsection{Review scope of the sources}

All six state the pin \gsfull{} (25 September 2026) and were written on 28 September 2026.
\begin{itemize}
\item Source~09 read this report's \texttt{article.tex} and README at the pin, and the collection's
reader's guide \path{Algebra/SurrealNumbers/docs/README.md}; it cites the labels
\path{univ:thm:gapcorrespondence}, \path{univ:cor:gapspectrum}, \path{univ:thm:classbf},
\path{univ:q:cutspectrum} and \path{univ:q:completeness}, all present.
\item Source~10 used the report at the fixed commit, especially Part~II, Part~V and
Questions~19.2, 19.4, 19.7 with the batch-32 status note.
\item Source~11 inspected the recursive tree, the surreal subproject guide, the collection guide, this
report's guide and the opening of its LaTeX source (targeted, not exhaustive; shipped
\path{11-finite-complements-SOURCE_AUDIT.md}).
\item Source~12 used the pinned article and \path{Algebra/SurrealNumbers/README.md} (shipped
\path{12-real-traces-SOURCES.md}).
\item Source~13 used this report's README, its statements on Cohen forcing and its Section~19, and
\path{docs/surcomplex/first-kappa-coefficients} for comparison (shipped
\path{13-coordinate-spectra-SOURCES.md}).
\item Source~14 used the article and README, with attention to the fresh-sign correspondence, the
single-Cohen example and Questions 19.2, 19.4, 19.7, and the neighbouring report overview.
\end{itemize}
Each read Fischer--Koelbing--Wohofsky \cite{FKW} in its author-hosted version; source~11 checked
Lévy--Solovay \cite{LevySolovay}, source~14 Shelah \cite{Shelah2000}.

\subsection{Repository statements checked against the current collection}

This report is unchanged between the pin and the placement commit \texttt{0e53d1034}, so no source
statement about it is stale in that sense. Checked:
\begin{enumerate}
\item The labels and numbers the sources cite exist and mean what they say: ``Theorem~1.1''
(\cref{univ:thm:main-summary}), ``Theorem~7.1'' (\cref{univ:thm:omega}), ``Theorems~8.4--8.5''
(\cref{univ:thm:gapcorrespondence,univ:thm:fillingcone}), ``Proposition~13.2''
(\cref{univ:prop:add}), ``Theorem~15.1'' (\cref{univ:thm:classbf}), Questions~19.2, 19.4, 19.5, 19.7
(\cref{univ:q:cutspectrum,univ:q:completeness,univ:q:length,univ:q:families}).
\item True at the pin: that \cref{univ:prop:add} computed only the first Cohen gap and left larger
gaps open (sources 12, 13, 14); that \cref{univ:q:families} asks for agreement on $\C^N$ (source~09);
that its batch-32 note separates single-character forms from multicharacter spectra and families of
old fields (source~10); that first-kappa-coefficients already treats singleton thresholds
(source~13); and that the report's own source pin (\shorthash) differs from the audit pin
(source~13). These gaps are now filled by Part VIII; the dated statuses record it.
\item Overstated: source~09's ``implication proposed as a possible classification principle in the
repository'' (\cref{univ:gs:rem:overstatement1}) and source~10's ``suggestion that the least gap size
\ldots\ could classify the fields'' (\cref{univ:gs:rem:overstatement2}); the report only asked.
\item Source~11's shipped audit mentions a ``neighboring critical-point-defects guide''; the file is
\path{docs/foundations-and-computation/large-cardinal-embeddings-and-normal-forms/07-critical-point-defects-PROOF_STATUS.md},
a proof-status file rather than a guide, and no claim from it is used.
\end{enumerate}

\subsection{Finite checks}

Every source shipped a standard-library Python program and a recorded run. There are no finite fresh
signs, so none tests freshness, forcing, cofinality, preservation, saturation or class arguments.
All six were rerun for the merge on a copy of this directory with Python~3.14.4, with explicit output
paths; each output equals its recorded file apart from line endings.
\begin{itemize}
\item \path{code/09-cohen-amplification-verify_finite.py}
(\path{data/09-cohen-amplification-verification_results.json}): 260{,}865 canonical-cone cases,
260{,}865 comparison-antisymmetry cases, 5{,}704 block round trips, 183{,}103 block-prefix cases and
three symbolic image examples (finite Easton, Prikry, Namba under CH).
\item \path{code/10-isolated-coordinates-check_finite.py}
(\path{data/10-isolated-coordinates-verification_results.json}): 5{,}461 delimiter-code round trips,
36{,}409 prefix checks, 3 rejected inputs, 16{,}065 cone equivalences, 511 profile checks and 87{,}381
reversed-inclusion pairs.
\item \path{code/11-finite-complements-finite_checks.py}
(\path{data/11-finite-complements-results.json}): 469{,}107 checks, including 65{,}025 order
comparisons, 65{,}025 cone equivalences and 330{,}496 first-disagreement checks.
\item \path{code/12-real-traces-verify_finite.py}
(\path{data/12-real-traces-verification_results.json}): 65{,}025 dyadic comparisons, 65{,}025
prefix-separator pairs, 5{,}698 block round trips, 129{,}400 prefix dependencies, 2{,}047 ternary
round trips, 11 injectivity checks, and the principal-down-set embedding for all labelled partial
orders on at most four points (1, 1, 3, 19, 219 of them; 3{,}688 pair checks).
\item \path{code/13-coordinate-spectra-verify_finite.py}
(\path{data/13-coordinate-spectra-verification_results.json}): 64{,}897 separator-cone pairs,
1{,}365 block round trips, 22{,}302 prefix-locality checks, 5 rejected encodings, 7{,}533
condition extensions and 16{,}384 coordinate-mask pairs.
\item \path{code/14-relative-cutoffs-verify_finite.py}
(\path{data/14-relative-cutoffs-finite_checks.json}): 619{,}026 assertions, including 260{,}865
separator-cone checks and the finite and eventually periodic index-set identities behind the
coordinate formulas.
\end{itemize}

\subsection{Dependency ledger}

{\small
\begin{longtable}{@{}P{0.27\textwidth}P{0.4\textwidth}P{0.25\textwidth}@{}}
\toprule
Result & Essential inputs & Status\\
\midrule
\endhead
Transfer (\cref{univ:gs:thm:transfer}) & Part~II; compression; block codes & written proof; related to FKW Props.~5.2--5.3\\[3pt]
Filter-base ceiling & truth lemma; cofinal pigeonhole & written proof\\[3pt]
Generalized Cohen & closure; universal occurrence; ceiling & written proof; four routes\\[3pt]
Square chain condition & bounded maximal antichain & FKW Prop.~2.2, reproved\\[3pt]
Easton realization & transfer; FKW Thm.~4.24, Prop.~4.10 & imported input\\[3pt]
Isolated coordinates, finite complements, cubes & localization; high closure; density; FKW preservation & written relative proofs\\[3pt]
Countable coordinate spectra & localization; intersections; density; FKW Cor.~4.21 & written proof; named import for $\lambda^+$\\[3pt]
Cutoff gap pairs & birthday bound; canonical chains; name counting & written proof\\[3pt]
Real trace, Boolean family & standard part; Cohen quotients; coordinate detection & written proof; ZFC\\[3pt]
Base change, amplification & atomic absoluteness; nice names; trace & written proof\\[3pt]
Prikry, Namba spectra & FKW Thms.~7.21, 7.23; transfer & imported inputs\\[3pt]
Cohen--Prikry tower & Lévy--Solovay; Prikry spectrum; tower bound; $\omega$ criterion & written proof; measurable\\[3pt]
Real-form families & class back-and-forth (Part~V); transport; interpretation & Global Choice, except the set version\\
\bottomrule
\end{longtable}}

\noindent The ledger records the checks made in assembling the arguments; it is not a
proof-assistant log.

\appendix
\section{Provenance and repository comparison}\label{univ:app:audit}

\subsection{Pinned review scope of the sources}

All four sources were written against
\[
 \repo,\quad\text{commit }\fullhash,
\]
(``Index the newly placed drafts without extending review claims''), a commit
that precedes the batch-20 rewrite of the automorphism report. They read the repository
through a connector and did not modify or build it.
\begin{itemize}
\item Source~01 read \path{docs/README.md}, the opening 180 lines of the
foundations report, the returned opening text of the automorphism report, and
the root README; its scoped code searches used the connector's default-branch
index. It states that this is not a file-by-file audit.
\item Source~03 read the main README, the catalogue, the foundations README, and
the sections of the automorphism report on the unrestricted field, class
back-and-forth and genuinely different real forms.
\item Source~04 inspected the recursive file tree, the root description, the
catalogue, and the foundations and automorphism sources, not every manuscript
in full; a repository-wide code search returned an incomplete result that it did
not treat as exhaustive.
\item Source~08 read the root README, the documentation inventory, the
automorphism report's README, and selected source sections including the class
back-and-forth, the involution conjugacy criterion and the countably cofinal
real form.
\end{itemize}

\subsection{Statements corrected against the current collection}

\begin{enumerate}
\item \emph{Searches for ``forcing'' and ``saturation''.} Source~01 reports that
searches for \texttt{forcing}, \texttt{saturation} and \texttt{distributivity}
``returned no matches''; source~08 that searches for \texttt{forcing} and
\texttt{saturation} ``returned no hits''. At the pin, and now, the foundations
report contains the phrase ``saturation for small cuts'' in its discussion of
completeness and a subsection ``ETR, class forcing, and two invalid moves''
(\path{found:sub:etr}). Neither contains the results of this report, and both
sources presented their searches as limited. Source~04's statement that the
foundations report's forcing references did not supply the cross-universe
spectrum is accurate.
\item \emph{Report counts.} Source~01 describes \path{docs/README.md} as listing
forty catalogued reports plus four newly placed drafts. That was the state at the
pin. The collection now has 44 report directories, including this report and
its sibling.
\item \emph{The automorphism report.} It was rewritten after the pin (commit
\texttt{fb182ea}, the batch-20 write), and its labels gained the prefix
\path{saut:}. The sources' descriptions of what they read in it (``opening
text'', ``the newly placed manuscript'', selected sections) refer to the pinned
version. This report cites its current numbers and labels: Theorem~9.1
(\path{saut:thm:acfhomogeneity}), Section~10 (\path{saut:sec:realforms}),
Theorem~10.2 (\path{saut:thm:inequivalentrealform}), Proposition~B.1
(\path{saut:prop:conjugacycriterion}) and Section~14. Source~04's remark
that the automorphism source contained no matches for \texttt{forcing},
\texttt{saturat} or \texttt{nonconjugate} is true of the pinned text; the
current text still has no \texttt{saturat} and no \texttt{nonconjugate}, and
uses ``forcing'' once only in the unrelated sense of forcing an equation.
\item \emph{A missed prior result.} Sources 03 and 08 correctly credit the
collection with class back-and-forth, the fixed-field conjugacy criterion and
the countably cofinal real form, but none of the four sources cites the
set-cut interpolation criterion (present at the pin as
\path{prop:conjugacycriterion}, now \path{saut:prop:conjugacycriterion}).
\Cref{univ:rem:whatisnew} records what it implies and what remains new.
\item \emph{Formalized degeneracy.} Source~01's statement that the internal
degeneracy of set-indexed convergence is ``discussed and formalized'' in the
repository is accurate: \path{docs/FORMALIZATION.md} marks the
subset, convergent-net and Cauchy-net clauses of \path{found:thm:discrete}
as proved, the last under \texttt{HasSmallCutFillers}
(\cref{univ:rem:foundations}).
\end{enumerate}

\subsection{Literature boundary}

Conway and Gonshor supply the construction and cut theory
\cite{Conway,Gonshor}; van den Dries--Ehrlich with their erratum, and
Bournez--Guilmant, the birthday-field background \cite{vdDE,vdDEerr,BournezGuilmant};
Marker the quantifier-elimination and saturation machinery \cite{Marker};
Ehrlich's absolutely saturated models the class-saturation background, whose
terminology concerns all sets of the ambient universe and is not contradicted by
the use of a larger universe \cite{Ehrlich1989}; van den Dries--Ehrlich's
homogeneous universal $H$-fields and Ehrlich--Kaplan's surreal ordered
exponential fields the richer structures \cite{vdDE2018,EhrlichKaplan}; Jech,
Prikry, Gitik and Benhamou the forcing facts \cite{Jech,Prikry,Gitik,Benhamou};
Fischer--Koelbing--Wohofsky the forcing notion of freshness \cite{FKW};
Berenbeim the reverse-simplicity question and the suggestion to study surreals
under forcing \cite{Berenbeim}; Holy and coauthors the class-forcing completion
theory \cite{HolyEtAl}. Rangel--Mariano and Bertram connect surreal
constructions with set-theoretic foundations; their themes and questions are not
claimed or resolved here \cite{RM,Bertram}. Each source's targeted literature
search did not locate the exact theorems; none is an exhaustive bibliography.

\section{Dependency and proof-obligation ledger}\label{univ:app:ledger}

{\small
\begin{longtable}{@{}P{0.27\textwidth}P{0.36\textwidth}P{0.29\textwidth}@{}}
\toprule
Result & Essential inputs & Boundary checked\\
\midrule
\endhead
Old arithmetic and elementarity & Absolute signs, prefix witnesses for old cuts, set-like Conway recursion; RCF/ACF QE & Field elementarity is not universe elementarity.\\[3pt]
Ground-set containment & Birthday bounds; common ordinals; old bounded sign sets & The containing set need not be small.\\[3pt]
One old side & Ground-set containment; ground cut theorem & Both sides external sets; one literally in $M$.\\[3pt]
Regular first length & Minimal new sequence; old coding maps & Cardinality and cofinality computed in $N$; $M$-regularity by persistence.\\[3pt]
Interval tree and decoder & Set-length recursion; explicit splitting; two cuts at limits & Decoder total, with a default value, in $M$.\\[3pt]
Absolute interval code & Per-input set recursion; absolute arithmetic and cuts & No set of all children; labels range over $\Ord$.\\[3pt]
Block codes & Two-bound cone identity; set-recursive decoding & Ordinal, not natural, sums; successor runs.\\[3pt]
Main spectrum & Minimal new sequence; decoder; one-side lemma; RCF QE & Saturation equivalence only for uncountable $\kappa$.\\[3pt]
First gap & Regularity; nested endpoints; one-side interpolation & Cofinalities in $N$ with set witnesses.\\[3pt]
Gap correspondence & Prefix convexity; simplest separators; constant-tail closure & Set-presented gaps only; no Global Choice.\\[3pt]
Least gap & Old-object coding; terminal-plus blocks & Sizes computed in $N$.\\[3pt]
First new birthday & Characteristic signs; minimality; ordinal coding in $M$ & $\beta$ is not a birthday of a given number.\\[3pt]
$\omega$-saturation & Countable syntax over finite old parameters; rational cuts & No new reals is not no new countable sequences.\\[3pt]
Pure surcomplex saturation & Ground envelopes; algebraic closures of set fields; ACF QE & No class-sized transcendence basis or global isomorphism.\\[3pt]
Conjugation spectrum & Fixed field; order by nonzero squares; coordinates & Naming $i$ adds one parameter.\\[3pt]
Bounded fragments & Birthday cut bound; sign-cone lemma & Regularity in $N$ separate from equality of short signs.\\[3pt]
External Cauchy rigidity & Old positive lower bounds & Fine surreal radii; no cut filling needed.\\[3pt]
Closed, nowhere dense & Open cones; old radii below new ones & Needs $M\ne N$ for nowhere-denseness.\\[3pt]
Class isomorphism, real forms & Global Choice; set-stage back-and-forth; gap invariance & No valuation, simplicity or exponential preserved.\\[3pt]
Set-complete Boolean algebra & Cylinder maps; common levels for set families & Least-level set pairs, not class objects.\\[3pt]
No class-complete algebra & Global Choice; proper-class antichain; diagonal subclass & Stronger completeness, extra axiom.\\
\bottomrule
\end{longtable}}

The ledger records the mathematical checks made in assembling the arguments. It
is not a proof-assistant log.

\section{Finite checks}\label{univ:app:checks}

Three sources shipped standard-library Python programs; source~01 shipped none.
(The checks of sources 09--14 are described in \cref{univ:gs:sec:audit}.)
They check finite sign comparisons and interval formulas. There are no finite
fresh signs, since a finite sequence of old entries is always old, so none of
them tests freshness, forcing, regularity, saturation or any class argument.

\begin{itemize}
\item \path{code/03-fresh-sign-gaps-check_finite_signs.py} checks
first-disagreement comparison, convexity of finite prefix cones and the
terminal-plus block code with its endpoint ordering, for value lists of length
at most four with entries in $\{0,1,2,3\}$ and cones among signs of length at
most six. Recorded run
(\path{data/03-fresh-sign-gaps-finite_checks.json}): 127 sign strings, 16{,}129
ordered-pair comparisons, 127 cone checks, 340 block-code cases and 1{,}744
nested endpoint pairs, all passing. For the list $(2,0,1)$ the code is
$\codeplus=+{+}{+}{-}\,{+}{-}\,{+}{+}{-}$, with endpoint lengths $(2,3)$, $(4,5)$,
$(7,8)$ and values $2<\tfrac52<\tfrac{43}{16}<\tfrac{173}{64}<\tfrac{87}{32}<\tfrac{11}4<3$.
\item \path{code/04-universe-saturation-verify_finite_models.py} checks, in exact
rational arithmetic, the successor interval formulas \eqref{univ:eq:childab},
the closed formulas of \cref{univ:prop:finiteformula}, separation of children
with finite labels, unique finite decoding, and the cut/prefix equivalence of
\cref{univ:lem:signcone} for finite strings. Recorded run
(\path{data/04-universe-saturation-finite_check_results.txt}): 1{,}365 interval
nodes, 20{,}475 child pairs, 1{,}364 decodings, 65{,}025 cut/prefix comparisons.
\item \path{code/08-forcing-saturation-verify_sign_code.py} checks sign order
against dyadic values, the two-bound identity of \cref{univ:lem:twobounds},
decoding, nesting and finite cone intersections for $\code(f)$. Recorded run
(\path{data/08-forcing-saturation-verification_results.txt}): 304{,}000 checks.
For $(2,0,1)$ the code is ${+}{+}{-}{+}\,{-}{+}\,{+}{-}{+}$ with value
$\tfrac{219}{128}$ and intervals
$(\tfrac32,2)\supset(\tfrac{13}8,\tfrac74)\supset(\tfrac{109}{64},\tfrac{55}{32})$.
\end{itemize}
All three suites were rerun for this merge on a copy of the directory, with
Python~3.14; each output agrees with its recorded file except for line endings
on Windows. Commands are given in the README.

\clearpage
\phantomsection
\addcontentsline{toc}{section}{References}
\begin{thebibliography}{99}

\bibitem{Conway}
J. H. Conway, \emph{On Numbers and Games}, 2nd ed., A K Peters, 2001 (first
edition, Academic Press, 1976).

\bibitem{Gonshor}
H. Gonshor, \emph{An Introduction to the Theory of Surreal Numbers}, London
Mathematical Society Lecture Note Series 110, Cambridge University Press, 1986.

\bibitem{vdDE}
L. van den Dries and P. Ehrlich, ``Fields of surreal numbers and
exponentiation,'' \emph{Fundamenta Mathematicae} \textbf{167} (2001), 173--188.
\url{https://doi.org/10.4064/fm167-2-3}.

\bibitem{vdDEerr}
L. van den Dries and P. Ehrlich, ``Erratum to `Fields of surreal numbers and
exponentiation','' \emph{Fundamenta Mathematicae} \textbf{168} (2001),
295--297. The corrected statements are also restated in
\cite[Theorems 1.1--1.2 and 3.3]{BournezGuilmant}.

\bibitem{BournezGuilmant}
O. Bournez and Q. Guilmant, ``Surreal fields stable under exponential and
logarithmic functions,'' arXiv:2201.08199, 2022.
\url{https://arxiv.org/abs/2201.08199}.

\bibitem{Ehrlich1989}
P. Ehrlich, ``Absolutely saturated models,'' \emph{Fundamenta Mathematicae}
\textbf{133} (1989), 39--46. \url{https://doi.org/10.4064/fm-133-1-39-46}.

\bibitem{vdDE2018}
L. van den Dries and P. Ehrlich, \emph{Homogeneous universal $H$-fields},
arXiv:1807.08861, 2018. \url{https://arxiv.org/abs/1807.08861}.

\bibitem{EhrlichKaplan}
P. Ehrlich and E. Kaplan, \emph{Surreal ordered exponential fields},
arXiv:2002.07739, version 3, 2021. \url{https://arxiv.org/abs/2002.07739}.

\bibitem{Marker}
D. Marker, \emph{Model Theory: An Introduction}, Graduate Texts in Mathematics
217, Springer, 2002. \url{https://doi.org/10.1007/b98860}.

\bibitem{Jech}
T. Jech, \emph{Set Theory: The Third Millennium Edition, Revised and Expanded},
Springer Monographs in Mathematics, Springer, 2003.
\url{https://doi.org/10.1007/3-540-44761-X}.

\bibitem{Prikry}
K. L. Prikry, ``Changing measurable into accessible cardinals,''
\emph{Dissertationes Mathematicae (Rozprawy Matematyczne)} \textbf{68} (1970),
55~pp.

\bibitem{Gitik}
M. Gitik, ``Prikry-type forcings,'' in M. Foreman and A. Kanamori (eds.),
\emph{Handbook of Set Theory}, Springer, 2010, 1351--1447.
\url{https://doi.org/10.1007/978-1-4020-5764-9_17}.

\bibitem{Benhamou}
T. Benhamou, ``Prikry forcing and tree Prikry forcing of various filters,''
arXiv:1801.04424, 2018. \url{https://arxiv.org/abs/1801.04424}.

\bibitem{FKW}
V. Fischer, M. Koelbing and W. Wohofsky, ``Fresh function spectra,''
\emph{Annals of Pure and Applied Logic} \textbf{174} (2023), no.~9, 103300.
\url{https://doi.org/10.1016/j.apal.2023.103300}.

\bibitem{Berenbeim}
A. Berenbeim, \emph{Review of Surreal Number Construction and Forcing
Construction}, research notes dated 8 April 2020, 35 pp. Question 2, p.~32;
proposed further directions, p.~34. Listed by the author as incomplete notes.
\url{https://www.berenbeim.com/papers}.

\bibitem{HolyEtAl}
P. Holy, R. Krapf, P. L\"ucke, A. Njegomir and P. Schlicht, ``Class forcing,
the forcing theorem and Boolean completions,'' arXiv:1710.10820, 2017.
\url{https://arxiv.org/abs/1710.10820}.

\bibitem{RM}
D. Rocha Rangel and H. L. Mariano, \emph{An algebraic (set) theory of surreal
numbers, I}, arXiv:1911.12726, 2019. \url{https://arxiv.org/abs/1911.12726}.

\bibitem{Bertram}
W. Bertram, \emph{On Conway's Numbers and Games, the Von Neumann Universe, and
Pure Set Theory}, arXiv:2501.04412, version 2, 2025.
\url{https://arxiv.org/abs/2501.04412v2}.

\bibitem{RepoAutomorphisms}
\emph{Automorphisms of the Surcomplex Numbers}, report in \repo,
\path{docs/surcomplex/surcomplex-field-automorphisms/}, label prefix
\path{saut:}. The sources read its pinned version at \shorthash.

\bibitem{EhrlichBSL}
P. Ehrlich, ``The absolute arithmetic continuum and the unification of all
numbers great and small,'' \emph{Bulletin of Symbolic Logic} \textbf{18} (2012),
no.~1, 1--45. \url{https://doi.org/10.2178/bsl/1327328438}. Cited by sources
09--14 of Part~VIII.

\bibitem{Hamkins}
J. D. Hamkins, ``Gap forcing,'' \emph{Israel Journal of Mathematics}
\textbf{125} (2001), 237--252. Preprint \url{https://arxiv.org/abs/math/9808011}.
Cited by source~09.

\bibitem{CodyGitman}
B. Cody and V. Gitman, ``Easton's theorem for Ramsey and strongly Ramsey
cardinals,'' \emph{Annals of Pure and Applied Logic} \textbf{166} (2015),
no.~9, 934--952. \url{https://arxiv.org/abs/1209.1133}; the form of Easton's
lemma used is Lemma~3.8 of \texttt{arXiv:1209.1133v1}. Cited by source~10.

\bibitem{LevySolovay}
A. L\'evy and R. M. Solovay, ``Measurable cardinals and the continuum
hypothesis,'' \emph{Israel Journal of Mathematics} \textbf{5} (1967), 234--248.
\url{https://doi.org/10.1007/BF02771612}. Cited by source~11.

\bibitem{Shelah2000}
S. Shelah, ``Embedding Cohen algebras using pcf theory,'' \emph{Fundamenta
Mathematicae} \textbf{166} (2000), no.~1--2, 83--86.
\url{https://doi.org/10.4064/fm-166-1-2-83-86}. Cited by source~14, through
\cite[Corollary~4.21]{FKW}.

\end{thebibliography}
\end{document}
